\documentclass[11pt,oneside,openany]{book}
\usepackage[T1]{fontenc}
\usepackage[utf8]{inputenc}
\usepackage{lmodern,microtype,amsmath,amssymb,amsthm,mathtools,mathrsfs,bm}
\usepackage[a4paper,margin=25mm,headheight=15pt]{geometry}
\usepackage{array,booktabs,longtable,tabularx,enumitem,xcolor,graphicx}
\renewcommand{\tabularxcolumn}[1]{>{\raggedright\arraybackslash}p{#1}}
\usepackage{xurl}
\usepackage[hypertexnames=false,colorlinks=true,linkcolor=blue!40!black,citecolor=blue!40!black,urlcolor=blue!40!black]{hyperref}
\usepackage{bookmark,fancyhdr}
\hypersetup{pdftitle={Shadow Theory and Quantum Measurement: Source Dynamics, Event Laws, and Physical Records},pdfauthor={Research programme of Jeremy Rodgers},pdfsubject={Integrated monograph, version 2: two constitutive completions}}
\pagestyle{fancy}\fancyhf{}\fancyhead[L]{\small Shadow Theory and Quantum Measurement}\fancyhead[R]{\small Integrated edition}\fancyfoot[C]{\thepage}
\setlength{\emergencystretch}{3em}
\setlist{itemsep=2pt,topsep=5pt}
\setcounter{tocdepth}{1}
\newtheorem{theorem}{Theorem}[chapter]
\newtheorem{lemma}[theorem]{Lemma}
\newtheorem{proposition}[theorem]{Proposition}
\newtheorem{corollary}[theorem]{Corollary}
\theoremstyle{definition}
\newtheorem{definition}[theorem]{Definition}
\newtheorem{assumption}[theorem]{Assumption}
\newtheorem{example}[theorem]{Example}
\newtheorem{counterexample}[theorem]{Counterexample}
\theoremstyle{remark}\newtheorem{remark}[theorem]{Remark}
\providecommand{\E}{\mathbb E}
\providecommand{\Prb}{\mathbb P}
\providecommand{\PP}{\mathbb P}
\providecommand{\QQ}{\mathbb Q}
\providecommand{\tr}{\operatorname{tr}}
\providecommand{\TV}{d_{\mathrm{TV}}}
\providecommand{\dd}{\,\mathrm d}
\providecommand{\id}{\operatorname{id}}
\providecommand{\ket}[1]{\lvert#1\rangle}
\providecommand{\bra}[1]{\langle#1\rvert}
\providecommand{\pos}[1]{[#1]_+}
\providecommand{\norm}[1]{\left\lVert#1\right\rVert}
\providecommand{\proj}[1]{\ket{#1}\bra{#1}}
\providecommand{\ii}{\mathrm i}
% Safe in the original book preamble. The additions also provide these
% defensively, so including this file does not create duplicate definitions.
\providecommand{\pLaw}{\operatorname{Law}}
\providecommand{\pVar}{\operatorname{Var}}
\providecommand{\pSupp}{\operatorname{supp}}
% Reuse the original \TV = d_{\mathrm{TV}} as a symbol: \TV(P,Q).
% No packages, theorem environments or counter changes are required.

% Load in the book preamble. No theorem environment is redefined.
\providecommand{\mcH}{\mathcal H}
\providecommand{\mcS}{\mathcal S}
\providecommand{\mcNcdf}{\mathsf F}
\providecommand{\mcNtail}{\overline{\mathsf F}}

\title{%
  \bfseries Shadow Theory and Quantum Measurement:\\
  Source Dynamics, Event Laws, and Physical Records\\[6mm]
  {\large Two Constitutive Completions}\\[2mm]
  {\normalsize A controlled Bell-path limit and a
  massive-configuration alternative}
}

\author{%
  Jeremy Rodgers\thanks{Website:
  \href{https://everythingequation.com}{everythingequation.com}}\\
  \small Independent Researcher
}

\date{%
  15 September 2026\\[2mm]
  \small Integrated monograph, version 2
}

\begin{document}
\frontmatter
\maketitle

\chapter*{Abstract}
\addcontentsline{toc}{chapter}{Abstract}
This monograph develops the quantum measurement programme associated with
Shadow Theory's distinction between source and readout. It preserves the
programme's detector, preparation, protection, record and continuation
theorems, together with the counterexamples that delimit their domains,
and establishes two distinct constitutive completions.

In the pilot-medium completion, a primitive canonical bond action produces
individual Hamiltonian currents. Conservative exporters create signed
packets; a deterministically moving, independently prepared spatial gas
supplies a complete marked-contact history; finite opposite-packet
recombination suppresses balanced service; and carrier population tracking
produces the residence denominator. Under the declared interaction
catalogue P1--P4 and preparation assumptions, these mechanisms converge in
total variation on complete tagged paths, in unchanged physical time, to
the minimal Bell process. The proof covers nodes, reversals and null
intervals on a fixed finite graph and horizon. A static finite clock
Hamiltonian incorporates source, fuel, pending and loss states, physical
copies, retained reset receivers, an inaccessible reference and
noncommuting continuation. Monomial copy cuts are crossed exactly once on
the first pass, giving actual sampled-history faithfulness. Finite members
have controlled deviations from the limiting event law, and the clock's
recurrence bounds the storage claim.

The massive-configuration completion instead postulates a continuous
configuration guidance law and complete initial equilibrium. Within its
smooth finite domain it provides conservative actual motion, an exact
massive pointer writer, retained null and loss branches, fixed historical
archives, transported-trap reset and spatial feedback. A finite autonomous
controller approximates the driven programme. Pointer derivative and
absolute surface-current estimates control historical corruption in
addition to complete retained-output error. This constitution supplies an
alternative measurement theory, not a derivation of discrete Bell jumps.

The resulting internal resolution is conditional: the pilot theory gives
controlled effective Bell dynamics and material records within its
enlarged constitution; the massive theory gives a separate operational
completion. Neither derives its interaction and preparation postulates
from source/readout incompleteness. Exact finite-resource Bell dynamics,
a smooth realization of the entire hybrid source, and broader preparation
and storage domains remain identified extensions. All comparisons specify
their output spaces, retained systems and positive-probability conditions.

\chapter*{Conventions, provenance and scope}
\addcontentsline{toc}{chapter}{Conventions, provenance and scope}
An \emph{internal resolution} means a mathematically closed
assumption-to-prediction chain inside an explicitly declared physical
constitution. In the pilot theory the Bell law is a controlled effective
limit, with the material programme included in the same ordinary
configuration graph. This usage makes no claim of empirical confirmation,
universal preparation of equilibrium, or deduction of P1--P4 from the older
source/readout premise. The massive alternative has its own postulated
guidance and equilibrium laws. No theorem exchanges axioms between them.

A constitutive assumption specifies an interaction, an ontology or a
statistical resource. A theorem derives consequences under named
assumptions. A comparison transfers only the outputs that its common
measurable space includes. Counts are random atomic measures; their
expectations and predictable compensators are different objects. Reduced
population evolution does not itself select an actual event generator.
At a coarse aperture the directed traffic is a sum of positive microscopic
currents; taking the positive part after summation requires a separate
directional-alignment condition.

Throughout the discrete chapters, $J_{nm}$ is current into configuration
$n$ from $m$, $w_m$ is its coherent weight, and $N_{nm}$ denotes an event
count when supplied with these indices. The scalar $N$ in the pilot
construction is the number of carriers. Normalized service intensities
are denoted by $\Phi$, while $\lambda$ denotes a conditional event
intensity. The incidence matrix $B$ is not an apparatus operator in those
chapters. In the massive chapters $q$ and $Q_t$ are continuous material
coordinates and their actual path; internal coherent keys are not
additional discrete actual occupancies. Symbols local to an apparatus are
defined again at that apparatus. The pilot clock's integer position and
the massive controller's continuous position are distinct variables in
distinct theories.

For probability laws, $\TV(P,Q)=\sup_A|P(A)-Q(A)|$; for density operators,
$D(\rho,\sigma)=\tfrac12\|\rho-\sigma\|_1$. Four conclusions must be kept
separate: complete state distance, endpoint distribution distance,
faithfulness to a specified actual past label, and total variation on
entire physical-time paths. The first two do not imply the latter two.
Whenever conditional laws are compared, both conditioning probabilities
are positive. In particular, $\TV(P,Q)<P(E)$ with $P(E)>0$ guarantees
$Q(E)>0$ and permits Lemma~\ref{found:errors}; taking a minimum with one
cannot define a conditional law on a null event.

An inaccessible reference remains in the joint state. Any memory,
classical key, loss product or reset receiver that can return belongs to
the complete experiment. The complete post-interaction joint state is
retained; an imperfect operation does not insert an ideal daughter. General
statements fix the graph, input class, programme and physical horizon
before resource limits are taken. Uniformity over growing programmes
requires additional estimates.

References M01--M30 and C01--C03 identify the inherited manuscripts and
checkpoints; F01--F07 identify their foundation sources. Their substantive
results retain their assumptions. The two resolution papers' essential
arguments are integrated below. The supplied AI-generated audit was used
as working material, after verifying its manuscript hashes, and supplies
no external validation authority. Established constructions are attributed
to primary literature; no historical-priority claim is made.

\chapter*{Roadmap and principal conclusions}
\addcontentsline{toc}{chapter}{Roadmap and principal conclusions}
The opening chapters distinguish information completion from statistical
selection and define complete experiments. Part~\ref{part:kinetic}
establishes the canonical action, conservative export and the conditional
signed-queue kinetic theorem. These are the comparison layer for the
physical construction in Part~\ref{part:pilot}; the kinetic proof is given
once, rather than repeated inside that construction.

The pilot argument has three successive comparisons. The spatial gas
approximates the full marked contact process; finite recombination
approximates excess-only signed service; and the kinetic theorem
approximates the complete tagged Bell path. Theorem~\ref{p:main}
combines them on one output space. The positive part comes from fast
recombination, the weight denominator from residence tracking, and the
limiting timing law from independent spatial preparation. The elementary
contact rules contain no Bell ratio. Chapter~\ref{p:chapter-records}
then uses that same complete path law for autonomous actual copies,
null/loss resources and retained continuation. Appendix~\ref{p:smooth}
gives a smooth realization of isolated contacts, with its asynchronous
source limitation explicit.

Parts~\ref{part:access} and \ref{part:statistics} explain why the enlarged
constitution is necessary to this conclusion. Readable pilot histories,
reciprocal writers, balanced traffic, cycle reassignment and non-Markov
timing remain valid countermodels when their respective interactions or
statistics are admitted. Variational selection, chamber stirring and
MPBT remain conditional routes. Their successful implications are
retained; their premises are not silently used to justify the pilot
interaction catalogue.

Parts~\ref{part:continuous}--\ref{part:records} develop preparation
consistency, finite detector instruments, protected complete transport,
configuration records and historical archives. Their event and null
constitutions are kept explicit. Part~\ref{part:massive} introduces the
separate massive theory, with complete definitions and proofs through
Theorem~\ref{mc:closure}. Its controller estimate includes derivatives
and absolute surface flux because ordinary wave distance alone cannot
prove an archive's history.

Part~\ref{part:preparation} preserves the conditional apparatus
preparation results and their access and returning-memory boundaries.
The finite discrete chain in Chapter~\ref{int:chain-chapter} remains a
useful limiting-model calculation. The final synthesis in
Chapter~\ref{syn:chapter} states exactly how the pilot implementation
supplies a physical approximation and how the massive alternative
provides its own chain. Stronger supplied-copy repairs, chamber limits,
and finite benchmarks remain in the appendices; their different resources
and constitutions are not merged.

\begin{center}\small
\begin{tabularx}{\textwidth}{@{}>{\raggedright\arraybackslash}p{.20\textwidth}XX@{}}\toprule
Question & Pilot medium & Massive configuration\\\midrule
Actual ontology & Predesignated carrier on a complete ordinary graph; additional retained pilot medium & Continuous material positions; internal spinor keys\\
Selected path & Minimal Bell process as a full-path TV limit & Exact postulated guidance flow\\
Records & Actual monomial copies on the finite clock's first pass & Massive pointer paths and held spatial archives\\
Autonomy & Exact finite ordinary Hamiltonian; hybrid pilot dynamics & Smooth semibounded controller approximating a fixed driven programme\\
Statistical input & Independent spatial gas and admitted carrier preparation & Complete initial squared-norm equilibrium\\
Unproved deeper implication & Interaction/access catalogue and preparation from weaker premises & Guidance and equilibrium from weaker premises\\\bottomrule
\end{tabularx}
\end{center}
\tableofcontents
\mainmatter


% Begin consolidated source: parts/00_foundations.tex
\chapter{Source, readout and the statistical target}
\label{found:chapter}

\section{What the source/readout premise supplies}
Let $\Omega$ be a space of candidate source states and let a declared equivalence remove gauge or descriptive redundancy. Write $S$ for the resulting reduced source space. A readout is a surjection $p:S\to T$. A source relation $a:S\to A$ is relevant to a stated question. This elementary specification is deliberately prior to a Hilbert space, a probability measure or a stochastic process. The applicable foundation results are the descent and target-completion theorems of \cite{F01,F02,F03}, and the dynamical closure theorems of \cite{F05,F06}. Their finite versions suffice here.

\begin{theorem}[Descent and the coarsest target completion]\label{found:descent}
There is a map $\bar a:T\to A$ with $a=\bar a\circ p$ if and only if $a$ is constant on every fiber of $p$. It is then unique. For any family $(a_\alpha)_{\alpha\in I}$, put
\[
 e(s)=(p(s),(a_\alpha(s))_{\alpha\in I}),\qquad E=e(S).
\]
Every surjective representation $f:S\to F$ from which $p$ and every $a_\alpha$ can be recovered factors uniquely through $e$: there is a unique surjection $h:F\to E$ with $e=h\circ f$.
\end{theorem}
\begin{proof}
If $a=\bar a\circ p$, points with the same readout have the same $a$ value. Conversely define $\bar a(t)=a(s)$ for any $s\in p^{-1}(t)$; fiber constancy makes the definition independent of the choice. Surjectivity gives uniqueness. In the second assertion write the recovery maps as $p=u\circ f$ and $a_\alpha=v_\alpha\circ f$. Then $h(z)=(u(z),(v_\alpha(z))_\alpha)$ lies in $E$ because $f$ is onto, and has the asserted factorization. It is onto because $e$ is onto and unique because $f$ is onto.
\end{proof}

This minimality orders retained information. It neither minimizes jump counts nor chooses a probability measure on $E$. For example, if $S=\{0,1\}\times\{a,b\}$, $p$ retains only the first coordinate and the nominated relation retains the second, the coarsest completion is all of $S$. Infinitely many transition generators can act on that same completed space. No representative-selection question distinguishes those generators.

\begin{proposition}[Dynamical closure]\label{found:closure}
For a deterministic source evolution $\Phi_t:S\to S$, an autonomous readout map $U_t$ satisfying $p\Phi_t=U_tp$ exists exactly when $p\Phi_t$ is constant on fibers of $p$. If this holds for all $t$ and $\Phi$ is a semigroup, then $U$ is a semigroup.
\end{proposition}
\begin{proof}
The existence assertion is Theorem~\ref{found:descent} applied to $a=p\Phi_t$. For the semigroup property, $U_{s+t}p=p\Phi_{s+t}=p\Phi_s\Phi_t=U_sU_tp$; surjectivity of $p$ gives $U_{s+t}=U_sU_t$.
\end{proof}

For a finite supplied Markov source with generator $L$, the corresponding criterion is equality of the total rates from any two source states in the same fiber to each destination fiber. It is a condition on a generator already given, not a recipe selecting that generator. Part~\ref{part:statistics} proves the version needed for event coarse-graining. The geometric realization theory \cite{F04} starts from a specified action and a declared reduction; its orbit space need not be a smooth manifold. A discrete readout therefore does not, by itself, exclude a continuous underlying source. The bulk/brane example in \cite{F07} similarly pushes forward an already supplied source-history law. These scope facts prevent geometric or informational incompleteness from silently becoming a probability postulate.

\section{Hamiltonian current and actual events}
Fix a finite orthogonal resolution $(P_n)_{n=1}^d$ of a coherent Hilbert space and a self-adjoint Hamiltonian $H(t)$. The sector spaces may be degenerate and may contain an inaccessible reference. For a normalized solution of $i\hbar\dot\Psi=H\Psi$, define
\begin{equation}\label{found:current}
 \Psi_n=P_n\Psi,\quad w_n=\|\Psi_n\|^2,\quad
 J_{nm}=\frac{2}{\hbar}\operatorname{Im}\langle\Psi_n,H\Psi_m\rangle.
\end{equation}
Throughout this text $J_{nm}$ is current \emph{into} $n$ from $m$; $H_{nm}=P_nHP_m$. Direct differentiation gives
\begin{equation}\label{found:continuity}
 J_{nm}=-J_{mn},\qquad \dot w_n=\sum_{m\ne n}J_{nm},\qquad
 |J_{nm}|\le\frac{2\|H\|}{\hbar}\sqrt{w_nw_m}.
\end{equation}
Indeed $\dot w_n=2\operatorname{Re}\langle\Psi_n,-iP_nH\Psi/\hbar\rangle$, and the $m=n$ term is real before multiplication by $-i$. Self-adjointness proves antisymmetry; Cauchy--Schwarz proves the bound.

The actual-event target is a law on right-continuous sector histories $X\in D([0,T],\{1,\ldots,d\})$. If $N_{nm}$ counts $m\to n$ events, a predictable intensity $\lambda_{nm}$ means
\begin{equation}\label{found:compensator}
 N_{nm}(t)-\int_0^t1_{\{X_{s-}=m\}}\lambda_{nm}(s\mid\mathcal F_{s-})\,ds
 \quad\text{is a local martingale.}
\end{equation}
The filtration $\mathcal F$ is part of the claim. Exposing a microscopic timer can enlarge it and change a conditional intensity even when the same natural-history law remains.

Bell's minimal generator is
\begin{equation}\label{found:bell}
 \lambda^{\rm B}_{nm}(t)=\frac{[J_{nm}(t)]_+}{w_m(t)}\quad(w_m(t)>0).
\end{equation}
This is an established construction \cite{BellQFT}; its physical selection is the present problem. Equation~\eqref{found:continuity} is not Equation~\eqref{found:compensator}. In particular a random counting measure is atomic whereas $J_{nm}(t)dt$ is an absolutely continuous signed measure. They are not equal path by path. Expected directed traffic $F_{nm}=\mathbb E[1_{X=m}\lambda_{nm}]$ may satisfy $F_{nm}-F_{mn}=J_{nm}$ under a statistical current-realization law. That law introduces probability-bearing content. Under it the surplus family $F_{nm}=[J_{nm}]_++s_{nm}$, $s_{nm}=s_{mn}\ge0$, remains; history freedom can remain even when $s=0$.

For time-dependent Markov rates, first-exit survival while the origin stays $m$ is
\[
 S_m(t_0,t)=\exp\left[-\int_{t_0}^t\sum_{n\ne m}\lambda_{nm}(u)\,du\right].
\]
A constant hazard gives an exponential waiting time in physical time. In general only the integrated hazard threshold is exponential. No substitution of exposure time for physical time is made when claiming \eqref{found:bell}.

\section{The predictive-current completion}\label{found:predictive-section}
A concrete use of the foundation is possible without choosing a jump law. Let $\mathcal H$ now be finite dimensional and let $H_a$, $a\in\mathcal A$, be a finite list of admitted constant control Hamiltonians. Set
\[
 C^a_{nm}=\frac{P_nH_aP_m-P_mH_aP_n}{i\hbar},\qquad
 \mathcal L_a^*O=\frac{i}{\hbar}[H_a,O].
\]
Let $\mathcal W$ be the smallest real subspace of Hermitian operators containing every $P_n,C^a_{nm}$ and invariant under every $\mathcal L_a^*$. It is computed by repeatedly adding the images of a basis under these maps; dimension is at most $(\dim\mathcal H)^2$. For a basis $O_1,\ldots,O_ell$ of $\mathcal W$, define $\xi(\rho)=(\operatorname{tr}\rho O_i)_i$.

\begin{theorem}[Exact predictive-current representation]\label{found:predictive}
Two density operators have the same $\xi$ if and only if all future sector weights and currents agree under every finite sequence of the admitted controls. The image of $\xi$ is the coarsest representation sufficient for that target, in the factorization sense of Theorem~\ref{found:descent}. It evolves autonomously under the admitted controls.
\end{theorem}
\begin{proof}
Heisenberg evolution under control $a$ is $e^{t\mathcal L_a^*}$. Invariance of the finite-dimensional space $\mathcal W$ implies invariance under this exponential, hence under products of such exponentials. Equality of $\xi$ therefore gives every nominated future expectation. Conversely equality of future expectations for all nonnegative pulse durations implies equality of all one-sided mixed derivatives at zero. These are expectations of the ordered words $\mathcal L_{a_1}^*\cdots\mathcal L_{a_k}^*P_n$ and of the corresponding words applied to $C^a_{nm}$; they span $\mathcal W$. Thus $\xi$ agrees. The factorization follows from the equality of target fibers. In a fixed basis, invariance writes $\mathcal L_a^*O_i=\sum_jA^a_{ij}O_j$, so $\dot\xi=A^a\xi$ during that control.
\end{proof}

For $H=\hbar g\sigma_x$, the waves $(|0\rangle\pm i|1\rangle)/\sqrt2$ have equal sector weights and admit the same current occupancy $X=0$, but their $J_{10}$ are $\mp g$. The source distinction lost by occupancy readout is active for the next-event question. Retaining the current or its predictive completion restores information needed to evaluate a proposed generator. It does not select the generator. Classical preparation keys that remain active are appended as separate variables, rather than erased by replacing a complete preparation with a density matrix of only one subsystem.

\chapter{Complete experiments and comparison conventions}
\label{found:experiments}

\section{Constitutions and permissible transfers}
The source in a theorem is its mathematical physical domain, not a claimed exhaustive description of source reality. The following six inherited domains organize the earlier constructions. The two added completions are compared in the roadmap and specified in Parts~\ref{part:pilot} and~\ref{part:massive}.
\begin{center}\small
\begin{tabularx}{\textwidth}{@{}>{\raggedright\arraybackslash}p{.23\textwidth}XX@{}}\toprule
Constitution & Complete state and statistical input & Event and null rule\\\midrule
Canonical reactive source & Coherent canonical field, action coordinates, finite packets and carriers; intrinsic reaction chemistry & Carrier jumps; unfinished queues survive zero-current holds\\
Classical--quantum reader & Classical provenance and coherent bank; native Wiener law or CPC reconstruction premises & Conditional ray evolution and continuous physical records\\
Absorbing extraction & Loaded coherent outlets; fundamental actual transition law & Alternatives removed at capture; held-null rule explicitly specified\\
Full-wave configuration & Surviving joint wave and actual configuration; guidance or an equivariant generator and preparation law & Configuration moves without global wave collapse; unused waves may return\\
Variational discrete source & Complete coherent wave and sector history; expected current matching and a path-entropy principle & Selected Markov history and its zero-background limit\\
Material contacts & Chamber geometry, conversion field or chiral reservoir with stated readiness & Model-specific reactions or first passage; transfer needs a theorem\\\bottomrule
\end{tabularx}
\end{center}
The variational finite-configuration source may supply the generator of a full-wave record model when the complete graph and preparations agree. A continuous configuration detector is not automatically that discrete model. Neither a reduced quantum-latch instrument nor a dissipative detector current is silently identified with the Hamiltonian process \eqref{found:bell}.

A complete finite experiment specifies the initial source and apparatus law, one unknown input and its inaccessible reference, a control programme, all physical records and retained quantum states, failures, nulls and exhaustion, and every later return interaction. An event record $z$ and its normalized daughter $\rho_z$ are compared together through the unnormalized output $\sigma_z=p_z\rho_z$. An approximate declaration retains its actual daughter; it is not replaced by an ideal state after the declaration.

\section{Distances and composition}
For probability measures $P,Q$ on the same measurable output space, $d_{\rm TV}(P,Q)=\sup_E|P(E)-Q(E)|$. For states, $D(\rho,\sigma)=\frac12\|\rho-\sigma\|_1$. A classical--quantum output has norm $\int\|\sigma(dz)\|_1$, with sums in the finite case. Half-diamond distance is used only for linear maps with the declared reference extension. A deterministic pushforward or a common Markov kernel contracts total variation. A quantum channel contracts trace distance even after an identity reference extension. Unknown nonlinear continuation has no such automatic property.

\begin{lemma}[Sequential replacement and rare conditioning]\label{found:errors}
If $m$ successive complete Markov kernels on the same retained state spaces differ uniformly by at most $\epsilon_j$ in total variation, their complete history laws differ by at most $\sum_j\epsilon_j$. The same conclusion holds for completely positive trace-preserving instruments in half-diamond distance, including adaptive classical control.

If $\|\sigma-\tau\|_1\le\delta$, $p=\operatorname{tr}\sigma>0$, $q=\operatorname{tr}\tau>0$, then
\[
 D(\sigma/p,\tau/q)\le\min\{1,\delta/p\}.
\]
If $d_{\rm TV}(P,Q)\le\epsilon$ and $P(E)=p>0$, $Q(E)>0$, then
\[
 d_{\rm TV}(P(\cdot\mid E),Q(\cdot\mid E))\le\min\{1,2\epsilon/p\}.
\]
\end{lemma}
\begin{proof}
Replace one stage at a time. The prefix supplies an admitted input and the remaining common suffix contracts the chosen distance. Retaining every prior classical record makes the same argument valid on the full history and for record-controlled kernels. For quantum instruments, attach the entire reference and retained history in each half-diamond bound. For conditioning,
\[
 \|\sigma/p-\tau/q\|_1\le\|\sigma-\tau\|_1/p+|q-p|/p\le2\delta/p.
\]
For measures, apply the analogous triangle argument to restrictions to $E$; both $|P(B\cap E)-Q(B\cap E)|$ and $|p-q|$ are at most $\epsilon$.
\end{proof}

The lemma requires uniform control on all actual suffix inputs. Calibrating a finite list of preparations is not such a bound. An archive discarded in one theorem cannot later return in its claimed domain. A failure branch may consume a resource or alter a source even if the displayed apparatus has reset. These conventions are used throughout the constructions rather than repaired only in the final comparison.

% End consolidated source: parts/00_foundations.tex

\part{Hamiltonian Current Production and Reactive Realization}\label{part:kinetic}

% Begin consolidated source: parts/01_kinetic.tex
\providecommand{\kinTV}{d_{\mathrm{TV}}}
\providecommand{\kinLaw}{\operatorname{Law}}
\providecommand{\kinVar}{\operatorname{Var}}

\chapter{Canonical current production and conservative export}
\label{kin:production}

This chapter consolidates the source action and conservative exporter of
\cite{M01}, retaining the distinction drawn in \cite{M02} between coherent
current, weighted candidate flux and actual event intensity. The source is
an explicit classical canonical field coupled to physical carrier and packet
variables. The action derives signed production calibration within its
declared class. The following chapter supplies an intrinsic reaction law
as a comparison model and proves its Bell limit. Chapter~\ref{p:chapter-medium}
retains that action and comparison theorem while replacing elementary Markov
clocks and instantaneous opposite-packet cancellation by a finite gas and
finite recombination. Its ordinary material interface is an additional
constitutive rule, not an identification of every earlier classical port
with an affine quantum instrument.

\section{Complete state, currents and physical time}
Let a finite graph have vertex set $V$, oriented edge set $E$ and incidence
matrix $B$, with $B_{re}=-1$, $B_{qe}=1$ for $e=(r,q)$. A fixed orthogonal
sector decomposition $\{P_r\}$ acts on a finite-dimensional carried space.
An inaccessible reference is included in the sector fibres: all admitted
blocks below act as $H_{qr}\otimes I_R$ when that is the physical experiment.
Fix a normalized coherent vector and a bounded piecewise $C^1$ Hamiltonian
programme on a physical horizon $[0,T]$. Write
\begin{equation}
 i\hbar\dot\Psi=H(t)\Psi,\qquad
 w_r=\|P_r\Psi\|^2,\qquad
 J_{qr}=\frac2\hbar\operatorname{Im}\langle\Psi_q,H_{qr}\Psi_r\rangle.
 \label{kin:current}
\end{equation}
Thus $\dot w=BJ$, $J_{qr}=-J_{rq}$, and
\begin{equation}
 |J_{qr}|\le\frac2\hbar\|H_{qr}\|\sqrt{w_qw_r}.
 \label{kin:nodal}
\end{equation}
The same time variable is used by wave evolution, reaction generators and
record clocks. No exposure-time change is made in the Bell limit.

For a candidate jump process with actual one-time law $\rho$, actual mean
flux is $\rho_r\lambda(q\mid r)$. The expression $w_r\lambda(q\mid r)$ is
only a coherent-weighted candidate flux until initialization and equivariance
are established. Relative to a complete history filtration, $X_{t-}$ is
already definite; an intensity specifies the compensator
\[
 N_{qr}(t)-\int_0^t1_{\{X_{s-}=r\}}\lambda_s(q\mid r)\,ds.
\]
It is not the conditional probability of the current occupied sector.

The complete finite-source state contains $\Psi$, connection--action pairs
$(\chi_e,\Pi_e)$, exporter residues, queues, $N$ carrier positions,
controllers, resource counters and all actual export, cancellation and
reaction records. Further coordinates are added explicitly when they become
active. The distinguished carrier supplies one projected incidence history.
The other carriers are physical source constituents, not supplied copies of
the unknown quantum input.

\section{A common canonical action}
The source field is a classical canonical complex amplitude. Introduce
\begin{align}
 \mathcal S&=\int\left[\frac{i\hbar}{2}
 (\Psi^\dagger\dot\Psi-\dot\Psi^\dagger\Psi)
 +\hbar\sum_e\Pi_e\dot\chi_e-h(\Psi,\chi,t)\right]dt,
 \label{kin:action}\\
 h&=\sum_r\langle\Psi_r,H_{rr}\Psi_r\rangle+
 \sum_{e=(r,q)}\left(e^{i\chi_e}
 \langle\Psi_q,H_{qr}\Psi_r\rangle+\mathrm{c.c.}\right).
 \label{kin:energy}
\end{align}
The absence of $\Pi$ from $h$ is exact passive charge ownership within this
source action. It is a physical postulate, not a consequence of incomplete
readout. Hamilton's equations give
\begin{equation}
 i\hbar\dot\Psi=H(\chi,t)\Psi,\qquad \dot\chi_e=0,\qquad
 \dot\Pi_e=-\hbar^{-1}\partial_{\chi_e}h
 =\frac2\hbar\operatorname{Im}
 \left(e^{i\chi_e}\langle\Psi_q,H_{qr}\Psi_r\rangle\right).
 \label{kin:torque}
\end{equation}
At the prepared connection $\chi=0$, this is $\dot\Pi=J$ for the same
Hamiltonian that moves the coherent field. In particular $w-B\Pi$ is
conserved. Bounded finite-dimensional $H$ gives a global solution on every
finite horizon. For a time-independent programme $h$ is conserved; an
external programme supplies work $\int\partial_t h\,dt$.

\begin{theorem}[Primitive connection ownership]\label{kin:ownership}
Suppose the source has the symplectic form in \eqref{kin:action}, common
canonical action unit $\hbar$, a real quadratic energy additive over
individual vertices and primitive binary bonds, and no independent
multi-bond interaction. Require physical endpoint covariance
\[
 \Psi_r\mapsto e^{i\alpha_r}\Psi_r,\qquad
 \chi_{(r,q)}\mapsto\chi_{(r,q)}+\alpha_q-\alpha_r,
\]
with bare $H_{qr}$ neutral and $h$ independent of $\Pi$. Fixing the coherent
block $H_{qr}$ at $\chi_e=0$ fixes the bond torque there to $J_{qr}$, up to
energy terms with zero bond torque. This class is nonempty.
\end{theorem}
\begin{proof}
A general binary cross term is
$\langle\Psi_q,T_e(\chi_e)\Psi_r\rangle+\mathrm{c.c.}$.
Covariance for every endpoint phase difference $\delta$ requires
$T_e(\chi+\delta)=e^{i\delta}T_e(\chi)$, hence
$T_e(\chi)=e^{i\chi}H_{qr}$. Endpoint-diagonal terms are connection
independent. Differentiation in the canonical equation yields
\eqref{kin:torque}. Every finite Hermitian block matrix realizes these
premises. No probability or carrier selector entered the calculation.
\end{proof}

\begin{counterexample}[A gauge-invariant force bypass]\label{kin:loop}
On an oriented triangle let
$\Theta=\chi_{12}+\chi_{23}+\chi_{31}$ and add
$h_{\rm loop}=-\hbar k\|\Psi\|^2\sin\Theta$.
This is gauge invariant. At $\chi=0$ it changes no coherent Hamiltonian,
but adds $k$ to each clockwise action current. Its divergence is zero.
Thus gauge invariance and vertex conservation alone do not identify the
original Hamiltonian edge current. The primitive binary-additivity premise
of Theorem~\ref{kin:ownership} is the equation-level exclusion.
\end{counterexample}

The normalization has physical content as well. Replacing the symplectic
term by $\hbar c_e\Pi_e\dot\chi_e$ gives $\dot\Pi_e=J_e/c_e$.
The true conjugate action is $P_e=c_e\Pi_e$ and obeys $\dot P_e=J_e$.
Counting fixed $P_e$ quanta retains calibration; counting raw $\Pi_e$
quanta changes the packet charge.

Genuine coherent circulation must not be removed with this bypass. For
\[
 \Psi=3^{-1/2}(1,1,1)^{\mathsf T},\qquad
 H=\frac{3\hbar j}{2}
 \begin{pmatrix}0&-i&i\\i&0&-i\\-i&i&0\end{pmatrix},\quad j>0,
\]
one has $H\Psi=0$, constant $w_r=1/3$, and
$J_{21}=J_{32}=J_{13}=j$. The meters advance on every bond. Minimizing
activity subject only to stationary vertex weights would incorrectly
replace this specified current by zero.

\section{Export without a production probability law}
Give each packet charge $1/N$. On each edge set $k_e(0)=u_e(0)=0$ and
\[
 u_e=N(\Pi_e-\Pi_e(0))-k_e.
\]
Between exports $\dot u_e=NJ_e$. At the first hit of $s\in\{-1,1\}$,
perform the hybrid rewrite
\begin{equation}
 (u_e,k_e,Z_e)\longmapsto(0,k_e+s,Z_e+s).
 \label{kin:exportrewrite}
\end{equation}
$Z_e$ is a signed queue: an opposite arriving packet cancels an unexposed
packet on that edge. Every lineage and cancellation is retained. The
canonical $\Pi_e$ does not jump. The rewrite changes the decomposition of
stored action into emitted charge and residue.

\begin{proposition}[Conservative export estimate]\label{kin:export}
Let $A_e^N=k_e/N$ and $L_e=\int_0^T|J_e|dt$. Then
\begin{equation}
 \sup_{t\le T}\left|A_e^N(t)-\int_0^tJ_e(s)ds\right|\le N^{-1},
 \qquad \#\{\text{exports on }e\}\le NL_e.
 \label{kin:exportbound}
\end{equation}
For an arbitrary initial $u_e(0)\in[-1,1]$, the discrepancy is at most
$2/N$ and the activity bound gains at most one export. Such a residue may
be correlated with an admitted initial history.
\end{proposition}
\begin{proof}
Conservation gives $A_e^N-\int J_e=-u_e/N$, or
$(u_e(0)-u_e(t))/N$ with a nonzero initial residue. Each completed
reset-to-reset excursion consumes at least $1/N$ of the total variation
of $\Pi_e$. The first excursion from a nonzero residue costs at most one
exception to this count. This also prevents export accumulation in finite
time.
\end{proof}
More generally, an exact conservation law
$k_e/N+v_e=v_e(0)+\int J_e$ with $|v_e|\le C/N$ gives discrepancy
$2C/N$. Under a uniform $O(N)$ elementary-activity budget, the kinetic
proof below applies to this larger exporter class. Fixed packet charge,
bounded storage and the hybrid rewrite are constitutive inputs; neither
a particular packet-birth distribution nor a statistical self-averaging
law is needed for signed production.

A growing compound packet batch cannot obey this same accounting. If
$\Pi$ is continuous and a compound export contains $K_N$ elementary
charges, the residual jump has magnitude $K_N/N$. Bounded storage forces
$K_N\le2C$. Thus the countermodels with $K_N=N/\mu_N\to\infty$ require a
larger cell or an additional impulsive production force. Mere matching of
conditional production means would not exclude them.

\chapter{Binary reactions and the physical-time Bell limit}
\label{kin:limit}

The theorem in this chapter is the intrinsic-chemistry comparison result
from \cite{M01}. It replaces the normalized allocation rule of \cite{M02}
by physical pair counting. The proof separates global flux tracking from
the later small-weight localization of a tagged path. Within this chapter,
a complete additive Markov reaction generator is supplied: neither the
action nor source non-reconstructibility selects it. The finite-gas and
recombination construction of Chapter~\ref{p:chapter-medium} derives this
comparison law in a controlled limit under a new microscopic constitution.

\section{Elementary reactions, null evolution and resources}
There are $N$ source carriers at positions $X_a\in V$ and
$x_r=N^{-1}\#\{a:X_a=r\}$. One predesignated carrier $a_*$ is observed.
A positive queued packet on $e=(r,q)$ can react with every carrier at $r$:
\[
 (P_e^+,C_a@r)\longrightarrow C_a@q+\text{record }(e,a,+).
\]
The negative packet uses the opposite direction. Every eligible pair has
rate $\kappa_e\mu_N/N$, where
$0<\kappa_-\le\kappa_e\le\kappa_+<\infty$ are fixed. Pair generators add,
there is no other carrier-changing reaction, and all carrier-specific
response has been excluded from this scalar class. These last statements
are completeness commitments, not mere relabeling symmetries.

For a function $F$ of the complete source state and archive, the generator is
\begin{equation}
 \mathcal G_NF=
 \sum_e\frac{\kappa_e\mu_N}{N}|Z_e|
 \sum_{a:X_a=o(e,Z_e)}[F(R_{ea}s)-F(s)],
 \label{kin:generator}
\end{equation}
where $R_{ea}$ consumes the signed packet, moves its owner and appends the
actual mark. Empty origins have no eligible pair: their packets remain
queued. Between stochastic events the wave, residues and declared clocks
continue their deterministic equations. Conditional survival along this
no-reaction flow is $\exp(-\int R_Ndt)$, with $R_N$ the sum in
\eqref{kin:generator}. A constant-rate exponential waiting time is asserted
only on a holding interval with constant total rate.

Let $m_e=\mu_N Z_e/N$. Direct counting gives the normalized bulk fluxes
and tag rates
\begin{align}
 \Phi_e^{N,+}&=\kappa_e x_r[m_e]_+,&
 \Phi_e^{N,-}&=\kappa_e x_q[-m_e]_+,\label{kin:flux}\\
 \lambda_{e,a_*}^{N,+}&=\kappa_e[m_e]_+\quad(X_{a_*}=r),&
 \lambda_{e,a_*}^{N,-}&=\kappa_e[-m_e]_+\quad(X_{a_*}=q).
 \label{kin:tagrate}
\end{align}
There are $|Z_e|n_r$ eligible pairs and $|Z_e|$ pairs containing a fixed
origin carrier. Consequently $\lambda_{qr}^N=\Phi_{qr}^N/x_r$ on an
occupied tag origin. This division is derived from composition of pairs;
it was not inserted as a normalized choice rule.

Put $L=\sum_eL_e$. Every reaction consumes a packet, so at most
$NL+O(1)$ native reactions occur. A source with capacity for that many
export and reaction entries has no overflow on the promised domain.
Smaller resources require an explicit exhaustion branch. Assigning
degenerate energy to packet, carrier and archive labels conserves the
stated source energy at rewrites; it does not derive material memory cost
or a cyclic reset.

\section{Global tracking, including empty origins}
Assume empty initial queues and calibrated initial populations with
$\mathbb E\|x^N(0)-w(0)\|_1\to0$. Deterministic census preparation is a special case. Randomized census preparation is also allowed; conditional on the initial field, it is independent of future comparison reaction clocks. Let $z^N=Z/N$ and
$e^N=A^N-\int Jdt$. The exact balance is
\begin{equation}
 x^N(t)+Bz^N(t)-w(t)=x^N(0)-w(0)+Be^N(t).
 \label{kin:balance}
\end{equation}
Write $\eta_N=\|x^N(0)-w(0)\|_\infty+C_B/N$ and define
\begin{align}
 \epsilon_{F,N}&=\sum_e\mathbb E\int_0^T
 \bigl(|\Phi_e^{N,+}-[J_e]_+|+
 |\Phi_e^{N,-}-[-J_e]_+|\bigr)dt,\label{kin:epsF}\\
 \epsilon_{x,N}&=\mathbb E\sup_{t\le T}\|x^N(t)-w(t)\|_1.
 \label{kin:epsx}
\end{align}

\begin{theorem}[Global mass-action tracking]\label{kin:tracking}
For the fixed finite programme above, bounded-variation currents and
\begin{equation}
 \mu_N\longrightarrow\infty,\qquad \mu_N/N\longrightarrow0,
 \label{kin:hierarchy}
\end{equation}
one has $\epsilon_{F,N}\to0$ and $\epsilon_{x,N}\to0$.
The theorem includes coherent cycles, current reversal, empty origins,
zero coherent weights and dark intervals. It is not uniform over a growing
sector graph or arbitrarily rapidly varying response coefficients.
\end{theorem}
\begin{proof}
We give the tracking and low-population steps separately. The same pathwise inventory and variation estimates hold after conditioning on the initial census; expectations below average that census as well as reaction clocks. Fix a cutoff
$\delta>0$ and, on each edge, run a companion signed queue from the same
empty state and the same exports with service function
\[
 \phi_t(u)=a_+(t)[u]_+-a_-(t)[-u]_+,\qquad
 a_+=\kappa_e\max(x_r(t-),\delta),\quad
 a_-=\kappa_e\max(x_q(t-),\delta).
\]
All divided-difference slopes lie between $a_0=\kappa_-\delta$ and
$a_1=\kappa_+$. The companion is driven by adapted coefficients; they are
not asserted independent of either process. Since each physical reaction
uses one of the $NL+O(1)$ packets, the sum of absolute population jumps is
at most $2L+O(N^{-1})$. Hence the variations of $a_\pm$ are bounded
independently of $N$ at fixed $\delta$.

For $m^*=\mu_NZ^*/N$, compensated reaction counting yields
\begin{equation}
 dm^*=\mu_N[J-\phi_t(m^*)]dt+\mu_Nde^N+dM,
 \qquad d\langle M\rangle_t=\frac{\mu_N^2}{N}|\phi_t(m^*)|dt.
 \label{kin:companion}
\end{equation}
Let $y$ solve the same adapted finite-variation equation without $M$, and
let $y_J$ omit both $M$ and $de^N$. Scalar monotonicity and
$\|e^N\|_\infty\le c_0/N$ give
\begin{equation}
 \|y-y_J\|_\infty\le2c_0\mu_N/N,\qquad
 \|y_J\|_\infty\le J_*/a_0,\qquad J_*:=\sup_{e,t}|J_e(t)|.
 \label{kin:dettracking}
\end{equation}
For the first inequality, write the difference equation with a measurable
divided-difference coefficient $a(t)\in[a_0,a_1]$. Its solution at $t$
is $\mu_N\int_0^t\exp[-\mu_N\int_s^t a(u)du]\,de^N(s)$.
This is a pathwise Stieltjes identity, not an anticipative stochastic
integral. Integration by parts bounds its absolute value by
$2\mu_N\|e^N\|_\infty$. The second inequality follows because the drift
points toward $[-J_*/a_0,J_*/a_0]$.

The instantaneous root
$f(t)=[J(t)]_+/a_+(t)-[-J(t)]_+/a_-(t)$ satisfies
\[
 \kinVar(f)\le \frac{\kinVar(J)}{a_0}
 +\frac{J_*}{a_0^2}[\kinVar(a_+)+\kinVar(a_-)].
\]
Contraction between jumps of $f$ and summation of its jumps imply
\[
 \int_0^T|y_J-f|dt\le\frac{|f(0)|+\kinVar(f)}{\mu_Na_0}.
\]
For $\xi=m^*-y$, production jumps cancel. The square-jump identity and
monotonicity, with $V=\mathbb E\xi^2$ and $\alpha=\mu_N/N$, give
\[
 V'\le-2\mu_Na_0V+
 \mu_N\alpha a_1(J_*/a_0+2c_0\alpha+\sqrt V).
\]
Young's inequality absorbs the square-root term into $\mu_Na_0V$ and
gives $\sup_tV\le C_\delta(\alpha+\alpha^2)$. Thus
\begin{equation}
 R_{\delta,N}:=\sum_e\mathbb E\int_0^T|\phi_t(m_e^*)-J_e|dt
 \le C_\delta\left(\mu_N^{-1}+\frac{\mu_N}{N}
 +\sqrt{\frac{\mu_N}{N}}\right).
 \label{kin:R}
\end{equation}
This controls directional flux because a signed queue exposes only one
orientation, and
\[
 |a_+[m^*]_+-[J]_+|+|a_-[-m^*]_+-[-J]_+|
 =|\phi_t(m^*)-J|.
\]

It remains to remove $\delta$ without assuming the desired population
closeness. At time $t$ let $K=\{r:x_r<\delta\}$, and let $I_K,O_K$ be
normalized queued charge directed into and out of $K$. Summing
\eqref{kin:balance} over $K$ gives
\[
 w_K+O_K\le |V|\delta+I_K+|V|\eta_N.
\]
An incoming queue has origin outside $K$, so its origin population is at
least $\delta$. Its expected service count therefore bounds
\[
 \mathbb E\int_0^TI_Kdt\le
 \frac{L+O(N^{-1})}{\kappa_-\mu_N\delta}.
\]
Consequently
\begin{equation}
 \mathbb E\int_0^Tw_Kdt\le |V|\delta T+
 \frac{L+O(N^{-1})}{\kappa_-\mu_N\delta}+|V|T\mathbb E\eta_N.
 \label{kin:smallmass}
\end{equation}
The target directional current whose physical origin lies in $K$ is
therefore bounded, using \eqref{kin:nodal} and Cauchy--Schwarz, by
\begin{equation}
 D_{\delta,N}\le C_H\sqrt{T\left(|V|\delta T+
 \frac{L+O(N^{-1})}{\kappa_-\mu_N\delta}+|V|T\mathbb E\eta_N\right)}.
 \label{kin:smallflux}
\end{equation}

Couple physical and companion queues with identical exports and
minimum-rate baseline services. View companion service as baseline
service with physical coefficients plus extra service where $x<\delta$.
Every unmatched baseline service contracts their absolute signed queue
difference by $1/N$; an extra companion service can enlarge it by at most
$1/N$. Starting from equality, the expected number of baseline mismatches
is at most the expected number of extra services. The latter normalized
count is at most $R_{\delta,N}+D_{\delta,N}$, since it is supported on
low-population origins. Comparing the two directional flux vectors thus
costs at most twice this count. Adding the direct companion error gives
\begin{equation}
 \epsilon_{F,N}\le3R_{\delta,N}+2D_{\delta,N}.
 \label{kin:fluxbridge}
\end{equation}
Take $N\to\infty$ at fixed $\delta$, then $\delta\downarrow0$, to prove
flux convergence. Finally the physical population process has drift
$B(\Phi^{N,+}-\Phi^{N,-})$ and an $O(N^{-1})$ quadratic-variation budget,
because it has at most $NL+O(1)$ jumps of size $1/N$. The martingale
maximal inequality, initial calibration and integrated flux convergence
prove \eqref{kin:epsx} tends to zero.
\end{proof}

\section{The complete tagged path and the nodal boundary}
Total variation means $\kinTV(P,Q)=\sup_A|P(A)-Q(A)|$ on the measurable
space $D([0,T],V)$ of finite-sector c\`adl\`ag paths. It controls event
ordering, exact event times, finite null windows and every common measurable
stopping or coarse record of the path. It does not by itself control an
additional archive absent from that output space.

\begin{lemma}[Finite-graph Bell existence through nodes]\label{kin:existence}
For \eqref{kin:current}, the minimal rates
$\lambda^B_{qr}=[J_{qr}]_+/w_r$ on positive-weight origins define a
nonexplosive inhomogeneous jump process starting at $w(0)$, with law
$w(t)$ at every time. If $\nu\le Cw(0)$, the same construction starts at
$\nu$, remains dominated by $Cw(t)$ and does not occupy a zero-weight
sector. Its conditional law is unique within this time-inhomogeneous
Markov class.
\end{lemma}
\begin{proof}
Each open set $\{t:w_r(t)>0\}$ is a countable union of intervals.
No finiteness of the nodal set is inferred from piecewise $C^1$ regularity.
On compact subintervals of these positive-weight components, standard
integrated-hazard first-jump construction is unique until explosion or a
nodal boundary; the increasing compact localization defines the minimal
process on their union. Killing at
such boundaries gives the minimal forward solution. The nonnegative
vector $w$ solves its forward balance equation, so successive first-jump
iteration, or positive Volterra iteration, bounds each partial-transition
sum by $w$; with initial $\nu$ the bound is $Cw$.
If a holding path remains at $r$ while $w_r$ tends to zero, write incoming
and outgoing positive currents as $I_r,O_r$. Then
$\dot w_r=I_r-O_r$ and
$\lambda^B_{\rm out}(r)=O_r/w_r\ge-\dot w_r/w_r$.
Integrating shows that the holding survival to that zero is zero.
The dominated killed law also gives
\[
 \mathbb E N_{[0,T]}\le
 C\int_0^T\sum_{q,r}[J_{qr}(t)]_+dt<\infty.
\]
Thus neither explosion nor nodal killing loses mass. For initial $w(0)$,
normalization and domination imply equality with $w(t)$. For general
$\nu$, normalization gives the asserted dominated process. The
first-jump construction determines its law uniquely. This is the
finite-graph existence argument underlying the standard Bell process
\cite{BellQFT}; it is not a selection of that process among all event laws.
\end{proof}

\begin{theorem}[Tagged physical-time path convergence]\label{kin:path}
Under Theorem~\ref{kin:tracking}, initialize the distinguished carrier
with fixed law $\nu\le Cw(0)$ while keeping calibrated total populations.
Then
\begin{equation}
 \kinTV\bigl(\kinLaw(X_{a_*}^N),\kinLaw(Q^B)\bigr)\longrightarrow0,
 \qquad \lambda^B(q\mid r;t)=\frac{[J_{qr}(t)]_+}{w_r(t)}.
 \label{kin:Belllimit}
\end{equation}
The target starts at $\nu$. The conclusion is unchanged by fixed positive
edge coefficients or uniformly bounded initial exporter residues.
\end{theorem}
\begin{proof}
Let $N_\varepsilon$ count crossings of the deterministic weights through
a regular level $\varepsilon$. One-dimensional coarea gives
$\int_0^1N_\varepsilon d\varepsilon\le\sum_r\kinVar(w_r)$.
There is a sequence $\varepsilon_k\downarrow0$ such that
$\varepsilon_kN_{\varepsilon_k}\to0$: otherwise $N_\varepsilon$ would
have a nonintegrable $c/\varepsilon$ lower bound near zero.
The probability that the Bell path visits a sector while its weight is
at most $\varepsilon$ is bounded by initial small-weight mass, jump
influx into those sectors, and deterministic downcrossings at which the
path is already in that sector. By Lemma~\ref{kin:existence} and
\eqref{kin:nodal}, a valid bound is
\begin{equation}
 b_C(\varepsilon)=C\bigl(|V|\varepsilon+
 C_0T\sqrt\varepsilon+\varepsilon N_\varepsilon\bigr).
 \label{kin:nodebudget}
\end{equation}
Jump influx uses the integrated current bound at a small destination;
each downcrossing contributes at most $C\varepsilon$.

Couple the target and tagged jumps at minimum conditional intensities
until they disagree, the target reaches the small-weight region, or
$\sup_t\|x^N-w\|_1>\varepsilon/2$. In the remaining states,
\[
 |\lambda^N_{qr}-\lambda^B_{qr}|
 \le\frac2\varepsilon|\Phi^N_{qr}-[J_{qr}]_+|
 +\frac{2J_*}{\varepsilon^2}|x_r^N-w_r|.
\]
The full microscopic marginal is preserved by this coupling although its
queues depend on the tag's past. The target marginal retains its own
Markov rates. A union and compensator bound gives
\begin{equation}
 \kinTV\le b_C(\varepsilon)
 +\frac{2\epsilon_{x,N}+2\epsilon_{F,N}}{\varepsilon}
 +\frac{C_GJ_*T}{\varepsilon^2}\epsilon_{x,N}.
 \label{kin:pathbound}
\end{equation}
Take $N\to\infty$ at each fixed $\varepsilon_k$, then $k\to\infty$.
The queue proof's cutoff $\delta$ was already removed; the two
localizations are not interchanged. This proves the whole-path claim.
\end{proof}

At a unique ready sector $w(0)$ is a point mass, so the initial tag is
certain. For general $w(0)$, population calibration and the tag's initial
law remain explicit preparation requirements. The limiting intensity is
relative to the natural tagged history conditional on the declared
coherent programme. Exposing all queues or microscopic random resources
is a different filtration, with no automatic Bell conditional law.

\section{Finite tails and the boundary of convergence}
For a Rabi pulse $\hbar=g=1$, $H=\sigma_x$ and initial state $(1,0)$,
\[
 w_0=\cos^2t,\quad w_1=\sin^2t,\quad J=\sin2t,\quad
 k(t)=\lfloor N\sin^2t\rfloor\quad(0\le t\le\pi/2).
\]
If $b$ carriers occupy sector $1$, then $Z=k-b$ and the exact rates are
\begin{equation}
 b\to b+1:\ \frac{\kappa\mu_N}{N}(k-b)(N-b),\qquad
 b\to b-1:\ \frac{\kappa\mu_N}{N}(b-k)b,
 \label{kin:binaryrates}
\end{equation}
using only the positive rate in its respective region. For $N=1$ the
first export is at $T=\pi/2$; a dark hold then has transfer probability
$1-e^{-\kappa\mu_Ns}$. A zero-current interval need not be a no-event
interval at finite resources.

At fixed $\nu=\kappa\mu$, the $N\to\infty$ fluid equation for the root
fraction is $x'=-\nu x(x-w_0)$, $x(0)=1$. Lipschitz drift comparison and
the $O(\nu T/N)$ martingale budget prove that limit. With
$A(t)=\int_0^tw_0(s)ds$, direct differentiation gives
\[
 \frac1{x(t)}=e^{-\nu A(t)}\left(1+\nu\int_0^te^{\nu A(s)}ds\right).
\]
Since $A(T)-A(T-s)=s^3/3+O(s^5)$, the substitution
$s=\nu^{-1/3}u$ and dominated Laplace asymptotics yield
\begin{equation}
 x(T)\sim\frac{3^{2/3}}{\Gamma(1/3)}\nu^{-2/3}.
 \label{kin:boundarylayer}
\end{equation}
This is an iterated limit, not a uniform joint $(N,\mu_N)$ error rate.
In a dark hold, the exact number $M$ of unfinished carriers decreases at
rate $\kappa\mu_NM^2/N$. From $M_0\ge1$ its mean clearance time is
\[
 \frac{N}{\kappa\mu_N}\sum_{m=1}^{M_0}\frac1{m^2}.
\]
It grows as $N/\mu_N$ whenever an unfinished carrier is present. A fixed
marked-carrier limit does not establish finite-time reset of the entire
source bank.

The scale and regularity assumptions are also substantive. A cell of size
$1/\mu$ releasing periodic bursts under constant current $j$ drives the
fast output $y(s)=e^{-s}/(1-e^{-1/j})$ on $0\le s<1/j$.
Its mean is $j$, but its mean absolute deviation from $j$ is positive.
Integrated hazards may converge while time densities fail to converge in
TV. The estimate \eqref{kin:R} needs $\mu\times\text{storage size}\to0$.
Likewise bounded coefficients $1+\epsilon\cos(\mu_Nt)$ need not have
uniformly bounded variation and are outside the fixed-response theorem.

\chapter{Physical equalization and complete-history response control}
\label{kin:neutrality}

The scalar reaction class can be supplied by an actual material
neutralization mechanism within a specified response law. The exact
finite-yield result comes from \cite{M03}; the stronger comparison of
complete marked histories and returning-work resources comes from
\cite{M21}. Neither theorem excludes a second physical input to the
reaction coefficient.

\section{Exact finite-yield neutrality}
Give each carrier a transferable valence $v_a$ and let its aperture have
length $F_e(v_a)>0$, with common bounded Lipschitz $F_e$. A homogeneous
intrinsic contact generator per unit length integrates to pair rate
$\kappa_e\mu_NF_e(v_a)/N$. Homogeneity, perfect transmission and
completeness of this aperture as the reaction input are disclosed laws.
They allow arbitrary initial unequal valences.

On a connected undirected valence graph take capacities $c_{ab}>0$ and
zero-sum admitted forcing $f$. Define
\begin{equation}
 \Phi(v)=\sum_{\{a,b\}}c_{ab}|v_a-v_b|,\qquad
 \dot v+\partial\Phi(v)\ni f(t),\qquad \sum_af_a=0.
 \label{kin:yield}
\end{equation}
Equivalently the antisymmetric edge transfer satisfies
$\ell_{ab}\in c_{ab}\operatorname{Sign}(v_a-v_b)$ and
$\dot v_a=f_a-\sum_b\ell_{ab}$. This is an ideal finite-yield constitutive
interaction. For bounded measurable forcing, implicit steps
\[
 v^{k+1}=\arg\min_v\left\{
 \frac{\|v-v^k-\Delta t f^k\|^2}{2\Delta t}+\Phi(v)\right\}
\]
are uniquely defined. Bounded subgradients supply equicontinuity of their
interpolants; the closed monotone graph identifies a limiting absolutely
continuous solution. Monotonicity gives
$\frac12\frac d{dt}\|v-u\|^2\le0$ for solutions with identical forcing,
so that solution is unique. Endogenous controllers must separately have a
well-posed causal update; arbitrary ill-posed feedback is not admitted.

\begin{theorem}[Cut capacity and finite-time locking]\label{kin:locking}
Let $C(S)=\sum_{\{a,b\}\in\partial S}c_{ab}$ and
$f(S)=\sum_{a\in S}f_a$. Consensus $v=\bar v\mathbf1$ can remain
stationary exactly when $|f(S)|\le C(S)$ for every proper nonempty $S$.
If uniformly $C(S)-|f(S)|\ge\eta>0$, the solution reaches consensus by
\begin{equation}
 T_{\rm lock}\le\frac{\sqrt{2NE(0)}}{\eta},\qquad
 E(0)=\tfrac12\sum_a(v_a(0)-\bar v)^2,
 \label{kin:locktime}
\end{equation}
and remains there while the non-strict cut condition holds. Afterwards
all eligible carrier coefficients are exactly
$\kappa_e\mu_NF_e(\bar v)/N$.
\end{theorem}
\begin{proof}
At consensus the allowable transfers form the box $|\ell_e|\le c_e$;
stationarity means its divergence equals $f$. Summing over cuts proves
necessity. Its divergence image is compact convex with support function
$\Phi(u)$. Sorting $u_{(1)}\le\cdots\le u_{(N)}$, with $S_k$ the indices
above the $k$th gap, yields
\[
 \Phi(u)=\sum_{k=1}^{N-1}(u_{(k+1)}-u_{(k)})C(S_k),\quad
 f\cdot u=\sum_{k=1}^{N-1}(u_{(k+1)}-u_{(k)})f(S_k).
\]
The cut inequalities imply $f\cdot u\le\Phi(u)$ for every $u$.
Finite-dimensional separation proves sufficiency. The mean is conserved.
One-homogeneity gives $\xi\cdot v=\Phi(v)$ for
$\xi\in\partial\Phi(v)$. Strict slack, the same coarea identity and
$D=\max v-\min v$ imply
\[
 \dot E=f\cdot v-\Phi(v)\le-\eta D
 \le-\eta\sqrt{2E/N}.
\]
Integration for $\sqrt E$ proves \eqref{kin:locktime}. Stationarity and
uniqueness prevent subsequent departure. Substitution in the actual
aperture length proves the response equality without an initial
valence-distribution assumption.
\end{proof}

For a complete graph with $c_{ab}\ge\rho/N$ and $|f_a|\le F<\rho/2$,
one has the sharper $D(t)\le[D(0)-(\rho-2F)t]_+$. Indeed, if tied maximum
and minimum groups have sizes $m,k$ and $D>0$, averaging their saturated
external transfers gives
$\dot D\le2F-\rho(2N-m-k)/N\le2F-\rho$.
At consensus the admissible flow $(f_a-f_b)/N$ balances the forcing.
The network has $O(N^2)$ physical links. Opening measurement at a fixed
preparation time beyond the uniform locking bound avoids leaking different
first-lock times; those histories remain in the source.

\section{Approximate equality with all old records retained}
A smooth alternative assigns material $h_a$ and response $g_e(h_a)$ with
$g_-\le g_e\le g_+$ and common Lipschitz constant $L_g$. Admit
\begin{equation}
 dh_a=-\nu(h_a-\bar h)dt+dK_a,\qquad
 \sum_a dK_a=0,\qquad
 V_N=\int_0^T\|dK\|_\infty,
 \label{kin:smooth}
\end{equation}
where every returning-memory write is represented by a well-posed causal
finite-variation $K$. The complete pair rate is
\[
 r_{ea}=\frac{\kappa_e\mu_N}{N}|Z_e|g_e(h_a)
 1_{\{X_a=o(e,Z_e)\}}.
\]
The comparison source has exactly the same material, records, initial
law and controller rules, but replaces $g_e(h_a)$ by $g_e(\bar h)$.
Write their complete marked path laws as $P_N,\bar P_N$.

\begin{theorem}[Complete-history equalization bound]\label{kin:equalization}
If $\sum_e|Z_e|\le B_N$ and
$D(t)=\max_a|h_a-\bar h|$, then
\begin{equation}
 \kinTV(P_N,\bar P_N)\le
 \min\left\{1,C_N\mathbb E\int_0^TD(t)dt\right\},\qquad
 C_N=\kappa_+\mu_NB_NL_g.
 \label{kin:equalbound}
\end{equation}
It holds for arbitrary common initial correlations, including retained
heat, controller and material histories. With uniform bounds $D_0,E_0$:
\begin{align}
 \kinTV(P_N,\bar P_N)&\le C_NTD_0e^{-\nu\tau}
 &&\text{after preload $\tau$ and no further writes},\label{kin:preload}\\
 \kinTV(P_N,\bar P_N)&\le C_N(D_0+V_N)/\nu
 &&\text{for live exchange},\label{kin:live}\\
 \kinTV(P_N,\bar P_N)&\le C_N
 \sqrt{T(E_0+W_{\max})/\nu}
 &&\text{for the work budget below}.
 \label{kin:workbound}
\end{align}
Each bound may be capped at one.
\end{theorem}
\begin{proof}
Until the first unmatched event, couple common channels at their minimum
rate and all excess rates separately. Every causal controller then agrees
between the two histories. At a common state,
\[
 \sum_{e,a}|r_{ea}-\bar r_{ea}|
 \le\frac{\kappa_+\mu_NL_g}{N}\sum_e|Z_e|
 \sum_{a:X_a=o}|h_a-\bar h|\le C_ND(t).
\]
The first-mismatch compensator and coupling inequality prove
\eqref{kin:equalbound}, including all shared record functions. Variation
of constants in \eqref{kin:smooth} yields
$d(t)=e^{-\nu t}d(0)+\int_{(0,t]}e^{-\nu(t-s)}dK(s)$ for
$d=h-\bar h\mathbf1$. Integrating its sup norm proves
\eqref{kin:preload} and \eqref{kin:live}.
For $E=\|d\|_2^2/2$, define the signed injected work
\[
 W_T=\int_{(0,T]}d(s-)\cdot dK(s)
 +\tfrac12\sum_{s\le T}\|\Delta K(s)\|_2^2.
\]
The finite-variation chain rule gives
$E(T)+\nu\int_0^T\|d\|_2^2dt=E(0)+W_T$.
If $W_T\le W_{\max}$, Cauchy--Schwarz proves \eqref{kin:workbound}.
No independence of returning writes from old records was used.
\end{proof}

A tag offset $D$ at conserved mean has
$\|d\|_2^2\ge D^2N/(N-1)$, with equality when all other offsets equal
$-D/(N-1)$. Maintaining it requires power at least
$\nu D^2N/(N-1)$; a finite initial stored energy may temporarily pay that
cost. The effective dissipative equation and its heat ledger are a new
material model, not a thermal-equilibrium derivation.

If the right side of \eqref{kin:equalbound} tends to zero and the decorated
scalar baseline satisfies Theorem~\ref{kin:path}, projection and triangle
inequality give the Bell path limit with additional error equal to that
right side. This removes exact initial equality within the specified
response class. It does not compare all retained records with an
input-independent archive.


\subsection{A sharper fixed-tag theorem for growing conductance networks}
The complete-history bound is deliberately stronger than a fixed-tag bound.
The latter can converge under a less demanding material scale. Suppose
\[
 \dot v_a=\sum_{b\ne a}g_{ab}(v_b-v_a)+f_a,
 \qquad g_{ab}=g_{ba}\in[g_-,g_+],\quad |f_a|\le F,
 \quad\sum_af_a=0,
\]
with fixed positive $g_-,g_+$, well-posed controls, and a preserved range on
which $h_-\le F_e(v)\le h_+$ and $F_e$ is Lipschitz. The physical network
has $O(N^2)$ links. Its diameter satisfies
\begin{equation}
 D(t)\le D_0e^{-Ng_-t}+\frac{2F}{Ng_-}(1-e^{-Ng_-t}).
 \label{kin:networkdiameter}
\end{equation}
Indeed, subtract the equations at an almost-everywhere maximum and minimum;
the common complete-graph contribution contracts their difference by
$Ng_-$, and the two forces contribute at most $2F$.

\begin{proposition}[Fixed-tag neutrality at the network scale]
\label{kin:networkneutrality}
With the canonical exporter, initial population calibration, and pair
coefficient $\kappa_e\mu_NF_e(v_a)/N$, the global flux/population conclusions
of Theorem~\ref{kin:tracking} hold. At a positive-weight path cutoff
$\varepsilon$, the additional tagged path error due to unequal material is
at most $C(D_0+FT)/(\varepsilon Ng_-)$, besides the canonical flux,
population and nodal errors. Thus the tagged Bell limit holds under
\eqref{kin:hierarchy} without demanding that the complete-history bound
\eqref{kin:equalbound} itself vanish.
\end{proposition}
\begin{proof}
Define weighted occupancies
$H_{e,r}^N=N^{-1}\sum_{a:X_a=r}F_e(v_a)$ and
$\bar h_{e,r}=H_{e,r}^N/x_r$ when $x_r>0$.
They satisfy $h_-x_r\le H_{e,r}^N\le h_+x_r$. Their total variation is
uniformly bounded: carrier moves contribute at most a constant times
$h_+L$, while
\[
 \frac1N\sum_a|\dot v_a|\le Ng_+D+F
\]
and \eqref{kin:networkdiameter} give a uniform integrated bound for
material variation. In the companion proof replace service slopes by
$\kappa_e\max(H_{e,r}^N,h_-\delta)$. The deterministic root tracking and
escape martingale estimates still apply. Extra companion service can occur
only when $x_r<\delta$, and an incoming queue to that cut has weighted
origin at least $h_-\delta$. Thus the low-population and mismatch estimates
hold with $\kappa_-$ replaced by $\kappa_-h_-$, proving global convergence.

The exact tag rate is
$\lambda_{e,*}^N=(F_e(v_*)/\bar h_{e,r})\Phi_e^N/x_r$.
The first ratio is bounded by $R=h_+/h_-$ and differs from one by at most
$\operatorname{Lip}(F_e)D/h_-$. On the stopped region $w_r\ge\varepsilon$,
$x_r\ge\varepsilon/2$, its difference from Bell's rate is bounded by
\[
 \frac{2R}{\varepsilon}|\Phi_e^N-[J_e]_+|
 +\frac{2RJ_*}{\varepsilon^2}|x_r-w_r|
 +\frac{J_*\operatorname{Lip}(F_e)}{\varepsilon h_-}D.
\]
Integrating \eqref{kin:networkdiameter} proves the added term. Apply the
same minimum-rate path coupling and the same order of nodal limits as in
Theorem~\ref{kin:path}. This controls a fixed tag; it does not assert
complete-source archive equivalence or uniformly bounded-degree resources.
\end{proof}

\section{Response bypasses and histories that remain active}
An exceptional carrier with pair multiplier $h_a>0$ is selected with
probability $h_a/\sum_{b:X_b=r}h_b$. One tag of fixed multiplier $h$ has
vanishing bulk fraction but limiting rate $h[J_{qr}]_+/w_r$.
In a monotone star with channel strengths $p_i$ and tag multipliers $h_i$,
\begin{equation}
 \mathbb P(i)=\frac{h_ip_i}{\sum_jh_jp_j},\qquad
 S(t)=(\cos^2\theta(t))^{\sum_jh_jp_j}.
 \label{kin:exceptional}
\end{equation}
This follows by integrating the hazards
$2h_ip_i\dot\theta\tan\theta$.
At $p=(3/4,1/4)$ and $h=(6/5,2/5)$ the labels are $(9/10,1/10)$ while
the unlabelled waiting law is unchanged. A second stress-dependent factor
$\zeta_a$ multiplying $g_e(h_a)$ restores this freedom even after exact
material equalization. It violates response completeness, not conservation.

Conservation and fast averaging also fail to remove event-associated
keys. If a rapidly alternating sign $\sigma(t)=\pm1$ gives intensity
$\ell(1+\alpha\sigma(t))$, the survival tends to $e^{-\ell t}$ as the
period shrinks. At an event print the instantaneous sign. Its limiting
positive-key probability by $T$ is
$(1+\alpha)(1-e^{-\ell T})/2$, whereas a scalar rate with the same
alternating key gives $(1-e^{-\ell T})/2$. This difference follows by
integrating the corresponding marked intensity over half periods.
The full-record gap is $\alpha(1-e^{-\ell T})/2$. Correct averaged timing
does not imply equality of complete marked histories.

A reversible clamp makes the same distinction physically explicit.
For conjugate material $p$ and $P_\perp=I-\mathbf1\mathbf1^{\mathsf T}/N$,
$H_{\rm eq}=-\nu h^{\mathsf T}P_\perp p$ gives
$d(t)=e^{-\nu t}d(0)$ but
$\pi(t)=e^{\nu t}\pi(0)$, $\pi=P_\perp p$.
After duration $\tau$, an admitted return
$H_{\rm ret}=\omega\|\pi\|^2/2$ for time $s$ gives
\[
 d_{\rm after}=e^{-\nu\tau}d(0)+\omega s e^{\nu\tau}\pi(0).
\]
The contracted information survives in conjugate stress. On any fixed
compact preparation tube a smooth Hamiltonian cutoff preserves these
finite trajectories, but the required stress range grows exponentially.
A bounded coupling coefficient is not a bounded-return-work condition.

\begin{remark}[Premises of the kinetic comparison]\label{kin:boundary}
The action removes a separately assigned signed-production calibration;
conservative export removes detailed production statistics; pair counting
removes an explicit normalized carrier selector; material exchange replaces
exact initial neutrality in its declared class. This comparison theorem
still assumes common packet charge, population calibration, complete Markov
chemistry and exhaustive response inputs. The pilot construction
replaces elementary stochastic clocks and instantaneous cancellation by
finite microscopic dynamics, while retaining charge, preparation and response
assumptions and specifying a new ordinary material coupling rule.
The countermodels in Part~\ref{part:access} remain valid for their stated classical
contacts: restrictive-looking access can preserve the native law while
exporting source information.
\end{remark}

% End consolidated source: parts/01_kinetic.tex

\part{Deterministic Pilot Medium and Autonomous Material Records}\label{part:pilot}

% Pilot-medium additions. Insert as a new Part after kin:boundary.
% No class, document environment or bibliography environment belongs here.
\providecommand{\pLaw}{\operatorname{Law}}
\providecommand{\pVar}{\operatorname{Var}}
\providecommand{\pSupp}{\operatorname{supp}}

\chapter{A deterministic pilot medium and its Bell limit}
\label{p:chapter-medium}

The preceding kinetic chapters identify a precise conditional bridge:
canonical action export, scalar signed-queue service and population
calibration imply the complete Bell path in physical time.  In
Remark~\ref{kin:boundary}, instantaneous directional cancellation,
intrinsic Markov chemistry and the admission of record contacts remain
physical commitments.  We now replace the first two by explicit finite
pilot dynamics and state a new common material interaction catalogue.
The signed queue survives as a comparison process, so
Theorems~\ref{kin:tracking} and \ref{kin:path} are used without repeating
their proofs.

The new ingredients have separate roles.  Independent initial positions
in a finite freely moving gas generate a controlled Poisson contact-law
limit.  Opposite activation species coexist and may both react at every
finite resource size; rapid binary recombination makes their complete
reaction histories close to those of the signed queue.  A designated
carrier represents the entire ordinary configuration.  The next chapter
puts source, apparatus, actual copies, reset recipients and continuation
inside one autonomous ordinary Hamiltonian and proves historical
faithfulness along that carrier's limiting path.

These are new constitutive and effective results.  They do not derive
the complete interaction catalogue from source/readout incompleteness,
nor identify it with the distinct entropy, chamber or MPBT constitutions
developed later.  The broader physical-access countermodels remain valid
when their extra couplings are admitted.

For a fixed ordinary configuration basis the target is
\begin{equation}\label{p:target}
 J_{YX}(t)=\frac{2}{\hbar}\operatorname{Im}
     \bigl(\Psi_Y(t)^*H_{YX}(t)\Psi_X(t)\bigr),\qquad
 \lambda^B_{Y\leftarrow X}(t)=\frac{[J_{YX}(t)]_+}{w_X(t)},
 \quad w_X(t)=|\Psi_X(t)|^2.
\end{equation}
The quotient is used at positive-weight occupied origins, with nodes
handled by Lemma~\ref{kin:existence}.  All total variations below use the
convention $\TV(P,Q)=\sup_A|P(A)-Q(A)|$ on a stated common output space.
The Bell output is the complete designated ordinary path.  The pilot
medium has additional retained coordinates; conditioning on all of their
initial values makes the finite process deterministic.

Each experiment has a fixed finite graph, a prescribed bounded piecewise
$C^1$ coherent programme with bounded-variation currents, and a fixed
horizon $[0,T]$.  A controller responding to records is included as an
ordinary factor in the same Hamiltonian; it is not an unanalysed adaptive
change to the deterministic source programme.  The autonomous
construction uses a static finite $H_F$.  Pilot resources grow while
the ordinary graph, $H$ and $T$ remain fixed.

\section{The microscopic constitution}
\subsection{Complete state and interaction catalogue}
Let $V$ be the finite configuration set of \emph{all ordinary material systems in the experiment}. A configuration specifies source, actuator, excitation, fuel, loss remnants, working display, archive, every reset receiver, control clock and any finite inaccessible reference. A reference with an actual basis coordinate is included in $V$; already-isolated reference experiments use blocks $H_{\rm local}\otimes I_R$. The explicit preparation variant below first entangles the reference, then proves its isolation beyond a one-way clock boundary. An internal reference fibre is an alternative declared sector convention, not an unnoticed coarse graining of a finer Bell process.

Fix an orientation $e=(r,q)$ for each off-diagonal bond in the union of the programme's nonzero supports and set $b_e=e_q-e_r$, where $(e_r)_{r\in V}$ is the free vertex basis. The incidence matrix $B$ has columns $b_e$. The complete microstate consists of
\begin{equation}
 \left(\Psi,(\chi_e,\Pi_e)_e,(u_e,k_e)_e,
       (X_a)_{a=1}^N,\mathcal P,\mathcal F,\mathcal G,\mathcal R\right).
 \label{p:eq-microstate}
\end{equation}
Here $\mathcal P$ is the finite bank of charged, unused and spent packet slots; $\mathcal F$ contains exporter blank/fuel cells; $\mathcal G$ contains every incoming and outgoing gas coordinate; and $\mathcal R$ contains all contact products and history receivers. One predesignated carrier, say $Q=X_1$, is the actual ordinary configuration. All carrier labels obey identical rules. The remaining carrier positions are pilot degrees of freedom on configuration space, not extra independently prepared copies of the ordinary quantum input.

The laws are the following complete inventory.
\begin{enumerate}[label=\textbf{P\arabic*.}]
\item The normalized canonical field and primitive bond connections obey the action in Section~\ref{p:sec-source}. All ordinary forces, including coherent feedback, enter its numerical Hermitian matrix $H$. An ordinary controller is another factor in that same matrix.
\item The pilot medium has the bounded action exporter, packet spectrum and contact reactions specified below. The reaction list is complete. Scalar carrier response and the common action unit are physical assumptions.
\item There is no additional force vertex $f(\Pi,\mathcal P,\mathcal G,\mathcal R,X_2,\ldots)B_{\rm ordinary}$ in the field energy or a direct classical rewrite of an ordinary display. The pilot medium influences ordinary configurations through the specified carrier motion only. All new ordinary apparatus must be represented in $H$ and obey the same rule.
\item Gas flight and collision contact are deterministic. The only randomness is a declared initial ensemble: independent spatial pilot-gas positions and marks, and an initial carrier ensemble. Conditional on the declared initial field, the joint law factors as gas ensemble times carrier ensemble. No future event times or desired Bell transition probabilities are used in preparation.
\end{enumerate}
P3 is a new interaction law. It is not an instruction to disregard a readable variable after permitting its read coupling. It asserts that such a coupling is absent from the equations. This differentiates the pilot gas from unrestricted ordinary classical matter. Quadratic ordinary energies form a closed algebra,
\[
 \{\Psi^\dagger A\Psi,\Psi^\dagger D\Psi\}
   =\Psi^\dagger[A,D]\Psi/(i\hbar),
\]
which motivates using the same coherent constitution for arbitrary composed apparatus. This algebraic observation does not derive P3. The force catalogue, special species and initial ensemble are explicit new physics.

\subsection{Energy, reversibility and finite resources}
The microscopic theory is a \emph{hybrid} theory: a canonical field, deterministic free flight and explicit deterministic contact/export rules. It is not advertised as a derivation of all these laws from one smooth Hamiltonian. The autonomous ordinary circuit is a single finite Hermitian Hamiltonian. Chapter~\ref{p:smooth} separately supplies a smooth positive kinetic Hamiltonian for a finite contact-permutation module; the main theorem uses the exact hybrid contact law.

Every contact map has a reversible finite lift with a blank receiver: pair the input $(s,\mathrm{blank})$ with $(F_c(s),\mathrm{record}(c,s))$ by a transposition, on a disjoint flagged output bank, and fix unused states. Here $s$ is the finite local logical input (packet species, carrier vertex and local flags), not the continuous complete field state. The outgoing receiver retains that local input. This proves finite reversible logic, not by itself smooth mechanical realizability. The incoming ensemble uses blank receivers, so inverse collisions require a different prepared incoming product. There are at most $M_N$ contacts and $\sum_e B_e$ exports on the promised horizon, hence finite preallocated capacity suffices. Spent particles and cells stay in \eqref{p:eq-microstate}.

For a total definition outside the promised resource horizon, an attempted export after the last unused slot sets a retained exhaustion flag, disables further exports, and leaves $\Pi$ and the now unbounded residue to their continuous evolution. Existing packets can still react. An exhausted contact-record bank sets its own retained flag and uses a fixed identity contact rule. Every exact tie is processed in a fixed order of channel and export labels. These branches make the finite device total; the stated budgets render them unreachable on the proved physical horizon. The final gas particle simply leaves the plane and remains in the outgoing inventory.

For a minimal energy assignment all pilot register states are degenerate and gas momentum is unchanged at ideal contact. The source action energy is conserved for static $H$. An export changes only the decomposition of a continuous stored action into a packet and a bounded remainder, not $\Psi,\Pi$ or this energy. Finite blank registers are consumed as low-entropy resources; a reset never produces them for free. Nondegenerate ordinary fuel and loss accounting is displayed in the material construction. A stronger universal smooth Hamiltonian or empirical material implementation is outside the constitutive claim.

\section{The canonical source interface and two-species export}
\label{p:sec-source}
Adopt the normalized canonical field and primitive connection action
\eqref{kin:action}--\eqref{kin:energy}, with the passive ownership and
binary-additivity hypotheses of Theorem~\ref{kin:ownership}.  At the
prepared connection $\chi(0)=0$ its equations are
\begin{equation}\label{p:eq-ownership}
 i\hbar\dot\Psi=H\Psi,\qquad \dot\chi=0,\qquad
 \dot\Pi_e=J_e,\qquad \dot w=BJ.
\end{equation}
Thus the source supplies each individual Hamiltonian edge current, not
only its divergence.  The loop-current counterexample
\ref{kin:loop} still applies when primitive binary additivity is dropped.
The conserved source moment map is $B\Pi-w$; none of this selects the
new physical contact law before it is specified.

The exporter below uses the same unit action bookkeeping as
Proposition~\ref{kin:export}, but it retains both packet orientations
instead of cancelling opposite arrivals immediately.
Take $k_e(0)=u_e(0)=0$ and $u_e=N(\Pi_e-\Pi_e(0))-k_e$. At a first hit $u_e=s\in\{-1,1\}$, put a packet of species $s$ in the next unused slot, mark its dedicated exporter blank/fuel cell spent with the retained sign and slot identifier, advance $k_e$ by $s$, and set $u_e$ to zero. Both species may remain simultaneously present. No cancellation is part of export. Tie events use a fixed ordering; gas ties with deterministic export times have probability zero. If $L_e\ge\int_0^T|J_e|\dd t$, then
\begin{equation}
 \left\|\frac{k_e}{N}-\int_0^\cdot J_e\dd t\right\|_\infty\le\frac1N,
 \qquad \#\mathrm{exports}_e\le NL_e.
 \label{p:eq-export}
\end{equation}
Indeed the first error is $-u_e/N$ and each full excursion consumes at least $1/N$ of action variation. Bounded nonzero initial residues give $2/N$ error and at most one extra birth. Choose $B_e=\lceil NL_e\rceil+1$ slots before the experiment. A bound from $H,T$ alone can be used, so the apparatus need not know an unknown input vector. The exporter uses finite increments of a canonical coordinate; it does not evaluate the Bell escape rate or supply a stochastic production clock.

\section{Binary species and contact geometry}
\subsection{Routing from a declared charge spectrum}
A carrier at $r$ has vector charge $e_r$. A positive packet on $e=(r,q)$ has charge $b_e=e_q-e_r$, a negative packet has $-b_e$, and spent slots and products are neutral. The elementary service consumes one packet and changes one carrier into one carrier. All other participants are neutral. Then
\[
 e_i+e_q-e_r=e_j
\]
forces $i=r,j=q$: otherwise the coefficient of $e_r$ on the left is negative. Thus the two possible services are
\begin{equation}
 P_e^++C_a@r\longrightarrow C_a@q+\mathrm{spent},\qquad
 P_e^-+C_a@q\longrightarrow C_a@r+\mathrm{spent}.
 \label{p:eq-service}
\end{equation}
Opposite packets can recombine to neutral products. This is routing from stoichiometry and charge, not from the sign of an instantaneous current. Charge does not derive completeness of the binary species list. Charged receivers, multipacket conversion and multicarrier moves would define other theories.

Writing $n_r=\#\{a:X_a=r\}$ and $Z_e=P_e^+-P_e^-$, the complete hybrid inventory
\begin{equation}
 \mathcal C=n+BZ+Bu-Nw
 \label{p:eq-charge}
\end{equation}
is conserved. Between events $\dot u=NJ$ and $\dot w=BJ$ cancel. A signed export changes $u_e$ by $-s$ and $Z_e$ by $s$; service changes $n$ by $sb_e$ and $Z_e$ by $-s$; recombination changes neither $n$ nor $Z$. Initially $\mathcal C=n(0)-Nw(0)$. This is an exact inventory identity, not an unsupported claim that the whole hybrid theory has a common Noether action.

\subsection{Deterministic candidate contacts}
Allocate a fixed channel for every potential pair $(e,b,a)$ of packet slot and carrier, with frequency parameter $r_{eba}=\kappa_e\mu_N/N$. Allocate a channel for every unordered pair of slots on edge $e$, with parameter $a_e>0$. A contact tests only the current species and carrier vertex. An eligible service implements \eqref{p:eq-service}; opposite slot species recombine and retain their signs and identities in the outgoing product; other contacts are null. The total candidate frequency is
\begin{equation}
 R_N=\sum_e\left(\kappa_e\mu_NB_e+a_e\binom{B_e}{2}\right).
 \label{p:eq-candidate}
\end{equation}
If the graph has no off-diagonal bonds, $R_N=0$, no beam is needed and the ordinary configuration is constant. Otherwise $R_N>0$. Partition a transverse cross-section into cells of relative areas $r_c/R_N$ for these channels. All frequencies and areas are fixed independently of $\Psi,w,J$ and the future material history. The scalar $\kappa_e$ expresses equal response for all carriers on that bond. There is no rule suppressing a minority packet's service.

Prepare $M_N$ distinguishable pilot particles with independent positions uniform on $[-L_N,0]$, equal speed $v>0$, and independent uniform transverse coordinates. Put $L_N=M_Nv/R_N>vT$. Each incoming particle freely crosses the contact plane once, at $t_i=-z_i/v$, and its transverse coordinate selects a channel. Its outgoing state and blank receiver are retained. The microdynamics from all initial coordinates is deterministic.

\paragraph{Which contact marks are compared.}
The primary input is the ordered contact history with channel marks
identifying the contacted edge, packet slot(s) and carrier.  The initialized
reaction device is independent of the gas conditional on the declared
coherent preparation.  Its output follows one common measurable causal
map with the stated tie and virtual-overflow conventions.  Neither all
initial gas positions nor conditioning on them belongs to this comparison.

If the physical identity of each incoming particle is also retained in
the output, decorate both count-conditioned laws identically.  Conditional
on $k\le M_N$ arrivals, their particle identities are a uniformly ordered
$k$-subset of the $M_N$ labels, independent of the ordered times and
channel marks.  For the Poisson comparison, draw one independent uniform
permutation of those labels, use its initial segment, and assign new
virtual receiver identities if $k>M_N$.  This is a common conditional
kernel; it does not make particle identities iid marks.  Channel marks
remain independent with the stated probabilities, and chemistry is
independent of the beam identity, so the reaction rates below are unchanged.

\section{Deriving complete timing from the spatial ensemble}
For comparison only, extend the reaction-map domain to a countably padded bank of virtual blank contact receivers. The first $M_N$ cells are the physical bank; the further cells have no counterpart in the physical finite beam, which cannot reach them. A Poisson contact sequence can use those mathematical extra cells, preserving its uncompromised reaction generator. Both contact histories are fed through this one total causal map. This convention avoids assuming that a Poisson count has a finite deterministic bound. It neither adds physical particles nor drops physical receivers. Accepted service and recombination events still have the finite stock bound below.

\begin{theorem}[Finite-gas marked-history bound]\label{p:gas}
Under the independent initialization and common output conventions just specified, on $[0,T]$ the complete marked contact process differs in total variation from a marked Poisson process of channel frequencies $r_c$ by at most
\begin{equation}
 \Delta_{\rm gas}=\frac{(R_NT)^2}{M_N}.
 \label{p:eq-gaserror}
\end{equation}
The same bound holds after any common measurable causal deterministic contact device, including its retained reaction state and exact event times. The compared output includes the local contact receivers, indexed in contact order; it does not include the incoming gas's unobserved continuous coordinates. Those remain in the physical microstate.
\end{theorem}
\begin{proof}
The number of beam crossings is $\operatorname{Bin}(M_N,p)$, $p=R_NT/M_N$. Conditional on that number, the ordered times are ordered independent uniforms on $[0,T]$ and marks have probabilities $r_c/R_N$. The Poisson process has exactly the same conditional distribution given its count, so contact-history total variation equals count total variation. A Bernoulli$(p)$ variable and a Poisson$(p)$ variable have distance $p(1-e^{-p})\le p^2$: their only excess Bernoulli mass is at one. Coupling $M_N$ independent such pairs and summing gives distance at most $M_Np^2$. On equal contact histories, the deterministic device evolves identically, including eligibility and all null records. Projection therefore contracts this bound.
\end{proof}
The finite gas has neither memoryless waits nor an inserted event tape. Given $K_t$ observed arrivals, its total conditional hazard is
\[
 \frac{M_N-K_{t-}}{M_N/R_N-t},
 \qquad
 \Prb(\text{no arrival in }(t,t+h]\mid\mathcal F_t)
  =\left(1-\frac{h}{M_N/R_N-t}\right)^{M_N-K_t}.
\]
The survival formula has $0\le h\le M_N/R_N-t$; the left limit in the hazard makes it predictable. These formulas follow by conditioning the remaining independent positions. In the Poisson comparison, independence of channel increments and predictable eligibility give, for a bounded state function $f$,
\begin{equation}
 \mathcal Gf(s)=\sum_c r_c\,[f(F_c(s))-f(s)]
 \label{p:eq-generator}
\end{equation}
as the jump part of the predictable compensator, almost everywhere in time. The full evolution also contains the specified deterministic flow and export rewrites. This is the derivation of both chemical clocks. It is stronger than a mean collision-frequency calculation. The spatial independence assumption is indispensable: equally spaced particles with a uniform global translation have the same uniform one-particle marginals but different complete waiting-time laws (Chapter~\ref{p:chapter-discrimination}).

\section{Recombination and the removal of surplus}
For the Poisson contact comparison, every eligible packet--carrier pair has derived rate $\kappa_e\mu_N/N$, and every opposite slot pair has derived rate $a_e$. Let $P_e^\pm$ be species counts, $m_e^{\rm mix}=\min(P_e^+,P_e^-)$, and $A_e,D_e$ the recombination and service counts. Exactly
\begin{equation}
 P_e^+(t)+P_e^-(t)+2A_e(t)+D_e(t)=\#\mathrm{births}_e(t)\le B_e.
 \label{p:eq-stock}
\end{equation}
The accepted reaction process is nonexplosive. Both service directions are active whenever both species and their origins are populated.

\begin{theorem}[Full reaction-state recombination comparison]\label{p:recombination}
Assume initially empty packet stock and the stated per-edge birth budgets. Let $S^a$ denote the reaction state of the Poisson comparison: field, actions, residues, all packet slots, carriers and reaction-product receivers, including null-contact records. It does not denote the original gas flight coordinates. The two-species reaction process has a lifted comparison process whose aggregate is exactly the instantaneous signed-queue model, with
\begin{equation}
 \TV\bigl(\pLaw(S^{a}),\pLaw(S^{\rm lift})\bigr)
 \le \Delta_{\rm ann}:=\sum_e\frac{\kappa_e\mu_NB_e}{2a_e}.
 \label{p:eq-annerror}
\end{equation}
Both processes retain slots, products and receivers. In particular this bound holds for their complete aggregate carrier paths and exact event times. It holds for any prescribed signed birth sequence of the stated budgets, including reversals and arbitrarily close births.
\end{theorem}
\begin{proof}
In each state match all minority packets with the same number of majority packets by a fixed ordering of their physical labels. This matching is a proof device. The lifted reference has the same births and all the same recombination channels, but services only the $|Z_e|$ unmatched excess packets. Let $\pi$ retain the field, actions, residues, net queues $Z$, all carrier coordinates and the common edge/carrier/direction service marks. It omits physical recombination times and does not identify their product archive with an instantaneous-cancellation archive. Recombination leaves $Z_e=P_e^+-P_e^-$ invariant; excess service changes it exactly as a signed-queue service. For each carrier at the excess-species origin there are $|Z_e|$ possible packet partners. Consequently the generator on functions of $\pi$ is exactly the signed-queue generator, independent of hidden matching and slot labels. This explicitly proves lumpability for this projection.

Couple all common transitions while full states agree. The physical process additionally services matched packets, with discrepancy hazard
\[
 h_e=\frac{\kappa_e\mu_N}{N}m_e^{\rm mix}(n_r+n_q)\le\kappa_e\mu_Nm_e^{\rm mix}.
\]
Continue correct marginals after the first such event. Since $m_e^{\rm mix}\le P_e^+P_e^-$ and the annihilation compensator gives
\[
 a_e\E\int_0^T P_e^+P_e^-\dd t=\E A_e(T)\le B_e/2,
\]
the stopping time $\tau$ of the first discrepancy obeys
\[
 \Prb(\tau\le T)\le
 \E\int_0^{T\wedge\tau}\sum_e h_e(t)\dd t
 \le\sum_e\kappa_e\mu_N\E\int_0^T P_e^+(t)P_e^-(t)\dd t
 \le\Delta_{\rm ann}.
\]
The stopped integral is bounded by the complete physical marginal
integral.  Probability of any discrepancy bounds entire-path total variation. The construction agrees on every retained receiver up to that discrepancy.
\end{proof}
The physical finite-$a$ aggregate is generally not Markov in $Z$ and carriers alone; its rates depend also on the mixed stock. No closure is assumed for that projection. An additional direct consequence is
\[
 \Prb(\text{any service while an edge is mixed})\le\Delta_{\rm ann},
\]
because, on $m_e^{\rm mix}>0$, its total service rate is at most $\kappa_e\mu_N\max(P_e^+,P_e^-)\le\kappa_e\mu_NP_e^+P_e^-$. Minority suppression follows from a competition of physical reaction speeds, not an imposed positive-part gate.

\section{The full Bell limit and its conditional history}\label{p:sec-mainlimit}
Define $\epsilon_{\rm kin}(N,T)$ to be the signed-queue tagged-path bound \eqref{kin:pathbound}, proved in Theorem~\ref{kin:path}. It tends to zero for fixed finite graph and programme if
\begin{equation}
 \mu_N\longrightarrow\infty,\qquad \mu_N/N\longrightarrow0,
 \qquad \E\|x^N(0)-w(0)\|_1\longrightarrow0.
 \label{p:eq-kineticscale}
\end{equation}
Initialize the tag from $\nu\le Cw(0)$ and the other $N-1$ carriers independently from $w(0)$, independently of the tag and gas conditional on the field. Equilibrium use takes $\nu=w(0)$, so all carriers are iid. The expected initial census error is at most $\sqrt{|V|/N}+2/N$, and the same tracking proof applies. This is an initial ensemble assumption, not a new draw at each event, a physical unknown-state sampling device, or a derivation of equilibrium preparation from dynamics. Deterministic counts with exchangeable labels instead give tag law $x^N(0)$, not exactly $w(0)$, and require adding their initial-law discrepancy if that variant is used.

\begin{theorem}[Deterministic microscopic selection of the Bell path]\label{p:main}
For the finite pilot-medium theory P1--P4, the entire designated ordinary-configuration path satisfies
\begin{equation}
 \TV\bigl(\pLaw(Q^N_{[0,T]}),\mathbb P^B_{H,\nu,[0,T]}\bigr)
 \le \Delta_N:=\frac{(R_NT)^2}{M_N}
       +\sum_e\frac{\kappa_e\mu_NB_e}{2a_e}
       +\epsilon_{\rm kin}(N,T).
 \label{p:eq-mainerror}
\end{equation}
In fixed physical units, $B_e=O(N)$, $\mu_N=N^{1/2}$, $a_e=N^2$, and $M_N=N^{10}$ give $R_N=O(N^4)$, gas error $O(N^{-2})$, recombination error $O(N^{-1/2})$, and $\Delta_N\to0$. No change of time accompanies this limit. The target is the nonexplosive minimal Bell process \eqref{p:target} with its complete natural conditional timing law.
\end{theorem}
\begin{proof}
Apply Theorem~\ref{p:gas} to the full causal source, exporter and contact maps. In the Poisson comparison these produce the per-slot mass-action reactions of Theorem~\ref{p:recombination}. Its lifted reference projects exactly to the signed queue of Theorems~\ref{kin:tracking} and \ref{kin:path}. Apply the latter theorem to the designated carrier and use the triangle inequality. The growth estimates follow from \eqref{p:eq-candidate}. Lemma~\ref{kin:existence} and the nodal coupling in Theorem~\ref{kin:path} identify the entire limiting law.
\end{proof}

For the limit, conditional on the declared initial field and programme, the full ordinary history through $t$ and current state $X$, the probability of holding at $X$ until $t+h$ is
\begin{equation}
 \exp\left[-\int_t^{t+h}\sum_{Y\ne X}\frac{[J_{YX}(s)]_+}{w_X(s)}\dd s\right]
 \label{p:eq-survival}
\end{equation}
on its positive-weight component. The first-event type has the corresponding competing-hazard density. Node localization supplies continuation when a component ends. The finite theory is history dependent through its unobserved gas and packets; the theorem does not claim pointwise convergence of conditional kernels on all rare pasts. It proves total variation of complete paths, which implies convergence for every measurable history event and every common stopped output. Conditional events of ideal probability $p>0$ have error at most $2\Delta_N/p$ when $\Delta_N<p$.

Individual edge ownership, absence of surplus, and timing thus have different proven origins. The primitive action supplies each $J_e$, not merely $BJ$; fast opposite-charge recombination removes excess physical directional service; independent incoming positions and scalar contact counting supply the complete chemical timing law; population tracking then supplies the residence denominator. In particular the denominator is not evaluated by a microscopic particle.

\paragraph{Complete configurations and projections.}
The theorem concerns $Q$, which already contains every ordinary source, archive, receiver, controller and reference coordinate of the declared experiment. It does not identify the Bell law with the full microstate \eqref{p:eq-microstate}. All pilot histories remain physically present under their different force law. A later programme that reads or returns pilot products is not covered by silently tracing them out. P3 rules out an extra ordinary read force; changes of pilot dynamics or resource recirculation require a new estimate. Coarse records contract \eqref{p:eq-mainerror}, but their jump rates need not be the positive part of a sum of fine currents.


\chapter{Autonomous material records in the pilot theory}\label{p:chapter-records}
\label{p:mat-section}

The pilot mechanism is applied to one finite graph containing the source,
apparatus, all receiving systems, inaccessible reference and a clock.  The
following construction specifies its material Hamiltonian and proves an
actual historical-record property of its Bell limit.  No continuum pointer
law or classical reader of a pilot coordinate is appended.  Circuit
Hamiltonians and engineered state-transfer chains provide useful context
\cite{pFeynman1986,pChristandl2004}; all facts used here are proved below.

\section{Exact autonomous propagation}

Let $U_0,\ldots,U_{\ell-1}$ be fixed unitaries on a finite material space
$\mathcal K$, including every resource and a reference $R$.  During the
portion declared to have an inaccessible reference, each gate acts as
the identity on $R$.  Proposition~\ref{p:prepared} also permits an
explicit earlier preparation stage involving $R$.  Set
\begin{equation}\label{p:mat-prefix}
 V_0=I,\qquad V_n=U_{n-1}\cdots U_0,\qquad
 c_n=\sqrt{(n+1)(\ell-n)}.
\end{equation}
The clock has basis $|0\rangle,\ldots,|\ell\rangle$.  For a frequency
$\Omega>0$, define the static matrix
\begin{equation}\label{p:mat-feynman}
 H_F=\hbar\Omega\sum_{n=0}^{\ell-1}c_n
 \left(|n+1\rangle\langle n|\otimes U_n+
 |n\rangle\langle n+1|\otimes U_n^\dagger\right).
\end{equation}
All its matrix edges are ordinary edges of the common canonical field.
In particular, the pilot link meters and reactions use the currents of
this complete $H_F$, rather than those of a source Hamiltonian with an
external ideal clock suppressed.

\begin{theorem}[Exact finite-clock programme]\label{p:mat-clock}
Starting with $|0\rangle\otimes\psi$, $\|\psi\|=1$, the field under
\eqref{p:mat-feynman} is
\begin{align}\label{p:mat-wave}
 \Psi(t)&=\sum_{n=0}^{\ell}\phi_n(t)|n\rangle\otimes V_n\psi,\\
 \phi_n(t)&=(-i)^n\sqrt{\binom\ell n}
   \cos^{\ell-n}(\Omega t)\sin^n(\Omega t).\nonumber
\end{align}
At the first transfer time
\begin{equation}\label{p:mat-time}
 T=\frac{\pi}{2\Omega},
\end{equation}
the state is $(-i)^\ell|\ell\rangle\otimes V_\ell\psi$.
Moreover $H_F+\hbar\Omega\ell I\ge0$; this shift changes no configuration
current.  The clock and all its correlations remain in the full model.
\end{theorem}
\begin{proof}
With $D=\sum_n|n\rangle\langle n|\otimes V_n$,
\[
 D^\dagger H_FD=H_C\otimes I,\qquad
 H_C=\hbar\Omega\sum_nc_n
       (|n+1\rangle\langle n|+|n\rangle\langle n+1|).
\]
On the permutation-symmetric subspace of $\ell$ qubits, $H_C$ is the
restriction of $\hbar\Omega\sum_{j=1}^\ell X_j$; the normalized state
with $n$ excitations has the displayed adjacent matrix element $c_n$.
Evolving $|0\rangle^{\otimes\ell}$ therefore gives
$(\cos(\Omega t)|0\rangle-i\sin(\Omega t)|1\rangle)^{\otimes\ell}$.
Its normalized symmetric coefficients prove \eqref{p:mat-wave} and
\eqref{p:mat-time}.  The spectrum of the qubit sum lies in
$[-\hbar\Omega\ell,\hbar\Omega\ell]$, proving the lower bound.
\end{proof}

\begin{proposition}[Input-independent shape of the clock weights]
\label{p:mat-levels}
For each complete material basis state $x$,
\begin{equation}\label{p:mat-weights}
 w_{n,x}(t)=\binom\ell n\cos^{2(\ell-n)}(\Omega t)
             \sin^{2n}(\Omega t)\,|(V_n\psi)_x|^2.
\end{equation}
On $[0,T]$, each positive regular level of each weight has at most two
crossings, independently of the unknown input.  On $[0,2T]$ it has at
most four.  Every nonzero coordinate has strictly positive weight in
the interior of the first pass; a vanishing coefficient
$(V_n\psi)_x$ gives an identically empty coordinate instead.
\end{proposition}
\begin{proof}
The factor depending on $\psi$ is a nonnegative constant.  For
$0<n<\ell$, logarithmic differentiation of the other factor gives
$2\Omega[n\cot(\Omega t)-(\ell-n)\tan(\Omega t)]$, which vanishes
once, at $\sin^2(\Omega t)=n/\ell$, and changes from positive to
negative.  For $n=0$ or $\ell$ the factor is monotone.  Reflection
around $T$ gives the second-pass count.  Positivity on $(0,T)$ follows
directly from \eqref{p:mat-weights}.
\end{proof}
The crossing count is a useful uniform input fact for a kinetic estimate.
By itself it is not a proof of every other uniform constant required by
that estimate.

\section{A faithful archive at a monomial clock cut}

A unitary is \emph{monomial} in the complete material basis when
\begin{equation}\label{p:mat-monomial}
 (U_m)_{yx}=e^{i\vartheta_x}\,1_{\{y=\pi(x)\}}
\end{equation}
for a permutation $\pi$.  Reversible copies and SWAPs are examples.

\begin{theorem}[Historical record at a monomial cut]\label{p:monomial}
For the Bell process of \eqref{p:mat-feynman} in initial equilibrium, a
monomial cut $m$ is crossed exactly once, from clock $m$ to clock $m+1$,
almost surely before $T$.  At that crossing the actual material
configuration is updated by $\pi$.  The material state immediately
before the crossing has law $|(V_m\psi)_x|^2$.

Suppose this permutation copies a working key $W$ into a blank archive
$A$, all earlier gates preserve its blank state, and all later gates
preserve the archive label.  Then the actual archive contains the
actual $W$ key at that crossing and stays unchanged for the rest of
the first pass.  A later monomial SWAP into a retained blank receiver
transfers the actual old working key into that receiver at its own
unique crossing.
\end{theorem}
\begin{proof}
Write $\phi_n=(-i)^n a_n$ with $a_n(t)>0$ on $(0,T)$, and put
$\xi_n=V_n\psi$.  The fine current at a complete edge across cut $m$ is
\begin{equation}\label{p:mat-finecurrent}
 J_{(m+1,y),(m,x)}(t)=2\Omega c_ma_{m+1}a_m
   \operatorname{Re}\left[(\xi_{m+1})_y^*(U_m)_{yx}(\xi_m)_x\right].
\end{equation}
For \eqref{p:mat-monomial}, the nonzero real factor is exactly
$|(\xi_m)_x|^2$.  Every current across that cut is therefore forward,
and its reverse Bell rate vanishes.  Since the clock initially lies
below the cut and finally lies above it with probability one, the
cut is crossed exactly once.  The permitted edge carries exactly
the permutation $\pi$.

Summing the forward current over $x$ gives
$2\Omega c_ma_{m+1}a_m$.  It is the time derivative of the field mass
strictly above the cut and integrates to one.  Integrating the
individual current thus gives $|(\xi_m)_x|^2$ for the pre-crossing
material state.  Once the process has crossed, it cannot return to
the earlier region.  All edges in the remaining region preserve $A$
by the later-gate hypothesis.  This proves the historical statement.
The same argument applies to the receiver SWAP.
\end{proof}

For completeness, the crossing-time density is
$2\Omega c_ma_{m+1}(t)a_m(t)$.  Under $z=\sin^2(\Omega t)$ it becomes
\[
 \frac{z^m(1-z)^{\ell-m-1}}{B(m+1,\ell-m)}\,\dd z.
\]
This is a derived clock-time law, not an additional random time draw.

\begin{remark}[Fine traffic is not coarse clock traffic]
\label{p:mat-finewarning}
For a general unitary $U_m$, the real factor in
\eqref{p:mat-finecurrent} can be negative.  The net clock flux can be
forward while some fine edges point backward.  For example, let a
Hadamard act on $S$ in the state
$\sqrt{2/3}|0,0_R\rangle+\sqrt{1/3}|1,+_R\rangle$.  At the fine edge
whose old and new source bits both equal one and whose reference bit
is zero, the real factor is $-1/12$.  Thus
Theorem~\ref{p:monomial} uses the monomial hypothesis essentially.
In particular, a later archive does not record every transient
excursion of an earlier nonmonomial resource gate.
\end{remark}

\section{Retained resources, reset and continuation}
\label{p:mat-resources}

\subsection{Finite production, fuel, capture, pending and loss states}

Let $S$ be a qubit with $P_a=|a\rangle\langle a|$, $a=0,1$.
The resource space $D$ has the orthonormal basis
\[
 |r\rangle,|p_0\rangle,|p_1\rangle,
 |c_0\rangle,|c_1\rangle,|l_0\rangle,|l_1\rangle.
\]
These are complete keys for the following local resources.  A unit of
production energy, fuel energy or pending/lost excitation has energy
$E>0$; a captured remnant has energy $2E$.
\begin{center}
\begin{tabular}{c|c c c c c}
Key & Production & Fuel & Site & Signal & Retained product\\\hline
$r$ & $E$ & $E$ & ready & none & none\\
$p_a$ & $0$ & $E$ & ready & pending $E$, mode $a$ & none\\
$c_a$ & $0$ & $0$ & spent, label $a$ & none & capture remnant $2E$\\
$l_a$ & $0$ & $E$ & ready & none & lost excitation $E$, mode $a$
\end{tabular}
\end{center}
Thus every listed complete resource configuration has total energy
$2E$.  The seven-dimensional space may be regarded as this sector of a
larger tensor inventory; the gates below extend as the identity on
its unused orthogonal complement.  The loss remnant is retained, not
traced away physically.

For $0<\eta<1$, define
\begin{align}\label{p:mat-generators}
 T_w&=\sum_{a=0}^1P_a\otimes
       (|p_a\rangle\langle r|-|r\rangle\langle p_a|),\\
 |b_a\rangle&=\sqrt\eta\,|c_a\rangle+
                   \sqrt{1-\eta}\,|l_a\rangle,\nonumber\\
 T_c&=\sum_{a=0}^1
       (|b_a\rangle\langle p_a|-|p_a\rangle\langle b_a|).
       \nonumber
\end{align}
They are real antisymmetric matrices.  The full finite unitaries are
\begin{equation}\label{p:mat-rotations}
 U_w(\theta)=I+\sin\theta\,T_w+(1-\cos\theta)T_w^2,
 \qquad
 U_c(\varphi)=I+\sin\varphi\,T_c+(1-\cos\varphi)T_c^2.
\end{equation}
Equivalently they are $e^{\theta T_w}$ and $e^{\varphi T_c}$; their
Hermitian generators are $iT_w$ and $iT_c$.  In particular
\[
 U_w|a,r\rangle=\cos\theta|a,r\rangle+
                       \sin\theta|a,p_a\rangle,
 \qquad
 U_c|p_a\rangle=\cos\varphi|p_a\rangle+
                       \sin\varphi|b_a\rangle.
\]
The orthogonal vector
$\sqrt{1-\eta}|c_a\rangle-\sqrt\eta|l_a\rangle$ is fixed by $U_c$.
These equations specify the gates on the full space, rather than an
isometry with unmentioned complementary states.

For any source/reference vector $\psi$, the first two gates give
\begin{equation}\label{p:mat-resourcevector}
 \cos\theta\,\psi|r\rangle+
 \sin\theta\sum_aP_a\psi
 \left[\cos\varphi|p_a\rangle+
 \sin\varphi\left(\sqrt\eta|c_a\rangle+
                       \sqrt{1-\eta}|l_a\rangle\right)\right].
\end{equation}
All coefficients retain their source and reference cofactors.  The
entire vector, including ready and pending sectors, is subsequently
used in the clock Hamiltonian.

\subsection{Nine exact gates with every receiving system retained}

Let $W,A,B_W$ be ternary registers, initially $0$, let $B_D$ be a
second seven-state resource initially $r$, and let $K$ be a qubit
initially $0$.  $W$ is the working status display, $A$ its archive,
$B_D,B_W$ are reset receivers, and $K$ is a later readout.  The full
ready state is
\begin{equation}\label{p:mat-ready}
 \psi_{SR}|r\rangle_D|0\rangle_W|0\rangle_A
 |r\rangle_{B_D}|0\rangle_{B_W}|0\rangle_K|0\rangle_C.
\end{equation}
The independent $B_D=r$ consumes another production/fuel supply of
energy $2E$.  Thus the two resource banks initially contain $4E$,
and their energy is conserved by the specified gates.  The register
labels and source/reference levels may be degenerate; the full
static clock-interaction energy is conserved separately.

Define
\begin{equation}\label{p:mat-status}
 \chi(r)=\chi(p_a)=\chi(l_a)=0,\qquad
 \chi(c_0)=1,\qquad\chi(c_1)=2.
\end{equation}
Ternary additions below are modulo three.  The nine gates are
\begin{align}\label{p:mat-circuit}
 U_0&=U_w(\theta),& U_1&=U_c(\varphi),\nonumber\\
 U_2: |d,w\rangle&\longmapsto|d,w+\chi(d)\rangle,
 &U_3: |w,a\rangle&\longmapsto|w,a+w\rangle,\nonumber\\
 U_4&=\operatorname{SWAP}_{D,B_D},
 &U_5&=\operatorname{SWAP}_{W,B_W},\\
 U_6&=|0\rangle\langle0|_A\otimes I_S
       +|1\rangle\langle1|_A\otimes I_S
       +|2\rangle\langle2|_A\otimes V,
 &V&=e^{-i\pi\sigma_y/6},\nonumber\\
 U_7&=H_S^{\rm Had},
 &U_8: |s,k\rangle&\longmapsto|s,k\mathbin{\oplus}s\rangle.
 \nonumber
\end{align}
Every unspecified factor is a spectator, and every gate acts as
$I_R$.  Here
$H^{\rm Had}=2^{-1/2}\left(\begin{smallmatrix}1&1\\1&-1\end{smallmatrix}\right)$.
Gate $U_6$ is coherent control by the archive operator; no actual
pilot log or actual classical archive value is inserted into the
field equation.

The clock now has $\ell=9$ and ten nodes.  For a reference qubit the
material space excluding the clock has dimension
$2\cdot2\cdot7\cdot3\cdot3\cdot7\cdot3\cdot2=10584$; the complete
clock/material graph has $105840$ vertices.  Its matrix is explicitly
\eqref{p:mat-feynman} with the gates \eqref{p:mat-circuit}.  Tensor and
permutation notation specify all entries without concealing a
postselected subspace.

\begin{proposition}[Actual facts retained by this circuit]
\label{p:mat-resourcehistory}
In its Bell limit, $W$ receives $\chi(D)$ at the unique crossing of
cut $2$.  This actual working value persists until its reset at
cut $5$.  At cut $3$, $A$ copies that value and retains it through
the end of the first pass.  At cut $4$, the actual old $D$ key moves
to $B_D$ and $D$ becomes ready.  At cut $5$, the actual old $W$ key
moves to $B_W$ and $W$ becomes blank.  The receiving keys remain
retained.  Gate $U_8$ similarly records the source computational key
after the coherent Hadamard continuation.
\end{proposition}
\begin{proof}
The cuts $2,3,4,5,8$ are monomial.  Before cut $2$, $W$ is blank;
the copying permutation leaves $D$ unchanged.  Gates $3$ and $4$
preserve $W$, and a return across cut $2$ is impossible.  Hence the
later archive copy records the same earlier actual status.
Before cut $3$, $A$ is blank, and every later gate preserves its
label.  Gates after $4$ preserve $B_D$; those after $5$ preserve
$B_W$.  Apply Theorem~\ref{p:monomial} at each cut.  The final cut has
no later gate.  These arguments identify actual past facts, not
merely correlations of final field weights.
\end{proof}
The status is a fact at its registered cut.  Preparation gates $0$
and $1$ can have provisional excursions; this proposition does not
claim that every such excursion was recorded.  A null archive value
means no captured key at the registered status cut.  Clock progress
distinguishes that completed null read from the initially blank
register.

\subsection{A null result with a reference-sensitive continuation}

Choose
\begin{equation}\label{p:mat-input}
 \psi_{SR}=\sqrt{\frac23}|0\rangle|0_R\rangle+
           \sqrt{\frac13}|1\rangle|+_R\rangle,
 \qquad |+_R\rangle=\frac{|0_R\rangle+|1_R\rangle}{\sqrt2},
\end{equation}
and $\theta=\pi/3$, $\varphi=\pi/4$, $\eta=2/3$.
The resource weights in \eqref{p:mat-resourcevector} are respectively
ready $1/4$, pending $3/8$, captured $1/4$ and loss $1/8$.
All are retained through the two SWAPs.  For example, after reset
the old pending or loss distinction resides in $B_D$ even though
the working $D$ has returned to $r$.

Let $\mathsf N=\{A=0\}$ be the null event.  Tracing the retained resource
factors only for this calculation gives the null subchannel
\begin{equation}\label{p:mat-nullmap}
 \mathcal N(\rho)=\frac14\rho+
          \frac12\sum_{a=0}^1P_a\rho P_a.
\end{equation}
This reduction is not the state used for the complete evolution.
The ready term retains source coherence; the distinct pending and
loss labels give the displayed dephased contribution.  The null
branch sees the identity in $U_6$, followed by a Hadamard and a
computational copy.  Writing $X+$ for $K=0$, one obtains
\begin{equation}\label{p:mat-nullnumbers}
 \Pr(\mathsf N)=\frac34,\qquad
 \Pr(\mathsf N,X+)=\frac{11}{24},\qquad
 \Pr(X+\mid\mathsf N)=\frac{11}{18}.
\end{equation}
The unnormalized inaccessible-reference state in that joint outcome is
\begin{equation}\label{p:mat-reference}
 \sigma_R^{\mathsf N,+}=\frac1{48}
       \begin{pmatrix}19&5\\5&3\end{pmatrix}.
\end{equation}
\begin{proof}
The source reduction of \eqref{p:mat-input} is
$\rho_S=\left(\begin{smallmatrix}2/3&1/3\\1/3&1/3\end{smallmatrix}\right)$.
Equation~\eqref{p:mat-nullmap} gives
$\mathcal N(\rho_S)=\left(\begin{smallmatrix}1/2&1/12\\1/12&1/4\end{smallmatrix}\right)$.
Its trace is $3/4$, and its $|+\rangle$ diagonal element is $11/24$.
For the reference calculation use
\[
 |r_0\rangle=\sqrt{2/3}|0_R\rangle,
 \qquad |r_1\rangle=\sqrt{1/3}|+_R\rangle.
\]
The ready contribution after the $+$ effect is
$(|r_0\rangle+|r_1\rangle)(\langle r_0|+\langle r_1|)/8$;
the other null contribution is
$(|r_0\rangle\langle r_0|+|r_1\rangle\langle r_1|)/4$.
Adding the two matrices gives \eqref{p:mat-reference}.
\end{proof}

The captured branches exercise noncommuting record-controlled
continuation.  Writing $a=0$ for $A=1$ and $a=1$ for $A=2$, the
four joint probabilities are
\begin{align}\label{p:mat-capturednumbers}
 \Pr(a=0,X+)&=\Pr(a=0,X-)=\frac1{12},\\
 \Pr(a=1,X+)&=\frac{2-\sqrt3}{48},&
 \Pr(a=1,X-)&=\frac{2+\sqrt3}{48}.\nonumber
\end{align}
Indeed capture has probability $1/4$ times the original source
population.  The $a=0$ daughter is $|0\rangle$ and gives equal $X$
probabilities.  The $a=1$ daughter becomes
$-\tfrac12|0\rangle+\tfrac{\sqrt3}{2}|1\rangle$ under $V$, yielding
the other two numbers.  They sum to $1/4$.  The complete reference
cofactors and both reset receivers remain attached throughout.

\subsection{What a finite material-path bound now proves}

\begin{proposition}[A basis-ready example with randomness only in the gas]
\label{p:prepared}
The concrete experiment above can start from one known complete ordinary
basis configuration, with every pilot carrier in that configuration.
Prepend to the nine gates \eqref{p:mat-circuit} the two gates
\begin{align}\label{p:mat-prepgates}
 U_{\rm prep}
 &=\left(P_0^S\otimes I_R+P_1^S\otimes H_R^{\rm Had}\right)
       \left(R_y^S(\alpha)\otimes I_R\right),\\
 \alpha&=2\arccos\sqrt{2/3},\qquad
 R_y(\alpha)=e^{-i\alpha\sigma_y/2},\qquad
 U_{\rm barrier}=I.\nonumber
\end{align}
Initialize $S,R$ in $|0,0_R\rangle$, and initialize all resource,
receiver, display and clock factors in the basis-ready states already
listed.  There are now eleven gates, twelve clock nodes and $127008$
complete ordinary configurations for a reference qubit.  The first-pass
duration is still $T=\pi/(2\Omega)$ for the new clock Hamiltonian.

In the Bell limit, the identity cut is crossed exactly once.  At that
crossing the source/reference configuration has law $|\psi_{SR}|^2$
for \eqref{p:mat-input}, with all remaining material resources still ready.
Thereafter the process cannot return to either preparation gate, and
every accessible later edge acts as the identity on $R$.  The resource,
archive, reset and continuation conclusions, including
\eqref{p:mat-nullnumbers}--\eqref{p:mat-capturednumbers}, are unchanged.
\end{proposition}
\begin{proof}
The sign convention in \eqref{p:mat-prepgates} gives
\[
 R_y(\alpha)|0\rangle=\sqrt{2/3}|0\rangle+
                                \sqrt{1/3}|1\rangle.
\]
The following controlled Hadamard therefore maps $|0,0_R\rangle$
exactly to \eqref{p:mat-input}.  The clock proof applies to these eleven
unitaries without modification: it does not require a spectator
reference during the explicitly designated preparation stage.
The new identity gate is the monomial cut $m=1$.
Theorem~\ref{p:monomial} gives its unique forward crossing and its
pre-crossing distribution, namely the modulus square of the prepared
vector.  No reverse current crosses this barrier on the first pass.
All later gates are precisely the earlier nine gates and are the
identity on $R$.  Their gate prefixes and final product act on the
same prepared vector as before.  The old monomial cuts are merely
shifted upward by two, so their historical proofs remain valid.
Finally $12\times10584=127008$.
\end{proof}

Here the initial field weight is a point mass at a known ordinary
configuration.  Taking all $N$ carriers there gives exact initial
census and tag equilibrium without a random carrier preparation.
Only the declared initial pilot-gas ensemble is random.  The preparation
uses the same static field Hamiltonian and pilot reactions as the rest
of the experiment; it does not resample a configuration at the barrier.
The former record cuts $2,3,4,5,8$ are $4,5,6,7,10$ in this variant.
The physical crossing-time densities use the new $\ell=11$ clock and
therefore change; the retained outcome probabilities remain the same.
The full Hamiltonian includes the earlier interaction with $R$, so the
reference is inaccessible only after the barrier, not throughout its
preparation.  This concrete example does not derive the general
unknown-input equilibrium postulate.  A finite pilot approximation
inherits the prepared crossing distribution and subsequent record
claims through its one full material-path error bound, including a
failure label if it has not reached the relevant cut by $T$.

\begin{corollary}[One error bound for the complete retained programme]
\label{p:programme}
Suppose the finite pilot construction for the complete static graph
\eqref{p:mat-feynman} obeys
\[
 \TV\bigl(\pLaw(Q^{\rm pilot}_{[0,T]}),
            \pLaw(Q^{\rm Bell}_{[0,T]})\bigr)\le\delta.
\]
Then every joint ordinary material record/history event in this
programme has probability error at most $\delta$.  In particular,
the probability that any historical assertion in
Proposition~\ref{p:mat-resourcehistory} fails is at most $\delta$;
the probabilities $\Pr(\mathsf N)$, $\Pr(\mathsf N,X+)$ and the captured joint outcomes
in \eqref{p:mat-capturednumbers} have that same error bound.  The
conditional value $\Pr(X+\mid\mathsf N)$ requires the branch-probability
correction stated below.
\end{corollary}
\begin{proof}
All the displayed records, copy crossings, receiver transfers and
later controlled outputs are measurable functions of the same full
material path.  Total variation contracts under their common
pushforward.  The Bell failure event has probability zero, so the
pilot failure event has probability at most $\delta$.
\end{proof}
This is a joint bound. A conditional comparison requires positive
probability under both compared laws; $\delta<p$ suffices when the
ideal branch probability is $p>0$. Lemma~\ref{found:errors} then gives
error at most $\min\{1,2\delta/p\}$. No uniform precision is asserted
on arbitrarily rare branches.  Nor does a classical-path TV estimate alone compare an
unobserved quantum density matrix in trace norm.  Reference-sensitive
operational tests are included by adjoining their admissible gates
to the same complete programme and applying the same argument.

All clocks, copies, fuel supplies, loss products and reset receivers
are part of the autonomous inventory.  Their field evolution is
exact; the kinetic, reaction and gas parameters enter only through
the separately proved material-path error.  The wave and currents
are analytic on the fixed finite graph, and
Proposition~\ref{p:mat-levels} provides the indicated input-uniform
level-crossing count where the kinetic bound needs it.

\paragraph{Exact time scope.}
Historical protection here is a first-pass statement on
$[0,\pi/(2\Omega)]$.  A finite reversible clock is not an absorbing
final state.  On its second pass the fine currents reverse, the
clock has nodes at the turnaround, and at $2T$ the circuit has
coherently undone itself.  The full pilot mechanism can be tested
against that extended Bell path, but the archive theorem does not
claim that a deliberately undone memory remains a record.  Every
required later feedback or hold operation must be included in the
declared first-pass programme.



\section{Uniformity and an explicit fixed-circuit resource rate}
\label{p:uniformity}
The qualitative limit in Theorem~\ref{p:main} uses the established
kinetic proof.  For the fixed clock circuit just constructed, its constants
can also be controlled uniformly over unknown inputs.  The following
quantitative refinement records the cutoff dependence explicitly; it is
not a new stochastic law or a rate for growing circuits.  The explicit
exponents follow bookkeeping isolated in the supplied working
audit~\cite{pIndependentAudit}; that citation records provenance, and the
cutoff proof is supplied here.

\begin{proposition}[Uniform fixed-circuit kinetic rate]\label{p:kinetic-rate}
Fix the complete finite graph, the clock Hamiltonian $H_F$, its first-pass
horizon $T$, and $0<\kappa_-\le\kappa_e\le\kappa_+$.  Suppose the
initial census has expected $\ell^1$ error $O(N^{-1/2})$ and the tag starts
from $\nu\le Cw(0)$ with fixed $C$.  The iid equilibrium and the known
basis-ready preparations satisfy this condition.  With
\[
 \mu_N=N^{1/2},\qquad \delta=N^{-1/7},\qquad
 \varepsilon=N^{-1/35},
\]
the signed-queue path error obeys
\begin{equation}\label{p:explicit-rate}
 \epsilon_{\rm kin}(N,T)=O(N^{-1/70}).
\end{equation}
For equilibrium use, the constants are uniform over all normalized inputs
on this fixed material space.  A reference included in the configuration
basis is part of the fixed graph; a declared spectator fibre does not
enlarge the operator-norm constants.  With the gas and recombination scales
of Theorem~\ref{p:main}, the full pilot path error is also $O(N^{-1/70})$.
\end{proposition}
\begin{proof}
Write $D=|V|$, $m=|E|$, $H_* =\|H_F\|$, and let $\pVar$ denote total
variation in physical time.  Each edge current is a bounded quadratic form
in the normalized input.  Safe input-independent bounds are
\[
 J_*\le 2H_*/\hbar,\qquad
 L\le 2mH_*T/\hbar,\qquad
 \sum_e\pVar(J_e)\le 4mH_*^2T/\hbar^2.
\]
The last inequality follows by differentiating each current expectation
and using $\|\dot\Psi\|\le H_*/\hbar$.  All source variation, birth and
census-jump budgets in Theorem~\ref{kin:tracking} are therefore uniform.

For $0<\delta\le1$ put $\alpha=\mu_N/N\le1$ and
$a_0=\kappa_-\delta$, $a_1=\kappa_+$.  In the proof of
Theorem~\ref{kin:tracking}, the instantaneous companion root has
$\pVar(f)=O(\delta^{-2})$ and $|f(0)|=O(\delta^{-1})$.
Its deterministic integrated tracking cost is consequently
$O((\mu_N\delta^3)^{-1})$.  The square-jump estimate there gives
\[
 V'\le-2\mu_Na_0V+
 \mu_N\alpha a_1(J_*/a_0+2c_0\alpha+\sqrt V).
\]
Using
$\alpha a_1\sqrt V\le a_0V+\alpha^2a_1^2/(4a_0)$ and $V(0)=0$
yields
\[
 \sup_{t\le T}V(t)\le
 \frac{\alpha a_1J_*}{a_0^2}
 +\frac{2c_0a_1\alpha^2}{a_0}
 +\frac{\alpha^2a_1^2}{4a_0^2}
 \le C\frac{\alpha+\alpha^2}{\delta^2}.
\]
The deterministic export discrepancy is $O(\alpha)$, so
\eqref{kin:R} and \eqref{kin:smallflux} have the explicit forms
\begin{align}\label{p:cutoff-rates}
 R_{\delta,N}&\le C\left[
 \frac{1}{\mu_N\delta^3}
 +\frac{\sqrt{\mu_N/N}+\mu_N/N}{\delta}\right],\\
 D_{\delta,N}&\le C\sqrt{\delta+(\mu_N\delta)^{-1}
                    +\E\eta_N}.\nonumber
\end{align}
Here and below constants depend on the fixed graph, programme and response
bounds, not on the input or the two cutoffs.

For iid equilibrium, direct multinomial variance gives
$\E\|x^N(0)-w(0)\|_1\le\sqrt{D/N}$.  For a separately initialized
dominated tag and $N-1$ independent equilibrium carriers, add at most
$2/N$.  The deterministic basis-ready census has zero initial error.
The tracking proof applies conditionally on the initial census and then
averages; Jensen's inequality handles the square root in the low-mass
estimate.  Its population martingale obeys
\[
 \epsilon_{x,N}\le\E\|x^N(0)-w(0)\|_1
   +C_B\epsilon_{F,N}+C\sqrt{(L+1)/N}.
\]
With $\mu_N=N^{1/2}$ and $\delta=N^{-1/7}$, the leading deterministic
term of \eqref{p:cutoff-rates} is $N^{-1/14}$, its square-root noise
term is $N^{-3/28}$, and $D_{\delta,N}=O(N^{-1/14})$.
Equation~\eqref{kin:fluxbridge} therefore gives
\[
 \epsilon_{F,N}+\epsilon_{x,N}=O(N^{-1/14}).
\]

For the clock family \eqref{p:mat-weights}, each fine weight is a
nonnegative input-dependent coefficient times a unimodal binomial
envelope.  Count entries into a sublevel set directly: each weight has
at most two boundary entries on $[0,T]$, including tangencies.  A
coefficient may put its maximum exactly at $\varepsilon$, so a common
regular value is not assumed.  Initial small-weight mass, jump influx
and these entries give the uniform node budget
\[
 b_C(\varepsilon)\le C\bigl(3D\varepsilon+C_0T\sqrt\varepsilon\bigr).
\]
Use this budget in \eqref{kin:pathbound}.  At
$\varepsilon=N^{-1/35}$ the node term and the
$\epsilon_{x,N}/\varepsilon^2$ term are $O(N^{-1/70})$; the
$ (\epsilon_{x,N}+\epsilon_{F,N})/\varepsilon$ term is
$O(N^{-3/70})$.  This proves \eqref{p:explicit-rate}.
The gas and recombination contributions are respectively $O(N^{-2})$
and $O(N^{-1/2})$, so they do not worsen this conservative rate.
\end{proof}

On $[0,2T]$ at most four level entries per weight replace the node bound
by $C(5D\varepsilon+2C_0T\sqrt\varepsilon)$; the same rate applies to
the reversed full path.  That second pass unwrites the records and is
not an extension of the first-pass archive theorem.  Merely assuming
$x^N(0)\to w(0)$ without a rate does not imply
\eqref{p:explicit-rate}.  Nor does this estimate apply uniformly to a
growing graph, increasing fine reference, changing Hamiltonian or unbounded
storage time.

\chapter{Discrimination, retained information and constitutive scope}\label{p:chapter-discrimination}\label{p:sec-rivals}
\section{A stationary coherent cycle tests individual edge ownership}
Take three vertices with $w_r=1/3$, $\Psi=(1,1,1)^T/\sqrt3$, and
\[
 H=\frac{3\hbar j}{2}
 \begin{pmatrix}0&-i&i\\i&0&-i\\-i&i&0\end{pmatrix},\qquad j>0.
\]
Then $H\Psi=0$ but $J_{21}=J_{32}=J_{13}=j$. The canonical exporters therefore produce clockwise packets even though every vertex population is stationary. Theorem~\ref{p:main} gives a clockwise Bell cycle at rate $3j$, with Poisson$(3jT)$ jump count and its full continuous event times. A generator which simply keeps the configuration fixed has the same one-time equilibrium and zero divergence, yet differs in path law by $1-e^{-3jT}$. It is excluded by individual primitive bond ownership, not by the continuity equation alone.

The nearest Markov surplus rival adds a constant $K>0$ to both directional equilibrium fluxes on each cycle bond. It also preserves $w$. Its clockwise rate is $3(j+K)$ and its counterclockwise rate is $3K$, independently of the occupied vertex. The probability of at least one backward jump is $1-e^{-3KT}$, whereas Bell's probability is zero. Thus full-path distance is at least that number. The microscopic theory suppresses this rival by the proved recombination hierarchy, rather than declaring two-way reactions impossible.

\section{The finite-recombination rival really has dark traffic}
On a two-vertex edge in a zero-current interval, put one packet of each species, no new births, and all $N$ carriers at the two endpoints. Total service rate is $\kappa\mu_N$ and recombination rate is $a$. The probability of at least one actual carrier event before $T$ is exactly
\begin{equation}
 \frac{\kappa\mu_N}{a+\kappa\mu_N}
       \left(1-e^{-(a+\kappa\mu_N)T}\right).
 \label{p:eq-darkrival}
\end{equation}
The first event is a race of the two derived contact mechanisms. Recombination first removes both packets; service first already makes the path nonconstant. The instantaneous net queue has $Z=0$ and never moves, so \eqref{p:eq-darkrival} is its exact carrier-path distance from this finite reaction experiment. A mixed stock can occur after opposite exports before recombination completes. This local test does not replace the empty-queue initialization of Theorem~\ref{p:main}. It shows that the new mechanism has finite-speed predictions and is not a renamed Jordan decomposition. If recombination is absent or too slow, surplus survives.

\section{The spatial ensemble selects timing beyond mean flux}
Let $\Delta=1/R$, $M\ge3$, and choose a single uniform $U\in[0,\Delta)$. Put the arrival positions at times $U+k\Delta$, $k=0,\ldots,M-1$, and randomly permute particle labels. Every individual arrival time is uniform on $[0,M/R)$, just as for the independent gas, and mean flux is $R$. Nevertheless on $[0,2/R]$ there are exactly two contacts with separation exactly $1/R$. The Poisson comparison assigns zero probability to that exact spacing. The complete contact-history distance is therefore one. Independence of initial positions, not their one-body density or mean pressure, excludes this rival. P4 makes that additional statistical content explicit.

\section{Recombination does not conceal a readable pilot ledger}
Let $E_\pm$ be signed export counts, $D_\pm$ signed consumption counts and $A$ the recombination count on one bond. For empty initial stock,
\[
 E_+=P^++D_++A,\qquad E_-=P^-+D_-+A.
\]
Consequently
\begin{equation}
 N(\Pi(t)-\Pi(0))=P^+(t)-P^-(t)+D_+(t)-D_-(t)+u(t).
 \label{p:eq-ledger}
\end{equation}
All retained recombination products cancel from the signed identity. Fast recombination removes surplus service but does not erase the source action from a joint ledger. A hypothetical ordinary classical reader of this ledger remains a countermodel to the unrestricted material source; the following physical-access part retains the corresponding explicit constructions. P3 replaces that unrestricted coupling constitution. Gauge invariance and finite work do not imply P3: $\Pi$ is gauge invariant and can be put in an additional invariant mixed energy if that extra force is permitted.

Within the new theory a reader is an ordinary device in the common $H$. Its probabilities in the Bell limit have the positive-effect form
\begin{equation}
 \Prb(M=m)=\|P_mU(\psi\otimes\alpha)\|^2
       =\langle\psi,E_m\psi\rangle,\qquad 0\le E_m\le I.
 \label{p:eq-effect}
\end{equation}
This follows from full-configuration equivariance and the actual joint unitary; it is not a separate measurement postulate. For a fixed compiled clock implementation, Proposition~\ref{p:kinetic-rate} supplies an input-uniform finite-resource error $\Delta_N$. An inaccessible reference remains in $U$ and is acted on by the identity. A direct rewrite of an actual ordinary source bit into a memory while leaving the joint field at $\psi\otimes|0_M\rangle$ would instead give positive actual probability to $M=1$ at zero coherent weight. It is not an allowed material interaction.

\begin{proposition}[A finite native-counter obstruction survives the new theory]
For $H=\hbar g\sigma_x$, $\tau=gT\in(0,\pi/4)$, no fixed positive record effect can reproduce the probability of an unchanged native $0\to1$ Bell event for all three inputs $|0\rangle$, $|1\rangle$ and $(|0\rangle-i|1\rangle)/\sqrt2$ with error smaller than
\begin{equation}
 \frac{\sin\tau(\cos\tau-\sin\tau)}{3}.
 \label{p:eq-nativegap}
\end{equation}
\end{proposition}
\begin{proof}
Direct two-state wave evolution gives event probabilities $\sin^2\tau,0,\sin\tau\cos\tau$. For the first and third preparations the current is forward throughout the interval, so there is at most one forward event and its probability is the loss of origin weight. For the second it is backward. If a positive effect has error at most $\varepsilon$ on all three, its trace is at most $\sin^2\tau+2\varepsilon$, whereas its expectation in the third vector is at least $\sin\tau\cos\tau-\varepsilon$. A positive expectation cannot exceed the trace. Rearrangement gives \eqref{p:eq-nativegap}.
\end{proof}
At $gT=\pi/8$ the gap is $(\sqrt2-1)/6$. A material apparatus with path error $\Delta_N$ cannot be a universal passive native counter with record error below this gap minus $\Delta_N$. The permitted coherent copy changes the complete source--memory field, so it does not claim to evade this obstruction. Its archive attests its own actual copy crossing. This explicitly distinguishes successful measurement from a fictitious passive record of every earlier native excursion.

\section{Fine and coarse paths remain distinct}
For any coarse boundary with microscopic currents $j_1,\ldots,j_k$, forward mean incidence is $\sum_i[j_i]_+$, not generally $[\sum_i j_i]_+$. Their discrepancy is
\[
 \frac12\left(\sum_i|j_i|-\left|\sum_i j_i\right|\right).
\]
The monomial-copy theorem computes every fine current and makes this discrepancy zero at that specific boundary. It does not infer the result from a coarse clock population. Arbitrary circuit gates can have negative fine forward factors. Reduced null maps are used only to compute probabilities after retaining every receiver in the joint model; they do not supply actual trajectory rates. Expected flux integrals likewise remain expectations of counts, not sample counting measures.

\section{Dependencies and the strength of the conclusion}
\subsection{Dependency map}
\begin{longtable}{>{\raggedright\arraybackslash}>{\raggedright\arraybackslash}p{0.25\textwidth}>{\raggedright\arraybackslash}>{\raggedright\arraybackslash}p{0.33\textwidth}>{\raggedright\arraybackslash}>{\raggedright\arraybackslash}p{0.34\textwidth}}
\toprule
Component & Prior status or premise & Present consequence\\
\midrule
\endhead
Canonical edge source & Primitive binary quadratic action; common $\hbar$; passive ownership & Retained. It fixes every individual Hamiltonian edge current, including cycles.\\
Packet production & Bounded conservative action exporter with finite elementary charge & Retained as explicit deterministic hybrid physics. No target waiting-time sampler.\\
Signed queue & Instant opposite cancellation was supplied & Replaced by actual two-species coexistence and finite recombination; new full-state error \eqref{p:eq-annerror}.\\
Reaction clocks & Additive Markov pair generator was supplied & Replaced by deterministic finite flight and independent spatial preparation; new complete-history error \eqref{p:eq-gaserror}.\\
Scalar response & Equal pair response and complete reaction list & Retained microscopic premises, with explicit charge routing and channel geometry.\\
Initial equilibrium & Calibrated population and tag law & General input premise retained. The explicit basis-ready preparation derives the example's distribution with randomness only in the pilot gas.\\
Ordinary admission & Unrestricted pilot-action/ledger readers obstruct common measurement & Replaced by P3: one coherent ordinary force algebra and an explicit absent mixed pilot force. Not derived from older mechanics.\\
Actuator and archive & Conditional physical record results & Explicit autonomous common $H_F$ with actual monomial crossing proofs, loss, pending, fuel and retained reset recipients.\\
Continuation & Complete-field conditional composition & Demonstrated through archive-controlled noncommuting source gates with an inaccessible reference and joint probabilities.\\
Entropy, chamber and MPBT routes & Their stated statistical, boundary and measure premises & Preserved as separate conditional results; not invoked to select this theory's event law.\\
\bottomrule
\end{longtable}

\subsection{What has been established}
With P1--P4 as its physical constitution, the finite theory has a fully specified state, deterministic event mechanism and finite resource inventory on every promised horizon. The controlled limit derives the \emph{original minimal Bell path}, not merely an operationally equivalent alternative, for the complete declared ordinary configuration. No positive part of $J$ or division by $w$ occurs in the microscopic service or contact law. Positive and negative packet species coexist at finite resources; the proved fast-recombination estimate explains their limiting directional selection. The gas preparation explains complete conditional timing rather than just its mean.

The autonomous ordinary experiment establishes more than Born endpoints. Fine-current signs prove that copies are actual historical records at their unique crossing boundaries and remain protected through the included reset and noncommuting continuation. The full material-path bound transfers this history statement to finite pilot resources with error at most $\Delta_N$. The example retains a nontrivial reference, null branch, pending and loss weight, fuel changes and both reset receivers. Its extended preparation starts all ordinary configurations and pilot carriers at one definite ready state, with no random carrier sampler; the same dynamics produces the entangled preparation before a one-way boundary isolates its reference.

This is verdict \textbf{(b)}, an explicit new constitutive internal completion, together with \textbf{(c)}, a controlled effective Bell closure. It is not verdict \textbf{(a)} relative to the old unrestricted source constitution. In particular neither the absence of a mixed pilot read force, the scalar reaction spectrum, nor the independent equilibrium/spatial ensemble is claimed to follow from gauge symmetry or generic mechanics. If any of those premises is refused, the corresponding countermodels above survive. That refusal distinguishes a different physical theory; it does not undo the conditional mathematics of this one.

\subsection{Limits of the result}
The theorem is for fixed finite ordinary programmes and growing finite pilot resources. It does not give an economical material substrate, a unique empirically selected new theory, a universal continuum-field limit, Bell rates conditional on every hidden pilot coordinate, or exact finite-resource equivariance. Its exact copy theorem protects records on the first pass of the retained clock. An unmodified later clock echo unwrites them, and a longer desired retention experiment must be included in a larger declared programme. The auxiliary smooth contact module does not prove a smooth common Hamiltonian for canonical export interleaved with all reactions; such an embedding is a stronger open mechanics problem, not a premise concealed in the present hybrid claim.

Thus the original unrestricted programme is not proved inevitable. A precisely stated replacement constitution now supplies a complete event-selection and measurement-chain realization with quantitative path errors. The massive continuous-configuration constitution developed separately in this monograph remains an alternative; its operational agreement must not be mislabeled as this Bell-path derivation.


\section{Connection to the remaining constitutions}
The physical-access results that follow allow couplings which P3 excludes.
They therefore remain precise tests of what the new law changes.  The
relative-entropy, MPBT and chamber arguments retain their own probability,
boundary and incidence premises.  They are neither needed in the
gas--recombination--kinetic proof nor invalidated by it.  A complete
experiment must choose one compatible source and material constitution,
as required in Chapter~\ref{found:experiments}.

The finite clock identity uses the engineered spin-chain matrix of
Christandl et al.\ \cite{pChristandl2004}; its gate-carrying conjugacy is
the Feynman construction \cite{pFeynman1986}.  The required matrices,
currents and histories were proved above rather than imported as an
unspecified computation theorem.  Bell-process existence retains the
provenance and proof in Lemma~\ref{kin:existence} and \cite{BellQFT}.
No historical-priority or empirical-confirmation claim follows from
these mathematical constructions.


\part{Physical Access, Material Contacts, and Compatibility Obstructions}\label{part:access}

% Begin consolidated source: parts/02_access.tex
\providecommand{\accTV}{d_{\mathrm{TV}}}
\providecommand{\accLaw}{\operatorname{Law}}
\providecommand{\accVar}{\operatorname{Var}}

\chapter{Readable histories and reciprocal source writers}
\label{acc:histories}

This part consolidates the finite compatibility witness of \cite{M11},
the reciprocal stored-action writer of \cite{M21}, the destructive and
native-product contacts of \cite{M22,M23}, and the accounted-work
counterexample of \cite{M25}. These are different levels of physical
restriction on access, tested on explicit complete experiments. They are
not a universal prohibition of Bell trajectories. Their common point is
that preserving a native path or accounting for a local disturbance does
not establish a common affine law for every additional record.

\section{One reusable detuned experiment}
Let physical sectors be the position factor of
$\mathbb C^2_{\rm pos}\otimes\mathbb C^2_{\rm int}\otimes\mathcal H_R$,
with $R$ inaccessible. Put $\hbar=1$ in this chapter and take
\begin{equation}
 H_D=g(\sqrt\eta\,\sigma_x+\sqrt{1-\eta}\,\sigma_z)_{\rm pos}
 \otimes\operatorname{diag}(1,2)_{\rm int}\otimes I_R,
 \quad 0<\eta<1,\quad T=\frac\pi{2g}.
 \label{acc:detunedH}
\end{equation}
Initially position and every native carrier are at $0$, the source action
and queues are empty, and the unknown normalized internal/reference vector
is $\Xi$. Set
$p=\|(\lvert0\rangle\langle0\rvert_{\rm int}\otimes I_R)\Xi\|^2$.
The two frequencies give
\begin{align}
 w_1(t)&=\eta\{p\sin^2(gt)+(1-p)\sin^2(2gt)\},\label{acc:weight}\\
 J_{10}(t)&=\eta g\sin(2gt)\{p+4(1-p)\cos(2gt)\}.
 \label{acc:current}
\end{align}
The bracket is decreasing on $[0,T]$, so the current reverses at most
once, from positive to negative. Write
$f_p(t)=p\cos^2(gt)+(1-p)\cos^2(2gt)$, so
$w_0=1-\eta+\eta f_p\ge1-\eta$.

\begin{proposition}[First-exit probability]\label{acc:firstexit}
The original Bell process initialized at the unique ready sector has
first-exit probability by $T$
\begin{equation}
 F(p)=\eta\left[1-\min_{0\le t\le T}f_p(t)\right].
 \label{acc:firstexitformula}
\end{equation}
It has $F(0)=F(1)=\eta$ and $F(2/3)=3\eta/4$.
\end{proposition}
\begin{proof}
Before first exit the integrated hazard is
$\int[J_{10}]_+/w_0$. On the sole increasing interval of $w_1$,
$J_{10}=-\dot w_0$, so survival equals the minimum attained $w_0$.
Subsequent negative current adds no outward first-exit hazard. For
$p=2/3$, set $x=\cos^2(gt)$; then
\[
 f_{2/3}=\tfrac23x+\tfrac13(2x-1)^2
 =\tfrac14+\tfrac43(x-\tfrac14)^2.
\]
Its minimum is $1/4$, at $gt=\pi/3$. The basis-frequency minima are zero.
\end{proof}

\begin{theorem}[Common complete-input affine floor]\label{acc:affine}
Suppose one complete represented ready preparation admits both the basis
ensemble with probabilities $(2/3,1/3)$ and the equal ensemble
\[
 \psi_\pm=\sqrt{2/3}\lvert0\rangle
 \pm\sqrt{1/3}\lvert1\rangle,
\]
without an active preparation tag distinguishing their methods. If their
complete represented matrix input is the same, no common affine record
law reproduces both original Bell first-exit histories. A common record
approximation with TV error at most $\epsilon$ on each preparation obeys
\begin{equation}
 \epsilon\ge\eta/8.
 \label{acc:affinefloor}
\end{equation}
This conclusion places no restriction on the possible event daughters.
\end{theorem}
\begin{proof}
Both ensembles have matrix $\operatorname{diag}(2/3,1/3)$ and the same
ready position. Their Bell first-exit probabilities are respectively
$\eta$ and $3\eta/4$. An affine map assigns their common complete input a
common event probability. The triangle inequality yields
$\eta/4\le2\epsilon$. The proof concerns this one record marginal, so
arbitrary subsequent daughter kernels cannot repair it.
\end{proof}

The theorem's complete-preparation qualification is essential. A canonical
source can regard the actual ray decomposition as physically different
data and predict different records consistently with its own ontology.
The theorem then shows that the proposed matrix quotient is insufficient;
it does not license discarding legitimate active provenance. Conversely,
if a common affine complete-input law is retained, the unmodified Bell
first-exit history cannot be made an arbitrarily accurate public record in
this comparison class. Terminal position weights do not show the conflict:
both ensembles end with position-one probability $2\eta/3$.

\section{A structural incompatibility of classical and quantum writers}
The difference between a classical canonical field and an affine
classical--quantum source is algebraic. Let
$\mathcal A=C^\infty(C,M_d)$ be the observable algebra of a classical
coordinate and finite quantum bank. A reversible infinitesimal
star-derivation $D$ satisfies the Leibniz rule.

\begin{proposition}[Center preservation]\label{acc:center}
Every derivation of $\mathcal A$ maps its center to its center. In
particular the rule $D(fI)=f'A$ with nonscalar Hermitian $A$ cannot be such
a reversible classical--quantum transport law.
\end{proposition}
\begin{proof}
For central $fI$ and arbitrary $B$, differentiate $[fI,B]=0$:
$[D(fI),B]+[fI,D(B)]=[D(fI),B]=0$.
Thus $D(fI)$ is central. Choosing $f'$ nonzero proves the final statement.
\end{proof}
This result does not invalidate the classical canonical action. It prevents
identifying that action's exact expectation writer with a reversible
transport of this different algebra without changing physical premises.
Dissipative instruments, their noise and their disturbance must be derived
or separately specified in the latter constitution.

\section{Reciprocity does not remove a delayed classical read}
Attach a canonical pointer $(Y,P)$ to the current action $\Pi_e$ of
Chapter~\ref{kin:production} through
\begin{equation}
 H_{\rm wire}(t)=k(t)\Pi_eP+\frac{P^2}{2M},\qquad M>0.
 \label{acc:wire}
\end{equation}
Its equations are
$\dot Y=k\Pi_e+P/M$, $\dot P=0$, $\dot\chi_e=kP$,
$\dot\Pi_e=-\partial_{\chi_e}H$.
The pointer can kick the connection; it is not a one-way expectation
assignment.

\begin{theorem}[Finite-width delayed action read]\label{acc:delayed}
Run the canonical source on $[0,T]$ with $k=0$, $\chi(0)=0$, and
$|Y_0|\le\epsilon_Y$, $|P_0|\le\delta$. During $[T,T+\tau]$ set the
coherent source Hamiltonian to zero and apply a pulse of area
$\beta=\int_T^{T+\tau}k(t)dt$. Then
\begin{equation}
 Y_{
\rm out}=Y_0+\frac{T+\tau}{M}P_0+\beta\Pi_e(T),\qquad
 \chi_{e,\rm out}=\beta P_0,\qquad \Pi_{e,\rm out}=\Pi_e(T).
 \label{acc:wireout}
\end{equation}
If the reaction generator has the scalar form \eqref{kin:generator}, or
the unchanged material response of Chapter~\ref{kin:neutrality}, the
native coherent field and all native reaction paths through read completion
are exactly unchanged. This includes unfinished-queue reactions in the
zero-current hold.
\end{theorem}
\begin{proof}
During production $k=0$. During reading the source Hamiltonian vanishes,
so $\Psi$ and $\Pi_e$ are constant; $P$ is constant throughout.
Integration gives \eqref{acc:wireout}. Queues, carriers and material retain
their original equations because their declared reaction coefficients
have no added $\chi$ dependence during the hold. The same primitive
reaction clocks therefore couple all native paths exactly. A later
connection kick cannot alter a record already secured in $Y$.
\end{proof}

In the detuned experiment, $\Pi(0)=0$ gives $\Pi(T)=w_1(T)=\eta p$.
If $\beta>0$ and a finite comparator has bounded error $e_{
\rm read}$,
\begin{equation}
 \epsilon_Y+(T+\tau)\delta/M+e_{
\rm read}<\beta\eta/4
 \label{acc:precision}
\end{equation}
guarantees that the bit $Y_{
\rm out}>3\beta\eta/4$ is one for $p=1$ and zero for $p=0,1/2$.
Equal $Z$ and equal $X$ ensembles then give bit probabilities $1/2$ and
$0$. No particular distribution inside the nonzero-width ready rectangle
is needed. For example, $\eta=1/2$, $\beta=M=1$, $g=1$, $\tau=1$,
$\epsilon_Y=e_{
\rm read}=1/64$, $\delta=10^{-3}$ satisfy the strict inequality.
A smooth comparator with constant response plateaux on the separated
bands can amplify the bit without a derivative force on those bands.

All pointer and connection variables remain represented. A later prescribed
Hamiltonian with
$\|H(\chi,t)-H(0,t)\|\le L_H(t)|\chi|$ obeys the wave bound
\[
 \|(U_\chi-U_0)\Psi_T\|\le|\beta|\delta\int L_H(t)dt.
\]
This controls the full vector including the reference. It is not a Bell
path bound near nodes and does not cover feedback through the new pointer.
An internal rotation noncommuting with $\operatorname{diag}(1,2)$ but
independent of $\chi$ is unaffected at the time it is applied.

\section{Equivalent encodings and noisy returns}
Deleting access to the symbol $\Pi$ is insufficient if its exact accounting
remains readable. For the monitored edge, let $R_e$ be net native reaction
count, $Z_e$ residual signed queue and $u_e$ exporter residue. Then
\begin{equation}
 N(\Pi_e-\Pi_e(0))=R_e+Z_e+u_e
 \label{acc:ledger}
\end{equation}
with pending and captured charges added when those resources are present.
The identity follows by adding every export and reaction rewrite.
Separate registers may reveal no useful scalar individually while their
joint return reconstructs it. A gauge-invariant difference of two such
actions is equally available to a canonical coupling $kA P$ if that
relational action is a physical observable.

\begin{proposition}[Finite scalar noise and retained keys]\label{acc:noise}
Let $N_0$ be any proper additive-noise variable independent of the source
value. The laws
$P_Z=\tfrac12\accLaw(N_0)+\tfrac12\accLaw(N_0+a)$ and
$P_X=\accLaw(N_0+a/2)$ cannot agree for $a\ne0$.
If $N_0$ has variance $\sigma^2<\infty$, then
\begin{equation}
 \accTV(P_Z,P_X)\ge\frac{a^2}{48(\sigma^2+a^2)}.
 \label{acc:noisebound}
\end{equation}
If the noise key itself is retained and returned jointly with the noisy
readout, the two joint laws have TV one.
\end{proposition}
\begin{proof}
Equality of characteristic functions would require
$\varphi(t)e^{iat/2}[\cos(at/2)-1]=0$ for every $t$.
Continuity and $\varphi(0)=1$ contradict this on a small punctured
neighborhood of zero. For the bound, center $N_0$ and test the bounded
cosine centered at $a/2$. Its expectation gap is
$(1-\cos(at/2))\operatorname{Re}\varphi(t)$ in absolute value.
Take $t=(\sigma^2+a^2)^{-1/2}$. Then
$\operatorname{Re}\varphi(t)\ge1-\sigma^2t^2/2\ge1/2$ and
$1-\cos(at/2)\ge a^2t^2/12$. A bounded test of absolute value at most
one has expectation gap at most $2\accTV$, giving \eqref{acc:noisebound}.
Finally $(Y,N_0)$ determines $Y-N_0$, whose values are in $\{0,a\}$ in
one preparation and $\{a/2\}$ in the other. Those supports are disjoint.
\end{proof}

\chapter{Destructive contacts, native products and reset attacks}
\label{acc:contacts}

The next constructions use the canonical action, packets and carriers of
Chapters~\ref{kin:production}--\ref{kin:limit}. They demonstrate why
packet consumption and ownership of the native move are different
restrictions, and why each still permits an informative record in an
explicit resource domain. They preserve the original Hamiltonian current.

\section{A single absorbing site with finite processing}
Choose one positive export direction and supply $B_N$ empty processing
pockets, one site $A_0+A_1=1$ initially empty, a clock and retained
processing records. Set $N>b>0$, $\Gamma_N>0$. Each positive packet has
\begin{align}
 P_i&\longrightarrow Q_e^+ &&\Gamma_N(1-b/N),\nonumber\\
 P_i+A_0&\longrightarrow A_1+\text{record}
 &&\Gamma_Nb/N,\label{acc:pocket}\\
 P_i+A_1&\longrightarrow Q_e^++A_1 &&\Gamma_Nb/N.\nonumber
\end{align}
Negative packets return to their native queue at rate $\Gamma_N$.
The total processing rate is $\Gamma_N$ regardless of site occupancy.
At most one packet of charge $1/N$ is retained in the site. All remaining
packets eventually return, including opposite-sign cancellations.

The ready pockets and site are independent of the unknown input and
primitive native reaction randomness. The new physical admission is that
an exported packet can contact this auxiliary site. No real threshold or
freely readable random seed is posited; exponential residences and
Bernoulli marks below are mathematical representations of the displayed
Markov chemistry. Degenerate labels conserve the stated energy, with
supplied finite storage and growing processing rates counted as resources.

\begin{proposition}[Exact finite-time null law]\label{acc:pocketnull}
For positive export times $t_i$,
\begin{equation}
 S_N(s)=\mathbb P(A_0(s)=1)=
 \prod_{t_i\le s}\left[1-\frac bN
 (1-e^{-\Gamma_N(s-t_i)})\right].
 \label{acc:nullproduct}
\end{equation}
After all $D$ positive packets complete, this is $(1-b/N)^D$.
If exports finish by $T$, reading at $T+\tau$ differs from the completed
binary law by at most $B_Ne^{-\Gamma_N\tau}$.
\end{proposition}
\begin{proof}
Every pocket has an independent rate-$\Gamma_N$ residence and a candidate
capture mark of probability $b/N$. The first completed candidate capture
occupies the only site. No capture by $s$ means every offered candidate is
unmarked or unfinished, giving the product. A union bound for the at most
$B_N$ uncompleted residences gives the finite-read estimate. It keeps,
rather than rejects, all unfinished branches.
\end{proof}

\begin{theorem}[Path stability under a one-packet diversion]\label{acc:tapstability}
Assume a label-symmetric native preparation and fixed programme with
$\sum_e|Z_e|\le B_N$. On $[0,T_h]$, the monitored and native tagged paths
obey
\begin{equation}
 \accTV\le\delta_{\rm tag,N}:=
 \min\left\{1,\frac{\kappa_+\mu_N}{N}
 e^{\kappa_+\mu_N(1+2B_N/N)T_h}
 \left(T_h+\frac{B_N}{\Gamma_N}\right)\right\}.
 \label{acc:tapbound}
\end{equation}
The same bound holds jointly with the same autonomous router history
adjoined to the unmodified comparator. It includes nodes and reversals.
\end{theorem}
\begin{proof}
Couple common edge/carrier reactions at minimum rates. Let
$Q=\sum_e|Z_e-Z'_e|$, let $D_c$ count disagreeing carrier locations and
let $M$ count unmatched native reactions. For agreeing carrier labels use
the queue difference, and for disagreeing labels the sum of rates. The
identity $|u_+-v_+|+|u_--v_-|=|u-v|$ yields
\[
 R_{
\rm mismatch}\le\kappa_+\mu_N\{Q+2(B_N/N)D_c\}.
\]
If $P(t)$ packets remain pending, signed conservation gives
$Q\le1+P+M$ and $D_c\le M$. This counts the one possible diversion.
Moreover $\int_0^{T_h}\mathbb EP(t)dt\le B_N/\Gamma_N$.
The unmatched-count compensator and Gronwall give
\[
 \mathbb EM(T_h)\le\kappa_+\mu_N
 e^{\kappa_+\mu_N(1+2B_N/N)T_h}(T_h+B_N/\Gamma_N).
\]
Permutation symmetry assigns a preselected label expected unmatched
count $\mathbb EM/N$. Its path can disagree only at such an event,
which proves \eqref{acc:tapbound}. For fixed coherent input, router
history depends only on deterministic exports and its own clocks.
The native comparison can be generated independently of these clocks;
the same coupling then preserves the decorated marginal. No division by
$w_r$ occurred in this estimate.
\end{proof}

\section{Actual capture times and simultaneous scales}
Let $A_N^+(t)$ count positive exports, $f=[J_e]_+$ on $[0,T]$, extended
by zero, and suppose its variation $V_f$ is finite. Define
\[
 d_N=\sup_{t\ge0}\left|\frac{A_N^+(t)}N-\int_0^tf(s)ds\right|.
\]
This positive-count discrepancy is stronger than signed residual control
when the number of reversals grows without bound.

\begin{theorem}[Latch-time TV estimate]\label{acc:latchtimes}
Let $\mathsf L$ be the first-arrival process of deterministic intensity
$bf(t)$ while ready and zero after its first event, including the null atom.
If $A_N^+(T)\le N\ell$, the actual latch path satisfies
\begin{equation}
 \accTV(\accLaw(L_N),\mathsf L)\le
 \delta_{\rm latch,N}:=\frac{b^2\ell}{N}
 +b\left[2(\Gamma_NT+1)d_N+\frac{V_f}{\Gamma_N}\right].
 \label{acc:latchbound}
\end{equation}
For a fixed input $\Xi$, the monitored tag and latch converge jointly to
the independent pair of the Bell path and $\mathsf L$ whenever
$\delta_{\rm tag,N}+\epsilon_{B,N}+\delta_{\rm latch,N}\to0$.
Here independence is conditional on $\Xi$, not on a density-matrix
preparation quotient or on all microscopic histories.
\end{theorem}
\begin{proof}
Each positive export at $t_i$ yields a candidate point with probability
$p_N=b/N$ and location density
$k_\Gamma(t-t_i)=\Gamma e^{-\Gamma(t-t_i)}1_{t\ge t_i}$.
A Bernoulli($p_N$) count and Poisson($p_N$) count have TV
$p_N(1-e^{-p_N})\le p_N^2$. Coupling all counts and locations bounds
candidate-point-process error by $b^2\ell/N$. The Poissonized intensity is
$a_N=(b/N)\sum_i k_\Gamma(\cdot-t_i)$.

Put $G=A_N^+/N-\int f$. Integration by parts gives
\[
 (k_\Gamma*dG)(t)=\Gamma G(t)
 -\Gamma^2\int_0^te^{-\Gamma(t-s)}G(s)ds.
\]
Its absolute value is at most $2\Gamma d_N$ before $T$ and decays
exponentially afterwards, so its $L^1$ norm is at most
$2(\Gamma T+1)d_N$. The BV translation bound
$\|f(\cdot-s)-f\|_1\le sV_f$ gives
$\|k_\Gamma*f-f\|_1\le V_f/\Gamma$.
Minimum-intensity Poisson coupling then costs at most the integral of
$|a_N-bf|$. Taking the first-point map, with its null value, contracts TV
and proves \eqref{acc:latchbound}. For the joint claim use the decorated
native comparison in Theorem~\ref{acc:tapstability}, the canonical Bell
limit, and independence of that comparator's tag from router clocks.
\end{proof}

With at most $K_+$ positive-current intervals, residual balance on each
interval gives $d_N\le2K_+/N$. For \eqref{acc:current} there is one initial
positive interval, the exporter starts at zero, and $d_N\le1/N$,
$V_f\le8\eta g$. Choose
\begin{equation}
 \Gamma_N=\Gamma_*\sqrt N,\qquad
 \mu_N=\sqrt{\log(N+1)},\qquad B_N/N\le\ell.
 \label{acc:simultaneous}
\end{equation}
Then \eqref{acc:tapbound} tends to zero, \eqref{acc:latchbound} is
$O(N^{-1/2})$, and the canonical hierarchy holds. All times are physical
and the duration remains fixed. The growing apparatus supplies $O(N)$
pockets, $O(\sqrt N)$ processing rate and unbounded total throughput.
Making $\Gamma_N$ arbitrarily faster is not an improvement: a narrow comb
of completion times near deterministic exports need not converge in TV
as a time density. The convolution estimate displays that restriction.

\begin{theorem}[Finite destructive-access gap]\label{acc:destructivegap}
For the two preparations of Theorem~\ref{acc:affine}, the completed
one-site contact has probability gap
\begin{equation}
 G_N=(1-b/N)^{\lfloor3N\eta/4\rfloor}
 -(1-b/N)^{\lfloor N\eta\rfloor}
 \longrightarrow G=e^{-3b\eta/4}-e^{-b\eta}>0.
 \label{acc:gap}
\end{equation}
It consumes at most one packet, and its tagged-path disturbance vanishes
under \eqref{acc:simultaneous}. The maximum limiting gap is $27/256$.
A finite error certificate is
\begin{equation}
 |G_N-G|\le\frac{2b}{N}
 +\frac{7\eta b^2}{8N(1-b/N)}.
 \label{acc:finitegap}
\end{equation}
At a finite read with binary-reader error at most $\epsilon_{
\rm bit}$
per preparation, any common affine approximation has error at least
\[
 \frac12[G_N-2B_Ne^{-\Gamma_N\tau}-2\epsilon_{
\rm bit}]_+.
\]
\end{theorem}
\begin{proof}
The maximum $w_1$ is $\eta$ for either basis input and $3\eta/4$ for
$p=2/3$. There is one positive rise, so the positive packet counts are
exactly the floors in \eqref{acc:gap}. Apply
Proposition~\ref{acc:pocketnull}. For $0\le v\le\eta$,
$0\le-\log(1-x)-x\le x^2/[2(1-x)]$ gives
\[
 |(1-b/N)^{\lfloor Nv\rfloor}-e^{-bv}|
 \le b/N+vb^2/[2N(1-b/N)].
\]
Sum this at $v=\eta,3\eta/4$. Clearing and reader error are probability
coupling bounds; an affine approximation must use a common probability,
so its two errors sum to at least the actual gap. Finally put $y=b\eta$.
The derivative of $e^{-3y/4}-e^{-y}$ vanishes at
$y=4\log(4/3)$, where the value is $27/256$.
\end{proof}

For a concrete finite binary calculation, $\eta=1/2$, $N=32$ and
$b=8\log(4/3)$ give $G_N=0.1053958545$. That number certifies the completed
latch statistic; it does not certify a small native path error at $N=32$.
The full programme runs the detuned pulse, retains pocket and queue
processing during a zero-current hold, reads the latch, and then applies
a fixed local internal noncommuting rotation. The old record, captured
packet, null branches and pending resources remain in the complete state.

\section{A reporter produced by the native reaction itself}
One may strengthen access by insisting that a record be made only when
its own packet performs its original native move. Prepare one cofactor
$A_0$, with bound states $A_{1,a}$ retaining the contacting carrier label.
At a positive native contact replace rate $r_{ea}$ by
\begin{align}
 P_e+C_{a,r}+A_0&\longrightarrow C_{a,q}+A_0
 &&(1-b/N)r_{ea},\nonumber\\
 P_e+C_{a,r}+A_0&\longrightarrow C_{a,q}+A_{1,a}
 &&(b/N)r_{ea}.
 \label{acc:nativemark}
\end{align}
After binding use the original native rate. Both branches consume the
same packet and obey $\Delta(n+BZ)=0$; neither diverts source charge.
The physical addition is a neutral retained reaction product.
Section~\ref{bench:history} gives a separate two-stage finite-delay
benchmark in which a first acquired marker is actually generated and
then retained through a noncommuting control. It distinguishes conditioning
the actual configuration from replacing the continuing wave.

\begin{theorem}[Exact native projection and reporter timing]\label{acc:reporter}
The projection onto the entire old native state and every native event time
is exactly unchanged by \eqref{acc:nativemark}. If $K_N(t)$ counts
eligible native reactions, then conditional on the complete native history,
\begin{equation}
 \mathbb P(A(t)=0\mid\mathcal H_{
\rm native})=(1-b/N)^{K_N(t)}.
 \label{acc:nativenull}
\end{equation}
The reporter path satisfies
\begin{equation}
 \accTV(\accLaw(L_N),\mathsf L)\le b\epsilon_{F,N}.
 \label{acc:nativelatch}
\end{equation}
For a predesignated carrier in the symmetric ready preparation, its joint
path and reporter converge to $\mathsf B_\Xi\otimes\mathsf L_
\Xi$ with bound
\begin{equation}
 \Xi_N(a)+b\epsilon_{F,N}+3b\ell/N,
 \label{acc:nativejoint}
\end{equation}
where $\Xi_N(a)$ is the right side of \eqref{kin:pathbound} at cutoff $a$.
\end{theorem}
\begin{proof}
For every function of the native state, the two generator terms sum to
$r_{ea}[F(R_{ea}s)-F(s)]$, independently of cofactor state. This exact
lumpability proves equality of native path laws, including any old
controller depending only on those native histories. Equivalently mark
independent eligible native events with probability $b/N$ until the first
mark. This representation has precisely the displayed chemistry and
proves \eqref{acc:nativenull}.

While ready, the reporter intensity is
$(b/N)\sum_ar_{ea}=b\Phi_e^N$. Minimum-rate coupling with intensity
$b[J_e]_+$ gives \eqref{acc:nativelatch}. For the joint comparison let $s$
be the tag's eligible intensity and $R=\sum_ar_{ea}$. Before binding,
the projected channels have rates
\[
 \text{tag only}:(1-b/N)s,\quad
 \text{latch only}:(b/N)(R-s),\quad
 \text{both}:(b/N)s.
\]
Against the independent target, their excess is at most
$|s-\lambda_B|+b|\Phi_e^N-[J_e]_+|+3bs/N$.
Use the localization in Theorem~\ref{kin:path} for the first term.
Exchangeability and the packet budget give $\mathbb E\int sdt\le\ell$,
proving \eqref{acc:nativejoint}. After binding only the native comparison
remains. The coupling preserves both the full microscopic marginal and
the independent target generator.
\end{proof}

This is stronger than the destructive-contact result: even the complete
native projection is undisturbed, without a limiting statement. The same
limiting gap \eqref{acc:gap} follows from the reporter law and kinetic flux
convergence. But the full native archive is not independent of the new
latch. The event that a non-null latch time equals one of its finitely
many native reaction times has actual probability $\mathbb P(A=1)$ and
probability zero for an independent continuously distributed latch time.
Thus
\begin{equation}
 \accTV(\accLaw(\mathcal H_{
\rm native},L_N),
 \accLaw(\mathcal H_{
\rm native})\otimes\accLaw(\widetilde L))
 \ge\mathbb P(A=1).
 \label{acc:fullhistoryfloor}
\end{equation}
The tagged limit does not survive every enlargement of the history
filtration as an independent-product claim.

\section{Permanent local damage and common-catalyst repair}
Suppose a binding event permanently changes the contacted carrier's pair
response by $\zeta\ne1$, leaving the others unchanged. Before binding the
models coincide, so the first reporter time and label are exactly those
of Theorem~\ref{acc:reporter}. After binding, a minimum-rate coupling has
mismatch intensity at most
\[
 \kappa_+\mu_N(1+2\ell)M+
 \kappa_+\mu_N|\zeta-1|\ell,
\]
where $M$ counts unmatched native reactions. Signed queue discrepancy and
carrier disagreement are at most $M$. Gronwall and symmetry give
\begin{equation}
 \accTV(\accLaw(X_1^{\rm damaged}),\accLaw(X_1^{\rm native}))
 \le\min\left\{1,
 \frac{|\zeta-1|\ell}{N(1+2\ell)}
 [e^{\kappa_+\mu_N(1+2\ell)T_h}-1]\right\}.
 \label{acc:damage}
\end{equation}
This tends to zero under \eqref{acc:simultaneous}. The bound also retains
the common reporter time and label. Selecting the damaged carrier later
from its retained label is a different experiment: for $\zeta=0$ it
cannot return under reversed current. No vanishing bound holds for that
actively selected carrier or the entire damaged microscopic state.

A genuinely shared interlock instead removes a common catalyst on capture.
Let $C=1$ mean active and let the programme age satisfy
\begin{equation}
 \dot a=C,\qquad H_C(t)=C H(a(t)),\qquad r_{ea}^C=C r_{ea}.
 \label{acc:gate}
\end{equation}
The controller, production and reactions pause together when $C=0$.
On a monotone one-jump interval with $J\ge0$, $w_1(0)=0$ and
$W=\int_0^TJdt<1$, the ungated Bell jump density is $J(t)$.
The rare-contact limit already proved gives an independent first-capture
intensity $bJ$ conditional on the fixed input.

\begin{proposition}[Interlock cost and retained-copy restoration]\label{acc:reset}
If the shared gate never resets, its native path distance is exactly
\begin{equation}
 \Delta_{
\rm gate}=W-\frac{1-e^{-bW}}{b}
 \ge\frac W2(1-e^{-bW}).
 \label{acc:gatecost}
\end{equation}
If a fuelled rate-$\rho$ reset copies the first capture into retained
memory and restores the catalyst, then the same first-capture law remains
and, for zero-extended BV $J$,
\begin{equation}
 \accTV(\text{native path with reset},\text{original native path})
 \le\frac{\accVar(J)+2\|J\|_\infty}{2\rho}.
 \label{acc:resetbound}
\end{equation}
The bound retains the same capture and reset variables in both compared
outputs and includes finite-window nulls.
\end{proposition}
\begin{proof}
Put $z(t)=\int_0^tJ$. With no reset the native density is
$J(t)e^{-bz(t)}$, of mass $(1-e^{-bW})/b$. The missing mass is moved to
the null atom, giving the equality. For $x=bW$, the inequality is
$x(1+e^{-x})-2(1-e^{-x})\ge0$, whose derivative is
$1-(1+x)e^{-x}\ge0$ and whose value at zero is zero.

For reset, condition on first capture $\sigma$ and pause
$D\sim\operatorname{Exp}(\rho)$, with $D=0$ if no capture. Jumps before
$\sigma$ remain fixed, and later jumps are delayed by $D$.
The affected subdensity is $f_\sigma(t)=J(t)1_{t>\sigma}$, of variation
at most $\accVar(J)+2\|J\|_\infty$. Translation gives
$\|f_\sigma(\cdot-D)-f_\sigma\|_1\le D\accVar(f_\sigma)$.
On the whole time line its total mass is unchanged, so TV is half this
$L^1$ distance. Average $\mathbb ED\le1/\rho$ and map late jumps to the
finite-horizon null. Keeping $\sigma,D$ throughout the conditioning proves
the decorated comparison.
\end{proof}

At finite $N$, gate every old drift and hazard and switch only after the
marked native move. For fixed $\sigma,D$, the resulting process is exactly
the old process under its measurable paused programme-age map. Adjoining
the independent reset-clock representation and applying that map to
\eqref{acc:nativejoint} preserves its error uniformly in $\rho$.
In the monotone domain, the full comparison therefore adds only
\eqref{acc:resetbound}. Letting $\rho\to\infty$ simultaneously with the
kinetic resources leaves the acquired record while suppressing the native
physical-time disturbance. A uniform bound on reset power or rate would
change the admitted resource domain; energy accounting alone supplies no
such bound.

\chapter{Accounted work, copied records and a complete return cycle}
\label{acc:work}

Canonical reciprocity excludes a frozen source-dependent recoil on an open
product phase space, but permits reciprocal replacements. This chapter
retains the strongest complete work-accounting counterexample from
\cite{M25}. Its actual memories use the common chamber medium of
\cite{M24}, while its work degree is classical canonical material. This
is a different hybrid constitution from the preceding canonical packet
source. The calculation is an architecture boundary, not a transfer of
its Bell theorem without proof.

\section{Canonical recoil is constrained but not neutral}
On a finite amplitude cone and a classical work pair $(Y,P)$ use
\[
\vartheta=\frac{i\hbar}{2}(z^\dagger dz-dz^\dagger z)+P\,dY,
 \qquad E_i(z,Y)=z^\dagger H_i z+Y.
\]
Here $Y$ has energy units and $P$ has time units. A frozen-source map
$(z,Y,P)\mapsto(z,Y+f(z),P)$ changes $d\vartheta$ by $dP\wedge df$.
It is canonical on an open product domain only if $df=0$. This proves a
real reciprocity constraint, but does not require equal source energies
on different contacts.

\begin{theorem}[Energy-matching reciprocal recoil]\label{acc:recoil}
For finite Hermitian $H_i,H_j$, define
\begin{align}
 U_{ji}(P)&=e^{iPH_j/\hbar}e^{-iPH_i/\hbar},\nonumber\\
 K_{ji}(P)&=H_i-U_{ji}(P)^\dagger H_jU_{ji}(P),\nonumber\\
 S_{ji}(z,Y,P)&=(U_{ji}(P)z,Y+z^\dagger K_{ji}(P)z,P).
 \label{acc:recoilmap}
\end{align}
This map is canonical, $E_j\circ S_{ji}=E_i$, and
$S_{kj}S_{ji}=S_{ki}$. Within maps of this form with $U(0)=I$ and the
stated phase/translation normalization, the two requirements determine
$U$ uniquely. Moreover
\begin{align}
 \|U_{ji}(P)-I\|&\le |P|\|H_j-H_i\|/\hbar,\nonumber\\
 \|K_{ji}(P)-(H_i-H_j)\|
 &\le |P|\|[H_i,H_j]\|/\hbar.
 \label{acc:recoilbound}
\end{align}
\end{theorem}
\begin{proof}
For any differentiable unitary $U(P)$ put $K=i\hbar U^\dagger U'$ and
$a=z^\dagger Kz$. Direct substitution gives
$S_U^*\vartheta-\vartheta=a,dP+P,da=d(Pa)$, proving canonicity.
Conversely its mixed source/work terms give the same $K$, up to the
fixed scalar translation. Energy matching requires
$K=H_i-U^\dagger H_jU$, hence
$i\hbar U'=UH_i-H_jU$. The displayed product is its unique solution.
The charts
$C_i(z,Y,P)=(e^{-iPH_i/\hbar}z,Y+z^\dagger H_i z,P)$ satisfy
$S_{ji}=C_j^{-1}C_i$, proving composition.
The first bound follows by Duhamel comparison of the two unitary groups;
the second follows by differentiating
$e^{iPH_i/\hbar}H_je^{-iPH_i/\hbar}$ and integrating its commutator norm.
Tensoring an inaccessible reference leaves the bounds unchanged.
\end{proof}

A pulse of $c_1(t)Pz^\dagger H_i z$ followed by a pulse of
$-c_2(t)Pz^\dagger H_jz$, each with coefficient integral one, implements
this map. This factorization alone does not account for switching work
at a native collision. Also an actual-event-dependent recoil changes
the source wave and invalidates direct use of the fixed-wave kinetic
proof. We next use an autonomous work cycle whose entire read, copy and
return can be calculated.

\section{The work rotor and physical memory medium}
Let the carried space include position, internal input, two memory bits
$B,C$, and an inaccessible reference $R$. Supply a canonical cylinder
$(\phi,Y)$, with time-valued angle of circumference $L$ and
$\{\phi,Y\}=1$. Its free energy is $Y$ on an active bounded energy band.
A smooth extension can be bounded below outside that band; the stated
experiment stays strictly inside it. This linear-dispersion work rotor
is a disclosed resource, not an ordinary quadratic kinetic pointer or
a finite-dimensional consequence of finite energy.

Put $A=I_{\rm pos}\otimes\lvert0\rangle\langle0\rvert_{\rm int}$ and
\[
 G_1=\sigma_y^B\otimes I_C,\qquad
 G_2=\lvert1\rangle\langle1\rvert_B\otimes\sigma_y^C,\qquad
 G_3=-\sigma_y^B\otimes\lvert1\rangle\langle1\rvert_C.
\]
Choose real fixed profiles $f,\ell_1,\ell_2,\ell_3$ and a fixed selector
$0\le\chi(Y)\le1$. The autonomous energy is
\begin{equation}
 \mathcal E=Y+f(\phi)\langle A\rangle+
 \hbar\ell_1(\phi)\chi(Y)\langle G_1\rangle+
 \hbar\ell_2(\phi)\langle G_2\rangle+
 \hbar\ell_3(\phi)\langle G_3\rangle.
 \label{acc:rotorenergy}
\end{equation}
The action has kinetic one-form
$i\hbar(\Psi^\dagger d\Psi-d\Psi^\dagger\Psi)/2+Y,d\phi$.
Its equations are
\begin{align}
 i\hbar\dot\Psi&=[fA+\hbar\ell_1\chi(Y)G_1+
 \hbar\ell_2G_2+\hbar\ell_3G_3]\Psi,\label{acc:rotorwave}\\
 \dot\phi&=1+\hbar\ell_1\chi'(Y)\langle G_1\rangle,\nonumber\\
 \dot Y&=-f'\langle A\rangle-\hbar\ell_1'\chi(Y)\langle G_1\rangle
 -\hbar\ell_2'\langle G_2\rangle-\hbar\ell_3'\langle G_3\rangle.
 \label{acc:rotorforce}
\end{align}
Thus the source-dependent work displacement is a reciprocal Hamiltonian
force, rather than a prescribed numerical impulse. Actual memory outcomes
are supplied by the following separately specified physical medium.

For every finite joint sector $q=(i,b,c)$, let
$\mathcal C_q=[0,w_q]\times[0,1]$, $w_q=\|P_q\Psi\|^2$.
An actual tracer has sector $q$, depth $d$ and transverse coordinate
$\zeta$. Set $f_{rq}=[J_{rq}]_+$,
$F_q=\sum_rf_{rq}$, $I_q=\sum_rf_{qr}$.
Between stirrings $\dot d=-F_q$ and $\zeta$ is fixed. The outgoing face
is divided into destination strips of relative width $f_{rq}/F_q$;
the incoming face of $r$ into strips of relative width $f_{rq}/I_r$.
At $d=0$ the tracer enters $r$ at $d'=w_r$ and its transverse coordinate
is rescaled affinely so that
$F_q,d\zeta=I_r,d\zeta'$. Zero-width strips are absent.
A fundamental rate-$\kappa$ stirring resets $(d,\zeta)$ in the current
sector to $(w_qU,V)$ for independent fresh uniforms $U,V$.

This geometry specifies squared-norm chamber volumes, rectified normalized
faces, perfect transmission and a Poisson uniform-refresh law. They are
not derived here. The initial unique ready chamber has unit density in
$(d,\zeta)$, with source-independent readiness. There is one unknown
carried input, not additional input specimens. Work readiness and future
stirring are independent of it.

\begin{lemma}[Density preservation in the finite memory medium]\label{acc:chamberdensity}
For a deterministic bounded piecewise-smooth coherent programme, the
chamber density identically one is preserved by the drift, portals and
stirring. Consequently actual sector probabilities are $w_q(t)$.
The process loses no probability at a finite-time explosion.
\end{lemma}
\begin{proof}
Interior drift has zero divergence in the chamber coordinates. The top
boundary moves at $\dot w_q=I_q-F_q$ while depth drift is $-F_q$,
so its relative incoming flux is $I_q$. The bottom outgoing flux is $F_q$.
The prescribed face rescaling equates each edge's boundary flux, and the
incoming and outgoing strip partitions are exhaustive. Hence the constant
density solves the transport equation with its moving-boundary conditions.
Stirring replaces a sector's density by its uniform conditional density;
it therefore also preserves density one. A killed construction has mass
bounded by this transport solution. Its expected portal count is bounded
by $\int\sum_qF_qdt<\infty$, and its expected stirring count by $\kappa T$.
Thus no finite-time explosion loses mass, and equality with the normalized
transport solution follows. Vanishing-volume chambers have zero occupied
probability. This is a preservation theorem under the supplied ready law,
not its preparation or statistical selection.
\end{proof}

\section{An exact read, genuine copy and complete work return}
Prepare $B,C$ in $\lvert00\rangle$ and a source vector $\Psi_S$ independent
of these memories. Choose $f=0$ at readiness, $f=-K$ during three disjoint
ordered pulses, and
\begin{equation}
 \int_0^Lf(\phi)d\phi=0,\qquad
 \int\ell_j(\phi(t))dt=\pi/2\quad(j=1,2,3).
 \label{acc:rotorprofiles}
\end{equation}

\begin{theorem}[Accounted-work acquisition with a retained copy]\label{acc:rotor}
For every fixed carried input and work-ready point, the solution is
\begin{equation}
 \phi_t=\phi_0+t\pmod L,\qquad
 Y_t=Y_0-f(\phi_t)p,\qquad
 \Psi_t=e^{-(i/\hbar)A\int_0^tf(\phi_u)du}\Psi_S\otimes m_t,
 \label{acc:rotorsolution}
\end{equation}
where $p=\langle A\rangle_{\Psi_S}$ is constant.
Its total energy is exactly $Y_0$. After a full rotor cycle,
\begin{equation}
(\Psi_S,\phi_0,Y_0,\lvert00\rangle)\longmapsto
 (\Psi_S,\phi_0,Y_0,
 \lvert0\rangle_B[\cos\beta\lvert0\rangle_C+
 \sin\beta\lvert1\rangle_C]),
 \quad \beta=\frac\pi2\chi(Y_0+Kp).
 \label{acc:rotorreturn}
\end{equation}
The source including inaccessible reference correlations, both work
coordinates and the display return exactly. The actual retained archive
has $\mathbb P(C=1)=\sin^2\beta$ at every finite $\kappa$.
\end{theorem}
\begin{proof}
For fixed complete initial data the Hamiltonian separates into source and
memory operators, so the wave remains a product across that partition.
Since $-iG_j$ are real and the initial memory amplitudes are real, $m_t$
stays real. Therefore every expectation $\langle G_j\rangle$ vanishes
at every time, not only at endpoints. The work equations reduce to
$\dot\phi=1$, $\dot Y=-f'p$. Also $A$ commutes with the entire contact,
so $p$ is constant. This proves \eqref{acc:rotorsolution} and
$\mathcal E=Y+fp=Y_0$.

On the plateau, $Y=Y_0+Kp$. The three exact memory gates give
\[
 |00\rangle\longmapsto
 \cos\beta|00\rangle+\sin\beta|10\rangle
 \longmapsto\cos\beta|00\rangle+\sin\beta|11\rangle
 \longmapsto|0\rangle_B(\cos\beta|0\rangle_C+\sin\beta|1\rangle_C).
\]
Thus $C$ is a physical copy of the display and is used to reset it.
The zero integral of $f$ restores the source unitary, periodicity restores
$\phi$, and $f=0$ restores $Y$. Lemma~\ref{acc:chamberdensity} makes the
final volume of $C=1$ its actual probability. No auxiliary Born draw or
ideal collapse was appended.
\end{proof}

The initial memory occupies a proper ready submanifold. On its reached
real-amplitude curve the Berry one-form vanishes while the metric line
element need not. There is no contradiction with a theorem about
information export by a source-fixed map on an open full product domain.
Hamiltonian reciprocity of the complete hybrid dynamics also does not give
one preparation-independent linear propagator after eliminating its
classical work trajectory.

\section{Finite positive-area resources and the separating experiment}
A concrete resource programme sets $K=\hbar g$, $L=16/g$ and
$u=g\phi\pmod{16}$. Let $s(x)=3x^2-2x^3$ on $[0,1]$. Define
\[
 \frac{f(u/g)}K=
 \begin{cases}
 -s(u-1),&1\le u\le2,\\
 -1,&2\le u\le7,\\
 -[1-s(u-7)],&7\le u\le8,\\
 0,&0\le u\le1\text{ or }15\le u\le16,
 \end{cases}
\]
and on $[8,15]$ use the negative of the preceding profile translated by
seven. Its two areas cancel. For $a_1=2.5,a_2=4,a_3=5.5$, set
\[
 \ell_j(\phi)=\pi g\sin^2(\pi(u-a_j))1_{[a_j,a_j+1]}(u).
\]
These profiles are $C^1$ and piecewise smooth, have the required pulse
areas, and use a finite four-dimensional memory.
Prepare independently of the source
\begin{equation}
 u_0\in[0,0.1],\qquad
 Y_0\in[10K-K/24,10K+K/24]
 \label{acc:workready}
\end{equation}
with the uniform rectangle law, of nonzero canonical area $\hbar/120$.
Its particular uniform density is not needed for the separated-band bit
below; support on this rectangle suffices.
Choose a fixed $C^1$ selector equal to one for
$|Y-(10+2/3)K|\le K/12$ and zero for
$|Y-(10+2/3)K|\ge K/6$. Then throughout the allowed ready bands,
\[
 \chi(Y_0+Kp)=1\quad(p=2/3),\qquad
 \chi(Y_0+Kp)=0\quad(p=0,1).
\]
All unknown inputs $p\in[0,1]$ obey $|Y_t-Y_0|\le K$,
$Y_t\in[8.95833K,11.04167K]$ and
$\|H_{
\rm field}\|\le(1+\pi)K$. The active band $[8K,12K]$ therefore has
strict headroom; no promised branch exhausts the work resource.
All couplings have finite duration and bounded strength on this domain.
Contacts engage and disengage where their profiles vanish, with zero
contact-energy switching cost. The actual read/copy/return is autonomous.

For a finite timing calculation let
$\Theta(v)=\pi v/2-\sin(2\pi v)/4$ on $[0,1]$, and put
$v_j=[gt+u_0-a_j]_0^1$, $x=\chi(Y_0+Kp)$.
The display and copied-archive capture CDFs during their respective pulses
are
\begin{equation}
 \mathbb P(L_B\le t)=\sin^2(x\Theta(v_1)),\qquad
 \mathbb P(L_C\le t)=\sin^2\beta\sin^2\Theta(v_2).
 \label{acc:rotortimes}
\end{equation}
The terminal null masses are $1-\sin^2\beta$.
Indeed the relevant currents are unidirectional, so crossing into the
daughter chambers is equivalent to having captured by that time;
Lemma~\ref{acc:chamberdensity} gives the CDFs. Differentiation gives the
ordinary time densities; for example the first is
$x\pi g\sin^2(\pi v_1)\sin(2x\Theta(v_1))$ on the pulse.
The reset has CDF $\sin^2\beta\sin^2\Theta(v_3)$ and leaves $C$
unchanged. Integrating these formulas over \eqref{acc:workready} gives
the full finite-read law. No conditioning on a successful branch is needed.

Run the detuned production \eqref{acc:detunedH} with $\hbar$ restored,
then a zero-current hold of duration $1/(2g)$, then the rotor programme.
Compare, using fixed orthogonal inaccessible reference vectors $r_0,r_1$,
\begin{align*}
 \mathcal E_{
\rm basis}:&\quad |0r_0\rangle\text{ with probability }2/3,
 \quad |1r_1\rangle\text{ with probability }1/3,\\
 \mathcal E_{
\rm sup}:&\quad
 \sqrt{2/3}|0r_0\rangle\pm\sqrt{1/3}|1r_1\rangle,
 \quad\text{each with probability }1/2.
\end{align*}
All declared apparatus readiness and active provenance agree; no ensemble
label is supplied to a controller. The averaged source/reference matrices
coincide. The actual raywise $p$ laws differ, and the exact returned
archive satisfies
\begin{equation}
 \mathbb P_{
\rm sup}(C=1)=1,\qquad
 \mathbb P_{
\rm basis}(C=1)=0.
 \label{acc:rotorgap}
\end{equation}
These are different complete source laws on rays; the result proves that
their common matrix quotient is insufficient under the hybrid coupling.
An actual retained preparation label would change the comparison and must
not be erased in applying this conclusion.

\section{A noncommuting continuation and native-path return}
After the rotor retain $C$ and apply
$H_P=\hbar h\sigma_y^{\rm pos}$ for $T_P=\pi/(4h)$.
This does not commute with the detuned position Hamiltonian.
At production end the source fibres are
\[
 \psi_0=-i\sqrt{1-\eta},P_0\Xi-P_1\Xi,\qquad
 \psi_1=-i\sqrt\eta,P_0\Xi.
\]
Their inner product is $p\sqrt{\eta(1-\eta)}$, so the final
position-one probability, in both read and clean experiments, is
$1/2+p\sqrt{\eta(1-\eta)}$. This equality alone is only a marginal
statement. The following stronger comparison controls the complete native
position path in the declared finite-stirring experiment.

\begin{theorem}[A retained record after arbitrarily small native-path change]
\label{acc:quietreturn}
For the three tested input classes $p=0,1,2/3$ with the preparation above,
compare the read experiment to the same production, hold, total port
duration and later probe, but with source/work and memory contacts off.
Both keep the same native stirring. If their final common quiet interval
has length $s$, then
\begin{equation}
 \accTV(\text{read native position path},
 \text{clean native position path})\le e^{-\kappa s},
 \qquad \text{retained archive gap}=1.
 \label{acc:quietbound}
\end{equation}
The explicit profiles provide $s=1/g$ on the total physical horizon
$\pi/(2g)+1/(2g)+16/g+\pi/(4h)$.
\end{theorem}
\begin{proof}
Couple initial tracer data and all stirring clocks. Production is identical.
There is no native position current during the hold or rotor, so native
position paths agree throughout that period. Memory transfers may change
the depth and transverse coordinates, which are not asserted returned.
By time $15/g$ the field, rotor and display have returned, all contacts
are off, and $C$ is certain for each tested class. If a stirring occurs
in the common final quiet interval, use the same uniform pair in the
occupied read and clean chambers, identified by their fixed memory label.
Their native weights agree, so depth and transverse coordinates then
coincide. They remain identical through the subsequent probe and all its
native crossings under the common coupling. Failure of this coupling is
confined to no stirring in an interval of length $s$, of probability
$e^{-\kappa s}$. The archive value is already certain and unchanged by
that stirring. This proves both statements, retaining exact native event
times and null intervals.
\end{proof}

The special deterministic-memory input classes are sufficient for the
separation. No uniform finite-$\kappa$ path estimate for every intermediate
$p$, arbitrary memory feedback or perturbation of this programme is
inferred from it. Increasing $\kappa$ is a supplied stirring-rate resource;
its Poisson count has no fixed finite upper bound. The ready uniform
coordinate and repeated uniform refresh remain statistical postulates.
This construction supplies a complete work and memory account, but no
new universal admission principle.

\section{Scope of the access conclusions}
The original packet source permits delayed reads of a classical action.
Destructive packet capture preserves a predesignated native path in a
controlled resource limit. A marked native product preserves the entire
native projection exactly. Local damage does not prevent that fixed-tag
limit, and a common gate can be reset while retaining its copied record.
The accounted-work model adds the stronger fact that reciprocal source
forces, autonomous bounded work excursions, genuine copying and exact
source/work/display return can coexist with a unit record gap in a
specific hybrid constitution.

Each result identifies an admitted coupling rather than asserting that
all source coordinates must be measurable. The constitutive class changes
where stated: the rotor is not silently the old packet source, and a
Bell law postulated or derived for one graph does not make all classical
resource extensions into a common affine instrument. These countermodels
show that conservation, reciprocity, consumption and finite local work
accounting alone do not select a compatible complete reaction and record
interface. The pilot theory supplies a different, explicit ordinary
material coupling inventory; its exclusion of the displayed classical
readers is new constitutive content. None of the countermodels proves that
every possible source completion fails, and the new inventory does not
invalidate a countermodel on its original permissive access domain.

% End consolidated source: parts/02_access.tex


% Begin consolidated source: parts/02_material.tex
\chapter[Material contact geometry and response]{Material contact geometry and the limits of common response}
\label{mat:geometry}

The constructions in this chapter replace or constrain the physical contact that couples a source to a work store. They do not identify ordinary classical work points with coherent material coordinates. That difference changes the state space, the allowed acquisition operations, and the predictions of a copied-return experiment. We first prove a positive construction in a complete material field, then retain a stronger counterexample in a distinct hybrid constitution. The latter shows why a shared origin of force and conversion is not sufficient for complete record compatibility \cite{M26,M28}.

Throughout, an unknown carried vector may be entangled with an inaccessible reference. Every source-dependent controller, old memory, and future-returning resource is part of the input. New ready packets and memory blanks are independent resources when independence is stated. An ensemble average does not erase active provenance.

\section{A complete material field}

Let $\mathcal K$ contain the finite source, reference, and internal apparatus indices. A material contact coordinate $y$ carries a section $\psi\in L^2(\mathbb R;\mathcal K)$. The section, including variations of its local orientation, is the primary source variable. A product $\psi(y)=\varphi(y)\Xi$ is an allowed preparation, not an invariant constraint during contact. At a finite regulator the real metric and symplectic form are
\begin{equation}
 g(u,v)=2\hbar\operatorname{Re}\langle u,v\rangle,
 \qquad \omega(u,v)=2\hbar\operatorname{Im}\langle u,v\rangle.
 \label{mat:metric}
\end{equation}
The material law requires smooth contact flows to preserve both forms, fix the zero field, and respect common phase. These are new constitutive assumptions. Positivity of energy, canonical reciprocity, and source/readout incompleteness do not by themselves imply them.

A nonempty massive example is the action
\begin{align}
 \mathcal I[\psi]&=\int dt\left\{\frac{i\hbar}{2}
 \int(\psi^\dagger\dot\psi-\dot\psi^\dagger\psi)dy-\mathcal E[\psi]\right\},\\
 \mathcal E[\psi]&=\int\left\{\frac{\hbar^2}{2M}\|\partial_y\psi\|^2
       +\psi^\dagger[H_S+V(y)A]\psi\right\}dy,
 \label{mat:action}
\end{align}
where $M>0$, $H_S,A$ are bounded Hermitian operators and $V$ is a bounded smooth real function. Variation gives
\[
 i\hbar\dot\psi=[-\hbar^2\partial_y^2/(2M)+H_S+V(y)A]\psi.
\]
The standard self-adjoint realization on $H^2$ supplies unitary evolution. This is a specified field dynamics; it is not a proof that all work substances have this form.

Where $\rho=\|\psi\|^2>0$, write $\psi=\sqrt\rho e^{iS/\hbar}z$, $z^\dagger z=1$, and $a=-i\hbar z^\dagger\partial_yz$. Direct differentiation gives
\begin{align}
 \frac{\hbar^2}{2M}\|\partial_y\psi\|^2
 &=\frac\rho{2M}(S'+a)^2+\frac{\hbar^2}{8M}\frac{(\rho')^2}\rho\\
 &\quad+\frac{\hbar^2\rho}{2M}
   (\|z'\|^2-|z^\dagger z'|^2).
 \label{mat:strain}
\end{align}
The last term is the squared norm of the orientation derivative normal to $z$ and is nonnegative. The changes $z\mapsto e^{i\theta(y)}z$, $S\mapsto S-\hbar\theta$ leave the section and every displayed energy term unchanged. Thus local orientation modes cannot generally be deleted without changing the physical variational problem.

\begin{theorem}[Linear transport from the stipulated field geometry]
\label{mat:linear}
On a connected finite-dimensional amplitude regulator, a twice differentiable contact vector field preserving \eqref{mat:metric}, fixing zero and respecting phase is $X\psi=-iH\psi/\hbar$ for a field-independent Hermitian $H$. On the test-function core $C_c^\infty(\mathbb R;\mathbb C^d)$, if the corresponding complex-linear local differential expression has order at most one, it is
\begin{equation}
 X\psi=-V(y,t)\partial_y\psi-\tfrac12\partial_yV(y,t)\psi
               -iB(y,t)\psi/\hbar,
 \quad V=V^\dagger,\quad B=B^\dagger.
 \label{mat:local}
\end{equation}
Conversely \eqref{mat:local} preserves the real metric on compactly supported test fields. A global unitary evolution additionally requires a self-adjoint realization of its generator.
\end{theorem}
\begin{proof}
In real coordinates metric preservation is the Killing equation
$\partial_bX_c+\partial_cX_b=0$. Differentiate in a third coordinate, add the first two cyclic identities and subtract the third. Equality of mixed derivatives gives $\partial_a\partial_bX_c=0$. Hence $X=B_0\psi+b_0$, with $B_0$ real antisymmetric. Zero fixation removes $b_0$. Symplectic preservation implies that $B_0$ commutes with the complex structure, so it is complex linear and anti-Hermitian. Therefore $B_0=-iH/\hbar$ with $H=H^\dagger$.

For the local statement write the field-independent expression as $L=-V\partial_y+C$. Integration by parts on the common core gives
\[
 L^\dagger=V^\dagger\partial_y+\partial_yV^\dagger+C^\dagger.
\]
Thus $L^\dagger=-L$ is equivalent to $V=V^\dagger$ and $C+C^\dagger=-V'$. The remaining anti-Hermitian part is $-iB/\hbar$. Reversing the integration proves the converse. Constant bounded $V$ and bounded smooth Hermitian $B$, with their declared domains, provide nonempty continuum realizations. No differentiability of an unbounded translation field on all of $L^2$ is asserted.
\end{proof}

The local continuity equation is
\begin{equation}
 \partial_t(\psi^\dagger\psi)+\partial_y(\psi^\dagger V\psi)=0.
 \label{mat:continuity}
\end{equation}
For $V=g(t)A$ and $B=0$, the spectral component with $A=a$ translates by $aK(t)$, where $K(t)=\int_0^tg(s)ds$. Translating the entire field at the expectation-dependent displacement $K\langle A\rangle$ violates the field-independent linearity in Theorem~\ref{mat:linear}. This is the specific mathematical restriction that changes the earlier classical writer. The theorem does not select the physical matrix $V$, prohibit a contact outside this material class, or assign actual coordinates a probability law.

\section{What an expectation writer leaves out}

Let $z,\varphi$ be normalized source and work vectors and let $A,B$ be self-adjoint, with $\varphi$ in the domain of $B$. Set $a=\langle z,Az\rangle$ and $b=\langle\varphi,B\varphi\rangle$. For $H_{\rm int}=gA\otimes B$,
\begin{align}
 Az\otimes B\varphi
 &=abz\otimes\varphi+b(A-a)z\otimes\varphi
   +az\otimes(B-b)\varphi\\
 &\quad +(A-a)z\otimes(B-b)\varphi.
 \label{mat:normaldecomp}
\end{align}
The last term is orthogonal to all tangent vectors of the product-state manifold. Therefore the exact squared normal speed is
\begin{equation}
 \|N\|^2=\frac{g^2}{\hbar^2}
          \operatorname{Var}_z(A)\operatorname{Var}_\varphi(B).
 \label{mat:normalspeed}
\end{equation}
All reference correlations can be included in $z$. On open preparation sets where both variances are positive, the product manifold is not invariant. Replacing the full dynamics by reciprocal equations on its tangent space removes a nonzero physical term.

For a projector $A$, a work momentum $P=-i\hbar\partial_y$ and $p=\|A\Xi\|^2$, the exact acquisition state is
\begin{equation}
 \psi_t(y)=(I-A)\Xi\varphi(y)+A\Xi\varphi(y-K(t)).
 \label{mat:translation}
\end{equation}
If $\varphi$ has finite position variance $s^2$, its resulting position density and variance are
\begin{equation}
 (1-p)|\varphi(y)|^2+p|\varphi(y-K)|^2,
 \qquad \operatorname{Var}(Y)=s^2+K^2p(1-p).
 \label{mat:variance}
\end{equation}
The mean shift is $Kp$, but the extra fluctuations and source/work correlation are determined by the same interaction. They have not been independently added to obscure a noiseless mean signal. The squared-norm interpretation of the material position density remains a preparation/configuration commitment, separately supplied in the relevant realization.

\section{A finite acquisition, genuine copy, and return}

Choose $\varphi\in C_c^\infty(\mathbb R)$ with unit norm and independent memory blanks $|00\rangle_{BC}$. Put $\Xi_0=(I-A)\Xi$, $\Xi_1=A\Xi$. During storage take the specified free work Hamiltonian to be zero; any other free motion belongs in a different complete pulse calculation. First implement \eqref{mat:translation}. Next apply a local read rotation
\[
 H_{\rm read}(t)=\hbar\dot b(t)\beta(y)\sigma_y^B,
 \qquad b:0\longrightarrow1,\quad 0\le\beta\le\pi/2.
\]
Use finite controlled rotations to copy $B$ to $C$ and reset $B$:
\[
 |00\rangle\mapsto\cos\beta|00\rangle+\sin\beta|10\rangle
 \mapsto\cos\beta|00\rangle+\sin\beta|11\rangle
 \mapsto|0\rangle_B(\cos\beta|0\rangle+\sin\beta|1\rangle)_C.
\]
Finally undo the source-controlled translation. With
$m(\theta)=\cos\theta|0\rangle+\sin\theta|1\rangle$, the complete final state is
\begin{equation}
 \psi_f(y)=\varphi(y)|0\rangle_B
       [\Xi_0m(\beta(y))_C+\Xi_1m(\beta(y+K))_C].
 \label{mat:echo}
\end{equation}
All rotations can use finite smooth pulses. The ideal prescribed schedule is a control resource; no autonomous microscopic clock has been derived by writing it down.

\begin{proposition}[Exact complete-state origin of the record/return tradeoff]
\label{mat:echoresult}
Under the squared-norm material configuration law, the actual terminal bit in \eqref{mat:echo} has probability
\begin{equation}
 R(p)=(1-p)r_0+pr_1,\qquad
 r_a=\int|\varphi(y)|^2\sin^2\beta(y+aK)dy.
 \label{mat:record}
\end{equation}
Let
\[
 D=\int|\varphi(y)|^2\cos[\beta(y+K)-\beta(y)]dy.
\]
Then the complete source/reference reduced state differs from its initial pure state by
\begin{equation}
 \tfrac12\|\rho_{SR,f}-|\Xi\rangle\langle\Xi|\|_1
       =\sqrt{p(1-p)}(1-D),\qquad
 |r_1-r_0|\le\sqrt{1-D^2}.
 \label{mat:echobound}
\end{equation}
In particular, exact source/reference return for one $0<p<1$ forces equal bit probabilities on the two source sectors.
\end{proposition}
\begin{proof}
Squaring the $C=1$ component of \eqref{mat:echo} and integrating gives \eqref{mat:record}, because $\Xi_0$ and $\Xi_1$ are orthogonal. The normalized work/archive branch vectors are $e_a(y)=\varphi(y)m(\beta(y+aK))$, with $\langle e_0,e_1\rangle=D$. Tracing these vectors multiplies the source off-diagonal block by $D$. In the span of the normalized nonzero $\Xi_0,\Xi_1$, the difference has eigenvalues $\pm\sqrt{p(1-p)}(1-D)$, proving the equality. Further,
\[
 |r_1-r_0|\le\int|\varphi|^2|\sin(\beta(y+K)-\beta(y))|dy
 \le\sqrt{1-D^2},
\]
by Cauchy--Schwarz and $\int|\varphi|^2\cos^2\theta\ge D^2$. Since $0\le D\le1$, exact return implies $D=1$ and therefore $r_0=r_1$.
\end{proof}

The retained archive is essential. For an arbitrary branch-preserving isometry
\[
 \Xi_a\longmapsto\Xi_a\otimes e_a,
 \qquad c=\langle e_0,e_1\rangle,
\]
including every returned key and environment in $e_a$, the same two-dimensional calculation gives
\begin{equation}
 \epsilon_p=\sqrt{p(1-p)}|1-c|,
 \quad d_{\rm arc}\le\sqrt{1-|c|^2},
 \quad \epsilon_p\ge\tfrac12\sqrt{p(1-p)}d_{\rm arc}^2.
 \label{mat:allarchives}
\end{equation}
The last inequality follows from $|1-c|\ge1-|c|\ge(1-|c|^2)/2$. A later common unitary on a returning archive preserves the complete-state comparison. An unread reduced channel alone would not supply this guarantee.

These elementary overlap inequalities are established Hilbert-space mathematics, recovered here inside an explicit contact. The new constitutive content is the full material field and its admitted transducers, not the algebraic inequality itself. A separately added ordinary classical writer of a nonlinear function of the source ray is excluded only by that material inventory premise.

\section{A distinct hybrid class: force supported by conversion}

The following construction retains an ordinary classical canonical pair $(R,P)$ and a finite coherent source. It must therefore be kept separate from the complete material field above. Let $H(P)=H(P)^\dagger$ and let conversion maps $L_c(P)$ have material coefficients independent of the unknown vector. Define
\begin{equation}
 K=\partial_PH,\qquad D=\sum_cL_c^\dagger L_c,
 \qquad F_\psi=\langle K\rangle_\psi,\quad h_\psi=\langle D\rangle_\psi.
 \label{mat:KD}
\end{equation}
The source-blind free coordinate velocity is omitted from $F$.

\begin{theorem}[Force--conversion domination]
\label{mat:domination}
For fixed $P$ and $c\ge0$, the following are equivalent:
\begin{align}
 |F_\psi|&\le c h_\psi\quad\hbox{for every unit ray},\\
 -cD&\le K\le cD,\\
 \ker D&\subseteq\ker K,\qquad
 \|D_{\rm supp}^{-1/2}K_{\rm supp}D_{\rm supp}^{-1/2}\|\le c.
 \label{mat:dominationeq}
\end{align}
The equivalence holds under arbitrary inaccessible reference extension. If $D\ge\kappa_0E$, $D=EDE$ and $K=EKE$, it holds with $c=\|K\|/\kappa_0$.
\end{theorem}
\begin{proof}
Taking the two signs of the quadratic-form inequality gives the operator inequalities. If $x\in\ker D$, the positive operators $cD\pm K$ have zero quadratic form at $x$ and hence annihilate $x$, so $Kx=0$. On the support, conjugation by $D^{-1/2}$ makes the two inequalities equivalent to spectrum in $[-c,c]$. Reversing this argument proves sufficiency. Tensoring identities preserves order, and the final assertion follows from $|\langle K\rangle|\le\|K\|\langle E\rangle$.
\end{proof}

Merely setting the dark--dark block of $K$ to zero is insufficient: a dark--reactive cross block has expectation of order $\sqrt{\langle E\rangle}$ near a dark ray. The linear bound excludes that block too. The equivalence is explicit; it is a characterization of a response restriction, not its physical necessity.

\section{A common binding level that still writes a null likelihood}

Let $A$ be a source projector, let $g,e,b$ label ground, reactive, and spent modes, and put
\[
 E=A\otimes|e\rangle\langle e|,
 \qquad X=A\otimes(|e\rangle\langle g|+|g\rangle\langle e|).
\]
The single contact has
\begin{equation}
 H(P)=\hbar[gX+(\Delta+\beta P)E],\qquad
 L=\sqrt\kappa A\otimes|b\rangle\langle e|,
 \quad D=\kappa E.
 \label{mat:hybridH}
\end{equation}
Its declared hybrid actualization law is
\begin{align}
 \dot\psi&=[-iH(P)/\hbar-\tfrac12(D-\langle D\rangle_\psi)]\psi,\\
 \dot R&=P/M+\hbar\beta\langle E\rangle_\psi,\quad\dot P=0,\\
 h_\psi&=\kappa\langle E\rangle_\psi,\qquad
 \psi^+=L\psi/\|L\psi\|.
 \label{mat:hybridlaw}
\end{align}
One capture consumes the channel; the accumulated coordinate remains. The normalized-history canonical force and stochastic extraction are primitive equations here. In particular, they have not been obtained from a unitary bath while silently treating its pointer as classical.

The same bound level supplies both responses:
\begin{equation}
 \dot R-P/M=(\hbar\beta/\kappa)h_\psi.
 \label{mat:commonresponse}
\end{equation}
This is a nonempty realization of Theorem~\ref{mat:domination}, with no direct incoming dark-sector tap. It has a finite exact resolvent self-energy
\[
 \Sigma(z,P)=\frac{g^2}{z-\Delta-\beta P+i\kappa/2},
 \quad -2\operatorname{Im}\Sigma(E,P)
   =\frac{g^2\kappa}{(E-\Delta-\beta P)^2+\kappa^2/4}.
\]
Both virtual dispersion and real loss remain present. A time-local elimination of the bound mode would require a further approximation; none is needed below.

For $p=\|A\Xi\|^2$, define exact bright amplitudes
\begin{equation}
 \binom{u_P(t)}{e_P(t)}=
 \exp\left[t\begin{pmatrix}0&-ig\\-ig&-\kappa/2-i(\Delta+\beta P)\end{pmatrix}\right]
 \binom10,
 \quad a_P=|u_P|^2+|e_P|^2.
 \label{mat:bright}
\end{equation}
Then $a_P'=-\kappa|e_P|^2$ and the unnormalized null vector and its mass are
\[
 \zeta_p=(I-A)\Xi\otimes g+A\Xi\otimes(u_Pg+e_Pe),
 \qquad N_p=1-p+pa_P.
\]
The capture density is $p\kappa|e_P(t)|^2dt$, with daughter $A\Xi/\sqrt p$ in the spent mode. The null daughter is $\zeta_p/\sqrt{N_p}$. Zero-mass daughters are never normalized.

\begin{proposition}[Null-likelihood writer]
\label{mat:nullwriter}
Before the first capture, the actual classical pointer in \eqref{mat:hybridlaw} satisfies
\begin{equation}
 R(t)=R(0)+Pt/M-\frac{\hbar\beta}{\kappa}\log N_p(t).
 \label{mat:logwriter}
\end{equation}
For $0<a_P(T)<1$ this gives a strictly increasing nonlinear response to $p$ on an actual null branch of positive mass.
\end{proposition}
\begin{proof}
The normalized reactive population is $p|e_P|^2/N_p$, so
$\partial_t\log N_p=-h_\psi$. Integration of \eqref{mat:commonresponse} gives \eqref{mat:logwriter}. The first and second derivatives of $-\log(1-p+pa)$ are $(1-a)/(1-p+pa)$ and $(1-a)^2/(1-p+pa)^2$, both positive when $a<1$.
\end{proof}

This countermodel is more restrictive than an independent force tap: it requires actual reactive support and a fixed ratio between force and hazard. Its failure cannot be repaired by mentioning reciprocal backaction without changing its equations.

\section{A complete finite counterexperiment}
\label{mat:counterexperiment}

Use three fresh binding labels on the same unknown carried vector, addressing $A_0=I$, $A_1=A$, $A_2=I-A$ in sequence. A first capture blocks further contacts. After two nulls read and copy a function of the first two pointers; the third contact is a later return test, not a condition for retaining the copy. At each exposure end switch off both exchange and conversion, retaining any excited residue. Take $\hbar=\beta=\kappa=M=1$, $\Delta=0$ and
\begin{equation}
 T=2\log2,\qquad g=\sqrt{1/16+\pi^2/(4\log^22)},
 \quad \omega=\sqrt{g^2-1/16}=\pi/T.
 \label{mat:parameters}
\end{equation}
At $P=0$,
\[
 u(t)=e^{-t/4}[\cos\omega t+\sin\omega t/(4\omega)],\quad
 e(t)=-ie^{-t/4}(g/\omega)\sin\omega t,
\]
so $(u(T),e(T))=(-1/\sqrt2,0)$ and $a=1/2$.

Compare the actual preparation laws
\begin{align}
 \mathcal E_B &: |0r_0\rangle\ (2/3),\quad |1r_1\rangle\ (1/3),\\
 \mathcal E_S &: \sqrt{2/3}|0r_0\rangle\pm\sqrt{1/3}|1r_1\rangle
                      \quad(1/2\text{ each}),
 \label{mat:ensembles}
\end{align}
with $r_0,r_1$ orthogonal inaccessible reference states and $A=|0\rangle\langle0|$. These have equal averaged density matrices but different complete ray laws. No active preparation label is supplied to the apparatus; if one is supplied it must stay represented.

With offsets and free drifts initially zero, the two-null pointer ratio is
\[
 z(p)=\frac{\log(1-p+p/2)}{\log(1/2)},
 \qquad z(0)=0,\ z(1)=1,\ z(2/3)=\log_2(3/2).
\]
A smooth finite plate equal to one on $[0.45,0.72]$, and zero near 0 and 1, is enabled only after two actual nulls. It is evaluated away from the denominator cutoff $R_0>1/2$. Its output is copied to an ordinary classical memory by the stipulated downstream mechanical contact.

\begin{theorem}[Persistent finite gap with all stopping branches retained]
\label{mat:finitegap}
The copied bit has probabilities $\Pr_B(C=1)=0$ and $\Pr_S(C=1)=1/3$. The four stopping branches for a fixed $p$ have masses
\begin{equation}
 \tfrac12,\qquad\tfrac p4,\qquad\tfrac{1-p}4,\qquad\tfrac14.
 \label{mat:fourbranches}
\end{equation}
On the triple-null branch the normalized original source/reference vector is exactly restored, while the already acquired copy survives.
\end{theorem}
\begin{proof}
For arbitrary bright survivals $a_0,a_1,a_2$, successive norm loss gives masses
\[
 1-a_0,\quad a_0p(1-a_1),\quad a_0(1-p)(1-a_2),
 \quad a_0[pa_1+(1-p)a_2].
\]
They sum to one for every $p$. Setting $a_j=1/2$ proves \eqref{mat:fourbranches}. The acquisition branch has mass $a_0(1-p+pa_1)$. Its value for $p=2/3$ is $1/3$, while the classifier is zero for both basis rays. For triple null, with $u=-1/\sqrt2$, the source maps are $uI$, $uA+(I-A)$, and $A+u(I-A)$. Their product is $u^2I=I/2$. Thus the return probability is $1/4$ for every ray, including every inaccessible reference extension.
\end{proof}

The continuous output is also fixed. Let
\[
 N_j(t)\Xi=(I-A_j)\Xi\otimes g_j+A_j\Xi\otimes v_{P_j}(t),
 \quad J_j(t)\Xi=\sqrt\kappa e_{P_j}(t)A_j\Xi\otimes b_j.
\]
The four unnormalized maps are
\begin{equation}
 J_0(t),\quad J_1(t)N_0(T),\quad
 J_2(t)N_1(T)N_0(T),\quad N_2(T)N_1(T)N_0(T).
 \label{mat:completekernel}
\end{equation}
Their squared norms are respectively time densities or the final null atom. For each history integrate \eqref{mat:hybridlaw} to its actual stopping time and retain the resulting pointers, clocks, flags, energy account, unused supplies, and copies. Although each source map in \eqref{mat:completekernel} is linear at fixed readiness, the full classical output depends nonlinearly on the input ray through \eqref{mat:logwriter}. Dropping that output is not a valid complete-instrument comparison.

The counterexample is stable at strictly positive phase volume. Independently prepare $|P_j|\le10^{-4}$ and $|R_j(0)|\le10^{-3}$. For total apparatus horizon $3T+12$,
\[
 \|v_P(t)-v_0(t)\|\le|P|t,\quad
 |a_P(T)-1/2|\le d:=2T10^{-4},\quad
 |\delta R_{\rm free}|\le10^{-3}+10^{-4}(3T+12).
\]
These follow by Duhamel for contraction propagators and direct integration of free velocity. Substitution into the logarithmic expressions bounds the attained ratio intervals by
\[
 \begin{array}{c|c}
 p& R_1/R_0\\\hline
 0&[-0.003792,0.003792]\\
 1&[0.99089,1.00920]\\
 2/3&[0.57814,0.59185]
 \end{array}
\]
and $R_0>0.6899$. An additional ratio error $0.02$ leaves all classifications fixed. If $\bar a$ denotes the exact average of $a_P(T)$ over the supplied momentum law, then
\begin{equation}
 \Pr_S(C=1)=\bar a(1/3+2\bar a/3),\qquad
 \Pr_S(NNN,C=1)=\bar a^2,
 \label{mat:positivewidth}
\end{equation}
while both basis probabilities remain zero. Triple-null unnormalized error is at most $\epsilon=3T10^{-4}$; its normalized complete source/reference/binding trace distance from the returned ray is at most $\epsilon/(1/2-\epsilon)<0.000833$. The finite-width estimates condition uniformly on each supported readiness, so postselection does not invalidate them.

The physical conclusion is limited and useful: reactive support plus common force/conversion coefficients does not force complete record affinity. A stochastic resource that compensates energy at capture is a supplied law in this hybrid model, not an independently derived bath. A separately admitted finite-stirring carrier realization \cite{M27} supplies its native paths; its renewal and chamber-volume premises are not consequences of the null-writer theorem.

\chapter[Directed transport and physical acquisition]{A common directed interaction for timing, continuation, and acquisition}
\label{mat:chiralchapter}

This chapter consolidates the directed-transport repair of the preceding null-writer model \cite{M29}. Its elementary decay clock is replaced by conservative spatial transport and a prepared actual coordinate. The same derivative that causes effective incoming loss also changes the dynamics after a physical position copy. It is a complete conditional detector constitution, not an embedding of the old canonical actual carrier. Its directed, non-semibounded generator also differs from the later massive guidance construction and finite pilot theory.

\section{Complete source and conservative Hamiltonian}

Let $\mathcal K=\mathcal S\otimes\mathcal R$, with $\mathcal R$ inaccessible, and let $A$ be a projector on $\mathcal S$. The cell space is
\begin{equation}
 \mathcal H_{\rm cell}=L^2(\mathbb R_-;\mathcal K)_g
            \oplus L^2(\mathbb R;A\mathcal K)_c.
 \label{mat:chiralspace}
\end{equation}
The $g$ mode is stationary and incoming. The $c$ mode moves right; its regions $x<0$ and $x>0$ are called reactive and spent. They are parts of one channel, not separate modes connected by a prescribed collapse. The incoming $g/e$ components remain coherent and share one actual position $X$.

For fixed $g>0$, $v>0$ and real $\Delta$, define
\begin{equation}
 H\binom GC=
 \binom{\hbar g A C|_{x<0}}
 {-i\hbar v\partial_x C+\hbar\Delta1_{x<0}C+\hbar g1_{x<0}AG},
 \quad D(H)=L^2(\mathbb R_-;\mathcal K)\oplus H^1(\mathbb R;A\mathcal K).
 \label{mat:chiralH}
\end{equation}
The full experiment retains all active and spent fields, actual coordinates, copies, clock phases, unused cells, and supplied classical histories. A conditional source slice is not the whole state when waves can return.

\begin{lemma}[Self-adjoint realization and current]
\label{mat:selfadjoint}
The operator \eqref{mat:chiralH} is self-adjoint. In the incoming region
\[
 \partial_t\rho+\partial_xj=0,
 \qquad\rho=\|G\|^2+\|C\|^2,\qquad j=v\|C\|^2,
\]
and on $x>0$, $\rho=\|C\|^2$ and $j=v\rho$.
\end{lemma}
\begin{proof}
The diagonal operator $0\oplus(-i\hbar v\partial_x)$ is self-adjoint on the stated domain. Restriction and extension across the half-line form mutually adjoint bounded exchange blocks; detuning is bounded and real. The bounded self-adjoint perturbation theorem gives the result. In the local norm derivative the exchange terms cancel, the real detuning contributes zero, and the remaining term is $-v\partial_x\|C\|^2$. The full-line $H^1$ trace matches current across zero.
\end{proof}

Prepare one unknown $\Psi\in\mathcal K$, $\|\Psi\|=1$, and
\begin{equation}
 (G_0,C_0)=(\chi\Psi,0),\qquad
 \chi(x)=\sqrt\alpha e^{\alpha x/2}\ (x<0),\qquad
 \Pr(X_0\in dx)=\alpha e^{\alpha x}dx,
 \label{mat:chiralready}
\end{equation}
with $\alpha>0$. The actual coordinate follows $\dot X=j/\rho$. Both the guidance law and the squared-norm readiness are explicit statistical/mechanical commitments. There is no additional Born measurement of the source. The ready vector belongs to $D(H)$, since the differentiated component is initially zero, and $\|H\Phi_0\|^2=\hbar^2g^2\|A\Psi\|^2$.

The ideal directed Hamiltonian is unbounded below. A finite ready energy variance does not make it a semibounded microscopic material model. No such bath realization is claimed. Prescribed switching and outgoing cam control remain resources.

\section{Exact event law and retained source}

Write $p=\|A\Psi\|^2$, $\kappa=\alpha v$, and define
\begin{equation}
 \dot u=-ige,\qquad
 \dot e=-igu-(\kappa/2+i\Delta)e,
 \qquad (u(0),e(0))=(1,0),\quad a=|u|^2+|e|^2.
 \label{mat:reducedODE}
\end{equation}
The next result derives this dissipative-looking equation from \eqref{mat:chiralH}.

\begin{theorem}[Common transport, waiting law, and continuation]
\label{mat:common}
For \eqref{mat:chiralH}--\eqref{mat:chiralready}, the exact field is
\begin{align}
 G(x,t)&=\chi(x)[(I-A)\Psi+u(t)A\Psi],&&x<0,\\
 C(x,t)&=\chi(x)e(t)A\Psi,&&x<0,\\
 C(x,t)&=\sqrt\alpha e(t-x/v)A\Psi,&&0<x<vt,
 \label{mat:exactfield}
\end{align}
with zero outgoing field for $x>vt$. If $\tau$ is the first crossing of zero and $N_p=1-p+pa$, then
\begin{align}
 \Pr(\tau>t)&=N_p(t),&\Pr(\tau\in dt)&=\kappa p|e(t)|^2dt,\\
 \Psi_\varnothing(t)&=\frac{(I-A)\Psi\otimes g+
          A\Psi\otimes[u(t)g+e(t)e]}{\sqrt{N_p(t)}},
 &\Psi_\tau&=A\Psi/\sqrt p,\\
 X_t-X_0&=-\alpha^{-1}\log N_p(t),&&t<\tau,\\
 \mathcal L(X_t\mid\tau>t)&=\alpha e^{\alpha x}1_{x<0}dx.
 \label{mat:commonlaw}
\end{align}
The daughter formula applies only when $p>0$. Outgoing temporal amplitudes remain part of the complete field.
\end{theorem}
\begin{proof}
Substitute the incoming ansatz into Schr\"odinger evolution. Since $\chi'=\alpha\chi/2$, the transport derivative supplies precisely $-\alpha ve/2$ in the time equation. This gives \eqref{mat:reducedODE}. Forward characteristics and continuity at zero give the outgoing field. The front at $x=vt$ is continuous because $e(0)=0$. Self-adjoint uniqueness identifies the solution.

The reduced equation gives $a'=-\kappa|e|^2$. The incoming norm is $N_p$ and the outgoing norm is $p\kappa\int_0^t|e(s)|^2ds=1-N_p$. Current is nonnegative; a trajectory crosses at most once. More explicitly, the incoming velocity is independent of $x$:
\[
 \dot X=vp|e|^2/N_p=-\alpha^{-1}\partial_t\log N_p.
\]
Thus $U=e^{\alpha X_0}$ is uniform and $\tau>t$ iff $U<N_p(t)$. This proves the waiting law directly from the prepared actual coordinate. The incoming factorization and outgoing source orientation give the conditional vectors. A surviving coordinate $x$ came from $x_0=x+\alpha^{-1}\log N_p$; change of variables gives joint surviving density $N_p\alpha e^{\alpha x}$. Dividing by $N_p$ proves the last line. At finite times $a>0$, since the two-dimensional propagator is invertible, so finite nulls have no normalization singularity.
\end{proof}

In the reaction-history filtration containing the occurrence time but no initial-coordinate record, the compensator is
\begin{equation}
 \int_0^{t\wedge\tau}\frac{\kappa p|e(s)|^2}{N_p(s)}ds.
 \label{mat:compensator}
\end{equation}
Indeed conditional survival from $s$ to $t$ is $N_p(t)/N_p(s)$; differentiation gives the intensity. In the larger filtration revealing $X_0$ and the complete field, crossing is predictable. Its compensator is the predictable boundary count, not \eqref{mat:compensator}. Prepared configuration randomness has replaced intrinsic chemical randomness; these are different complete constitutions.

An actual finite output cam can be included on $0<x<\ell$ by
\[
 H_{\rm cam}=v[p_x+c f(x)p_y],\qquad
 f\in C_c^\infty(0,\ell),\quad\int_0^\ell f(x)dx=1.
\]
Its divergence-free characteristics have $\dot x=v$, $\dot y=vcf(x)$; each completed passage shifts the latch by $c$. A ready latch packet narrower than $c/3$ has disjoint unactuated and completed regions. The record completion time is $\tau+\ell/v$. At a finite deadline, incomplete cams remain pending outputs. A clearance interval completes them without inventing a completed record at the earlier cut.

\section{A physical copy changes the reaction through the same derivative}

Prepare a memory packet $\eta(y)$ of positive width and the actual joint coordinate law $|\chi(x)|^2|\eta(y)|^2dx\,dy$, conditionally independent of the complete prior source and actual past. This is an additional ready-resource law, not a consequence of the product wave alone. Freeze the binding contact, then apply $H_{\rm copy}(t)=a_c(t)x p_y$ with $\int a_cdt=1$. It is a source-blind coordinate acquisition. The complete new ready field is
\begin{equation}
 G_0(x,y)=\chi(x)\eta(y-x)\Psi,\qquad C_0=0.
 \label{mat:copiedready}
\end{equation}
During binding the memory coordinate is stationary and retained. In the Fourier convention $\eta(z)=(2\pi)^{-1/2}\int e^{ikz}\widehat\eta(k)dk$, put
\begin{equation}
 U_t=\mathcal F^{-1}[\widehat\eta(k)u_{\Delta-vk}(t)],\qquad
 V_t=\mathcal F^{-1}[\widehat\eta(k)e_{\Delta-vk}(t)].
 \label{mat:copiedamplitudes}
\end{equation}
Subscripts denote detuning in \eqref{mat:reducedODE}.

\begin{theorem}[Complete acquired-record and reaction law]
\label{mat:copytheorem}
The exact incoming field after the copy is
\begin{equation}
 \Phi_-(x,y,t)=\chi(x)\{\eta(y-x)(I-A)\Psi\otimes g+
       A\Psi\otimes[U_t(y-x)g+V_t(y-x)e]\}.
 \label{mat:copiedfield}
\end{equation}
Define
\begin{align}
 m(y)&=\int_{-\infty}^0|\chi(x)|^2|\eta(y-x)|^2dx,\\
 b(y,t)&=\int_{-\infty}^0|\chi(x)|^2
                 (|U_t(y-x)|^2+|V_t(y-x)|^2)dx,\\
 s_p(y,t)&=(1-p)m(y)+pb(y,t).
\end{align}
Then
\begin{align}
 \Pr(y\in dy,\tau>t)&=s_p(y,t)dy,\\
 \Pr(y\in dy,\tau\in dt)&=\kappa p|V_t(y)|^2dy\,dt,\\
 \Pr(\tau>t\mid y)&=s_p(y,t)/m(y),\\
 h_y(t)&=1_{t<\tau}\kappa p|V_t(y)|^2/s_p(y,t).
 \label{mat:copylaw}
\end{align}
After an acquired $y$ and a null with $s_p(y,t)>0$, the retained incoming vector is $\Phi_-(\cdot,y,t)/\sqrt{s_p(y,t)}$. Conditional formulas are used only for $m(y)>0$ and positive survivor density, almost everywhere in the actual record law. Its spatial degree cannot be dropped if it can subsequently return.
\end{theorem}
\begin{proof}
Fourier transforming in $y$ writes the incoming profile as $\chi(x)e^{-ikx}\widehat\eta(k)$. Applying $-iv\partial_x$ contributes $-i\kappa/2-vk$, so the bright detuning becomes $\Delta-vk$. Inverse Fourier transformation yields \eqref{mat:copiedfield}. Orthogonality of $A\Psi$ and $(I-A)\Psi$, including the reference, gives $s_p$. The boundary flux at $x=0$ is $v\alpha p|V_t(y)|^2$, while integrated continuity gives $\partial_tb=-\kappa|V_t(y)|^2$. Conditioning on the initial memory density $m$ proves the final two formulas.
\end{proof}

This is the connecting equation: copying position changes the spatial dependence, and the unchanged transport derivative changes the reaction law and continuation. The copy does not leave an independent classical displacement available for the old log-likelihood experiment. Discarding $y$ gives only the coarser survivor
\[
 1-p+p\overline a(t),\qquad
 \overline a(t)=\int|\widehat\eta(k)|^2a_{\Delta-vk}(t)dk,
\]
which cannot replace the record-conditioned law or its complex retained amplitudes.

\begin{proposition}[Finite position resolution can suppress conversion]
\label{mat:copybound}
For the Gaussian $\eta_\sigma(y)=(2\pi\sigma^2)^{-1/4}e^{-y^2/(4\sigma^2)}$, if $4g^2>\kappa^2/4$, then
\begin{equation}
 \Pr(\tau\le T)\le
 \kappa pT\frac{\sigma\sqrt{2/\pi}}v
 \frac{4\pi g^2}{\sqrt{4g^2-\kappa^2/4}}.
 \label{mat:resolutionbound}
\end{equation}
The bound is uniform in the fixed detuning $\Delta$.
\end{proposition}
\begin{proof}
At detuning $\delta$, the two dissipative eigenvalues satisfy $\operatorname{Re}\lambda_\pm\le0$ and
\[
 e_\delta(t)=-ig\frac{e^{\lambda_+t}-e^{\lambda_-t}}{\lambda_+-\lambda_-},
 \quad |\lambda_+-\lambda_-|^2
 =|(\kappa/2+i\delta)^2-4g^2|
 \ge\delta^2+4g^2-\kappa^2/4.
\]
Therefore $|e_\delta(t)|^2\le4g^2/(\delta^2+4g^2-\kappa^2/4)$. The Gaussian momentum density is bounded by $\sigma\sqrt{2/\pi}$. Insert this bound into the integrated event density, change variable $\delta=\Delta-vk$, and integrate the Lorentzian. Integration over $[0,T]$ proves \eqref{mat:resolutionbound}.
\end{proof}

At \eqref{mat:parameters} with $\kappa=v=1$, the coefficient is $15.941283\,p\sigma$. Thus $p=2/3$, $\sigma=10^{-3}$ gives conversion below $0.010629$, whereas the uncopied contact converts with probability $1/3$. Resolution improves at momentum cost $\operatorname{Var}(p_y)=\hbar^2/(4\sigma^2)$. This is a finite acquisition/backaction consequence, not a universal tradeoff for all contacts. An active controller trying to undo the correlation must retain every copy in the new Hamiltonian calculation.

\section{Null reuse and the preparation boundary}

At \eqref{mat:parameters}, a source phase correction turns the endpoint bright amplitude $-1/\sqrt2$ into $1/\sqrt2$ without changing the incoming spatial factor. The null coordinate distribution in Theorem~\ref{mat:common} is again the ready exponential. Consequently the same unmeasured cell, after a null and a finite source phase correction of duration $d$, admits $m$ attempts with
\begin{equation}
 \Pr(\hbox{all }m\hbox{ null})=1-p+p2^{-m},\qquad
 f_k(s)=p2^{-(k-1)}\kappa|e(s)|^2,\quad 0<s<T.
 \label{mat:reuse}
\end{equation}
Here $f_k(s)ds$ is the unconditional first-arrival mass at physical time $(k-1)(T+d)+s$. The proof is multiplication of the actual null filter; there are no reactions while $g=v=0$ in the phase windows. This is conditional renewal of an unmeasured exponential resource, not independence of successive nulls. After a position copy the factorization fails and Theorem~\ref{mat:copytheorem} replaces \eqref{mat:reuse}. After capture the cell is spent.

\begin{proposition}[Invariant incoming span selects a shape, not a probability law]
\label{mat:profile}
For constant nonzero $g$ and $v>0$, a one-dimensional incoming spatial span $\{\chi(x)\xi\}$ is invariant for every binding/source vector under the incoming differential expression if and only if
\[
 \chi'=z\chi,\qquad
 \chi(x)=\sqrt{2\operatorname{Re}z}\,e^{zx},
 \quad\operatorname{Re}z>0,
\]
up to a constant phase. The effective decay coefficient is $2v\operatorname{Re}z$ and detuning is $\Delta+v\operatorname{Im}z$.
\end{proposition}
\begin{proof}
Applying the incoming generator to an $e$ component requires $\chi'$ to lie in the same one-dimensional span. The weak equation $\chi'=z\chi$ has only exponential solutions. Square integrability on the negative half-line requires positive real exponent and fixes normalization. Conversely substitution proves invariance and the displayed coefficients. With only initial $g$ support, the nonzero exchange creates an $e$ component and its next derivative imposes the same condition.
\end{proof}

This characterizes a resource property. It does not prepare that resource from a broader class, select squared-norm actual coordinates, or establish universality of directed material transport.

The old three-contact experiment has a useful regression test. If only the present surviving clock coordinate is acquired after a null, its conditional law is the source-independent ready exponential. For $R_j=-c\alpha X_j$, the two surviving clocks are independent $\operatorname{Exp}(1/c)$ variables conditional on the double-null source branch. Applying the same finite plate as before gives
\[
 \theta=\int_{c/2}^{\infty}\frac{e^{-r/c}}c
       \left(e^{-0.45r/c}-e^{-0.72r/c}\right)dr,
\]
independently of $p$. The four source stopping masses remain \eqref{mat:fourbranches}; hence both ensembles in \eqref{mat:ensembles} have $\Pr(C=1)=\theta/3$ and $\Pr(NNN,C=1)=\theta/4$. A physical initial-coordinate copy instead obeys \eqref{mat:copylaw}, with its changed timing and daughters. This distinction prevents using the unmeasured null displacement and a freely known initial position in the same experiment.

The supported chain for this directed cell is therefore
\[
 \begin{gathered}
 \text{specified coherent field and interface Hamiltonian}\\
 +\ \text{guidance and independent exponential squared-norm readiness}\\
 \Longrightarrow\ \text{actual first passage, null continuation, product field}\\
 \Longrightarrow\ \text{copy-dependent detuning and complete acquired-record law}.
 \end{gathered}
\]
The original canonical Hamiltonian current and its actual carrier are absent from this cell's complete state. Adding Bell rates on an enlarged configuration would be a benchmark, not a derivation of that embedding. A conversion boundary flux is not the original Hamiltonian edge current, and the reaction-history intensity is not a law in the filtration revealing the actual ready coordinate. Complete material-field admission, guidance, readiness, and the ideal directed channel remain physical commitments of this cell. Its theorems do not close the canonical event-law bridge; the pilot construction addresses that bridge with a different complete state and microscopic mechanism, while the massive construction establishes a different operational event law.

% End consolidated source: parts/02_material.tex

\part{Statistical Selection of Event Histories}\label{part:statistics}

% Begin consolidated source: parts/03_statistics.tex
\chapter{Current, incidence, and conditional timing}
\label{stat:current-chapter}

The Hamiltonian continuity equation does not by itself specify a stochastic
history. This part consolidates a variational current-to-history selection
argument and the countermodels that delimit it. The principal result is a
finite-horizon theorem over a class that initially permits history-dependent
reaction rates. A proposed relative-entropy law selects an explicit
finite-background process, and its zero-background limit is the Hamiltonian
Bell process, with a total-variation estimate on complete paths. The
probability-bearing premises remain visible. This result selects timing
within its stated class, but does not deduce a stochastic law from
source/readout incompleteness alone \cite{M30,C03,BellQFT}. The later pilot
construction uses microscopic gas and reaction dynamics instead of this
path-entropy premise.

\section{The wave programme and the actual history}

Let $\mathcal H$ be finite dimensional and let $P_1,\ldots,P_d$ be a fixed
orthogonal resolution of its identity. A normalized source wave satisfies
\begin{equation}\label{stat:schrodinger}
 i\hbar\dot\Psi_t=H(t)\Psi_t,
 \qquad w_m(t)=\|P_m\Psi_t\|^2,
 \qquad
 J_{nm}(t)=\frac{2}{\hbar}\operatorname{Im}
       \langle\Psi_t,P_nH(t)P_m\Psi_t\rangle.
\end{equation}
The time parameter is physical time. It is not replaced by accumulated
exposure. Throughout the main theorem $H$ has finitely many constant,
self-adjoint segments on a fixed interval $[0,T]$. This gives an analytic
extension of the wave across each closed segment considered separately.
No finite-node assertion is inferred merely from analyticity on an open
interval with a possibly singular endpoint.

The source state includes $\Psi_t$, an actual sector $X_t$, and all
future-active apparatus and memory factors. For a fixed complete preparation
and an open-loop joint Hamiltonian, $\Psi_t$ is deterministic. An event changes
$X_t$; it does not collapse or reset the full wave. A null interval advances
the same wave while leaving the actual sector unchanged. A record is a
specified function of the joint configuration or of a physically acquired
archive. A freely readable copy of the unperturbed history is not supplied
by this definition.

Let $N_{nm}(t)$ count transitions $m\to n$ up to time $t$. A candidate law
$P$ is carried by the finite-jump paths in
$D([0,T],\{1,\ldots,d\})$. Its natural filtration is denoted
$\mathcal F_t$. Predictable intensities are written
$\lambda_{nm}(t\mid\mathcal F_{t-})$ on the event $X_{t-}=m$; they may initially
depend on the entire observed native history. The compensator convention is
\begin{equation}\label{stat:compensator}
 A_{nm}(t)=\int_0^t 1_{\{X_{s-}=m\}}
       \lambda_{nm}(s\mid\mathcal F_{s-})\,ds,
 \qquad N_{nm}-A_{nm}\text{ is a martingale}.
\end{equation}
All candidates used below have finite expected total jump count. Define their
directed expected flows by
\begin{equation}\label{stat:mean-flow}
 F_{nm}(t)=\mathbb E_P\!\left[
   1_{\{X_{t-}=m\}}\lambda_{nm}(t\mid\mathcal F_{t-})\right].
\end{equation}
Thus $\mathbb E_P N_{nm}(B)=\int_B F_{nm}(t)\,dt$ for every measurable time
set $B$.

\begin{assumption}[Statistical current realization]\label{stat:matching}
For each complete unordered sector pair, the actual expected signed event
flow equals the specified Hamiltonian current:
\begin{equation}\label{stat:pair-current}
 F_{nm}(t)-F_{mn}(t)=J_{nm}(t)
 \quad\text{for almost every }t.
\end{equation}
\end{assumption}

This is an identification of a source current with an expectation of actual
counts. Assigning a conserved nonprobabilistic charge the density $w$ and
flux $J$ does not already prove it. Likewise, the random atomic measure
$dN_{nm}$ cannot be set equal pathwise to the usually absolutely continuous
measure $[J_{nm}]_+dt$. Even the signed difference of the two pathwise count
measures is atomic. Equation~\eqref{stat:pair-current} concerns mean measures,
and \eqref{stat:compensator} concerns conditional expectations.

A permanently inaccessible reference may be included inside the sector
fibers. For a finite carried system with an arbitrary reference and local
controls $H_{SA}\otimes I_R$, the initial Schmidt support in $R$ is finite and
is preserved by all such controls. The finite-dimensional theorems therefore
apply on that invariant support. This does not grant control of the reference
or extra copies of the unknown input.

\section{Three different restrictions on an event law}

\begin{proposition}[Continuity, pair matching, and remaining traffic]
\label{stat:traffic-family}
The wave data in \eqref{stat:schrodinger} satisfy
\begin{equation}\label{stat:continuity}
 J_{nm}=-J_{mn},\qquad \dot w_n=\sum_{m\ne n}J_{nm}.
\end{equation}
Under Assumption~\ref{stat:matching}, every candidate satisfies
\begin{equation}\label{stat:offset}
 \pi_n(t)=w_n(t)+c_n,
 \qquad \pi_n(t)=P(X_t=n),\quad c_n=\pi_n(0)-w_n(0).
\end{equation}
At any time the nonnegative flows realizing a fixed pair current have
exactly the form
\begin{equation}\label{stat:surplus}
 F_{nm}=[J_{nm}]_++K_{nm},\qquad
 K_{nm}=K_{mn}\geq0.
\end{equation}
Matching only the divergences $\sum_m(F_{nm}-F_{mn})=\dot w_n$ is weaker:
it permits antisymmetric current reassignments with zero divergence in
addition to the symmetric surplus in \eqref{stat:surplus}.
\end{proposition}
\begin{proof}
Differentiate $\langle\Psi,P_n\Psi\rangle$ using Schr\"odinger evolution and
self-adjointness. Inserting the sector resolution gives
\eqref{stat:continuity}. On every finite-jump path,
\[
 1_{\{X_t=n\}}-1_{\{X_0=n\}}
   =\sum_{m\ne n}\bigl(N_{nm}(t)-N_{mn}(t)\bigr).
\]
Take expectations, substitute \eqref{stat:pair-current}, and integrate
\eqref{stat:continuity}. This proves \eqref{stat:offset}. For a pair with
$J_{nm}\geq0$, set $K_{nm}=F_{mn}$. Then $F_{nm}=J_{nm}+K_{nm}$. Reversing
the orientation covers the other sign and proves uniqueness of the surplus.
Finally, for a current array $\widetilde J$ with the same divergence as $J$,
$C=\widetilde J-J$ is antisymmetric and has zero row sums. Conversely, any
such $C$ preserves the divergence. Taking
$F_{nm}=[J_{nm}+C_{nm}]_++K_{nm}$ realizes that divergence with nonnegative
flows. On graphs with cycles nonzero choices of $C$ exist.
\end{proof}

Consequently equivariance, pairwise Hamiltonian-current matching, minimal
directed traffic, and full conditional timing are different mathematical
claims. Initial equality $\pi(0)=w(0)$ plus pair matching implies equivariance;
it does not remove $K$, and it does not make conditional rates constant
across histories with the same present sector.

\begin{example}[A stationary Hamiltonian circulation]\label{stat:ring}
On three sectors let $S|m\rangle=|m+1\bmod3\rangle$ and
\begin{equation}\label{stat:ring-hamiltonian}
 H=\frac{i\hbar a}{2}(S-S^\dagger),\qquad
 \Psi=\frac{|0\rangle+|1\rangle+|2\rangle}{\sqrt3},\qquad a>0.
\end{equation}
Here $H\Psi=0$, each weight is $1/3$, and the clockwise current on each
edge is $a/3$. Thus zero actual traffic preserves the weights but fails pair
matching. Within the matching class, the Markov rates
\begin{equation}\label{stat:ring-surplus}
 k_{m+1,m}=(1+c)a,\qquad k_{m,m+1}=ca,\qquad c\geq0,
\end{equation}
all preserve the weights and currents. They use the same declared source
coordinates $(\Psi,X)$; an additional quotient-fiber representative is not
needed. Their first-event survival is $e^{-(1+2c)at}$, so the surplus changes
finite event timing.
\end{example}

This example separates informational completion from a physical preference
for fewer events. It also shows why minimizing total activity subject only
to population continuity would erase a stationary Hamiltonian circulation:
zero activity is then admissible. An event-law variational principle must
state whether it constrains each Hamiltonian edge or only its divergence.

\begin{corollary}[Reduced reservoir rates do not bound actual turnover]
\label{stat:reservoir-turnover}
Suppose a reduced population calculation supplies
\[
 \dot p_1=k_{10}(t)p_0-k_{01}(t)p_1,\qquad p_0+p_1=1,
\]
with bounded nonnegative rates on a finite interval. For every bounded
nonnegative function $b(t)$, the time-inhomogeneous Markov rates
\[
 \widehat\lambda_{10}=k_{10}+b p_1,\qquad
 \widehat\lambda_{01}=k_{01}+b p_0
\]
give the same population solution from the same initial law. Their extra
expected traffic is $2\int b p_0p_1\,dt$.
\end{corollary}
\begin{proof}
The extra directed flows both equal $B=b p_0p_1$, so they cancel in the
forward equation. Uniqueness of the finite-state forward equation gives
the asserted marginals, and summing the two extra flows gives the count.
The rates remain bounded even at a node: this realization of the formal
addition $B/p_i$ needs no singular division there.
\end{proof}
For example, take $p_0=p_1=1/2$, $k_{10}=k_{01}=\kappa$ and constant
$b>0$. The first-event survival changes from $e^{-\kappa T}$ to
$e^{-(\kappa+b/2)T}$, while every population is unchanged. Even when a
weak-coupling calculation gives $\kappa=O(|g|^2)\to0$, the additional
turnover need not vanish. This is a counterexample to inference from a
reduced population equation, not an assertion that these rates arise
from the same fully specified microscopic Hamiltonian. A microscopic
record-reliability argument must identify its actual generator or bound
surplus traffic separately \cite{C03}; compare
Equation~\eqref{cfg:traffic} and the product-field calculation in
Appendix~\ref{bench:chapter}.

\section{Minimal mean incidence does not determine a history}

\begin{proposition}[Non-Markov rivals with exactly minimal mean flow]
\label{stat:renewal}
For the fixed wave programme of Example~\ref{stat:ring}, there are stationary
clockwise processes having weights $1/3$, mean edge flow $a/3$, and no
counterclockwise events, whose complete path laws differ from the clockwise
rate-$a$ Bell process. One such competitor has bounded natural-filtration
intensities and finite relative entropy to every strictly positive
constant-rate Markov reference on a finite horizon.
\end{proposition}
\begin{proof}
First let $\Theta$ be uniform on $[0,3)$ and set
\[
 X_t=\lfloor\Theta+at\rfloor\bmod3.
\]
The translating phase is uniform modulo $3$ at each time, so sector weights
are $1/3$. Each edge is crossed with stationary intensity $a/3$ in the mean.
The first waiting time is uniform on $[0,1/a]$ and all subsequent interevent
times equal $1/a$. Hence the no-event probability is $(1-at)_+$, differing
from $e^{-at}$ by $e^{-1}$ at $t=1/a$. This competitor has a singular clock
and is not used in the finite-entropy minimization below.

For a regular competitor take a stationary renewal process whose interevent
law is Erlang with shape $2$ and rate $2a$:
\[
 f_L(u)=4a^2u e^{-2au},\qquad
 S_L(u)=e^{-2au}(1+2au),\qquad \mathbb E L=1/a.
\]
Such a process on the line can be constructed by taking the interval
containing zero from the size-biased law $a\ell f_L(\ell)d\ell$, placing zero
uniformly in that interval, and using independent copies of $L$ to either
side. Give the sector at zero an independent uniform label and advance it
clockwise at each renewal. Stationary event intensity is $a$, and uniform
translation of the independent sector label gives mean flow $a/3$ on every
edge. There are no reverse events.

The residual first-wait survival is
\begin{equation}\label{stat:erlang-survival}
 S_{\rm res}(t)=a\int_t^\infty S_L(u)\,du
              =e^{-2at}(1+at).
\end{equation}
Its hazard before the first observed event is
$a(1+2at)/(1+at)$, bounded by $2a$. After an observed event, the hazard at
age $u$ is
\begin{equation}\label{stat:erlang-hazard}
 h_L(u)=\frac{4a^2u}{1+2au},
\end{equation}
also bounded by $2a$ and not constant. Thus the natural history changes the
conditional law while the current and mean directed incidence remain
unchanged. The process has finitely many expected jumps, and its bounded
intensities give finite path relative entropy against any reference whose
finite set of rates is strictly positive. The entropy identity proved in
Lemma~\ref{stat:entropy-identity} makes this last assertion explicit: its
integrand is bounded on $0\leq\lambda\leq2a$, including $\lambda=0$.
\end{proof}

No assumption of complete current information rules out this example by
itself. One can encode its age in a hidden source coordinate or leave it in a
history-dependent kernel. Requiring $(\Psi,X)$ to be statistically sufficient
for the future rules it out, but that requirement supplies a stochastic
closure law. In a time-inhomogeneous Markov theory the first-wait survival
from $m$ at $s$ is
\begin{equation}\label{stat:survival}
 S_m(s,t)=\exp\left(-\int_s^t\sum_{n\ne m} k_{nm}(u)\,du\right).
\end{equation}
It is an ordinary exponential in elapsed time only when the total rate is
constant. The usual unit exponential variable is a threshold in integrated
hazard, not a claim that every physical waiting time has constant rate.

\section{Calibration from a control-stable statistical premise}

\begin{theorem}[Coherent-evacuation calibration]\label{stat:calibration}
Suppose a single finite wave programme admits an actual finite-mean-count
history satisfying Assumption~\ref{stat:matching}. Suppose that for each
sector $n$ there is a time $t_n\in[0,T]$ at which $w_n(t_n)=0$. Then
$\pi(t)=w(t)$ throughout that programme, without separately stipulating the
initial occupancy distribution.
\end{theorem}
\begin{proof}
By \eqref{stat:offset}, $\pi_n=w_n+c_n$. At $t_n$, positivity implies
$c_n=\pi_n(t_n)\geq0$. Both distributions have total mass one, so
$\sum_n c_n=0$. Therefore every $c_n$ vanishes.
\end{proof}

A fixed ancillary flip supplies a nonempty programme with the required
zeros. Let an unknown source $S$ be entangled with an inaccessible reference
$R$, let $\{Q_k\}$ be a chosen resolution on $S$, and prepare the \emph{wave}
of a two-level ancilla $A$ as $|0\rangle$. Declare the joint sectors
$Q_k\otimes I_R\otimes|b\rangle\langle b|$ and use
\begin{equation}\label{stat:ancilla-flip}
 H_A=\hbar g I_{SR}\otimes\sigma_x,\qquad
 0\leq t\leq\pi/(2g).
\end{equation}
With $v_k=\|(Q_k\otimes I_R)\Psi\|^2$, direct evolution gives
\begin{equation}\label{stat:flip-current}
 w_{k0}=v_k\cos^2(gt),\qquad w_{k1}=v_k\sin^2(gt),
 \qquad J_{k1,k0}=gv_k\sin(2gt).
\end{equation}
Every $(k,1)$ sector is empty initially, and every $(k,0)$ sector is empty
finally. No actual ancillary distribution was used in this calculation.
The theorem forces the complete occupancy weights if current realization
holds in every $k$ block. Merely matching the summed ancillary current does
not calibrate the source-sector weights. The flip leaves the unknown
$SR$ wave and its reference correlations untouched.

This is a consistency theorem for an admitted class of preparations and
controls, not a relaxation procedure for arbitrary nonequilibrium. An
initial $\pi\ne w$ simply cannot satisfy the same statistical-current
premise throughout that programme. To infer the law for an experiment that
chooses another later control, the allowed preparations must obey causal
stability: the already prepared distribution cannot depend on which later
allowed programme is selected. Preparation of the ancillary wave and
validity of blockwise current realization under its control remain physical
premises. A reversed flip does not make its retained history a fresh
independent resource.

\begin{corollary}[Finite calibration defect]\label{stat:calibration-defect}
Suppose instead that
$\pi_n(t)=w_n(t)+c_n+e_n(t)$, where $e_n(0)=0$ and
$|e_n(t)|\leq\delta_n$. If every sector has a time with
$w_n(t_n)\leq\zeta_n$, then, with total variation normalized as one half of
$\ell^1$ distance,
\begin{equation}\label{stat:calibration-bound}
 d_{\rm TV}(\pi(0),w(0))\leq\sum_n(\zeta_n+\delta_n).
\end{equation}
\end{corollary}
\begin{proof}
Positivity at $t_n$ gives $c_n\geq-\zeta_n-\delta_n$. Since $\sum_n c_n=0$,
$d_{\rm TV}=\sum_{c_n<0}|c_n|$, proving the bound. For a current defect
$r_{nm}=F_{nm}-F_{mn}-J_{nm}$, the counting identity supplies
$e_n(t)=\int_0^t\sum_m r_{nm}(s)ds$, so the hypothesis can be checked from an
integrated statistical-current error rather than presumed pointwise.
\end{proof}

\chapter{Variational selection of complete event paths}
\label{stat:entropy-chapter}

\section{The proposed law and the initial candidate class}

Fix positive symmetric constants $a_{nm}=a_{mn}$ for $n\ne m$, with
$a_{nm}\leq a_{\max}$, independent of the unknown wave. Let
$R_\varepsilon$ be the Markov path law having rates $\varepsilon a_{nm}$ and
uniform initial sector distribution, for $\varepsilon>0$. It is a reference
probability measure. No physical equilibrium bath or uniform random
controller is thereby prepared, but the reference carries explicit Poisson
statistical structure.

\begin{assumption}[Relative-entropy reaction principle]\label{stat:entropy-law}
For a prescribed complete wave programme and fixed initial occupancy law,
the actual finite-background path law minimizes
$D(P\Vert R_\varepsilon)$ among laws realizing each expected Hamiltonian edge
current. The zero-background constitution is the finite-horizon limit of
these minimizers as $\varepsilon\downarrow0$.
\end{assumption}

This principle can be stated before the Bell formula: neutral reaction
statistics are changed by the least relative entropy compatible with a
specified signed transport. Its status is a proposed physical law on path
probabilities, not an inference from ignorance or from a quotient theorem.
Both the choice of reference and the extremization are additional
statistical input. The principle has a finite-background prediction
containing two-way traffic; it is not simply the assertion that every rate
already equals the Bell ratio. The connection of this proposed law to an
independently derived material interaction remains open \cite{M30,BFG}.

For the theorem, the candidate class consists of all laws with initial
$w(0)$, finite expected jump count, predictable natural-history intensities,
finite relative entropy to $R_\varepsilon$, and
Assumption~\ref{stat:matching}. Initial $w(0)$ is either supplied, or follows
from Theorem~\ref{stat:calibration} on its stronger control-stable domain.
The competitors are not required to be Markov. Their wave programme is the
same deterministic one. All finite classical provenance is retained or
conditioned upon; different active preparations are not identified by
averaging their density matrices.

Define, for $y>0$,
\begin{equation}\label{stat:entropy-density}
 \ell(x;y)=x\log(x/y)-x+y,\qquad x\geq0,
\end{equation}
with $0\log0=0$. Also set $\ell(0;0)=0$ and $\ell(x;0)=+\infty$ for $x>0$.

\section{Existence and normalization at nodes}

The construction below is needed before the optimizer is identified with
an actual probability law. Rates can diverge as a wave sector vanishes;
bounded rates are not silently assumed.

\begin{lemma}[Nonexplosive construction for the selected fluxes]
\label{stat:existence}
For $\varepsilon\geq0$, define directed fluxes
\begin{align}
 F^\varepsilon_{nm}(t)
 &=\frac{\sqrt{J_{nm}(t)^2+
            4\varepsilon^2a_{nm}^2w_m(t)w_n(t)}+J_{nm}(t)}{2},
                                                        \label{stat:flux}\\
 F^0_{nm}(t)&=[J_{nm}(t)]_+.
\end{align}
On $w_m(t)>0$ set $k^\varepsilon_{nm}=F^\varepsilon_{nm}/w_m$.
For the finite piecewise-constant Hamiltonian domain, these rates define a
unique nonexplosive Markov law from initial $w(0)$ with marginals $w(t)$.
Values assigned to outgoing rates at isolated zero-weight times have no
effect on this law and may be zero. Identically empty sectors are never
occupied. For $\varepsilon>0$, the law has finite entropy relative to
$R_\varepsilon$.
\end{lemma}
\begin{proof}
Write $h=\max_t\|H(t)\|$. Cauchy--Schwarz gives
\begin{equation}\label{stat:current-bound}
 |J_{nm}(t)|\leq\frac{2h}{\hbar}\sqrt{w_m(t)w_n(t)}.
\end{equation}
In particular, a current vanishes if either endpoint weight is zero.
The fluxes satisfy $F^\varepsilon_{nm}-F^\varepsilon_{mn}=J_{nm}$ and
\begin{equation}\label{stat:excess-bound}
 F^\varepsilon_{nm}=[J_{nm}]_++s^\varepsilon_{nm},\qquad
 0\leq s^\varepsilon_{nm}=s^\varepsilon_{mn}
       \leq\varepsilon a_{nm}\sqrt{w_mw_n}.
\end{equation}
The last inequality follows from
$\sqrt{x^2+4b^2}\leq|x|+2b$. Therefore
\begin{equation}\label{stat:finite-flow}
 \int_0^T\sum_{n\ne m}F^\varepsilon_{nm}(t)\,dt<\infty.
\end{equation}

On each constant-Hamiltonian segment, a nonidentically zero weight has only
finitely many zeros, because it is analytic on a neighborhood of the closed
segment. Subdivide at all such zeros and at the Hamiltonian switches. In the
interior of one resulting interval $(a,b)$ the occupied-sector set is fixed,
and its rates are continuous and locally bounded. Put
$q_m=\sum_{n\ne m}k^\varepsilon_{nm}$ and use the survival function
\eqref{stat:survival}. The usual first-jump construction on compact
subintervals gives a minimal process, killed only if infinitely many jumps
accumulate or it reaches an endpoint through a state whose escape integral
diverges.

The wave weights solve the forward equation for these fluxes. On a
node-free interval, variation of constants gives
\begin{equation}\label{stat:duhamel}
 w_n(t)=w_n(a)S_n(a,t)+
   \int_a^t \sum_{m\ne n}k^\varepsilon_{nm}(u)w_m(u)S_n(u,t)\,du.
\end{equation}
At a left endpoint with $w_n(a)=0$, derive the formula first from
$a+\delta$. The boundary term is at most $w_n(a+\delta)$, which tends to zero;
the nonnegative integral has the required limit. At a positive-weight left
endpoint the ordinary limit applies. Identically zero sectors have zero
incoming flow.

Iteration of the positive integral equation counts paths by their number
of jumps. Its minimal subprobability solution $\pi$ is bounded componentwise
by the nonnegative solution $w$. Consequently its stopped counting
compensator has expectation bounded by
\eqref{stat:finite-flow}. Infinite-jump accumulation has probability zero:
for the stopping time of the $L$th jump, the probability that it occurs
before $T$ is at most the bound in \eqref{stat:finite-flow} divided by $L$.

There is no separate loss of mass into a vanishing sector at $b$. Indeed,
on any positive-weight portion,
$q_m\geq-\dot w_m/w_m$, and thus
\begin{equation}\label{stat:node-survival}
 S_m(s,t)\leq w_m(t)/w_m(s).
\end{equation}
If $w_m(b)=0$, the right side tends to zero as $t\uparrow b$. A path entering
$m$ at any earlier time therefore cannot remain there up to $b$. A path
with infinitely many entries was already excluded. With no explosion and
no node loss, the minimal subprobability has total mass one. The domination
$\pi\leq w$ then forces $\pi=w$. At the finitely many boundaries the process
occupies only positive-weight sectors almost surely, so the next interval
can be joined with no discretionary reset. The first-jump construction is
unique on every compact node-free interval; these joins give uniqueness
from the stated initial law.

For completeness, finite entropy is not lost at the nodes. Set
$b_{nm}=\varepsilon a_{nm}\sqrt{w_mw_n}$. By
\eqref{stat:current-bound}, when both weights are positive,
\[
 F^\varepsilon_{nm}=\sqrt{w_mw_n}\,r_{nm}(t),
\]
where for fixed $\varepsilon>0$ all $r_{nm}$ lie between strictly positive
finite constants: this follows directly from \eqref{stat:flux} with
$|J_{nm}|/\sqrt{w_mw_n}\leq2h/\hbar$ and the finitely many positive
$a_{nm}$. Hence
\[
 \left|\log\frac{F^\varepsilon_{nm}}{\varepsilon a_{nm}w_m}\right|
 \leq C_\varepsilon+\tfrac12|\log w_n|+\tfrac12|\log w_m|.
\]
Each nonzero analytic weight has a finite-order zero, so its logarithm grows
at most as a constant times $1+|\log|t-t_0||$ near a node. Multiplication by
the bounded flux in \eqref{stat:finite-flow} is locally integrable. Empty
intervals contribute zero candidate flow and at most the finite reference
escape term. The entropy identity below therefore has a finite right side,
including the finite initial entropy against the uniform distribution.
\end{proof}

\section{Entropy on complete histories}

\begin{lemma}[Likelihood identity and history penalty]
\label{stat:entropy-identity}
For a candidate law $P$ in the specified class, let $u_m=1/d$. Then
\begin{align}\label{stat:kl}
 D(P\Vert R_\varepsilon)
  =D(w(0)\Vert u)+
   \int_0^T\sum_{m}\mathbb E_P\!\left[
    1_{\{X_{t-}=m\}}\sum_{n\ne m}
    \ell(\lambda_{nm}(t\mid\mathcal F_{t-});
                    \varepsilon a_{nm})\right]dt.
\end{align}
Every candidate has marginal $w(t)$, and consequently
\begin{equation}\label{stat:jensen}
 D(P\Vert R_\varepsilon)\geq D(w(0)\Vert u)
 +\int_0^T\sum_{n\ne m}
       \ell(F_{nm}(t);\varepsilon a_{nm}w_m(t))\,dt.
\end{equation}
Equality holds precisely when, for almost every time and every occupied
origin, the natural-history intensity is almost surely the deterministic
quantity $F_{nm}(t)/w_m(t)$.
\end{lemma}
\begin{proof}
For a finite jump history with successive states $x_0,\ldots,x_L$ and jump
times $t_1<\cdots<t_L$, its likelihood factors into initial probability,
conditional survival factors, and the intensities of its realized jumps.
Relative to the strictly positive reference the log likelihood is
\[
 \log\frac{w_{x_0}(0)}{u_{x_0}}
 +\sum_{i=1}^L\log
       \frac{\lambda_{x_i x_{i-1}}(t_i\mid\mathcal F_{t_i-})}
            {\varepsilon a_{x_i x_{i-1}}}
 -\int_0^T\sum_{n\ne X_{t-}}
        (\lambda_{nX_{t-}}-\varepsilon a_{nX_{t-}})\,dt.
\]
Taking expectations of the jump sum with its compensator gives
\eqref{stat:kl}. For unbounded intensities one first stops when the total
count or integrated intensity exceeds a finite bound and clips the
logarithm. The negative part of $x\log(x/y)$ is bounded by $y/e$, and the
reference rates are bounded. Thus this negative part is uniformly
integrable on the finite horizon. The positive part converges by monotone
truncation; the integrated linear terms converge by the finite-mean-count
assumption. Finite relative entropy, or the value $+\infty$ before
restriction to the candidate class, yields the same identity in the limit.

Proposition~\ref{stat:traffic-family} with initial $w(0)$ gives marginal
$w(t)$. Conditional on $X_{t-}=m$, the mean intensity is $F_{nm}/w_m$.
Strict convexity of $x\mapsto\ell(x;y)$ and Jensen's inequality give
\[
 w_m\,\mathbb E[\ell(\lambda_{nm};\varepsilon a_{nm})\mid X_{t-}=m]
 \geq w_m\ell(F_{nm}/w_m;\varepsilon a_{nm})
 =\ell(F_{nm};\varepsilon a_{nm}w_m).
\]
At zero origin weight the candidate contributes zero. Equality in a
strictly convex Jensen inequality requires an almost surely constant
conditional intensity, including the zero-mean case by nonnegativity.
Integrating its nonnegative defect proves the stated equality criterion.
\end{proof}

This is the connecting step that an optimization of mean traffic lacks.
Histories with the same current sector are allowed different responses at
the outset, as in Proposition~\ref{stat:renewal}. The specified reference
and entropy principle penalize that variation. Markovity is a consequence
of this particular statistical choice, not of conserving the mean current.

\begin{theorem}[Selection of rates and physical-time path law]
\label{stat:selection}
Within the preceding finite piecewise-constant wave domain, the
relative-entropy reaction principle has a unique minimizing path law
$P_\varepsilon$, up to null histories. It is the Markov law of
Lemma~\ref{stat:existence}, with rates
\begin{equation}\label{stat:selected-rate}
 k^\varepsilon_{nm}(t)=
 \frac{\sqrt{J_{nm}(t)^2+
       4\varepsilon^2a_{nm}^2w_m(t)w_n(t)}+J_{nm}(t)}{2w_m(t)}
 \quad(w_m(t)>0).
\end{equation}
As $\varepsilon\downarrow0$ its complete path law converges to the Bell law
$P_{\rm B}$ with conditional rates
\begin{equation}\label{stat:bell-rate}
 k^{\rm B}_{nm}(t)=[J_{nm}(t)]_+/w_m(t),
\end{equation}
and
\begin{equation}\label{stat:path-bound}
 d_{\rm TV}(P_\varepsilon,P_{\rm B})
 \leq\int_0^T\sum_{n\ne m}s^\varepsilon_{nm}(t)\,dt
 \leq\varepsilon a_{\max}(d-1)T.
\end{equation}
The selected law is consistent under restriction to earlier intervals and
does not use future controls to determine an earlier intensity.
\end{theorem}
\begin{proof}
The Jensen lower bound reduces the problem to independent pairwise convex
optimizations at each time. On an unordered edge write $x=F_{nm}$,
$y=F_{mn}$, $x-y=J_{nm}$, $c=\varepsilon a_{nm}w_m$, and
$d'=\varepsilon a_{nm}w_n$. When both weights are positive, minimize
$\ell(x;c)+\ell(y;d')$ along $x-y=J_{nm}$. In the interior, differentiation
along that constraint gives
\[
 \log(x/c)+\log(y/d')=0,
 \qquad xy=cd'=\varepsilon^2a_{nm}^2w_mw_n.
\]
The derivative tends to minus infinity at the admissible boundary where
one flow is zero, and the objective is strictly convex and coercive.
Consequently its unique minimum is the positive solution of these two
equations, namely \eqref{stat:flux}. If an endpoint weight is zero, its
outgoing flow must vanish for finite entropy; current matching and zero
Hamiltonian current then force the other flow to vanish too.

Lemma~\ref{stat:existence} realizes all the minimizing fluxes by a
nonexplosive Markov law of finite entropy. It attains the Jensen bound.
Any other minimizer must attain both the pointwise pair minima and the
strict conditional Jensen equality. Its intensities therefore agree with
\eqref{stat:selected-rate} on all occupied histories for almost every
time. The unique first-jump construction proves equality of the full laws,
not merely of their one-time marginals.

To prove \eqref{stat:path-bound}, start both processes at the same sampled
initial sector and couple them while their states coincide. Use common
jumps at the Bell rates and additional jumps for the finite-background
process at rates $s^\varepsilon_{nm}/w_m$. Upon separation continue with
any coupling of the correct marginals. On compact node-free intervals this
is an ordinary finite-state jump construction; the nonexplosive limits
from Lemma~\ref{stat:existence} join it across the node times. The chance of
separation is at most the expected number of extra jumps before separation.
The probability that the coupled pair is still together at $m$ is at most
$P_\varepsilon(X_t=m)=w_m(t)$. Hence that expectation is at most
$\int\sum_{n\ne m}s^\varepsilon_{nm}dt$. By
\eqref{stat:excess-bound} and Cauchy--Schwarz,
\[
 \sum_{n\ne m}\sqrt{w_nw_m}
       =\left(\sum_m\sqrt{w_m}\right)^2-1\leq d-1.
\]
The coupling inequality yields \eqref{stat:path-bound}. It controls the
entire finite path, including null intervals and event times.

Finally, the optimizer at time $t$ depends only on the current $w,J$ and
reference coefficients. Optimizing on an earlier interval produces the
same rates there. More generally, restarting at a deterministic time with
its induced marginal gives the same restricted Markov kernel. Thus the
finite-horizon variational definition is projectively consistent on the
fixed wave programme, and coincident prefixes of two allowed open-loop
programmes have coincident selected prefix laws.
\end{proof}

The rates supply an executable process: draw competing jumps using the
integrated survival \eqref{stat:survival}, update the actual sector at a
jump, and continue the full Schr\"odinger wave. The exponential thresholds
used to simulate that derived kernel are a representation of the selected
path measure. They are not newly postulated physical threshold coordinates
whose readability has been established.

\section{The exact limit and a finite-background discriminator}

\begin{proposition}[Activity and timing are separate levels of the law]
\label{stat:lexicographic}
For every admissible $P$ with the common marginal $w$, the exact identity is
\begin{align}\label{stat:entropy-expansion}
 D(P\Vert R_\varepsilon)
 ={}&\mathbb E_PN_T\log(1/\varepsilon)+D(P\Vert R_1)\\
 &+(\varepsilon-1)\int_0^T\sum_m w_m(t)
                  \sum_{n\ne m}a_{nm}\,dt.\nonumber
\end{align}
The Bell law is also obtained by first minimizing expected total activity
under pair-current matching and then minimizing $D(P\Vert R_1)$ among
those activity minimizers.
\end{proposition}
\begin{proof}
Expand \eqref{stat:entropy-density} to obtain
$\ell(x;\varepsilon y)=\ell(x;y)+x\log(1/\varepsilon)
 +(\varepsilon-1)y$, and integrate the identity
\eqref{stat:kl}. Its linear count term is $\mathbb E_PN_T$ by the
compensator. The final term is fixed across competitors.
For each pair, \eqref{stat:surplus} gives activity $|J|+2K$.
Thus the smallest integrated activity forces $K=0$ almost everywhere.
With these fixed minimal flows, Lemma~\ref{stat:entropy-identity} selects
the Markov intensities $[J]_+/w$ uniquely. Their finite entropy against
$R_1$ follows by the same logarithmic node estimate as in
Lemma~\ref{stat:existence}; alternatively positive-part fluxes obey
\eqref{stat:current-bound} and have no worse integrability. This proves
the second assertion independently of any formal exchange of minimization
and limit.
\end{proof}

One must not minimize directly against the zero-rate reference. A path with
a jump is singular to that reference. The zero-background limit is a
stated constitutive prescription supported by Theorem~\ref{stat:selection}.
It is not a supplied finite reservoir with a demonstrated preparation
procedure or a finite-cost realization of the limit.

At finite background the selected opposite flows satisfy
\begin{equation}\label{stat:traffic-product}
 F^\varepsilon_{nm}F^\varepsilon_{mn}
       =\varepsilon^2a_{nm}^2w_mw_n.
\end{equation}
This relation and the response to zero current are consequences beyond
terminal Born weights. For $H=0$, uniform $w$, and all $a_{nm}=a$, the
selected chain jumps at rate $\varepsilon a$ toward each other sector,
whereas the Bell path remains constant. Therefore
\begin{equation}\label{stat:sharpness}
 d_{\rm TV}(P_\varepsilon,P_{\rm B})
       =1-e^{-\varepsilon a(d-1)T}.
\end{equation}
Indeed $P_{\rm B}$ is concentrated on constant paths; the finite-background
law assigns them total mass $e^{-\varepsilon a(d-1)T}$, with the same
uniform initial-state proportions. Its remainder is supported on paths
having at least one jump. This proves equality and first-order sharpness
of \eqref{stat:path-bound}. It is a native-path discriminator, not an
admission theorem for a passive recorder of those paths.

The elementary pair optimization is related to established relative-entropy
contractions for Markov empirical flows and currents \cite{BFG}. The
Bell positive-current construction is established independently
\cite{BellQFT}. Here the consolidated result is the finite-horizon
history selection, node-safe realization, uniform background path bound,
and explicit scope of their source-law interpretation. No claim of
historical priority for the square-root formula is needed.

\section{Complete outputs, memories, and conditioning}

\begin{corollary}[Operational consequences on a common experiment]
\label{stat:output-bound}
Fix a complete finite apparatus programme satisfying the theorem, including
all its returning memories and control factors, and let $\mathcal O$ be
any common measurable map of its actual path and deterministic wave into
an output space. Then
\begin{equation}\label{stat:pushforward}
 d_{\rm TV}(\mathcal O_*P_\varepsilon,\mathcal O_*P_{\rm B})
 \leq\delta,
 \qquad \delta=\varepsilon a_{\max}(d-1)T.
\end{equation}
For a common event $E$ with $P_{\rm B}(E)\geq r>0$, the conditional output
laws, whenever both exist, differ in total variation by at most
$\min(1,2\delta/r)$. If the output depends only on the final actual sector
and the common deterministic wave, its two laws are exactly equal.
\end{corollary}
\begin{proof}
A measurable pushforward cannot increase total variation: the preimage of
each output event is a path event. For conditional laws, write $p=P_{\rm B}$
and $q=P_\varepsilon$. For any subevent $B\subset E$,
\[
 \left|\frac{q(B)}{q(E)}-\frac{p(B)}{p(E)}\right|
 \leq \frac{|q(B)-p(B)|}{p(E)}
       +\frac{q(B)}{q(E)}\frac{|q(E)-p(E)|}{p(E)}
 \leq\frac{2\delta}{r}.
\]
Applying this to preimages proves the conditional claim. Finally both
endpoint sector distributions are $w(T)$, and the wave is identical, so
any common endpoint map has exactly equal distributions.
\end{proof}

The output statement preserves actual classical flags, copied keys, null
outcomes, and event-time stopping if they are defined by the same complete
experiment. It invokes no contraction theorem for an unknown nonlinear
state update. If an output includes a normalized relative wave, use the
same sector projection in both models and assign a common arbitrary value
at zero-weight sectors, which occur with zero endpoint probability.

Endpoint equality is insufficient for a claim that an archive faithfully
reports an actual earlier sector. At $\varepsilon>0$, the reference in this
chapter has strictly positive rates on every pair. It can therefore induce
surplus jumps between occupied archive sectors even when the Hamiltonian
has no matrix element between them. For two such sectors of weights
$1/2$, rate $\varepsilon a$ in either direction gives an archive-label flip
probability
\begin{equation}\label{stat:archive-background}
 \frac{1-e^{-2\varepsilon aT}}{2}.
\end{equation}
The endpoint label remains uniform, but its correlation with the initial
actual label decays. Exact archive support in the configuration chapters
therefore applies to the Bell limit, or requires the approximate path
budget above at finite background. Replacing this reference by one with
zero rates on archive-changing pairs changes the hypothesis and requires
rechecking feasibility and entropy support; it is not done silently.

The filtration is part of the theorem. Uniqueness of a law on $X$ does not
select every joint extension by a hidden variable. For example, if a
selected first event has time $\tau$ with $0<P(\tau\leq s)<1$, the extension
$Z=1_{\{\tau\leq s\}}$ leaves the entire marginal path law unchanged but
makes $P(\tau\leq s\mid Z)=Z$. Its one-bit archive is informative if a
physical interaction actually makes $Z$ available. Minimizing entropy on
$X$ alone imposes no constraint on that joint extension. A claimed theorem
for such an enlarged source must include the variable and its acquisition
interaction in the compared complete laws, or state the conditional
independence premise that excludes the extension. This is an admission
boundary, not a reason to treat all internal variables as readable.

For growing apparatus banks the bound is finite-programme control, not
dimension-free closure. A simultaneous limit requires
$\varepsilon a_{\max}(d-1)T\to0$ for the \emph{complete} retained sector
count $d$ and physical horizon $T$. Classical mixtures can be treated by
conditioning on their retained provenance and averaging the bound. Random
feedback is covered when represented by a finite retained coherent
controller in one deterministic joint programme. Substituting a random
branch-dependent wave into the unconditional fixed-wave proof requires a
new conditional formulation; no such substitution is implicit here.

\chapter{Minimal traces, projected histories, and the constitutive boundary}
\label{stat:boundary-chapter}

\section{What the Jordan theorem does and does not select}

The successive names ``no-surplus incidence,'' ``single-channel exclusion,''
``Jordan representation,'' and ``minimal positive boundary trace'' describe
closely related conditions in the checkpoint corpus. They are consolidated
here into one equivalence theorem, with the statistical objects made
explicit \cite{C03,M30}. This prevents the same mathematical condition from
being counted repeatedly as an independent derivation.

\begin{theorem}[Minimal positive mean incidence]\label{stat:jordan}
Fix a complete edge and an integrable prescribed signed current $J(t)$.
Let $\mu(B)=\int_B J(t)dt$ and let finite nonnegative measures $\alpha,\beta$
satisfy $\alpha-\beta=\mu$. There is a unique finite nonnegative measure
$\rho$ such that
\begin{equation}\label{stat:jordan-family}
 \alpha=\mu^++\rho,\qquad \beta=\mu^-+\rho,
 \qquad \alpha+\beta=|\mu|+2\rho,
\end{equation}
where $\mu^+(dt)=[J(t)]_+dt$ and $\mu^-(dt)=[-J(t)]_+dt$.
The following conditions are equivalent:
\begin{enumerate}
 \item $\alpha+\beta$ is the least possible total incidence measure;
 \item its total mass is the least possible among representations of $\mu$;
 \item $\alpha$ and $\beta$ are mutually singular;
 \item $\alpha=\mu^+$ and $\beta=\mu^-$.
\end{enumerate}
If $\alpha(dt)=F_{nm}(t)dt$ and $\beta(dt)=F_{mn}(t)dt$, these are also
equivalent to $F_{nm}F_{mn}=0$ almost everywhere and to
$F_{nm}+F_{mn}=|J_{nm}|$ almost everywhere.
\end{theorem}
\begin{proof}
Let $A_+$ and $A_-$ be a Hahn decomposition for $\mu$. For every measurable
$B$,
\[
 \mu^+(B)=\mu(B\cap A_+)
          =\alpha(B\cap A_+)-\beta(B\cap A_+)\leq\alpha(B).
\]
Thus $\rho=\alpha-\mu^+$ is nonnegative. The identity
$\alpha-\beta=\mu^+-\mu^-$ gives also $\rho=\beta-\mu^-$, proving
\eqref{stat:jordan-family} and uniqueness. Its least sum in measure order
is $|\mu|$, and its least total mass is $|\mu|([0,T])$; either is attained
exactly when $\rho=0$. If $\alpha$ and $\beta$ are mutually singular, their
common submeasure $\rho$ must vanish. Conversely the Jordan parts are
mutually singular. In the absolutely continuous case,
$\rho(dt)=K(t)dt$ with $K\geq0$, and the final equivalences follow by
checking the two signs of $J$.
\end{proof}

The theorem supplies a unique \emph{minimal representation} of a signed
mean measure. It does not say that an actual source must choose this
representation. In particular, opposite pathwise counts are already
supported at distinct times in an ordinary jump process with no simultaneous
events, even when $K>0$. Their pathwise mutual singularity is therefore not
the mean-measure minimality appearing in the theorem.

\begin{corollary}[The exact scope of the MPBT rate statement]
\label{stat:mpbt}
Suppose a source has initial occupancy $w(0)$, realizes each expected pair
current, satisfies any equivalent minimal-incidence condition in
Theorem~\ref{stat:jordan}, and is admitted to be a time-inhomogeneous Markov
process with rates depending on the deterministic wave programme and the
present sector. Then its rates on occupied sectors are
$[J_{nm}]_+/w_m$. Conversely, the nonexplosive Bell law in
Lemma~\ref{stat:existence} has those minimal mean incidence measures.
\end{corollary}
\begin{proof}
Pair matching and initial equality imply $\pi=w$. Markovity gives
$F_{nm}=w_m k_{nm}$, while the theorem gives $F_{nm}=[J_{nm}]_+$.
Division is valid for $w_m>0$. The converse follows by taking expectations
of the Bell compensator with marginal $w$.
\end{proof}

Without Markov admission, the same minimal mean measures give only
\[
 \mathbb E\left[1_{\{X_{t-}=m\}}
       \lambda_{nm}(t\mid\mathcal F_{t-})\right]=[J_{nm}(t)]_+,
\]
and Proposition~\ref{stat:renewal} disproves uniqueness of timing.
Without initial equality, the means generated by an assigned Bell rate are
$\pi_m[J_{nm}]_+/w_m$, generally not $[J_{nm}]_+$. These premises cannot be
removed by calling a probability-weighted flux a mass flux.

One equivalent cost formulation has an immediate additional meaning. For a
real sector observable $f$, define its quadratic jump variation by
$[f(X)]_T=\sum_{t\leq T}(f(X_t)-f(X_{t-}))^2$. Under pair matching,
\begin{equation}\label{stat:quadratic-activity}
 \mathbb E\bigl([f(X)]_T\bigr)
 =\int_0^T\sum_{\{n,m\}}
       (|J_{nm}|+2K_{nm})(f_n-f_m)^2\,dt.
\end{equation}
The compensator proves the identity. Requiring this quantity to be minimal
for every $f$ is equivalent to $K=0$ almost everywhere: choose sector
indicator functions to detect any positive surplus on an incident edge.
This gives an independently interpretable fluctuation cost, but its
universal minimization remains mathematically equivalent to the same
minimal-incidence assumption. It is not a weaker theorem forcing it.

\section{What isolation and a one-use reaction resource can select}

The older incidence manuscript contains two narrower physical selection
arguments worth preserving separately from the statistical variational
law \cite{M02}. They operate on an explicitly specified reaction class.

\begin{proposition}[Isolation and stochastic modularity]
\label{stat:isolated-bond}
Suppose a Markov candidate assigns rates $k_{nm}$ to primitive binary
bonds. Assume that each such bond can physically be isolated while holding
its instantaneous wave and Hamiltonian block fixed; that disabling other
bonds leaves the two rates on the retained bond unchanged; that an absent
bond has no transition; and that in each isolated experiment the
coherent-weighted candidate flux $f_{nm}=w_mk_{nm}$ obeys the coherent
continuity equation. Then $f_{nm}-f_{mn}=J_{nm}$ on every positive-weight
bond in the full experiment.
\end{proposition}
\begin{proof}
In the isolated experiment the destination continuity equation has exactly
one intersector term, namely $f_{nm}-f_{mn}$. Its Schr\"odinger value is
$J_{nm}$ for the unchanged wave and block. Stochastic modularity transfers
this equality to the original experiment.
\end{proof}

The conclusion eliminates divergence-free current reassignment, but
$f_{nm}=[J_{nm}]_++K_{nm}$ with $K_{nm}=K_{mn}\geq0$ still survives.
The assumed compatibility of a \emph{coherent-weighted candidate} flux is
not automatically an empirical statement about $\pi_mk_{nm}$; calibrated
initialization and equivariance establish their equality. Likewise,
Hamiltonian locality alone does not imply stochastic modularity or the
physical ability to isolate every joint-configuration bond of an
entangled preparation.

\begin{proposition}[Restricted exact selector with a spent resource]
\label{stat:licensed}
In addition to Proposition~\ref{stat:isolated-bond}, suppose the complete
actual reaction list on a monotone contact consists only of
\[
 (r,\mathrm{live},a)\longrightarrow(i,\mathrm{spent}_i,a\star i),
\]
where $a\star i$ retains an immutable certificate, no inverse reaction
restores the live resource, and $J_{ir}\geq0$ throughout the allowed
window. Then every positive-weight forward rate is
$k_{ir}=J_{ir}/w_r$ and every reverse rate is zero. From calibrated initial
weights these are the actual minimal Bell rates on that window.
\end{proposition}
\begin{proof}
The full reaction list makes $f_{ri}=0$. Proposition~\ref{stat:isolated-bond}
gives $f_{ir}-f_{ri}=J_{ir}$, so $w_rk_{ir}=J_{ir}$. From initial $w$, the
resulting forward equation preserves $w$; existence on the stated finite
wave domain follows from Lemma~\ref{stat:existence} restricted to its
nonzero currents. Thus coherent-weighted and actual mean fluxes agree.
\end{proof}

A nonempty example has one live state $r$, spent states $i$, normalized
coefficients $\sum_i|c_i|^2=1$, and
\[
 H=i\hbar g\sum_i(c_i|i\rangle\langle r|
                  -\overline c_i|r\rangle\langle i|),
 \qquad \Psi_0=|r\rangle.
\]
For $0\leq t\leq\pi/(2g)$,
$w_r=\cos^2(gt)$, $w_i=|c_i|^2\sin^2(gt)$, and
$J_{ir}=2g|c_i|^2\sin(gt)\cos(gt)$. The selected forward rates are
$2g|c_i|^2\tan(gt)$. In a finite window $s<\pi/(2g)$, the null probability
is $\cos^2(gs)$ and the probability of spent record $i$ is
$|c_i|^2\sin^2(gs)$. This exact selector has a concrete irreversible
reaction-resource premise. Continuing the same Hamiltonian past current
reversal contradicts its inverse-free reaction list and pair matching;
the restricted theorem does not extend unchanged.

Nor does directed certificate growth replace bond ownership. On a diamond
$r\to u\to f$ and $r\to v\to f$, take positive candidate flows $j$ on all
four edges. Add $k$ along the $u$ route and subtract $k$ along the $v$ route,
with $0<k<j$. Vertex divergences are unchanged and every edge remains
forward, but the next-route ratio at $r$ becomes $(j+k):(j-k)$. At $k=j/2$
it is $3:1$ instead of $1:1$. Isolation and stochastic modularity are the
assumptions that disallow this independent reassignment; acyclic histories
and spent certificates alone do not.

\section{Projection can create visible countertraffic}

A declared fine process and a grouped readout must not be assigned
independent minimal laws without checking compatibility. Let a partition
$B_1,\ldots,B_r$ group the fine sectors and let $Y_t=a$ when $X_t\in B_a$.
Set
\begin{equation}\label{stat:coarse-data}
 W_a=\sum_{i\in B_a}w_i,\qquad
 J^{\rm c}_{ba}=\sum_{i\in B_a,j\in B_b}J_{ji}.
\end{equation}
The actual projected flow is a sum of directed fine flows, not the positive
part of their signed sum.

\begin{theorem}[Sign alignment and strong lumpability]
\label{stat:coarse}
For a fine Bell process with marginal $w$, the mean directed flow of its
projected readout is
\begin{equation}\label{stat:coarse-flow}
 F^{\rm proj}_{ba}=[J^{\rm c}_{ba}]_++C_{ab},\qquad
 C_{ab}=\frac{\sum_{i\in B_a,j\in B_b}|J_{ji}|
                    -|J^{\rm c}_{ba}|}{2}\geq0.
\end{equation}
The excess vanishes precisely when all nonzero currents across that block
pair have the same orientation. On intervals with positive fine weights,
the generator-level condition
\begin{equation}\label{stat:lumpability}
 \sum_{j\in B_b}\frac{[J_{ji}]_+}{w_i}=\kappa_{ba}(t)
 \quad\text{independently of }i\in B_a
\end{equation}
is sufficient for the projected process to be Markov in its own history.
It is necessary if the same projected transition kernel is required for
every fine starting state in a block at every starting time. Under this
condition,
\begin{equation}\label{stat:coarse-rate}
 \kappa_{ba}=\frac{[J^{\rm c}_{ba}]_++C_{ab}}{W_a}.
\end{equation}
Thus a common projected Bell law requires both this closure and sign
alignment. These are not necessary conditions for all endpoint quantum
measurement statistics.
\end{theorem}
\begin{proof}
Put $A=\sum_{\rm cross}[J_{ji}]_+$ and
$B=\sum_{\rm cross}[-J_{ji}]_+$. Then
$A-B=J^{\rm c}_{ba}$ and $A+B=\sum_{\rm cross}|J_{ji}|$. Since
$A=[A-B]_++\min(A,B)$, \eqref{stat:coarse-flow} follows, with
$C_{ab}=\min(A,B)$. It is zero precisely when positive and negative
cross-block currents do not both occur.

Conditional on the entire projected past, the current fine state has some
posterior supported in its present block. Under
\eqref{stat:lumpability}, averaging the fine exit intensity into another
block gives $\kappa_{ba}$ for every such posterior. The projected counting
compensators are therefore those of the deterministic Markov generator
$\kappa$; uniqueness of its first-jump construction gives the Markov law.
Necessity at the claimed generator-level scope follows by starting in two
fine states $i,i'$ in one block and comparing the first-order probability
of entering $B_b$ during $[t,t+dt]$. Equality of the common projected
kernel requires equality of their block-rate sums. This is strong
lumpability, not a necessary condition for a single exceptional initial
mixture to exhibit Markov projection \cite{F06}.
Multiplying \eqref{stat:lumpability} by $w_i$ and summing over $B_a$ gives
$W_a\kappa_{ba}=F^{\rm proj}_{ba}$, proving \eqref{stat:coarse-rate}.
\end{proof}

The projection theorem is stated locally on positive-weight intervals to
avoid assigning rates to unoccupied fine states. When its identities hold
on all such intervals of the piecewise-constant domain, the node-safe law
of Lemma~\ref{stat:existence} supplies the joins. No freely chosen
post-node source distribution is introduced.

\begin{example}[Minimal microscopic motion with zero coarse current]
\label{stat:parity}
On four cyclic sectors use
$H=i\hbar a(S-S^\dagger)$ and the uniform wave
$(1,1,1,1)/2$. Every fine clockwise current is $a/2$, so every clockwise
rate is $2a$. Group even sites into one block and odd sites into the other.
The opposite fine routes cancel in $J^{\rm c}$, giving $J^{\rm c}=0$, while
$F^{\rm proj}_{10}=F^{\rm proj}_{01}=a$. Each fine state has a single exit
of rate $2a$ into the other block, so the projection is strongly lumpable
and its parity flips at rate $2a$. Therefore
\[
 P(Y_t\ne Y_0)=\frac{1-e^{-4at}}2.
\]
A process constructed anew from the coarse Bell formula with zero current
would never change parity. Both have constant coarse weights $1/2$, but
they are different histories of the same readout. Microscopic minimality
therefore does not mean minimality after every grouping.
\end{example}

\begin{example}[A sufficient hidden-factor realization]\label{stat:factor}
Let $\mathcal H=\mathcal H_C\otimes\mathcal H_E$ with rank-one product
sectors $(c,\eta)$, product wave $\psi\otimes\xi$, and Hamiltonian
$H_C\otimes I+I\otimes H_E$. The factorization is preserved. A fine edge
changing $c$ while keeping $\eta$ fixed has
\[
 J_{(c',\eta),(c,\eta)}=|\xi_\eta|^2J^C_{c'c},\qquad
 w_{c\eta}=|\psi_c|^2|\xi_\eta|^2.
\]
Its Bell rate, on occupied sectors, is $[J^C_{c'c}]_+/|\psi_c|^2$,
independent of $\eta$. Crossedges with both labels changing are absent.
Thus the $c$ readout has aligned fluxes and is strongly lumpable; its
projected law is the coarse Bell law. The conclusion relies on this
factorization and interaction structure. A correlated returning factor
with new couplings need not satisfy either identity.
\end{example}

This is the precise meaning of transfer to a coarse aperture in the
checkpoint argument \cite{C03}. Erasing route labels does not erase their
crossing counts. An integrated current is an expected signed count, not
an outcome probability unless the protocol also guarantees the relevant
single-event and survival conditions.

\section{Finite directional response: a completed counterconstruction}

The checkpoint proposed a finite-response sign gate as a possible physical
implementation of directional exclusion. Its correct partial result can
be stated as a theorem, while retaining the distinction between a flux
ansatz and an actual rate law.

\begin{proposition}[Directional exclusion with lag is not current matching]
\label{stat:gate}
Let $d(t)\in[-1,1]$ and define proposed directional fluxes on a two-sector
cut by
\begin{equation}\label{stat:gate-flux}
 q_f=|J|[d]_+,\qquad q_r=|J|[-d]_+.
\end{equation}
They have no simultaneous opposite traffic, but their signed flux is
$|J|d$. They realize the prescribed current exactly if and only if
$d=\operatorname{sgn}J$ wherever $J\ne0$.
For $t\geq0$, take
\[
 J(t)=-J_0\tanh(t/T_J),\qquad
 \tau\dot d=-1-d,\qquad d(0)=1,
 \qquad J_0,T_J,\tau>0.
\]
On any finite interval on which the probabilities below stay strictly
between zero and one, the completion that prescribes \emph{actual}
fluxes $q_f,q_r$ has an explicit Markov realization with occupation $p_1$
and
\begin{align}
 w_1(t)&=w_*-J_0T_J\log\cosh(t/T_J),\label{stat:gate-wave}\\
 p_1(t)-w_1(t)
    &=2J_0\int_0^t e^{-s/\tau}\tanh(s/T_J)\,ds,
                                                  \label{stat:gate-error}\\
 0<p_1(t)-w_1(t)&\leq2J_0\tau^2/T_J\qquad(t>0).
                                                  \label{stat:gate-bound}
\end{align}
The target current is negative after zero, while the gate selects the old
positive direction until $t=\tau\log2$. This is a finite preparation and
timing counterexample to equating directional exclusion with Bell selection.
\end{proposition}
\begin{proof}
Equation~\eqref{stat:gate-flux} gives
$q_fq_r=0$, $q_f-q_r=|J|d$, and $q_f+q_r=|J||d|$.
The gate solution is $d=-1+2e^{-t/\tau}$. Integrating $\dot w_1=J$ gives
\eqref{stat:gate-wave}. Define
$p_1(t)=w_*+\int_0^t|J(s)|d(s)ds$. Subtraction gives
\eqref{stat:gate-error}. Positivity of the integrand proves strictness,
and $\tanh(s/T_J)\leq s/T_J$ with
$\int_0^\infty s e^{-s/\tau}ds=\tau^2$ proves
\eqref{stat:gate-bound}.

For example, choose the horizon so that
\[
 J_0T_J\log\cosh(T/T_J)<w_*,\qquad
 w_*+2J_0\tau^2/T_J<1.
\]
Then both $w_1$ and $p_1$ remain inside $(0,1)$. Set
\begin{equation}\label{stat:gate-rates}
 k_{10}(t)=\frac{q_f(t)}{1-p_1(t)},\qquad
 k_{01}(t)=\frac{q_r(t)}{p_1(t)}.
\end{equation}
These are bounded on the finite closed interval, and $p$ solves their
master equation with the chosen initial distribution. Uniqueness of the
finite-state forward equation proves that $p$ is the actual occupation,
and hence that the assigned $q$ are its actual directed fluxes.
The underlying target wave current is nonempty as well: take
$\Psi_t=(\sqrt{1-w_1(t)},\sqrt{w_1(t)})$ and
\[
 H(t)=i\hbar\Omega(t)(|1\rangle\langle0|-|0\rangle\langle1|),
 \qquad
 \Omega(t)=\frac{J(t)}{2\sqrt{w_1(t)(1-w_1(t))}}.
\]
Direct differentiation gives Schr\"odinger evolution and current $J$.
This smooth finite programme is used only for the gate counterexample;
it is not an extension of the piecewise-constant selection theorem.
\end{proof}

The rates in \eqref{stat:gate-rates} depend on the prepared ensemble law
through $p$. The proposition is an executable time-dependent statistical
completion for the given preparation, not a derived preparation-independent
local material actuator. If instead one divides the same $q$ by the wave
weights, the actual occupation $\nu$ obeys
\begin{equation}\label{stat:gate-wrong-denominator}
 \dot\nu_1=\frac{\nu_0}{w_0}q_f-\frac{\nu_1}{w_1}q_r,
\end{equation}
and \eqref{stat:gate-error} no longer follows. A correct mean-flux estimate
cannot be transferred across those different completions. Nor does a gate
bound on total traffic alone control signed current: the gate can permit
almost the right amount of traffic in the wrong direction.

The scope of the useful estimate is nevertheless explicit. In the
prescribed-actual-flux completion, its expected total count is
$\int_0^T|J||d|dt\leq\int_0^T|J|dt$, and
$P(N_T\geq1)\leq\min(1,\mathbb E N_T)$. Its occupation discrepancy is
quadratic in response time for this smooth reversal. These are completed
finite calculations. They neither prove the original Bell law at finite
response time nor establish a passive apparatus measuring that native
traffic without changing the experiment.

\section{One chain, two possible constitutive choices}

The variational selection chain can now be given precisely.
For a fixed finite complete source and control domain:
\begin{enumerate}
 \item Supply linear coherent evolution, one actual sector, the declared
       fundamental resolution, and every future-active memory.
 \item Impose expected pairwise Hamiltonian-current realization.
       Supply initial $w$, or impose its control-stable coherent-evacuation
       domain so that Theorem~\ref{stat:calibration} forces it.
 \item Impose the relative-entropy reaction principle with its neutral
       Markov reference and zero-background prescription.
       Theorem~\ref{stat:selection} then selects the conditional Bell path
       law in physical time, including nodes and null intervals.
 \item Admit explicit coherent apparatus contacts in that same complete
       source. Apply their joint law and physical record maps, retaining
       returning factors. Corollary~\ref{stat:output-bound} transfers the
       path estimate to those complete experiments; it does not supply
       universal admission of all possible contacts.
\end{enumerate}
A second internally specified route replaces the entropy premise by an
explicit Markov premise and minimal positive mean incidence. It yields the
same Bell kernel by Corollary~\ref{stat:mpbt}, but it does not explain timing
within a history-dependent class. These are two constitutive presentations,
not successive deductions from increasingly weak old assumptions.

The earlier kinetic construction supplies a different conditional route:
Hamiltonian current production, conservative packet export, scalar pair
chemistry, and a residence/path limit \cite{M01}. Its comparison proof
retains the intrinsic chemistry and full participation hypotheses.
Chapter~\ref{p:chapter-medium} provides their new finite-gas realization
and a finite-recombination comparison. The earlier direct incidence model
already allocated normalized escape responses \cite{M02}; it cannot be
credited with deriving that allocation a second time. The variational law
above can instead replace the kinetic reaction premise as a statistical
constitution, but it does not derive that chemistry from the canonical
action. A full-state record interface requires its own common physical
inventory; equality of mean currents alone never establishes that embedding.

The predictive-current quotient developed in the foundations chapter
identifies the source distinctions needed to determine every future
$w,J$ under the admitted coherent controls. Once a conditional generator
has been selected, that quotient together with $X$ determines its kernel.
The direction of this implication matters: the quotient organizes the
required source information; it does not manufacture the probability law.

\begin{center}
\small
\begin{tabular}{>{\raggedright\arraybackslash}p{0.27\textwidth}>{\raggedright\arraybackslash}p{0.35\textwidth}>{\raggedright\arraybackslash}p{0.29\textwidth}}
\hline
Ingredient & Consolidated consequence & Input retained\\
\hline
Hamiltonian continuity & Antisymmetric currents and weight evolution & Wave,
Hamiltonian, resolution\\
Expected pair-current matching & $\pi-w$ constant; conditional calibration &
Statistical event-current identification\\
Minimal mean incidence & Positive-current numerator & Physical no-surplus
premise; no timing selection alone\\
Path entropy minimization & Finite-background rate and Markov timing &
Neutral reference and extremal statistical law\\
Zero-background limit & Bell paths with \eqref{stat:path-bound} & Limit
prescription, fixed finite domain\\
Physical record projection & Common-output error and coarse-law test &
Admitted joint contacts; sign alignment/closure where claimed\\
\hline
\end{tabular}
\end{center}

The event-law block is mathematically selected within an explicit
statistical constitution. Its entropy principle and expected edge-current
constraint are not consequences of the older information-completion
premises. Chapter~\ref{p:chapter-medium} addresses the same Bell target by
an independently specified microscopic model, and
Chapter~\ref{p:chapter-records} supplies ordinary physical records under
that model's additional coupling rule. This is constitutive internal
completion with controlled effective errors, not a derivation of the
entropy principle or a universal necessity theorem for Bell dynamics.

% End consolidated source: parts/03_statistics.tex

\part{Preparation Consistency and Continuous Record Laws}\label{part:continuous}

% Begin consolidated source: parts/04_cpc.tex
\chapter{Complete preparations and finite causal admission}
\label{cpc:chap:admission}

The results in this part concern an operational question upstream of the finite instrument construction: which aspects of a preparation can affect an actual record together with everything retained after it? Three distinct answers survive the corpus. Complete preparation consistency directly requires descent through a matrix-valued preparation readout. A finite absorbing tag supplies a weaker diagonal identity under a different physical constitution. Calibrated remote experiments, causal separation and uniform stability reconstruct instruments for an initially unknown writer. Their assumptions are stated separately; they are not three successive eliminations of the same premise \cite{M08,M09,M10,M13,M14}.

\section{The complete bank and its proposed readout}

Fix a finite coherent bank $\mathcal H_S$, of dimension $d$, containing every quantum memory that the declared programme can return. Classical data $c\in C$ contain actual clocks, controller states, provenance, allocations and future-active randomizer keys. An inaccessible reference $R$ is retained mathematically and receives no control. A source preparation is an actual probability measure $\nu(dc,d\psi)$ on classical states and normalized rays of the complete coherent bank. Its proposed readout is the positive matrix-valued measure
\begin{equation}
 \pi(\nu)(B)=\int 1_{\{c\in B\}}|\psi\rangle\langle\psi|\,\nu(dc,d\psi).
 \label{cpc:eq:readout}
\end{equation}
Ordinary randomization obeys the law of total probability. Equation~\eqref{cpc:eq:readout} is a definition of a proposed readout, not a proof that it is sufficient. In particular, different ensembles with the same barycenter may still be distinct source preparations.

An unknown writer produces an actual finite mark $a$ and a normalized retained ray $\xi$, including unsuccessful, null and exhausted branches. At a fixed complete classical input define
\begin{equation}
 A_a(P)=\mathbb E_P[1_{\{a\}}|\xi\rangle\langle\xi|],\qquad
 P=|\psi\rangle\langle\psi|.
 \label{cpc:eq:unknownwriter}
\end{equation}
These are positive matrices with $\sum_a\operatorname{tr}A_a(P)=1$. Their dependence on $P$ is initially unrestricted. No instrument, ensemble affinity or completely positive extension is implied by writing an algebraic average of actual daughters.

\begin{definition}[Complete preparation consistency]
\label{cpc:def:cpc}
For a specified apparatus and preparation domain, CPC requires that every finite declared event and its full unnormalized matrix continuation depend on $\nu$ only through~\eqref{cpc:eq:readout}. Reference-compatible CPC additionally requires local event maps to extend as $\mathcal T_E\otimes\operatorname{id}_R$, positively on every admitted joint input and finite reference.
\end{definition}

A preparation key which can later return belongs in $c$ or the coherent bank. Thus comparing a keyed eigenstate lottery with an unkeyed coherent preparation is not an application of CPC. Conversely, a hidden coordinate cannot be omitted merely because it is not displayed at the present cut.

\begin{example}[A marginally invisible returning bit]
Let $C$ be a fair classical sign and $H$ a second sign. Preparations $H=C$ and $H=-C$ have the same separate marginals and the same fixed quantum ray. A later admitted writer $Z=CH$ distinguishes them perfectly. Keeping the joint distribution of $(C,H)$ repairs this particular omission; keeping only the two marginals does not. A finite visible bank is complete only relative to a proved interaction or causal-domain restriction.
\end{example}

\section{A reversible algebra law and its exact boundary}

The corpus supplies a nonempty reversible constitution on the trivial algebra bundle
$\mathfrak A=C^\infty(C,M_d(\mathbb C))$. The matrix product, classical center and preparation independence of transport are physical assumptions. They permit a useful exact statement \cite{M10}.

\begin{theorem}[Complete reversible transport]
\label{cpc:thm:algebra}
Every regular preparation-independent $*$-derivation of $\mathfrak A$ has the form
\begin{equation}
 \mathcal D X(c)=v(c)\cdot\nabla X(c)+i[H(c),X(c)],\qquad H(c)=H(c)^\dagger.
 \label{cpc:eq:derivation}
\end{equation}
Here $v$ is a classical vector field; $H$ can be chosen trace zero. In particular a reversible classical velocity cannot depend independently on the unknown ray.
\end{theorem}
\begin{proof}
For every scalar $f$ and matrix section $X$, $[fI,X]=0$. Applying the derivation rule gives $[\mathcal D(fI),X]=0$, so $\mathcal D$ preserves the center. A regular derivation of smooth scalar functions is a vector field $v$. Subtracting $v\cdot\nabla$ leaves a $C^\infty(C)$-linear derivation on each full matrix fiber. Such derivations are inner. Adjoint preservation makes the inner generator $iH$ with $H$ Hermitian; the scalar ambiguity is removed by choosing trace zero. The trivial bundle and regularity make this choice a regular section.
\end{proof}

For a concrete realization solve $\dot\chi_t(c)=v(\chi_t(c))$ and $i\dot U_t(c)=H(\chi_t(c))U_t(c)$. Then
$\alpha_tX(c)=U_t(c)^\dagger X(\chi_t(c))U_t(c)$
preserves products and adjoints. Moving clocks and classically controlled noncommuting gates are allowed.

This excludes a precise competing interaction. A reciprocal hybrid Hamiltonian
$h(\psi,\theta,p)=g\sin p\langle\psi|A|\psi\rangle$
gives
\begin{equation}
 i\dot\psi=g\sin p\,A\psi,\qquad
 \dot\theta=g\cos p\langle A\rangle_\psi,\qquad \dot p=0.
 \label{cpc:eq:hybrid}
\end{equation}
At $p=0$ it writes an expectation into a classical clock while leaving the ray unchanged. Its operator lift would require $\mathcal D(\sin\theta I)=gA$ at $\theta=p=0$, contradicting center preservation when $A$ is nonscalar. Using the scalar $\langle A\rangle$ instead makes transport preparation dependent. The rival is well defined in a different hybrid constitution. The theorem excludes it by a specific algebraic law, not by asserting that it would spoil CPC. Nor does the theorem select an irreversible stochastic writer.

\section{Finite remote calibration of an unknown writer}

We now develop the causal route without assuming that~\eqref{cpc:eq:readout} is sufficient for the unknown writer. The following admissions include an already selected finite native monitor. This route therefore reconstructs additional writers and their compatibility; it is not an independent derivation of the monitor used for calibration \cite{M09,M10}.

For a bounded score in $[0,1]$ of the complete local writer and its later diagnostic, write $F_R(\Psi;\zeta)$ for its actual expectation on a joint ray. The exterior data $\zeta$ retain every relevant preparation and apparatus history. On a detached product input write $f(\psi)$ for the local expectation. Impose:
\begin{enumerate}
\item The coherent preparations and remote gates below are physically admitted with the same complete local clock, resources and programme. The finite reference monitor has its stated actual likelihood and continuation.
\item The two remote settings, with their outputs unread locally and no transmitted or returning influence during the test, change the unconditioned local score by at most $\nu$.
\item Uniform stable separation holds:
\begin{equation}
 |F_R(\Psi;\zeta)-f(\psi)|\le L\sqrt{1-|\langle\Psi,\psi\otimes r\rangle|^2}
 \label{cpc:eq:stability}
\end{equation}
for every relevant product comparison, including the actual conditional exterior histories and resource preparations. A common modulus can replace $Lu$.
\item For the exact limiting conclusion, arbitrarily accurate finite reference monitors are available while the compared local physical age and programme stay fixed. A fixed resource ceiling supports only the finite-error conclusion.
\end{enumerate}
The stability assumption allows nonlinear source dynamics. For example, synchronous coupling of globally Lipschitz source SDEs gives
$\mathbb E\|X_t-X'_t\|^2\le e^{(2L_b+L_\sigma^2)t}\|X_0-X'_0\|^2$;
a Lipschitz final score therefore has a uniform continuity bound on a controlled preparation class. Bounds uniform only at each fixed hidden gain are insufficient.

Let $\rho=\sum_i p_i|\psi_i\rangle\langle\psi_i|=\sum_a\lambda_a|e_a\rangle\langle e_a|$, with $\lambda_a>0$, and prepare
$\Omega_\rho=\sum_a\sqrt{\lambda_a}|e_a\rangle|a\rangle$.
The columns
\begin{equation}
 U_{ia}=\frac{\sqrt{p_i}\langle e_a|\psi_i\rangle}{\sqrt{\lambda_a}}
 \label{cpc:eq:ensembleunitary}
\end{equation}
are orthonormal: this follows by evaluating $\rho$ between the eigenvectors. They extend to a unitary on a sufficiently padded reference and give
$(I\otimes U)\Omega_\rho=\sum_i\sqrt{p_i}\psi_i\otimes|i\rangle$.
This is a vector identity and a physical gate target. The unknown writer has not been assigned Born probabilities by this identity.

Monitor the reference observable $B=\sum_i(i-1)\Delta|i\rangle\langle i|$ for time $T$ with gain $k$. The admitted native monitor, whose explicit construction appears in Chapter~\ref{cpc:chap:selection}, supplies densities
\begin{equation}
 g(y)=\sum_i p_i g_i(y),\quad g_i=N(2k(i-1)\Delta T,T),\qquad
 \Psi_y=\sum_i\sqrt{\frac{p_i g_i(y)}{g(y)}}\,\psi_i\otimes|i\rangle.
 \label{cpc:eq:finitecalibration}
\end{equation}
Let $c(y)$ be the nearest-mean classifier and $\epsilon=2\Phi(-k\Delta\sqrt T)$, where $\Phi$ is the standard normal distribution function. Direct Gaussian integration and orthogonality of the reference labels give
\begin{equation}
 \operatorname{TV}(\mathcal L(c(Y)),p)\le\epsilon,\qquad
 \mathbb E\bigl[1-|\langle\Psi_Y,\psi_{c(Y)}\otimes c(Y)\rangle|^2\bigr]\le\epsilon.
 \label{cpc:eq:calibrationerror}
\end{equation}
Indeed the second integrand integrates to the actual classification error; each interior Gaussian contributes two tails, each endpoint one. The mixture index used to bound the first inequality is only a coupling variable, not an additional source selector.

\begin{theorem}[Finite causal affinity]
\label{cpc:thm:affinity}
Two eligible ensembles of the same matrix, compared through the above common preparation, satisfy
\begin{equation}
 \left|\sum_i p_i f(\psi_i)-\sum_j q_j f(\varphi_j)\right|
 \le \nu+2(L\sqrt\epsilon+\epsilon).
 \label{cpc:eq:affinitydefect}
\end{equation}
Under exact causality, uniform stable separation and a finite calibration ladder with $\epsilon\to0$, every such bounded score has the form
$f(\psi)=\langle\psi|E|\psi\rangle$, $0\le E\le I$. A calibration reference of dimension $2d$ suffices for this exact score statement.
\end{theorem}
\begin{proof}
Insert $f(\psi_{c(Y)})$ between the actual remote-prepared score and the ensemble average. The first error is at most $L\sqrt\epsilon$ by~\eqref{cpc:eq:stability},~\eqref{cpc:eq:calibrationerror} and Jensen's inequality. The classifier-distribution error costs at most $\epsilon$ because $0\le f\le1$. Compare the two remote gates on the same purification and use the causal discrepancy $\nu$. This proves~\eqref{cpc:eq:affinitydefect}.

For the limit, define $F(\rho)$ by any spectral ensemble, with at most $d$ members. The vanishing comparison defect makes the definition independent of the spectral choice. The union of spectral ensembles for two density matrices has at most $2d$ members and is comparable with a spectral ensemble of their mixture in the same padded reference. Thus $F$ is affine. Finite-dimensional duality represents it as $\operatorname{tr}(E\rho)$; its range supplies $0\le E\le I$.
\end{proof}

For a general separation modulus $\omega$, replace $L\sqrt\epsilon+\epsilon$ by
$\epsilon+\inf_{0<u\le1}\{\omega(u)+\epsilon/u^2\}$, using Markov's inequality in~\eqref{cpc:eq:calibrationerror}. The exact limit is a theorem about a family of finite experiments. It does not claim exact sharp preparation at finite gain or grant an unbounded calibration ladder at fixed cost.

\section{Continuation, references and constructive finite repair}

A finite output probe can recover continuation coordinates without imposing an ideal projection on the unknown daughter. For a projector $Q$ use a native diagnostic with gain $h$ and duration $\tau$. The smooth record score $s_c(y)=\Phi(c(y-h\tau))$, $c>0$, obeys
\begin{equation}
 \mathbb E[s_c(Y_\tau)\mid\xi]=e+v\langle\xi|Q|\xi\rangle,\quad
 e=\Phi\!\left(-\frac{ch\tau}{\sqrt{1+c^2\tau}}\right),\quad v=1-2e>0.
 \label{cpc:eq:deconvolution}
\end{equation}
This is the Gaussian convolution identity for the two eigenrecord laws. It uses no random postprocessing device. Calibration must hold on the actual histories produced by the unknown writer, not merely on an unrelated calibration ensemble.

Apply Theorem~\ref{cpc:thm:affinity} to the mark indicator and mark times probe score. Equation~\eqref{cpc:eq:deconvolution} then recovers each matrix coordinate of $A_a$. For a Polish raw-record space, use bounded continuous functions of the earlier record; finite regular matrix-valued measures are determined by those functions. Discontinuous bins require boundary control when only approximate continuity is available.

\begin{theorem}[Exact operational reconstruction]
\label{cpc:thm:exactcp}
Suppose the exact hypotheses of Theorem~\ref{cpc:thm:affinity} hold for all required complete writer-plus-probe scores, including retained references. Suppose further that pure product inputs have product continuation locality,
$A_a^R(P\otimes Q)=A_a(P)\otimes Q$.
Then there is a normalized CP instrument $\mathcal I_a$ with
$A_a(P)=\mathcal I_a(P)$ and actual reference extension
$A_a^R(P)= (\mathcal I_a\otimes\operatorname{id}_R)(P)$.
\end{theorem}
\begin{proof}
Probe deconvolution makes every coordinate affine. It extends uniquely to a Hermiticity-preserving linear map $\mathcal I_a$, positive by spectral decomposition of a positive input. Actual normalization gives $\sum_a\mathcal I_a^*(I)=I$. Repeating the causal argument on the entire input-plus-reference bank, with a further calibration reference, gives a positive linear map $\mathcal G_{a,R}$ for its actual output. Product locality identifies it with $\mathcal I_a\otimes\operatorname{id}_R$ on product density matrices, which span the Hermitian tensor space. Positivity for a reference of input dimension implies complete positivity. This also identifies the actual extension, rather than merely constructing one possible CP extension.
\end{proof}

The following finite theorem removes the need to claim exact reconstruction from finite-accuracy comparisons. Let the local output dimension be $n$ and the number of marks $m$. Suppose equal-density ensemble comparisons have trace-norm defect at most $D_0$ for each $A_a$ and $D_1$ for each actual $A_a^R$, where $\dim R=d$. Comparisons involving at most $d^2+1$ and $d^4+1$ pure states respectively suffice. Retain pure-product continuation locality, but impose no linear tensor extension on entangled inputs. Set
\begin{align}
 c_j&=1+\sqrt2(j-1),&\eta_0&=c_dD_0,&\eta_1&=c_{d^2}D_1,\nonumber\\
 K_d&=8d-7,&\zeta&=\eta_1+K_d(\eta_0+\eta_1),&
 \tau_0&=m\eta_0+md\zeta.
 \label{cpc:eq:repairconstants}
\end{align}

\begin{theorem}[Finite instrument repair]
\label{cpc:thm:repair}
If $\tau_0<1$, a normalized CP instrument $\{\mathcal I_a\}$ satisfies, uniformly on pure inputs,
\begin{align}
 \sum_a\|A_a(P)-\mathcal I_a(P)\|_1
 &\le\delta_0:=\frac{2\tau_0}{1-\tau_0},\label{cpc:eq:repairlocal}\\
 \sum_a\|A_a^R(P)-(\mathcal I_a\otimes\operatorname{id})(P)\|_1
 &\le\delta_R:=m\zeta+md\zeta+
 \frac{(1+\tau_0)\tau_0}{1-\tau_0}.
 \label{cpc:eq:repairref}
\end{align}
The same inequalities hold after averaging actual preparation distributions. The theorem asserts operational approximation on the stated complete bank, not microscopic linearity.
\end{theorem}
\begin{proof}
Use the $d^2$ projectors onto $|i\rangle$, $(|i\rangle+|j\rangle)/\sqrt2$ and $(|i\rangle+i|j\rangle)/\sqrt2$ as a real Hermitian basis. A pure projector $P$ has a signed expansion $P=\sum_b t_bB_b$, with $\sum_b t_b=1$. Its off-diagonal coefficients are $2\operatorname{Re}P_{ij}$ and $-2\operatorname{Im}P_{ij}$. Their total absolute sum is at most $\sqrt2(d-1)$ because $\sum_{i<j}|P_{ij}|\le(d-1)/2$. The diagonal corrections add no more than the same amount to the original diagonal sum one. Hence
$\sum_b|t_b|\le1+2\sqrt2(d-1)$, and the positive coefficient sum is at most $c_d$.

Interpolate a Hermiticity-preserving linear $L_a$ by $L_a(B_b)=A_a(B_b)$. Moving the negative coefficients to the other side of the projector identity gives equal-density convex ensembles after division by the common mass, at most $c_d$. Therefore $\|A_a(P)-L_a(P)\|_1\le\eta_0$. In dimension $d^2$ obtain a linear $G_a$ within $\eta_1$ of $A_a^R$. Product locality gives
\[
 \|[G_a-L_a\otimes\operatorname{id}](P\otimes Q)\|_1\le\eta_0+\eta_1.
\]
A Schmidt projector has a signed expansion into product pure projectors with coefficient absolute sum at most $K_d$. To verify the stated conservative constant, write its off-diagonal pair as
$\sqrt{\lambda_i\lambda_j}(X_{ij}\otimes X_{ij}-Y_{ij}\otimes Y_{ij})/2$.
Each Hermitian $X_{ij}$ or $Y_{ij}$ has a pure-projector expansion with absolute coefficient sum at most four. Including diagonal products gives
$1+16\sum_{i<j}\sqrt{\lambda_i\lambda_j}\le1+8(d-1)=K_d$.
The Schmidt bases on the two factors may differ. Thus $L_a\otimes\operatorname{id}$ is within $\zeta$ of the actual positive output on every pure joint input.

Let $|\Omega\rangle=d^{-1/2}\sum_i|ii\rangle$ and $J_a=d(L_a\otimes\operatorname{id})(|\Omega\rangle\langle\Omega|)$. Its distance from a positive matrix is at most $d\zeta$, so $\operatorname{tr}J_a^-\le d\zeta$. Replace $J_a$ by its positive part, obtaining a CP map $L_a^+$. The correction is CP and has diamond norm at most $\operatorname{tr}J_a^-\le d\zeta$: the diamond norm of a CP map is $\|\Lambda^*(I)\|$, bounded by its Choi trace.

Set $T=\sum_a(L_a^+)^*(I)$. Actual normalization and the local interpolation bound give $\|T-I\|\le m\eta_0+md\zeta=\tau_0<1$. Thus
\[
 \mathcal I_a(X)=L_a^+(T^{-1/2}XT^{-1/2})
\]
is defined, CP and normalized. For $S=T^{-1/2}$,
\[
 \|S(\cdot)S-\operatorname{id}\|_\diamond
 \le\|S-I\|(\|S\|+1)\le\frac{\tau_0}{1-\tau_0}.
\]
The direct-sum map $(L_a^+)_a$ has diamond norm $\|T\|\le1+\tau_0$. Adding interpolation, Choi correction and normalization yields
$\tau_0+(1+\tau_0)\tau_0/(1-\tau_0)=2\tau_0/(1-\tau_0)$ locally. On the reference bank the initial interpolation term is $m\zeta$, giving~\eqref{cpc:eq:repairref}.
\end{proof}

Finite probes make the defect assumptions quantitative. If every relevant causal score has defect $\Delta=\nu+2(L\sqrt\epsilon+\epsilon)$, then~\eqref{cpc:eq:deconvolution}, applied also to the mark probability, bounds each rank-one matrix coordinate by $(1+e)\Delta/v$. The Hermitian operator norm is the supremum over such coordinates, so
\begin{equation}
 D_0\le\frac{n(1+e)}{v}\Delta_0,\qquad
 D_1\le\frac{nd(1+e_R)}{v_R}\Delta_1.
 \label{cpc:eq:finiteprobeerrors}
\end{equation}
These are uniform constitutive bounds over the admitted projector family. Finitely many observed frequencies do not certify them for an arbitrary unknown nonlinear writer.

\begin{corollary}[Finite adaptive record trees]
\label{cpc:cor:tree}
At each cut of a finite $N$-stage programme retain its complete bank and suppose~\eqref{cpc:eq:repairref} holds uniformly over the actual conditional preparation class, with error $\delta_{R,j}$. Define the comparison CP suffixes on every branch, including branches with zero actual probability. Then the declared finite record trees obey
\begin{equation}
 \operatorname{TV}(P_{\rm actual},P_{\rm CP})\le\frac12\sum_{j=1}^N\delta_{R,j}.
 \label{cpc:eq:treeerror}
\end{equation}
\end{corollary}
\begin{proof}
Replace the last writer first and continue backwards. At each replacement the prefix remains an actual source experiment; the suffix is a normalized CP programme and hence a contractive effect on the common retained bank. The uniform complete output trace-norm error changes every subsequent event probability by at most $\delta_{R,j}/2$. Summing the replacement errors proves the claim. No contractivity of an unknown nonlinear suffix is used.
\end{proof}

This is a finite record-tree theorem, not continuous-path total variation from finitely many matrix fits. For a declared event of actual probability $p>0$, conditioning a joint-law error $\delta$ can cost up to $2\delta/p$ when the comparison probability is positive. Rare branches require their own control.

Two tests protect the assumptions. First, transpose is positive locally but its tensor extension sends a Bell projector to a matrix negative on the antisymmetric ray; product locality and complete-reference affinity cannot be replaced by marginal no-signalling. Second, a rank-sensitive response which agrees with the required law for every nonproduct input but changes discontinuously at products defeats a finite calibration ladder without uniform stable separation. Neither test licenses imposing causality separately at inaccessible hidden values.

\chapter{Selection of a shared-current reader}
\label{cpc:chap:selection}

The statistical target in this chapter is a held Gaussian diagnostic, not the Hamiltonian Bell incidence process. Within that diagnostic class, the population drift, population diffusion and measured signal can be selected together. Gaussian noise, the shared-noise architecture and phase lifting retain their own physical content \cite{M13,M14,Filtering}.

\section{The broader trial class and the exact diagonal identity}

Fix orthogonal nonzero projectors $P_1,\ldots,P_n$ summing to $I$, and let $p_j=\|P_j\psi\|^2$. A normalized continuous source-ray process has, in its declared complete filtration,
\begin{equation}
 dp_j=b_j(\psi)\,dt+A_j(\psi)\,dW,\qquad
 dY=g(\psi)\,dt+dW,\quad Y_0=0.
 \label{cpc:eq:trial}
\end{equation}
The Wiener process is primitive. Coefficients are bounded and continuous, with local regularity sufficient for the indicated It\^o equations and initial moment derivatives. The simplex is preserved, so $\sum_j b_j=\sum_j A_j=0$. Every ray wholly within $P_j\mathcal H$ stays there and has the same calibrated record $Y_t=\beta_jt+W_t$, independently of the ray's direction inside that sector. Coefficients may initially depend on phases and other represented source data. There are no additional noises or jumps in this trial class.

\begin{lemma}[Finite diagonal support from CPC]
\label{cpc:lem:diagonal}
If the trial law obeys CPC, then for every finite-time path event $E$,
\begin{equation}
 \mathbb E_\psi[1_Ep_j(t)]=p_j(0)\mu_{j,t}(E),
 \label{cpc:eq:diagonal}
\end{equation}
where $\mu_{j,t}$ is Wiener path law with constant drift $\beta_j$.
\end{lemma}
\begin{proof}
CPC and actual randomization make
$F_{E,j}(\rho)=\operatorname{tr}[P_j\mathcal T_t(E;\rho)]$
a positive affine functional. Extend it homogeneously to positive matrices. Finite-dimensional duality gives $F_{E,j}(\rho)=\operatorname{tr}(H_{E,j}\rho)$ with $H_{E,j}\ge0$. Any vector in a different sector has zero value. Positivity implies that $H_{E,j}^{1/2}$ annihilates every such vector, so $H_{E,j}=P_jH_{E,j}P_j$. On every unit vector in sector $j$, calibration gives the value $\mu_{j,t}(E)$. Polarization then gives $H_{E,j}=\mu_{j,t}(E)P_j$, proving the identity. No measurement of $P_j$ is performed: it is a coordinate of the assumed continuation measure.
\end{proof}

\begin{theorem}[Population and signal selection]
\label{cpc:thm:selection}
For the regular calibrated shared-current class, identity~\eqref{cpc:eq:diagonal} forces
\begin{equation}
 b_j=0,\qquad g=\overline\beta:=\sum_k\beta_kp_k,\qquad
 A_j=(\beta_j-\overline\beta)p_j.
 \label{cpc:eq:selected}
\end{equation}
Only the mass and endpoint-moment consequences of the identity to first order are needed.
\end{theorem}
\begin{proof}
Integrating~\eqref{cpc:eq:diagonal} against one and against the path endpoint gives
$\mathbb E p_j(t)=p_j(0)$ and
$\mathbb E[Y_tp_j(t)]=\beta_jp_j(0)t$.
The bounded signal gives integrability; clipped endpoint tests justify the latter identity. Initial differentiation of the first gives $b_j=0$. It\^o's rule gives
\[
 d(Yp_j)=(Yb_j+gp_j+A_j)\,dt+(YA_j+p_j)\,dW.
\]
At $Y_0=0$ the second derivative identity therefore gives $gp_j+A_j=\beta_jp_j$. Sum over $j$ to obtain $g=\overline\beta$, and substitute. All initial rays are eligible; continuity includes boundary points.
\end{proof}

\section{A finite physical replacement for full CPC}

A tag need only establish~\eqref{cpc:eq:diagonal}; it need not establish all matrix-valued ensemble equivalences first. State the replacement independently. Attach a ready qubit by the admitted coherent isometry
\begin{equation}
 C_j\psi=P_j\psi\otimes|1\rangle+(I-P_j)\psi\otimes|0\rangle.
 \label{cpc:eq:tag}
\end{equation}
It is realized by the unitary $P_j\otimes X+(I-P_j)\otimes I$ on a $|0\rangle$ tag. Assume the target dynamics preserves $\operatorname{Ran}C_j$ and has exactly the same population/current law after tagging, including on entangled references. This tagging invariance is stronger than ordinary locality.

A physically specified tag reader has, by deadline $s$, a common factor $r_s>0$ such that tag-$1$ click probability is $r_s\|Q_1\Psi\|^2$ and that click captures the source into the corresponding sector. All opposite clicks, nulls and times remain actual outcomes. Assume disjoint-record interchange: for the same two isolated physical ports, clocks and preparations, the two permitted evaluation orders give the same joint law of their actual classical records. This is not an assertion that arbitrary sequential laboratory operations commute.

\begin{theorem}[Finite-tag diagonal bridge]
\label{cpc:thm:tagbridge}
The calibrated trial class, physical tagging invariance, the finite tag law and disjoint-record interchange imply~\eqref{cpc:eq:diagonal}, including an arbitrary inaccessible entangled reference. Exact phase-faithful tag capture is not needed for this scalar implication; capture into the indicated sector suffices.
\end{theorem}
\begin{proof}
Let $K$ denote tag-$1$ click by $s$. Evaluate the target first. The tag load at its later evaluation is $p_j(t)$ by copied-subspace preservation and tagging invariance, so
$\Pr_{A\to B}(E,K)=r_s\mathbb E[1_Ep_j(t)]$.
Evaluate the tag first. Event $K$ has probability $r_sp_j(0)$ and leaves a sector-$j$ input. Eigen-calibration gives
$\Pr_{B\to A}(E,K)=r_sp_j(0)\mu_{j,t}(E)$.
Interchange equates these probabilities and $r_s>0$ permits cancellation. Neither experiment discards its null or opposite-click branches.
\end{proof}

One nonempty tag constitution uses fresh independent thresholds $Z_j\sim\operatorname{Exp}(1)$, quadratic coupling loads $e_j=\|Q_j\Psi\|^2$, a held unchanged null vector and countdowns $\dot R_j=-h(e_j)$. Load subdivision at fixed receiver sensitivity equates
$e^{-t h(e_1+e_2)}$ with $e^{-t[h(e_1)+h(e_2)]}$; continuity and $h(0)=0$ force $h(e)=\gamma e$. The first-zero race then has
\begin{equation}
 \Pr(T\in du,J=j)=\gamma e^{-\gamma u}e_j\,du,\qquad
 \Pr(T>s)=e^{-\gamma s},\qquad r_s=1-e^{-\gamma s}.
 \label{cpc:eq:tagrace}
\end{equation}
This derives the scalar race from that kinetic constitution. Quadratic load, fresh exponential readiness, first-zero actualization and capture remain physical assumptions. Faithful capture gives multipliers
$N_j(u)=\sqrt\gamma e^{-\gamma u/2}Q_j$ and
$N_\varnothing=e^{-\gamma s/2}I$; these describe a sharp absorber, not a bounded finite diffusive reader.

The interface must be explicit. At a winner $j$ and time $u$, the change of variables $Z_j=\gamma e_ju$, $Z_k=R_k+\gamma e_ku$ gives density
\[
 \gamma e_je^{-\gamma u}\,du\prod_{k\ne j}e^{-R_k}\,dR_k.
\]
Thus present stopped residuals factor independently, conditional on the hit and time. At a null the analogous density is $e^{-\gamma s}\prod_ke^{-R_k}dR_k$. But storing the original winning threshold exposes $e_j=Z_j/(\gamma T)$. Disarmed residual exposure followed by load-sensitive reuse also fails the same preparation readout. Threshold retirement and prohibition of live numerical copying are therefore assumptions of this particular realization, not consequences of memorylessness. Later detector constitutions address that separate access problem; the diagonal bridge does not resolve it.

\section{Quantitative tags and the cost of small success probability}

Suppose $|h(x)-\gamma x|\le Cx^{1+\alpha}$ uniformly on $[0,1]$, $\alpha>0$. Split each sector into $m$ identical cells of load $e_j/m$, keeping the coherent microcell register, winning cell and stopped resources in both actual and comparison outputs. Then
\begin{equation}
 \sum_j|m h(e_j/m)-\gamma e_j|\le C m^{-\alpha}.
 \label{cpc:eq:refinerates}
\end{equation}

\begin{proposition}[Complete refined-tag estimate]
\label{cpc:prop:tagerror}
For the held vector, common capture and retirement rules, the complete refined tag differs from its extensive comparator by at most
\begin{equation}
 \eta_m(s)=\min\left\{1,\frac{C}{\gamma m^\alpha}(1-e^{-\gamma s})\right\}
 \label{cpc:eq:tagerror}
\end{equation}
in half trace distance on the record and retained bank, uniformly on input and reference. If each actual joint ordering changes by at most $\eta_m(s)$ on replacing its tag, and their record-interchange defect is $\chi_{t,s}$, then
\begin{equation}
 \sup_E\left|\mathbb E[1_Ep_j(t)]-p_j(0)\mu_{j,t}(E)\right|
 \le d_j(t):=\frac{\chi_{t,s}+2\eta_m(s)}{r_s}.
 \label{cpc:eq:diagdefect}
\end{equation}
\end{proposition}
\begin{proof}
Couple corresponding microcell clocks at the lesser of their rates, with separate excess clocks. The total discrepancy rate is at most $Cm^{-\alpha}$; the union rate is at least the comparator total $\gamma$. Thus the probability that the first union event by $s$ is discrepant is at most $Cm^{-\alpha}\int_0^se^{-\gamma u}du$. On all other branches the actual time, winning cell, captured or null vector agree. Couple stopped residuals by their identical exponential conditional kernels. This is a coincidence coupling of complete outputs; no contraction of an unknown nonlinear trial suffix is used. The triangle inequality between two orderings and the exact-tag comparison gives~\eqref{cpc:eq:diagdefect} after division by $r_s$.
\end{proof}

A small finite-time event discrepancy alone does not determine a generator. The next estimate states the necessary uniformity explicitly. Let $x$ denote complete source coordinates, $H_j(x)=g(x)p_j(x)+A_j(x)$, and suppose $|b_j|\le B$, $|g|,|\beta_j|\le G$, $\mathbb E d(x_u,x_0)\le K\sqrt u$, with Lipschitz constants $L_b,L_H$ for $b_j,H_j$. For $a\ge0$ put $M=Gt+a\sqrt t$ and let $\phi$ be the standard normal density. Then
\begin{align}
 |b_j(x_0)|&\le\frac{d_j(t)}t+\frac23L_bK\sqrt t,\label{cpc:eq:approxdrift}\\
 |H_j(x_0)-\beta_jp_j(x_0)|&\le z_j(t),\nonumber\\
 z_j(t)&=\frac{2Md_j(t)+4\sqrt t\phi(a)}t+
 \frac23(B+L_HK)\sqrt t+\frac{BG}2t.
 \label{cpc:eq:approxcross}
\end{align}
To prove the first, integrate $b_j(x_u)$, subtract $tb_j(x_0)$ and use the all-event mass defect. For the second, the signed endpoint measure has total variation norm at most $2d_j(t)$, so clipping at $M$ costs $2Md_j(t)$. In both compared laws $|Y_t|\le Gt+|W_t|$, giving two combined tail errors at most $4\sqrt t\phi(a)$. Integrate $\mathbb E[Y_ub_j(x_u)+H_j(x_u)]$, using $\mathbb E|Y_u|\le\sqrt u+Gu$ and the Lipschitz estimate. This proves~\eqref{cpc:eq:approxcross}. Summing gives
$|g-\overline\beta|\le\sum_jz_j$ and
$|A_j-(\beta_j-\overline\beta)p_j|\le z_j+p_j\sum_kz_k$.
For a uniform defect $d_j\le\epsilon$, choose $t=\epsilon^{2/3}$ and $a=\sqrt{4\log(1/\epsilon)}$; every displayed error vanishes if the regularity constants stay fixed. Coherent square-root lifting near a node requires additional weighted control.

For comparison, let $p_t^*$ be the selected binary process and define
$p_t=\operatorname{logistic}(\operatorname{logit}p_t^*+a\sin(2\pi t/T))$,
with continuous endpoint extension. Keep $Y$ unchanged. At $0,T$ the transformation is identity, so the complete final ray and current path agree with the selected model. Nevertheless the initial population drift is $2\pi a p(1-p)/T$. This regular clock-dependent rival is distinguished by an earlier stop and defeats inference from one deadline alone.

\section{Positive lift, physical likelihood and generated closure}

Impose positive sector lifting: each $P_j\psi_t$ is a positive scalar multiple of $P_j\psi_0$, with its internal vector fixed. For $L=\tfrac12\sum_j\beta_jP_j$ and $\ell=\langle L\rangle_\psi$, applying It\^o's formula to $\sqrt{p_j}$ in~\eqref{cpc:eq:selected} gives
\begin{equation}
 d\psi=-\tfrac12(L-\ell)^2\psi\,dt+(L-\ell)\psi\,dW,
 \qquad dY=2\ell\,dt+dW.
 \label{cpc:eq:native}
\end{equation}
Smooth coefficients on the compact unit sphere, and direct norm preservation, give a unique global strong process after a local Lipschitz extension around the sphere. This is a nonempty source realization with one actual Wiener-driven record.

\begin{theorem}[Likelihood and complete continuation]
\label{cpc:thm:likelihood}
For a held reader define
\begin{equation}
 M_t(y)=\sum_j\exp\left(\frac{\beta_jy}{2}-\frac{\beta_j^2t}{4}\right)P_j.
 \label{cpc:eq:multiplier}
\end{equation}
The physical process~\eqref{cpc:eq:native} has
\begin{equation}
 \Pr_\psi(dY)=\|M_t(Y_t)\psi\|^2\mathbb Q(dY),\qquad
 \psi_t=\frac{M_t(Y_t)\psi}{\|M_t(Y_t)\psi\|},
 \label{cpc:eq:likelihood}
\end{equation}
where $\mathbb Q$ is reference Wiener measure. Hence the finite event maps
$\mathcal I_t(E)(\rho)=\int_E M_t\rho M_t^\dagger\,d\mathbb Q$
are normalized CP maps with the displayed actual continuation and arbitrary inaccessible references.
\end{theorem}
\begin{proof}
Under $\mathbb Q$, the explicit held multiplier solves $dM=-L^2M\,dt/2+LM\,dY$. Gaussian integration gives
\[
 \mathbb E_{\mathbb Q}M_t^\dagger M_t
 =\sum_j\mathbb E_{\mathbb Q}e^{\beta_jY_t-\beta_j^2t/2}P_j=I.
\]
The bounded-coefficient likelihood argument proved in Theorem~\ref{det:likelihood} identifies this normalized multiplier law with~\eqref{cpc:eq:native}, including its reference extension. Here its held diagonal specialization supplies the explicit Gaussian solution; that later theorem treats the general Hamiltonian and predictable-control propagator. Positivity and complete positivity follow directly from the displayed event-map integral.
\end{proof}

The change of measure does not select its own physical measure: selection occurred through CPC or the tag bridge, with the remaining shared-Gaussian and positive-lift premises. The Gaussian-mixture representation of the final pointer density introduces no ontic eigenlabel. Installing such a label in the physical filtration could change the Wiener property required by~\eqref{cpc:eq:trial}.

If the calibrations are distinct, each bounded martingale $p_j$ converges and the limit is a vertex with probability $p_j(0)$. Indeed, for $V=\sum_j\beta_j^2p_j-\overline\beta^2$, It\^o gives $d\overline\beta=VdW$ and $dV=\mu_3dW-V^2dt$. Thus $\mathbb E\int_0^\infty V^2dt\le V(0)$; convergence of $p$ forces $V\to0$. The limiting support contains one calibration, and bounded convergence gives its probability. Equal calibrations remain an unresolved coherent sector. Every initially positive population remains positive at finite time because~\eqref{cpc:eq:multiplier} is invertible.

For a finite sequence of held readers, admitted coherent gates, fresh ready resources and record-dependent controls, multiply the full retained-bank factors. The joint density is $\|K_\omega\Psi\|^2$ against input-independent reference kernels; normalized primitive laws give normalization by successive integration. Tensor identities retain inaccessible references. In the tag realization, the disjoint target and tag factors commute, and their stopped resource kernels are input independent conditional on the record. Therefore both evaluation orders have the same density $\|N M\Psi\|^2$, proving that the tag assumptions have a common nonempty realization. This supplies CPC for the generated library. It does not prove universal admission of every extra source port or select simultaneous arbitrary noncommuting feedback SDEs.

\section{Decisive alternatives to overstrong selection claims}

\begin{example}[An unread affine channel with the wrong joint law]
For a qubit let
$dp=c p(1-p)dW$ and $dY=\beta p\,dt+dW$, with positive lift. With $f(p)=\sqrt{p(1-p)}$, $f''=-1/(4f^3)$ and the generator sends $f$ to $-c^2f/8$. Thus the unread channel is CP dephasing with coherence $e^{-c^2t/8}$ for every $c$. The equal eigenstate mixture and equal $|+\rangle,|-\rangle$ mixture both have matrix $I/2$, but the initial derivatives of $\mathbb E[Y_tp_t]$ differ by $(c-\beta)/4$. For a finite tag with success $r_s$, the exact order discrepancy is
\[
 \mathbb E_{A\to B}[Y_t1_K]-\mathbb E_{B\to A}[Y_t1_K]
 =r_s(c-\beta)\int_0^t\mathbb E[p_u(1-p_u)]\,du.
\]
The integral is positive for an interior input. Clipping $Y_t$ at a sufficiently large finite level preserves a nonzero bounded-record distinction. Unread affinity does not replace the joint event-and-continuation identity.
\end{example}

\begin{example}[Terminal Born weights with wrong continuation]
The absorbed Wright--Fisher process $dp=\sqrt{2\kappa p(1-p)}dW$ has terminal probability $p$ and mean absorption time
$-[p\log p+(1-p)\log(1-p)]/\kappa$.
These follow respectively from stopped martingales and the boundary-value equation $\kappa p(1-p)u''=-1$. With fixed phase, however, the generator sends $f(p)=\sqrt{p(1-p)}$ to $-\kappa/(4f)$.
Compare the equal mixture of positive-phase rays with $p=1/4,3/4$ with the mixture of $|+\rangle,|-\rangle$ of probabilities $(2+\sqrt3)/4,(2-\sqrt3)/4$. Both initial off-diagonal entries are $\sqrt3/4$. Their derivatives are $-\kappa/\sqrt3$ and $-\kappa\sqrt3/4$. Hence their unread states have trace distance $\kappa t/(4\sqrt3)+o(t)$. Stopping inside small neighborhoods of the initial interior points justifies this expansion. A later finite noncommuting reader of visibility $1-2e_X>0$ detects the difference. Correct terminal weights do not supply coherent continuation.
\end{example}

\begin{example}[Phase backaction survives CPC]
Put $a_j=\beta_j/2$ and choose arbitrary real $\theta_j$. The multipliers
\begin{equation}
 M_t^\theta(y)=\sum_j e^{a_jy-a_j^2t}
 e^{i\theta_j(y-a_jt)}P_j
 \label{cpc:eq:phasefamily}
\end{equation}
have exactly the same likelihood and populations as~\eqref{cpc:eq:multiplier}, and form normalized CP instruments. Their unread off-diagonal magnitudes acquire the additional factor
$e^{-(\theta_j-\theta_k)^2t/2}$, as direct Gaussian integration shows. A later noncommuting probe distinguishes them unless the recorded phase is compensated. Positive lifting selects the phase-free member; CPC and the diagonal bridge alone do not.
\end{example}

\chapter{Physical writing, innovation access and retained information}
\label{cpc:chap:access}

There are two complementary access theorems. The first uses the complete CP instrument already derived or independently admitted and its efficient record. The second uses an explicit linear Gaussian coupling class to calculate the disturbance of a source-blind correlated auxiliary output. Neither theorem establishes universal source admission from source/readout incompleteness \cite{M08,M09,M13}.

\section{Passive refinement of a complete efficient record}

Let $\mathcal I(dy)(\rho)=M_y\rho M_y^\dagger\mathbb Q(dy)$ retain the complete efficient record and every future-active quantum resource. An additional CP port $\mathcal J(dy,dz)$ is passive only if
$\int\mathcal J(dy,dz)=\mathcal I(dy)$
on every input and inaccessible reference. This is equality of the full record-and-continuation instrument, not merely equality of its unread channel.

\begin{theorem}[Efficient passive refinement]
\label{cpc:thm:refinement}
On standard Borel output spaces there is a state-independent classical Markov kernel $k(dz\mid y)$ such that
\begin{equation}
 \mathcal J(dy,dz)(\rho)=k(dz\mid y)M_y\rho M_y^\dagger\mathbb Q(dy).
 \label{cpc:eq:passivekernel}
\end{equation}
\end{theorem}
\begin{proof}
The positive Choi measure of $\mathcal J$ has marginal $|M_y\rangle\!\rangle\langle\!\langle M_y|\mathbb Q(dy)$, using matrix vectorization. Disintegrate its finite scalar trace measure over $y$. The conditional positive matrices average to a rank-one matrix. Every quadratic form orthogonal to its range is nonnegative with integral zero, hence vanishes almost everywhere. Each conditional Choi matrix is therefore a nonnegative scalar multiple of that rank-one matrix. Normalization supplies the kernel. Null $y$ sets can be assigned arbitrarily. Invertibility of $M_y$ is unnecessary.
\end{proof}

A classical copy, coarse-graining or independent randomized processing of $Y$ realizes the kernel. The theorem fails if one first forgets part of the efficient record or traces a returning archive: the resulting instrument can have higher Kraus rank. It also does not apply to a previously written extra memory unless its writing was included in the same instrument.

\begin{corollary}[No information from preserving a whole pure preparation class]
\label{cpc:pure-class}
Let $\mathcal I_a$ be a fixed finite-dimensional CP outcome map on a
subspace $\mathcal K$, retaining that same subspace. If for every unit
$\psi\in\mathcal K$ its output is a nonnegative multiple of
$|\psi\rangle\langle\psi|$, then
\[
 \mathcal I_a(\rho)=p_a\rho\quad\hbox{on }\mathcal K
\]
for an input-independent constant $p_a\geq0$. The same formula holds
after tensoring with an inaccessible reference.
\end{corollary}
\begin{proof}
For any Kraus decomposition, positivity and the one-dimensional output
support imply $A_{a\ell}\psi\in\mathbb C\psi$ for every $\psi$.
Apply this to a basis and to each sum of two basis vectors: every
$A_{a\ell}$ has the same eigenvalue on all basis vectors, hence equals
$c_{a\ell}I$ on $\mathcal K$. Summing gives
$p_a=\sum_\ell|c_{a\ell}|^2$ and proves the reference extension.
\end{proof}
This is the identity-channel case of efficient passive refinement and
preserves the checkpoint's whole-class qualification \cite{C01}.
It assumes CP outcome maps; it does not derive their admission or rule
out fitting a single isolated probability table.

For a held qubit reader $L=\lambda_0P_0+\lambda_1P_1$, put $\delta=|\lambda_1-\lambda_0|>0$. On eigeninput $j$,
$Y_T\sim N(2\lambda_jT,T)$ and the native innovation endpoint is $W_T=Y_T-2\lambda_jT$. The two full path laws are mutually absolutely continuous.

\begin{proposition}[Passive innovation precision]
\label{cpc:prop:precision}
With equal prior probabilities of the two unlabelled eigeninputs, every passive estimator $\widehat W_T$ and tolerance $0\le r<\delta T$ satisfy
\begin{equation}
 \Pr(|\widehat W_T-W_T|>r)\ge\Phi(-\delta\sqrt T).
 \label{cpc:eq:precision}
\end{equation}
The bound is attained by postprocessing $Y_T$.
\end{proposition}
\begin{proof}
For each observed path endpoint $y$, the two possible innovation centers differ by $2\delta T$. Success within radius $r$ implies correct nearest-center classification of the input. The equal-prior Bayes error of the two Gaussians is $\Phi(-\delta\sqrt T)$; the endpoint is sufficient for the constant-drift Brownian path. A midpoint classifier followed by $\widehat W_T=y-2\lambda_{\widehat j}T$ attains the bound: correct decisions have zero error and wrong decisions error $2\delta T$.
\end{proof}

An active supplied eigenlabel changes the complete input and invalidates this comparison. Conversely, a claimed exact innovation output identifies the eigeninput from $(Y_T-W_T)/(2T)$ and must change the full native instrument.

\begin{theorem}[Sharp complete disturbance for exact extra discrimination]
\label{cpc:thm:disturbance}
Suppose a CP enlarged instrument has a classical decoder which perfectly identifies the two eigeninputs. After forgetting the extra output but retaining the original $Y$ and original continuing bank, denote it by $\overline{\mathcal J}$. Then
\begin{equation}
 \frac12\|\overline{\mathcal J}-\mathcal I\|_\diamond
 \ge\frac12e^{-\delta^2T/2}.
 \label{cpc:eq:disturbance}
\end{equation}
This constant is attained in the wider class of CP instruments.
\end{theorem}
\begin{proof}
Perfect eigeninput decoding makes its two effects exactly $P_0,P_1$. Every Kraus operator associated with label $j$ therefore annihilates the opposite eigenvector. Thus $\overline{\mathcal J}$ kills $|0\rangle\langle1|$ and gives identical outputs on $|+\rangle$ and $|-\rangle$. Write $m_j(y)=e^{\lambda_jy-\lambda_j^2T}$. Their native output difference is $m_0m_1\sigma_x\mathbb Q(dy)$ and has trace distance
$\int m_0m_1d\mathbb Q=e^{-\delta^2T/2}=:v_T$.
The triangle inequality forces one error to be at least $v_T/2$.

For attainment use $\mathcal J(dy,j)(\rho)=M_yP_j\rho P_jM_y^\dagger\mathbb Q(dy)$. This is native reading after complete dephasing. On an arbitrary reference input its difference from native reading has trace norm $2v_T\|\rho_{01}\|_1$. Positivity implies $\|\rho_{01}\|_1\le\sqrt{\operatorname{tr}\rho_{00}\operatorname{tr}\rho_{11}}\le1/2$. Hence the half-diamond distance is at most $v_T/2$, with equality on $|+\rangle$.
\end{proof}

The attaining sharp instrument is a benchmark outside bounded finite diffusion. A finite invasive alternative adds an independent commuting native read of exposure $S=\mu^2\tau$. The optimal equal-prior error becomes $\Phi(-\sqrt{\delta^2T+S})$. On $|+\rangle$, its later unread coherence changes a noncommuting plus probability by
\begin{equation}
 \Delta p_X=\frac{e^{-\delta^2T/2}}2(1-e^{-S/2}).
 \label{cpc:eq:finiteinvasive}
\end{equation}
Both formulas follow by adding Gaussian log-likelihood information and multiplying coherence factors. A finite final $X$ reader multiplies the probability gap by its visibility $1-2e_X$. This provides a complete finite separating experiment, not an exact auxiliary projective measurement.


\section{Common amplitude writing selects multivariate response}

A more specific physical writing law links attenuation to the response of several correlated currents. Let commuting Hermitian loads $L_1,\ldots,L_m$ act on the complete finite bank, with joint eigenvalue vectors affinely spanning $\mathbb R^m$. Assume the common Stratonovich amplitude equation
\begin{equation}
 d\phi=\sum_i L_i\phi\circ dX_i-F\phi\,dt,\qquad F=F^\dagger,
 \quad dX=b([\psi])\,dt+dW,\quad d[W_i,W_j]=C_{ij}\,dt,
 \label{cpc:eq:multiwriting}
\end{equation}
where $\psi=\phi/\|\phi\|$, $C\ge0$ is fixed, $F$ is preparation independent and $b$ is finite and continuous on all rays. Impose conditional balance of every held load mean: $\mathbb E[d\langle L_i\rangle\mid\mathcal F_t]=0$. This last condition is statistical constitutive physics, stronger than conservation of an unread ensemble mean. It is not asserted to follow from reversibility or source incompleteness.

\begin{theorem}[Multivariate response and attenuation selection]
\label{cpc:thm:multivariate}
For the stated class there are fixed $c\in\mathbb R^m,d_0\in\mathbb R$ such that
\begin{equation}
 F=L^{\mathsf T}CL+c\cdot L+d_0I,\qquad
 b(\psi)=2C\langle L\rangle_\psi+c.
 \label{cpc:eq:multiselected}
\end{equation}
Conversely these coefficients define a regular norm-preserving finite-dimensional interaction after normalization, including singular $C$ and degenerate joint eigenspaces.
\end{theorem}
\begin{proof}
Set $a_i=\langle L_i\rangle$, $V_{ij}=\langle L_iL_j\rangle-a_ia_j$ and $G=F-L^{\mathsf T}CL$. Applying It\^o's quotient rule to $\langle\phi,L_i\phi\rangle/\|\phi\|^2$ yields
\begin{equation}
 da=2V\,dW+2\{V(b-2Ca)-\operatorname{Cov}(L,G)\}\,dt,
 \label{cpc:eq:balanceidentity}
\end{equation}
where $\operatorname{Cov}(K,G)=\tfrac12\langle KG+GK\rangle-\langle K\rangle\langle G\rangle$. Thus with $\beta=b-2Ca$, balance says $V\beta=\operatorname{Cov}(L,G)$.

Choose unit vectors $u_\lambda,v_\mu$ from different joint eigenspaces, put $\delta=\mu-\lambda$ and $\psi_p=\sqrt{1-p}u_\lambda+e^{i\theta}\sqrt p\,v_\mu$. Let $g_{\lambda\mu}=\langle u_\lambda,Gv_\mu\rangle$. The balance equation becomes
\[
 \delta\cdot\beta(\psi_p)=g_{\mu\mu}-g_{\lambda\lambda}
 +\frac{1-2p}{\sqrt{p(1-p)}}\operatorname{Re}(e^{i\theta}g_{\lambda\mu}).
\]
Bounded continuity as $p\downarrow0$, for every phase, forces $g_{\lambda\mu}=0$. Holding $u_\lambda$ fixed and varying $v_\mu$ in its degenerate eigenspace shows that $G$ has a constant quadratic form on that block, hence is scalar there; reversing the roles covers every block. Write the scalar values $g_\lambda$. At a fixed eigenray $u_0$ let $c=\beta(u_0)$; the same limit gives $g_\mu=c\cdot\mu+d_0$ on every joint eigenvalue. Therefore $G=c\cdot L+d_0I$.

Equation~\eqref{cpc:eq:balanceidentity} now reads $V(\beta-c)=0$. Rays whose occupied spectral points affinely span $\mathbb R^m$ are dense and have positive-definite covariance $V$, so $\beta=c$ there and everywhere by continuity. Removing the fixed drift offset and scalar amplitude gauge gives
\[
 d\psi=\sum_i(L_i-a_i)\psi\,dW_i
 -\frac12\sum_{ij}C_{ij}(L_i-a_i)(L_j-a_j)\psi\,dt.
\]
Its smooth sphere coefficients preserve norm; a square root of $C$ gives a strong global realization. Substitution verifies balance and the converse.
\end{proof}

The affine-span assumption identifies a precise access freedom. For $L=(kA,0)$ the second drift is unconstrained by charge balance; taking correlated covariance $C_{12}=r$ and setting $b_2=0$ is a surviving source-blind tap. The following additional physical loading family removes that freedom:
$L_1=kA\otimes I$, $L_2=\epsilon I\otimes B$ with $A,B$ nonscalar, common covariance, calibrated offsets, and response and attenuation continuous as $\epsilon\to0$. For every $\epsilon\ne0$ the joint spectrum spans a rectangle, so the theorem gives
\begin{align}
 F_\epsilon&=k^2A^2\otimes I+2rk\epsilon A\otimes B+
 \epsilon^2I\otimes B^2,\nonumber\\
 b_1^\epsilon&=2k\langle A\rangle+2r\epsilon\langle B\rangle,
 &b_2^\epsilon&=2rk\langle A\rangle+2\epsilon\langle B\rangle.
 \label{cpc:eq:loadedresponse}
\end{align}
The zero-load limit therefore has $b_2^0=2rk\langle A\rangle$, giving a signal copy rather than a source-blind innovation copy. This is a theorem about a common continuously loadable amplitude port. Merely attaching an arbitrary mechanical displacement sensor does not establish that it belongs to this family.

\begin{proposition}[Finite-load repair and its record error]
\label{cpc:prop:weakload}
In the same two-load amplitude class put $G_\epsilon=F_\epsilon-L^{\mathsf T}CL$. On a product calibrator with charge extrema $\pm h$ in equal superposition assume uniformly
\[
 |\mathcal D_\epsilon(B)|\le\eta_B(\epsilon),\quad
 \inf_g\|G_\epsilon-gI\|\le\delta_F(\epsilon),\quad
 |b_{2,\epsilon}-b_{2,0}|\le\omega_{\rm load}(\epsilon).
\]
Here $\mathcal D_\epsilon(B)$ is the conditional drift of $\langle B\rangle$ and the zero-load auxiliary is an unchanged spectator. Then
\begin{equation}
 |b_{2,0}-2rk\langle A\rangle|
 \le\inf_{0<\epsilon\le\epsilon_{\max}}
 \left\{\omega_{\rm load}(\epsilon)+\frac{\eta_B(\epsilon)}{2\epsilon h^2}
 +\frac{\delta_F(\epsilon)}{\epsilon h}\right\}.
 \label{cpc:eq:weakload}
\end{equation}
If the right side is $D$, the primary source and $Y$ continuation remain exactly unchanged, and $|r|<1$, comparison with the reciprocal copy over time $T$ gives
\begin{equation}
 D_{\rm KL}(P\|P_*)\le\frac{D^2T}{2(1-r^2)},\qquad
 \operatorname{TV}(P,P_*)\le\frac{D\sqrt T}{2\sqrt{1-r^2}}.
 \label{cpc:eq:weakloadpath}
\end{equation}
\end{proposition}
\begin{proof}
On a product input,~\eqref{cpc:eq:balanceidentity} gives
\[
 \mathcal D_\epsilon(B)=2\epsilon\operatorname{Var}(B)
 [b_{2,\epsilon}-2rk\langle A\rangle-2\epsilon\langle B\rangle]
 -2\operatorname{Cov}(B,G_\epsilon).
\]
The centered calibrator has $\langle B\rangle=0$, variance $h^2$ and covariance magnitude at most $h\inf_g\|G_\epsilon-gI\|$. Solve for the bracket and compare the loaded and unloaded drifts. This proves~\eqref{cpc:eq:weakload}. For the path estimate, the common primary continuation leaves only an auxiliary drift difference of covariance norm squared at most $D^2/(1-r^2)$. Bounded-drift Girsanov gives its relative-entropy cost one half of the time integral; Pinsker gives the displayed total-variation bound.
\end{proof}

Exact response follows if the load modulus vanishes and $\eta_B,\delta_F=o(\epsilon)$. If the measured drift is that of $\epsilon B$, the needed error is $o(\epsilon^2)$. At finite error, arbitrarily small load can be worse: for $\omega=L\epsilon^\alpha$ and constant $A_0=\eta_B/(2h^2)+\delta_F/h$, the optimizing load is $(A_0/(\alpha L))^{1/(1+\alpha)}$ when admissible. A smooth counterfamily $b_{2,\epsilon}=q[1-e^{-(\epsilon/\ell)^2}]+2\epsilon\langle B\rangle$, $q=2rk\langle A\rangle$, keeps an order-one zero-load suppression while its balance error is uniformly $O(\ell)$. Its load derivative diverges as $\ell^{-1}$, violating the required common modulus. Thus small unscaled calibration errors cannot replace the stated uniform hypotheses.

\section{Explicit linear Gaussian contacts and the source-blind cost}

Consider commuting source charge $A$ and independent selected Gaussian writing channels with real couplings $\lambda_\alpha A$. Their observed displacement vector is $X=HR$, where $R$ has unit independent noise. Put
\begin{equation}
 C=HH^{\mathsf T},\qquad v=2H\lambda.
 \label{cpc:eq:gaussianclass}
\end{equation}
On charge eigenvalue $a$, $X_T$ has mean $vaT$ and covariance $TC$. The selected source law gives unread coherence-decay rate
$\Gamma_{ab}=\tfrac12(a-b)^2\|\lambda\|^2$; further unread channels or stochastic phase kicks can only add their nonnegative dephasing contribution. The class assumes the common linear Gaussian amplitude-writing law and its stochastic signature. It does not derive these premises from the following optimization.

\begin{theorem}[Gaussian access--disturbance bound]
\label{cpc:thm:gaussianbound}
For the class~\eqref{cpc:eq:gaussianclass}, with $C^+$ the Moore--Penrose inverse,
\begin{equation}
 \Gamma_{ab}\ge\frac{(a-b)^2}{8}v^{\mathsf T}C^+v.
 \label{cpc:eq:gaussianbound}
\end{equation}
The complete held-record Fisher information for parameter $a$ is $T v^{\mathsf T}C^+v$. If $v\in\operatorname{Ran}C$, equality is attained with $\lambda=H^{\mathsf T}C^+v/2$ and no extra unread coupling.
\end{theorem}
\begin{proof}
$H^{\mathsf T}(HH^{\mathsf T})^+H$ is the orthogonal projection onto $\operatorname{Ran}H^{\mathsf T}$. Therefore
$\|\lambda\|^2\ge\lambda^{\mathsf T}H^{\mathsf T}C^+H\lambda=v^{\mathsf T}C^+v/4$,
which gives the bound and equality condition. On the support of the Gaussian covariance the endpoint log-likelihood derivative is $v^{\mathsf T}C^+(X_T-vaT)$; its variance is $Tv^{\mathsf T}C^+v$. The endpoint is sufficient for the held Brownian location family, so this is also the full-path information. Here Fisher information refers to this specified Gaussian location family; it does not posit a passive meter of an arbitrary unknown quantum expectation.
\end{proof}

In particular demand the output pair
\begin{equation}
 dY=2k\langle A\rangle\,dt+dW,\qquad
 dZ=r\,dW+\sqrt{1-r^2}\,dV,
 \quad |r|<1,
 \label{cpc:eq:sourceblind}
\end{equation}
where $V$ is independent and $Z$ is source-blind: it has no signal drift. Then
$C=\left(\begin{smallmatrix}1&r\\r&1\end{smallmatrix}\right)$ and $v=(2k,0)^{\mathsf T}$, so
\begin{equation}
 \Gamma_{ab}^{\min}=\frac{k^2(a-b)^2}{2(1-r^2)}
 =\frac{\Gamma_{ab}^{0}}{1-r^2}.
 \label{cpc:eq:blindcost}
\end{equation}
At $|r|=1$ the required $v$ lies outside $\operatorname{Ran}C$, so no finite coupling in this class realizes it. This penalty does \emph{not} apply to an ordinary signal copy $Z=rY+\sqrt{1-r^2}V$. That copy has $v=2k(1,r)^{\mathsf T}$ and $v^{\mathsf T}C^{-1}v=4k^2$, hence no necessary extra dephasing.

There is an explicit amplitude realization of the bound, which also locates its quantum preparation assumptions. For one finite exposure $\delta$, prepare a two-coordinate pure Gaussian meter
\[
 f_C(x)=(\det(2\pi C))^{-1/4}\exp(-x^{\mathsf T}C^{-1}x/4),\qquad C>0.
\]
On charge eigenvalue $a$, apply the controlled translation $x\mapsto x-\sqrt\delta\,va$ generated by the self-adjoint coupling $\sqrt\delta A\otimes v^{\mathsf T}P$. The retained meter state is $f_C(x-\sqrt\delta va)$. Its density has the required Gaussian displacement and two eigenmeter states have overlap
\begin{equation}
 \langle f_{C,b},f_{C,a}\rangle
 =\exp[-\delta(a-b)^2v^{\mathsf T}C^{-1}v/8].
 \label{cpc:eq:gaussianoverlap}
\end{equation}
Completing the square proves both statements. The interaction is unitary on source plus meter, preserves inaccessible references, and fixes the exported coherence loss. To obtain an actual classical record one must additionally supply the admitted native coordinate diagnostic or a declared configuration-record law; the unitary translation alone does not actualize a value. The Gaussian ready amplitude and fresh supply are also resources, not a derived equilibrium reservoir. Used meter states are retained; reset by SWAP exports their correlations to an archive and consumes a ready state.

The ordering of a copy matters physically. A canonical copying gate $e^{-ibQ P_R}$ sends $R\mapsto R+bQ$ and $P_Q\mapsto P_Q-bP_R$. Copying the source-written position afterward exports its signal and noise together. Copying an independent precursor position before the source interaction retains correlated noise but imparts a conjugate momentum kick which participates in the later source coupling. The complete Gaussian amplitude calculation~\eqref{cpc:eq:gaussianoverlap} quantifies that difference. A late erasure of the displayed copy does not undo a conjugate kick or an earlier exported memory.

\section{Delayed innovation writers and preparation provenance}

The abstract and Gaussian results must withstand a directly printed innovation. Suppose the selected reader is left unchanged but an extra physical writer prints $W_t$. Continue the same reader for another duration $s$, recording $Z=Y_{t+s}-Y_t$. Compare a phase-randomized coherent preparation with fixed interior populations against its eigenstate lottery; their preparation keys must be absent from all future-active resources. Since $\overline\beta$ is a bounded martingale and $d\overline\beta=VdW$,
\begin{equation}
 \mathbb E[Z\mid\mathcal F_t]=s\overline\beta_t,\qquad
 \mathbb E[W_tZ]=s\,\mathbb E\int_0^t V_u\,du>0
 \label{cpc:eq:innovationwitness}
\end{equation}
on the interior coherent preparation, whereas the eigenstate lottery gives zero. The equality follows by It\^o covariance; finite-time positivity makes it strict when distinct calibrations are occupied. The product is integrable, so a sufficiently large finite clipping yields a nonzero bounded-record distinction. Thus the extra output would make the source ensemble observable beyond~\eqref{cpc:eq:readout}. Printing it is a consistent stochastic extension if CPC is abandoned; the theorem identifies its consequence rather than declaring it meaningless.

The same issue can be delayed through a coherent qubit memory. A per-path memory unitary controlled by the mathematical innovation is an expectation-dependent source interaction when rewritten in physical $Y$ coordinates. Replacing its compensating scalar $\langle A\rangle I$ torque by the actual operator $A$ is a lawful common-linear-writing repair, but changes source/memory continuation. One cannot infer admission of the initial writer from the fact that its final diagnostic is an admitted native reader. Every earlier interaction that imprinted the returning memory must satisfy the claimed constitution.

Uniformity at weak load is equally important. A drift defect divided by auxiliary load $\epsilon$ may remain finite even if the unscaled balance defect vanishes. At the elementary level, $\min(\ell h,1)\to0$ for each fixed susceptibility $h$, but the family $h=1/\ell$ retains response one. A resource bound $\mathbb E h\le M$ instead gives the uniform error at most $\ell M$. The corresponding Gaussian auxiliary-load selection needs uniform balance and attenuation errors of order $o(\epsilon)$ in the unscaled convention; pointwise continuity at each preparation does not supply that premise \cite{M08,M09}.

\section{Literal source erasure and the cost of returning keys}

A possible alternative to restricting access is to scramble preparation ensembles physically. The following exact obstruction and attaining construction specify what that proposal costs \cite{M14}.

\begin{theorem}[Source-ensemble erasure tradeoff]
\label{cpc:thm:erasure}
Let $K(\psi,d\varphi)$ be a preparation-independent Markov operation on pure rays. If its output barycenter is $|\psi\rangle\langle\psi|$ for every pure input, then it fixes every ray almost surely. For qubit inputs define
\[
 \begin{gathered}
 \mu_z=(K_0+K_1)/2,\quad\mu_x=(K_++K_-)/2,\quad
 d_*=\operatorname{TV}(\mu_z,\mu_x),\\
 F=\frac14\sum_{v\in\{0,1,+,-\}}\int|\langle v,\varphi\rangle|^2K_v(d\varphi).
\end{gathered}
\]
With $\epsilon_*=(2-\sqrt2)/4$,
\begin{equation}
 F\le1-\epsilon_*+\epsilon_*d_*.
 \label{cpc:eq:erasure}
\end{equation}
Every point on this boundary is attained by a finite random-unitary interaction whose rotation key is inactive.
\end{theorem}
\begin{proof}
For the first statement, the expectation of the squared overlap with every vector orthogonal to $\psi$ is zero. Nonnegativity forces every output onto the ray of $\psi$.

Let $r(\varphi)$ be the output unit Bloch vector. Positivity of the four source measures implies
\[
 F\le\frac12+\frac14\left(\int|r_z|\,d\mu_z+\int|r_x|\,d\mu_x\right).
\]
The common submeasure of $\mu_z,\mu_x$ has mass $1-d_*$. On it $|r_x|+|r_z|\le\sqrt2$; each remaining measure has mass $d_*$ and its coordinate is at most one. Therefore
$F\le1/2+[\sqrt2(1-d_*)+2d_*]/4$, giving~\eqref{cpc:eq:erasure}.

With probability $u$ apply equiprobably $U_\pm=e^{\mp i\pi\sigma_y/8}$ and otherwise apply identity. On the rotation branch both input lotteries become the same four bisector rays; on the identity branch their supports are disjoint. Hence $d_*=1-u$ and $F=1-u\epsilon_*$. The averaged channel is $(1-u\epsilon_*)\operatorname{id}+u\epsilon_*\operatorname{Ad}_{\sigma_y}$; its half-diamond distance from identity is $u\epsilon_*$, attained on an $x$ or $z$ eigeninput and bounded above by convexity.
\end{proof}

The independent classical randomizer is a stated resource. If its key returns, an inverse conditional rotation restores the original source distinctions; the reduced erasure conclusion no longer concerns the complete output. This theorem obstructs literal pure-ray scrambling at zero disturbance. It does not obstruct operational CPC for a source that retains inaccessible distinctions but constrains their future interactions.

\section{The shared dependency boundary}

The exact assumption-to-consequence chain in this part is the following. Either CPC plus calibrated shared-current dynamics, or the independently specified finite tag plus tagging invariance and disjoint-record interchange, yields the diagonal identity. That identity fixes $b,A,g$. Positive lifting then supplies the particular native reader and its complete finite likelihood. Products of the admitted factors yield a reference-compatible library; finite instrument implementation is developed separately. Once a complete efficient instrument has been derived or independently admitted, passive refinements reduce to record postprocessing and invasive access has calculable costs.

The causal reconstruction chain is different: already calibrated finite native probes, physical preparation access, causal separation, stable response and product continuation locality constrain an unknown writer, including its daughters and reference extension. Its finite repair theorem gives a nearby instrument with explicit complete-bank error. It cannot be used to justify the monitor that supplied its own calibration.

None of these statements selects the original Hamiltonian Bell incidence law, primitive Gaussian noise, universal source contact admission or apparatus equilibrium. The mathematical role of Shadow source/readout structure is to identify the lost preparation information and require its fate to be checked through complete records, actual histories and returning banks. The stronger algebra, stochastic, calibration and admission laws remain explicit constitutive commitments.

% End consolidated source: parts/04_cpc.tex

\part{Finite Detectors and Measurement Instruments}\label{part:detectors}

% Begin consolidated source: parts/05_detectors.tex
\chapter{Joint capture and the state retained after an event}
\label{det:chapter-capture}

This chapter consolidates the strongest detector-selection argument of
\cite{M15,M16}. It begins with genuinely unknown stochastic daughters. A
positive-response test of absent joint channels determines their old-bank
component; finite click/null and double-click comparisons then determine a
held detector's scalar response. A further archive test is necessary to
determine the complete retained output. These implications hold within a
specified source constitution. In particular, preservation of all joint
zeros characterizes faithful extraction in the admitted probe class; it is
not a consequence established here from source/readout incompleteness.

Two subsequent chapters give different realizations. The conserved
converter with an actual diffusive reader has a finite, nonsharp conditional
state on every finite record. The finite receptor with an intrinsic latch
has exact selected branches but a nontrivial no-record receptor state.
Neither realization may inherit the other's null dynamics by a change of
notation.

\section{The initially open event kernel}

At a fixed complete classical preparation $c$, write a loaded binary outlet
as
\begin{equation}
 \Phi=|0\rangle_D v_0+|1\rangle_D v_1,
 \qquad v_j\in\mathcal B,
 \qquad p_j=\|v_j\|^2,
 \qquad p_0+p_1=1.                         \label{det:rows}
\end{equation}
The old coherent bank $\mathcal B$ contains the carried system, every
incoming quantum memory that can return, and an arbitrary inaccessible
reference. The classical variable $c$ retains active preparation labels,
clocks, reaction resources and controls. Different values of these data
are different complete preparations. A spent detector produces a blocked
record; its retained bank is unchanged unless a separately specified
replenishment interaction acts.

Before imposing the selection premises, an event in outlet $j$ at time $u$
has an arbitrary normalized kernel
\begin{equation}
 \kappa_{j,u}(\Phi;d\chi,dz),
 \qquad [\chi]\in\mathbb P(\mathcal B\otimes\mathcal M).       \label{det:kernel}
\end{equation}
The new bank $\mathcal M$ and classical archive $z$ remain available for
later experiments. Thus the old-bank daughter can be stochastic, mixed
after marginalization, nonlinear in the input, dephased or rotated.
The source kinematics is a kernel over full rays; branch-dependent finite
banks can be embedded in one declared direct sum. No CP instrument or
rank-one daughter is imposed at this stage.

\begin{assumption}[Held responsive detector class]\label{det:held-assumptions}
On an active binary detector the complete-history intensities are
$\lambda_j=h(p_j)$, where $h:[0,1]\to[0,\infty)$ is continuous,
$h(0)=0$, $h(1)=\gamma>0$, and $h(p)>0$ for $p>0$. A held null leaves the
coherent bank unchanged. The same function applies after the admitted
coherent loading, spectator embedding and previous disjoint event.
There are no additional unlisted gains or history variables in this
response. Intrinsic event randomness and these complete-history
intensities define the process; future simulation seeds are not present
physical coordinates.
\end{assumption}

For positive finite windows, this class gives
\begin{equation}
 \Lambda(p)=h(p)+h(1-p),\qquad
 F_t(p)=h(p)\int_0^t e^{-\Lambda(p)u}\,du,
 \qquad N_t(p)=e^{-\Lambda(p)t}.                \label{det:held-laws}
\end{equation}
$F_t(p)$ is an actual click probability, not an auxiliary Born
measurement. In particular $F_t(p)>0$ precisely when $p>0$.
The population-only response and frozen null are important restrictions;
later drive- or history-dependent extensions require additional premises.

\begin{assumption}[Joint zero and record-order tests]\label{det:zero-assumptions}
An admitted coherent probe $P$ on the old bank has the same responsive
binary response when its population is $\|(P\otimes I)\chi\|^2$.
If $Pv_j=0$ before separation, the disjoint two-port experiment cannot
produce both the original record $j$ and the probe's record $1$.
For the specified disjoint experiments, fixed local settings, resources
and windows, the two admitted orders of evaluation give the same joint
actual-record law. Neither port reads the other's newly created flag
during isolation.
\end{assumption}

The first condition is a source support law. The second is a finite
record-order consistency law, stronger than equality of unread marginals.
They do not require arbitrary noncommuting experiments to commute. A
probe that physically acts on a region participating in the first
reaction is not automatically disjoint.

\section{Positive responses determine the old-bank daughter}

\begin{theorem}[Faithful old-bank continuation]\label{det:dark-theorem}
Under Assumptions~\ref{det:held-assumptions}--\ref{det:zero-assumptions},
suppose $p_j>0$ and the probe
$P=I-|\widehat v_j\rangle\langle\widehat v_j|$, where
$\widehat v_j=v_j/\|v_j\|$, is admitted. Then
\begin{equation}
 \chi=\widehat v_j\otimes\eta
 \quad\kappa_{j,u}\text{-almost surely for almost every actual }u.
                                                        \label{det:faithful}
\end{equation}
The law of $\eta,z$ is not yet determined.
\end{theorem}
\begin{proof}
Set $q(\chi)=\|(P\otimes I_{\mathcal M})\chi\|^2$. Evaluate the
two-port experiment by the original event first. Its formerly dark
joint record has probability
\begin{equation}
 \int_0^t h(p_j)e^{-\Lambda(p_j)u}
 \int F_s(q(\chi))\,
       \kappa_{j,u}(\Phi;d\chi,dz)\,du.          \label{det:dark-integral}
\end{equation}
It vanishes by the source support law. The integrand is nonnegative,
the outer density is strictly positive and $F_s(q)>0$ for every $q>0$.
Consequently $q=0$ for kernel-almost every daughter, for almost every
$u$. The null space of $P\otimes I_{\mathcal M}$ is
$\operatorname{span}(v_j)\otimes\mathcal M$, proving the factorization.
The second probe's unknown daughter does not enter this argument.
Equivalently, finitely many positive rank-one probes spanning
$v_j^\perp$ give the same conclusion by intersecting their null spaces.
\end{proof}

\begin{proposition}[Extension to an inaccessible reference]
\label{det:spectator}
The theorem extends from owned binary calibration inputs to the old bank
in \eqref{det:rows}, without performing a probe on the unknown reference,
if the reaction kernel obeys isometric spectator compatibility:
\begin{equation}
 \kappa^{\mathcal B}_{j,u}((I_D\otimes W)\Phi_0)
  =(W\otimes I_{\mathcal M},\operatorname{id}_z)_*
       \kappa^{\mathcal B_0}_{j,u}(\Phi_0)         \label{det:isometry}
\end{equation}
for every isometry $W:\mathcal B_0\to\mathcal B$, including the
admitted loading and probe contexts. The local record law is unchanged
and no output is created outside the embedded old support.
\end{proposition}
\begin{proof}
The span of $v_0,v_1$ has dimension at most two. Choose an owned
two-dimensional bank $\mathcal B_0$ and an isometry $W$ mapping its
corresponding two vectors to those rows. Theorem~\ref{det:dark-theorem}
applies to this calibration input. Push its almost-sure conclusion
forward by \eqref{det:isometry}. This is a comparison of reaction laws,
not the physical application of $W$ or a reference control to the
unknown input. An $n$-outlet version needs owned calibration on a span
of dimension at most $n$.
\end{proof}

Spectator compatibility is additional physical content. Nonlinear
Schmidt reweighting, for example, can respect spectator isometries while
violating joint-zero preservation. Conversely, faithful extraction
preserves every tested zero. Thus, within the responsive complete-probe
class, the zero condition and the old-bank conclusion are equivalent.
This equivalence fixes the logical strength of the result.

\begin{counterexample}[Storage instead of extraction]\label{det:store-rival}
Charge conservation and record-order consistency alone allow arbitrary
$h$. Give each detector a fresh neutral memory $M_D$ and use the
event-dependent isometry
\[
 |a\rangle_D|0\rangle_{M_D}
 \longmapsto |\mathrm{vac}\rangle_D|a\rangle_{M_D}
                  |\mathrm{flag}\ j\rangle.
\]
The readiness excitation supplies the charged flag. The entire unknown
input, including its entanglement, is stored. Remote populations are
unchanged in each branch; the two local isometries commute. Their joint
double-event density is the product of the original local densities,
for every continuous positive $h$, including $h(p)=\gamma p^2$.
On $(|00\rangle+|11\rangle)/\sqrt2$, the initially absent pair
$(1,0)$ occurs with probability $F_t(1/2)F_s(1/2)>0$.
It violates the joint-zero premise, not conservation or resource
accounting. Calling its displayed label an extraction would be false.
\end{counterexample}

\section{Finite records determine the held scalar rate}

\begin{theorem}[Joint held-rate and continuation selection]
\label{det:joint-selection}
The assumptions of Theorem~\ref{det:dark-theorem}, with the spectator
extension where needed, and finite record-order consistency imply
\begin{equation}
 h(p)=\gamma p,\qquad
 \chi_{\mathcal B}=v_j/\|v_j\|\quad\text{at record }j.        \label{det:joint-law}
\end{equation}
It suffices to compare one pair of strictly positive finite windows
over all the owned preparations used in the proof.
\end{theorem}
\begin{proof}
Prepare the known triangular input
\[
 \Phi_{x,y}=\sqrt y|11\rangle+\sqrt{x-y}|10\rangle
                      +\sqrt{1-x}|00\rangle,
 \qquad 0<y<x<1.
\]
The initial populations of the two $1$ ports are $x,y$. The
already-proved daughter theorem makes the second population $y/x$
after the first port clicks $1$; in the reverse order the remaining
first population is $1$. These are consequences about carried vectors.

Compare the finite event ``first port clicks $1$ by $t$, second port
is null by $s$.'' Its two probabilities are
\begin{equation}
 F_t(x)e^{-\Lambda(y/x)s},\qquad
 e^{-\Lambda(y)s}F_t(x).                         \label{det:null-comparison}
\end{equation}
Choose $q\in(0,1)$ with $h(q)>0$, set $x=q,y=qa$, and cancel the
positive $F_t(q)$. Exponential injectivity yields
$\Lambda(a)=\Lambda(qa)$. Iterate and use continuity at zero to obtain
$\Lambda(a)=\Lambda(0)=\gamma$.

For the double-click event, put $k_t=\int_0^t e^{-\gamma u}du>0$.
Interchange gives $k_tk_s h(x)h(y/x)=k_tk_s\gamma h(y)$.
Thus $f=h/\gamma$ obeys $f(x)f(z)=f(xz)$ and
$f(u)+f(1-u)=1$. For $u,v>0$ with $u+v<1$,
\[
 f(u)+f(v)=f(u+v)
 \left[f\left(\frac{u}{u+v}\right)+
       f\left(\frac{v}{u+v}\right)\right]=f(u+v).
\]
Continuity supplies the boundary. Nonnegative additivity makes $f$
monotone, fixes all rationals in $[0,1]$ and, by rational bracketing,
fixes every real argument. Hence $f(p)=p$.
\end{proof}

This theorem selects a rate in a declared intrinsic reaction class,
not the Hamiltonian Bell current. It contains no calculation identifying
the fixed $\gamma$ with $[J_{qr}]_+/w_r$. Its positive finite-window
algebra is conditional on the source support, spectator and
record-order principles. The causal probability antecedent and its
projection premise are recorded in \cite{M15}; here the preceding
positive-integrand theorem supplies the conditional populations.

For a finite separating experiment, $h(p)=\gamma p^2$ and
$x=2/3,y=1/3,\ \gamma t=\gamma s=1$ give the difference
\begin{equation}
 \Delta=\frac45(1-e^{-5/9})(e^{-1/2}-e^{-5/9})
        =0.01117694879\ldots.                    \label{det:rate-witness}
\end{equation}
It is a finite click/null bit, so an exact real-valued time readout is
unnecessary. If the two complete readout implementations each have
total-variation error at most $\epsilon$, their separating gap remains
at least $\Delta-2\epsilon$.

\section{New archives are a separate obligation}

Equation~\eqref{det:faithful} only implies that the mean complete event
state has the form
\begin{equation}
 |\widehat v_j\rangle\langle\widehat v_j|
                         \otimes\tau_{j,u}(\Phi).             \label{det:archive-open}
\end{equation}
The new CQ archive may still depend on the unknown input. Positivity
forces factorization from a pure old marginal, but does not force
input independence of the other factor.

\begin{theorem}[Complete archive selection]\label{det:archive-selection}
After Theorem~\ref{det:joint-selection}, suppose the following further
experiments are admitted. Known source vectors can be coherently tagged
before separation; product spectators do not affect the local event
kernel; a finite spanning set of quantum projectors and a separating
family of classical events can probe the new archive after the event;
and record-order consistency holds for those tests. Probe settings are
chosen after the original event, so they test one fixed prior event
kernel. Then, at fixed actual local $c$,
\[
 \tau_{j,u}(\Phi)=\tau_{j,u,c}
\]
on inputs with $p_j>0$, for almost every resolved event time, in the
following measure-theoretic sense: the event/archive measures agree
with this input-independent conditional kernel. The null remains the
complete frozen held null.
\end{theorem}
\begin{proof}
Take two known input vectors $\psi,\varphi$ whose $j$ rows are nonzero.
Prepare $\Xi=\sqrt a\,\psi|0\rangle_C+
\sqrt{1-a}\,\varphi|1\rangle_C$, with $0<a<1$, before separating
the tag. This is an owned coherent calibration input, not a coherent
superposition of classical histories or a copy of the unknown input.

Let $P$ be one setting in the spanning family, with a strictly positive
input-independent probability $\mu_P$ chosen after the original event.
It can be generated by fresh owned calibration pointers using the
detector just selected. Retain setting nulls, failures and new archives;
they terminate or branch the test according to the declared programme.
They are not discarded from its normalization.

Compare the records: original click $j$, chosen setting $P$, new-archive
click $1$ in classical event $Z$, and tag click $0$. In either order,
the probability has the common positive factor
$r_A r_C\mu_P r_P a\|E_j\psi\|^2$, where
$r_X=1-e^{-\gamma s_X}$ and $E_j=\langle j|_D$ denotes the derived
row operation. After removing that factor the two terms are
\[
 \operatorname{tr}[P\tau_j(\Xi)(Z)],\qquad
 \operatorname{tr}[P\tau_j(\psi|0\rangle_C)(Z)].
\]
The original-first evaluation uses the factorization from the entire
old pure branch, which includes $C$. The new-archive probe cannot then
change the old tag. Reverse evaluation selects the $\psi$ branch first.
Interchange equates these expressions. The tag-$1$ comparison equates
the same first term with the corresponding expression for $\varphi$.
Product-spectator independence removes the idle pure tag. A spanning
set of $P$ determines the matrix for each $Z$; the separating classical
events determine its CQ measure. This proves input independence.

For a continuous original time, perform the argument on every finite
time bin. Equality of finite operator-valued measures gives equality
of conditional densities almost everywhere. No conditional assertion
at a prescribed zero-probability timestamp is required. If a single
pointwise kernel version over an uncountable input family is desired,
regularity in the input or a jointly measurable version must be added;
the actual instrument conclusion only needs equality of the measures.
\end{proof}

\begin{corollary}[Calculated event instrument]\label{det:intrinsic-instrument}
The generated held detector has complete maps, with deterministic
register embeddings understood,
\begin{align}
 \mathcal I_j(du)(\rho)
  &=\gamma e^{-\gamma u}E_j\rho E_j^\dagger
                   \otimes\tau_{j,u,c}\,du,\nonumber\\
 \mathcal I_{\varnothing}(\rho)&=e^{-\gamma s}\rho.           \label{det:intrinsic-maps}
\end{align}
These maps are affine and CP on every inaccessible reference extension,
and finite generated programmes close on the complete CQ preparation
measure.
\end{corollary}
\begin{proof}
Multiply the proved event density $\gamma e^{-\gamma u}p_j$ by the
proved daughter projector and fixed archive. The factor $p_j$ cancels.
Because $\sum_j E_j^\dagger E_j=I$, the integral of the click traces
plus the null trace is one. The explicit products $E_j\rho E_j^\dagger$
remain positive after tensoring with an arbitrary identity; attachment
of a fixed positive archive has the same property. Composition with
the specified coherent and classical controls is therefore linear and
positive on the complete referenced bank. Ordinary classical mixing
gives dependence only on the complete CQ measure in this generated
class. Universal admission of additional writers is not inferred.
\end{proof}

\section{Loading, early failures and the domain of a driven extension}

For orthogonal projectors $Q_0+Q_1=I$, an admitted loader on ready
$|r\rangle_D$ has
\begin{equation}
 H_{\mathrm{load}}=\omega\sum_j Q_j\otimes
 (|j\rangle\langle r|+|r\rangle\langle j|),
 \quad
 U_t(\Psi|r\rangle)=\cos(\omega t)\Psi|r\rangle
  -i\sin(\omega t)\sum_jQ_j\Psi|j\rangle.       \label{det:loader}
\end{equation}
The equality tensors with any inaccessible reference. Loading with
actualization disabled until $\pi/(2\omega)$ supplies the rows used
above. Conservation is of the stated one-excitation charge; it does
not by itself establish conservation of laboratory energy or another
source charge.

If the same instantaneous population response is additionally
postulated during a drive, all live modes, including ready and unused
failure modes, must be registered. Then $\sum_aE_a^\dagger E_a=I$
on the live space and
\begin{equation}
 K_a(u)=\sqrt\gamma e^{-\gamma u/2}E_aU_u,
 \qquad K_{\varnothing}=e^{-\gamma s/2}U_s.       \label{det:driven}
\end{equation}
Normalization follows by summing $K_a^\dagger K_a$ and integrating.
During \eqref{det:loader}, intended outlet $j$ contributes
$\gamma e^{-\gamma u}\sin^2(\omega u)Q_j\rho Q_j\,du$,
whereas premature ready capture contributes
$\gamma e^{-\gamma u}\cos^2(\omega u)\rho\,du$ and a failure flag.
The null is the actual $U_s\rho U_s^\dagger$, with its scalar survival.
Suppressing the ready-failure term would destroy normalization.

The effect of a finite record, with its time unread, is consequently
\begin{equation}
 F_a(T)=\int_0^T K_a(t)^\dagger K_a(t)\,dt
       =\int_0^T\gamma e^{-\gamma t}
                    U_t^\dagger E_a^\dagger E_aU_t\,dt.
 \label{det:driven-effect}
\end{equation}
In particular, for $E_i^\dagger E_i=P_i$ and a fixed bounded
self-adjoint source Hamiltonian $H_S$, restore $\hbar$ and take
$U_t=e^{-iH_St/\hbar}$. Then
\begin{equation}
 \begin{aligned}
 F_i(\infty)&=P_i+\frac{i}{\hbar\gamma}[H_S,P_i]+R_{i,\gamma},\\
 \|R_{i,\gamma}\|&\le
 \frac{\|[H_S,[H_S,P_i]]\|}{\hbar^2\gamma^2}
 \le\frac{4\|H_S\|^2}{\hbar^2\gamma^2}.
 \end{aligned}
 \label{det:driven-effect-remainder}
\end{equation}
Indeed, twice differentiating $U_t^\dagger P_iU_t$ gives a second
derivative of norm at most
$\|[H_S,[H_S,P_i]]\|/\hbar^2$. Its integral Taylor remainder is
bounded by half this constant times $t^2$; integrating against
$\gamma e^{-\gamma t}$ gives the displayed bound. Thus finite response
generally measures a time-averaged effect, even for exact projective
outlet channels. This calculation still assumes the driven-response
premise of \eqref{det:driven}.

With loading stopped at $\tau=\pi/(2\omega)$ and $s\ge\tau$,
the integrated loading failure is
\[
 f_\tau=\frac{1-e^{-\gamma\tau}}2+
 \frac{\gamma^2(1+e^{-\gamma\tau})}{2(\gamma^2+4\omega^2)},
 \qquad f_\tau\le\frac{\pi\gamma}{4\omega}.
\]
The remaining success has common factor
$r_s=1-f_\tau-e^{-\gamma s}>0$ and label law
$r_s\|Q_j\Psi\|^2$. This is the finite tag required by the scalar
selection bridge elsewhere in the monograph. The held theorem alone
does not derive the driven-response premise. Vacuum/live coherent
superpositions also require a different extension, because the sum of
the registered effects is then the live projector, not the identity.

\chapter{Conserved converters and finite measurement programmes}
\label{det:chapter-instruments}

The construction in this chapter uses the actual diffusive law of
\cite{M12}, rather than the sharp intrinsic primitive of the previous
chapter. The event law of that native reader remains a statistical
input. Once it is supplied, one coherent converter determines both
transferred packet amplitudes and continuing daughter amplitudes.
The results concern one unknown input with an arbitrary inaccessible
reference, all future-active archives and a finite physical programme.

\section{Native law and the exact instrument it generates}

Let $H_t,L_t$ be bounded Hermitian operators, preparation independent
at fixed classical inputs and actual record history. Coefficients and
feedback are sufficiently regular for finite-horizon strong solutions;
piecewise held controls suffice. The primitive equations are
\begin{align}
 dY_t&=2\ell_tdt+dW_t,\quad
        \ell_t=\langle\psi_t,L_t\psi_t\rangle,\nonumber\\
 d\psi_t&=\left[-iH_t-\tfrac12(L_t-\ell_t)^2\right]\psi_tdt
                  +(L_t-\ell_t)\psi_tdW_t.       \label{det:native}
\end{align}
$W$ is Wiener in the declared complete physical filtration. The
admitted record is $Y$; a separately readable innovation is not included
by this constitution. Actualization and this output signature are not
deduced from the calculations below. Independent ports have the joint
noise law specified by their wiring.

\begin{lemma}[Complete likelihood and continuation]\label{det:likelihood}
Under a reference Wiener measure for the coordinate $Y$, let
\[
 dM_t=(-iH_t-\tfrac12L_t^2)M_tdt+L_tM_tdY_t,
 \qquad M_0=I.
\]
Then
\begin{equation}
 \mathbb E_0 M_T^\dagger M_T=I,
 \quad \mathbb P_\psi(dY)=\|M_T(Y)\psi\|^2\mathbb P_0(dY),
 \quad \psi_T=\frac{M_T(Y)\psi}{\|M_T(Y)\psi\|}.             \label{det:likelihood-eq}
\end{equation}
The same formulas with $M\otimes I_R$ describe inaccessible references
and yield normalized CP instruments on the complete coherent bank.
\end{lemma}
\begin{proof}
It\^o multiplication gives
$d(M^\dagger M)=2M^\dagger LM\,dY$. Bounded coefficients on a finite
horizon give the required mean-one norm martingales. The stochastic
logarithm of $\|M_t\psi\|^2$ has coefficient $2\ell_t$; changing measure
therefore makes $W=Y-\int2\ell_tdt$ Wiener. Applying It\^o's quotient
rule to $M\psi/\|M\psi\|$ gives exactly \eqref{det:native}.
Pathwise uniqueness identifies this with the proposed actual process.
For an event $E$, integration of $M_Y\rho M_Y^\dagger$ over its
reference paths gives its unnormalized output. Normalization follows
from the first equality, and reference positivity follows from the
operator formula. State-independent actual classical control kernels
can be appended to the same complete-path integral.
\end{proof}

This is the standard filtering likelihood calculation \cite{Filtering}.
Its reference measure is a mathematical coordinate, not an equilibrium
bath that has been physically prepared. It does not independently
select the stochastic law \eqref{det:native}.

\section{A single charge-preserving transfer}

Let $C$ be a $d$-dimensional carrier, $F$ have ready state $|r\rangle$
and terminal labels $1,\ldots,m$, and $E$ have packet vacuum $|0\rangle$
and the same terminal labels. The apparatus supplies matrices
$D_i:C\to C$ with $\sum_iD_i^\dagger D_i=I$. Rectangular maps can be
placed in a declared direct-sum carrier. Define
\begin{equation}
 G=\sum_i\left(D_i\otimes|i,i\rangle\langle r,0|
              +D_i^\dagger\otimes|r,0\rangle\langle i,i|\right).
                                                        \label{det:converter}
\end{equation}

\begin{theorem}[Conserved converter and finite synthesis]
\label{det:converter-theorem}
The operator $G$ is Hermitian, has norm one and obeys
\begin{equation}
 e^{-i\theta G}\psi|r,0\rangle
  =\cos\theta\,\psi|r,0\rangle
      -i\sin\theta\sum_iD_i\psi|i,i\rangle.      \label{det:transfer}
\end{equation}
It conserves $Q_F+N_E$, where $Q_F=|r\rangle\langle r|$ and
$N_E=\sum_i|i\rangle\langle i|$. A connected finite library of
equal-charge two-level exchanges, phases and a fixed anchor exchange
synthesizes $G$ on its whole local bank. No input copies or reference
controls are needed.
\end{theorem}
\begin{proof}
The map $V\psi=\sum_iD_i\psi|i,i\rangle$ is an isometry from the
ready subspace to the paired destination subspace. On their direct
sum, $G$ exchanges $\psi|r,0\rangle$ with $V\psi$ and has square
the identity. It vanishes on the orthogonal complement. Its exponential
therefore gives \eqref{det:transfer}. Both exchanged spaces have
charge one, giving the commutator zero.

Complete the isometry $\psi|1,1\rangle\mapsto V\psi$ to a unitary
$W$ on the $dm$-dimensional paired destination space; extend it by the
identity elsewhere, including the ready subspace. For
\[
 E_*=I_C\otimes(|1,1\rangle\langle r,0|+
                         |r,0\rangle\langle1,1|)
\]
one has $G=WE_*W^\dagger$ on every subspace, hence
$e^{-i\theta G}=We^{-i\theta E_*}W^\dagger$.
Complex Givens elimination factors $W$ into at most
$dm(dm-1)/2$ two-level rotations and $dm$ phases. Swaps along the
connected admitted graph implement each required pair. For a path
graph this yields a crude $O((dm)^3)$ gate count. All coefficients
come from the known apparatus matrices. Additional conserved charges
require connectivity inside their joint allowed blocks, and are not
implied by the readiness-charge calculation.
\end{proof}

If the synthesis of $W$ uses $L$ gates with operator error at most
$\epsilon_g$ each, and the anchor angle error is $\epsilon_\theta$,
unitary telescoping gives
\begin{equation}
 \|\widetilde U-U\|\le2L\epsilon_g+\epsilon_\theta.            \label{det:gate-error}
\end{equation}
The corresponding half-diamond channel error is no larger. Packing
fixed-angle compiler gates into duration $\delta$ can require controls
of size $O(\delta^{-1})$, even when a directly available weak $G$ pulse
would cost only $O(\delta^{-1/2})$. The synthesis premise, control
bandwidth and stocked pure resources remain explicit physical inputs.

\section{Full-output comparison for a finite reader}

Monitor a retired packet for time $\tau$ with
$L=\kappa B$, $B=\sum_j b_jP_j$, and $H=0$, where different $b_j$
are separated by at least one. Put $\mathcal R=\kappa^2\tau$.
The exact multiplier and its reference measure are
\begin{equation}
 M_y=\sum_jm_j(y)P_j,
 \quad m_j(y)=e^{\kappa b_jy-\kappa^2b_j^2\tau},
 \quad \mathbb P_0(dy)=\mathcal N(0,\tau)(dy).               \label{det:read-multiplier}
\end{equation}
On eigenlabel $j$, the actual $y$ is normal with mean
$2\kappa b_j\tau$ and variance $\tau$. A nearest-mean decoder stores
$z(y)$ in a classical register. The complete actual state following a
converter is $M_yU\psi/\|M_yU\psi\|$, including reference, flag and
all packet components; a wrong declaration does not substitute a
desired sharp daughter.

\begin{theorem}[Reference-uniform complete reader error]
\label{det:reader-theorem}
Let $e$ be the largest eigenlabel classification error. Then
\begin{equation}
 e\le2\Phi(-\sqrt{\mathcal R})\le e^{-\mathcal R/2},
 \qquad \tfrac12\|\mathcal R_{\mathrm{native}}
                   -\mathcal R_{\mathrm{sharp}}\|_\diamond
                  \le\sqrt{2e}.                            \label{det:reader-error}
\end{equation}
Here the sharp comparison retains the same coordinate $y$, the whole
packet bank and a displayed eigenlabel; conditional on that eigenlabel
it draws $y$ with density $m_j^2$ relative to $\mathbb P_0$.
The bound survives every common admitted finite continuation, including
a return of the packet to the carrier.
\end{theorem}
\begin{proof}
Half the separation of neighboring normal means is at least
$\kappa\tau$. Gaussian tails give the first inequality, including
the two-sided interior labels; the usual bound
$2\Phi(-x)\le e^{-x^2/2}$ gives the second.

Before dephasing the classical coordinates, use the common direct-integral
output space with isometries
\[
 V\psi=\int^{\oplus}\!|z(y)\rangle\otimes M_y\psi\,
                      \mathbb P_0(dy)^{1/2},\qquad
 W\psi=\int^{\oplus}\!\sum_j|j\rangle\otimes m_j(y)P_j\psi\,
                      \mathbb P_0(dy)^{1/2}.
\]
The direct-integral coordinate denotes the same retained $y$ on both
sides. Orthogonality of the $P_j$ implies
\[
 V^\dagger W=\sum_jP_j\int_{z(y)=j}m_j(y)^2\mathbb P_0(dy)
                 \ge(1-e)I.
\]
This operator is real positive, so
$(V-W)^\dagger(V-W)\le2eI$. For any purified input, the trace
distance of the two output pure states is at most their vector
distance, hence at most $\sqrt{2e}$. Taking all references proves the
half-diamond bound. Common dephasing of actual classical records and
any common subsequent channel contract the distance. No packet
component has been silently discarded in this comparison.
\end{proof}

\begin{theorem}[One-pulse finite instrument realization]
\label{det:one-pulse}
For an admitted normalized finite instrument
$\mathcal E_a(\rho)=\sum_jD_{aj}\rho D_{aj}^\dagger$, use terminal
labels $(a,j)$, angle $\theta=\pi/2$ and a degenerate packet reader
$B=\sum_{a,j}b_a|a,j\rangle\langle a,j|$. Then its exact finite
actual output differs from the specified complete sharp extension by
half-diamond distance at most
\begin{equation}
 \sqrt2\,e^{-\mathcal R/4}.                    \label{det:one-pulse-bound}
\end{equation}
The reference is arbitrary and inaccessible. The multiplicity $j$
remains coherently retained rather than being measured.
\end{theorem}
\begin{proof}
Full transfer gives $-i\sum_aW_a\psi$, where
$W_a\psi=\sum_jD_{aj}\psi|a,j\rangle_F|a,j\rangle_E$.
The native output is proportional to $\sum_am_a(y)W_a\psi$.
The sharp comparison has probability $\|W_a\psi\|^2$ and, on branches
of positive probability, complete conditional vector $W_a\psi/\|W_a\psi\|$. Tracing multiplicity,
when it is truly unavailable for return, yields $\mathcal E_a$ on
the carrier. Apply Theorem~\ref{det:reader-theorem} to the eigenspaces
indexed by $a$. The ready amplitude is exactly zero after the ideal
pulse. Define the decoder also on all off-subspace/error inputs so
that gate errors, bad readiness and every false declaration still
have an actual continuation. No collision or population limit is used.
\end{proof}

Thus exposure $\mathcal R\ge4\log(\sqrt2/\epsilon)$ suffices for
reader error $\epsilon$ before preparation and gate errors. This is a
scheduled finite measurement, not an assertion of an exponential
physical detection time or an exact finite sharp effect.

\section{Exact absorbing grid calibration}

For a timed module, a new vacuum packet is supplied to each collision.
If the flag is already terminal, the converter annihilates its tensor
product with a new vacuum packet. This establishes absorption for the
directed circuit, without declaring every old packet incapable of
return. Write $c=\cos\theta,s=\sin\theta$. With $N$ cells and no
early stopping, its pre-read wave is
\begin{equation}
 c^N\psi|r,\mathrm{vac}\rangle
 -is\sum_{k=1}^Nc^{k-1}\sum_iD_i\psi|i,e_{ki}\rangle,        \label{det:absorbing-wave}
\end{equation}
where $e_{ki}$ is a single packet in cell $k$. Its exact native
record density is
\begin{equation}
 c^{2N}\prod_\ell\varphi_0(y_\ell)
 +s^2\sum_{k,i}c^{2k-2}q_i\varphi_i(y_k)
                       \prod_{\ell\ne k}\varphi_0(y_\ell),   \label{det:absorbing-density}
\end{equation}
where $q_i=\|(D_i\otimes I_R)\psi\|^2$ and $\varphi_i$ is the
corresponding normal density. The actual conditional wave is obtained
by multiplying each component of \eqref{det:absorbing-wave} by its
native multipliers and normalizing. The mixture notation in
\eqref{det:absorbing-density} does not create an actual hidden finite
label $(k,i)$; multiple alternatives can remain coherent.

\begin{theorem}[Exact reduced semigroup and rounded times]
\label{det:grid}
For a bin of length $\delta$, choose
$\cos\theta_\delta=e^{-\nu\delta/2}$. Let
$P=I_C\otimes|r\rangle\langle r|$, $S=I-P$ and
$T_i=D_i\otimes|i\rangle\langle r|$. After discarding this
reader output and packet for this reduced comparison only, a collision
is exactly $e^{\delta\mathcal L}$, where
\[
 \mathcal L(\rho)=\nu\sum_iT_i\rho T_i^\dagger
                         -\tfrac\nu2\{P,\rho\}.
\]
The sharp stopped instrument from a ready flag has
\begin{equation}
 \mathbb P(K=k,I=i)=e^{-\nu(k-1)\delta}
           (1-e^{-\nu\delta})q_i,
 \qquad \mathbb P(\varnothing)=e^{-\nu N\delta}.             \label{det:geometric}
\end{equation}
\end{theorem}
\begin{proof}
Tracing the fresh packet gives Kraus operators
$K_0=S+cP$, $K_i=-isT_i$. The native packet reader does not change
this reduced channel when both its outcome and packet are traced.
Since $\sum_iT_i^\dagger T_i=P$ and $T_iT_j=0$, the master equation
has block solution
\[
 \rho_{PP}(t)=e^{-\nu t}\rho_{PP}(0),\quad
 \rho_{PS}(t)=e^{-\nu t/2}\rho_{PS}(0),\quad
 \rho_{SS}(t)=\rho_{SS}(0)+(1-e^{-\nu t})
                    \sum_iT_i\rho_{PP}(0)T_i^\dagger.
\]
These are exactly the Kraus blocks at $t=\delta$, including ready/sink
coherence and any reference. Multiplying the quiet factors gives
\eqref{det:geometric}. A rate-$\nu$ exponential variable rounded by
$T_\delta=\delta\lceil T/\delta\rceil$ has these same probabilities.
This is a coupling of calculated comparison laws, not a physical
pre-event threshold inserted into the source.
\end{proof}

Grid and continuous unrounded timestamps have total-variation distance
one. Any continuous-time comparison must state the rounding map or a
weaker metric. The exact reduced semigroup also leaves the raw native
records and growing coherent archives unspecified; it cannot be used
alone to predict their returns. Its current is the constructed
dissipative flux, not the original Hamiltonian current in the Bell
target.

\section{Complete finite networks, readiness and stopping}

For normalized states put $D(\rho,\sigma)=\tfrac12\|\rho-\sigma\|_1$.
For actual classical records plus retained quantum output, the same
notation means the trace distance on their CQ direct integral. The
record total variation is bounded by this distance.

\begin{theorem}[Finite adaptive complete-output bound]
\label{det:network}
Consider a finite programme of the admitted converters, native readers,
state-independent coherent controls and classical feedback. Include all
future-active memories, clocks and inaccessible references. Give every
timeout, false declaration, blocked request and failed preparation an
actual continuation. On every reachable complete input assume uniform
gate errors $u_k$, reader errors $e_j$, initial trace error
$\epsilon_{\mathrm{prep}}$, and any justified bank-truncation error
$\epsilon_{\mathrm{cut}}$. Compared with the same-clock sharp
extensions specified above, the final complete output error is at most
\begin{equation}
 \min\left\{1,\epsilon_{\mathrm{prep}}+
        \sum_k u_k+\sum_j\sqrt{2e_j}+
        \epsilon_{\mathrm{cut}}+p_{\mathrm{ex}}\right\}.     \label{det:network-error}
\end{equation}
$p_{\mathrm{ex}}$ is needed only when an uncapped comparator can exceed
the physical stock and the protocols agree until exhaustion. If both
have the same cap and blocked continuation it is zero.
\end{theorem}
\begin{proof}
Pad a stopped control tree by identity channels with deterministic dummy
outputs up to its uniform maximum slot count. This does not advance
physical clocks; both outputs are evaluated at the declared common
cut, with the same specified storage dynamics. Replace one gate or
reader at a time. Its local complete error is bounded uniformly even
on entangled inputs containing earlier archives. The common suffix
is a normalized CQ channel and therefore contracts trace distance.
The triangle inequality gives the two sums, and contracts the initial
error. A proved complete bank approximation contributes its stated
error. Couple capped and uncapped implementations until the first
exhaustion; only that event's probability can differ. The same
reasoning works with continuous-valued records provided depth,
coefficients and local errors are uniform over the admitted controls.
No contraction of an unknown nonlinear continuation is used: these
channels were calculated from the stated native law.
\end{proof}

\begin{lemma}[Conditioning and supplied readiness]\label{det:ready-rare}
For complete outputs at distance at most $\epsilon$, a common event
has probabilities $p,q$ with $|p-q|\le\epsilon$. If $p,q>0$, its
conditional outputs obey
\begin{equation}
 D(\rho_E/p,\sigma_E/q)
      \le\min\{1,\epsilon/\max(p,q)\}.          \label{det:rare}
\end{equation}
For commuting rank-one ready projectors $P_j$ on dedicated resource
cells, set $s=\sum_j[1-\operatorname{tr}(\omega P_j)]$. Then for
arbitrary incoming correlations
\begin{equation}
 D(\omega_{XB},\omega_X\otimes P_{\mathrm{ready}})
                \le\min\{1,2\sqrt s\}.         \label{det:readiness}
\end{equation}
\end{lemma}
\begin{proof}
For the first claim, event/complement block decomposition gives
$\|\rho_E-\sigma_E\|_1+|p-q|\le2\epsilon$.
If $p\ge q$, normalizing both differences with $p$ yields the bound
$\epsilon/p$; reverse the roles for $q\ge p$.
For readiness, the commuting-projector union bound makes the total
ready probability at least $1-s$. Project a purification onto the
joint ready subspace and normalize. Its trace distance from the
original purification is at most $\sqrt s$. The projected state
factorizes the rank-one ready bank. Its $X$ marginal differs from the
original marginal by at most $\sqrt s$ by contraction. A triangle
inequality proves \eqref{det:readiness}.
\end{proof}

A physical reset is a swap with a supplied same-dimensional ready flag.
For $H_{\mathrm{sw}}=(\pi/(2\tau))\operatorname{SWAP}_{FF'}$,
the unitary is $-i\operatorname{SWAP}$ and preserves $Q_F+Q_{F'}$.
It sends
$\rho_{CFAR}\otimes|r\rangle\langle r|_{F'}$
to $|r\rangle\langle r|_F\otimes\rho_{CF'AR}$.
The old correlations are archived, not erased. This works also after
false positives and timeouts. Each reset consumes a disclosed fresh
resource and does not generate an independently distributed innovation.

A concrete return test shows why complete outputs matter. For one
$D=I$ half-transfer, the conditional two-component state is
$(|r,0\rangle-im|1,1\rangle)/\sqrt{1+m^2}$, where
$m=e^{\kappa y-\kappa^2\tau}$. Returning the same cell through the
inverse converter gives ready probability
\[
 q_{\mathrm{return}}(m)=\frac{(1+m)^2}{2(1+m^2)}.
\]
The two positive coarse declarations with $m=2$ and $m=10$ give
$0.9$ and $0.5990099\ldots$. Small intervals around those values have
positive actual probability. A coarse display alone is not sufficient
to predict the return; the full native multiplier is.

\section{Finite diffusion cannot produce an exact sharp effect}

\begin{theorem}[Finite diffusive sharpness obstruction]
\label{det:nonsharp}
Fix the same complete classical preparation and a bounded finite
programme of the native diffusion, state-independent unitary gates,
pure ready embeddings and bounded stopping times. Retain every quantum
discard in a mathematical dilation. Every record event of positive
physical probability has a positive-definite effect on the finite
input space. In particular a nontrivial exact projective effect is
impossible in this primitive class at finite bounded resources.
\end{theorem}
\begin{proof}
For several independent native coordinates, the linear fundamental
matrix satisfies
\[
 dM=(-iH-\tfrac12\sum_jL_j^2)Mdt+\sum_jL_jM\,dY_j.
\]
Its determinant is the nonzero stochastic exponential
\[
 \det M_t=\exp\left\{\int_0^t
     \operatorname{tr}(-iH-\sum_jL_j^2)ds
       +\sum_j\int_0^t\operatorname{tr}L_j\,dY_j\right\}.
\]
This follows from the matrix It\^o rule, or from solving the inverse
linear SDE. Bounded stopping does not create a zero determinant.
Unitaries and pure resource embeddings preserve injectivity; fixed
mixed resources can be purified. Thus, for almost every complete
reference history $h$, the map $K_h$ from the input to the enlarged
output is injective. Input-independent classical kernels can be
included in its common reference measure.

For every nonzero input $v$ and event $E$ of positive reference measure,
$\langle v,F_Ev\rangle=\int_E\|K_hv\|^2\mathbb P_0(dh)>0$.
A positive-probability physical event has positive reference measure.
In finite input dimension this makes $F_E$ positive definite. A
nontrivial projector instead has a nonzero kernel. Tracing an output
resource does not alter the event probability, so cannot evade the
effect conclusion.
\end{proof}

The theorem permits exact identification of different actual classical
tags, because it fixes their complete preparation. It excludes neither
new sharp/jump primitives nor limits of diverging exposure. A useful
quantitative constraint is obtained from
\[
 \mathcal E=\sup_h\int_0^T\sum_j
 [\lambda_{\max}(L_j)-\lambda_{\min}(L_j)]^2dt.
\]
The drifts for two inputs at the same coordinate history differ by at
most twice these spectral diameters. Girsanov's formula gives
$D_{\mathrm{KL}}(P_\psi\|P_\phi)\le2\mathcal E$.
If a classical event discriminates the two inputs with both errors at
most $\epsilon<1/2$, binary data processing and the monotonicity of
binary relative entropy give
\begin{equation}
 \mathcal E\ge\frac{1-2\epsilon}{2}
                       \log\frac{1-\epsilon}{\epsilon}.      \label{det:exposure-lower}
\end{equation}
This is an exposure requirement, not a laboratory energy theorem.
It explains the divergence required by ideal sharpness within this
specific diffusive source constitution.

\chapter{Finite receptor response, retained nulls and delayed records}
\label{det:chapter-response}

The first checkpoint \cite{C01} contains a distinct finite detector
whose coherent excitation is followed by an assumed intrinsic latch.
This chapter expands its exact null calculation and its controlled
fresh-cell limit into complete theorems. The comparison includes
classical event times, outcomes, the source and inaccessible references.
It excludes the future return of receptors that have been discarded.
That restriction is essential: the exact complete null contains
source--receptor correlations which a reduced source instrument omits.

\section{The finite receptor and its stochastic premise}

Let $P_i$ be a finite projective resolution on the carried system.
The apparatus has a ready vector $|A_0\rangle$, orthogonal excited
vectors $|A_i^*\rangle$ and mutually distinguished terminal
recorded/spent states $|R_i\rangle$. During an exposure set the source
Hamiltonian to zero and use
\begin{equation}
 H_{\mathrm{int}}=\sum_i g_iP_i\otimes
       (|A_i^*\rangle\langle A_0|+|A_0\rangle\langle A_i^*|),
 \qquad C_i=\sqrt{\Gamma_i}I_S\otimes
                     |R_i\rangle\langle A_i^*|.             \label{det:receptor}
\end{equation}
Units have $\hbar=1$. The $g_i,\Gamma_i$ are positive apparatus
parameters. $H_{\mathrm{int}}$ acts identically within each possibly
degenerate $P_i$ sector and trivially on the reference.

\begin{assumption}[Intrinsic apparatus latch]\label{det:latch-postulate}
An actual latch $i$ has conditional intensity $\|C_i\Psi\|^2$,
normalized daughter $C_i\Psi/\|C_i\Psi\|$, and the associated
no-event evolution generated by
$H_{\mathrm{eff}}=H_{\mathrm{int}}-\tfrac i2\sum_iC_i^\dagger C_i$.
An event creates its actual record and consumes this cell's readiness;
there is at most one event per cell. Future event randomness is the
specified Markov latch law. No separate physical sampling tape is
available. All terminal products are specified, and failed readiness
produces its declared failed or blocked branch.
\end{assumption}

This is a disclosed squared-norm statistical primitive. The following
theorems derive the effective detector response from it and the
coherent interaction. They do not derive this primitive from the
Hamiltonian Bell current. The same-charge assignments can account for
one readiness unit passing through excitation into the terminal flag;
an energetic reservoir model would require additional dynamics.

\begin{theorem}[Exact fresh-cell event and complete null]
\label{det:finite-receptor}
Prepare the cell independently in $|A_0\rangle$ and allow one
uninterrupted exposure. Define
\begin{equation}
 \dot a_i=-ig_ib_i,\qquad
 \dot b_i=-ig_ia_i-\tfrac{\Gamma_i}{2}b_i,
 \quad a_i(0)=1,\quad b_i(0)=0.                 \label{det:amplitudes}
\end{equation}
For $\rho_{ij}=(P_i\otimes I_R)\rho_{SR}(P_j\otimes I_R)$ and
$\eta_i(t)=a_i(t)|A_0\rangle+b_i(t)|A_i^*\rangle$, the complete
unnormalized null is
\begin{equation}
 \widetilde\rho_{SRA}(t)=\sum_{i,j}\rho_{ij}\otimes
                           |\eta_i(t)\rangle\langle\eta_j(t)|.             \label{det:full-null}
\end{equation}
The timed event map on source and reference is
\begin{equation}
 \mathcal J_i(dt)(\rho)=\Gamma_i|b_i(t)|^2\rho_{ii}\,dt,
 \quad S(t)=\sum_iw_i(|a_i(t)|^2+|b_i(t)|^2),
 \quad w_i=\operatorname{tr}\rho_{ii}.          \label{det:finite-events}
\end{equation}
For a nonzero event the retained source/reference state is
$\rho_{ii}/w_i$. If all channels have the same $g,\Gamma$, then
the receptor-discarded null is exactly
\begin{equation}
 \mathcal N_\Gamma^T(\rho)=|a(T)|^2\rho+
                         |b(T)|^2\mathcal D(\rho),
 \qquad \mathcal D(\rho)=\sum_i\rho_{ii}.       \label{det:dephased-null}
\end{equation}
It is generally different from $S_\Gamma(T)\rho$.
\end{theorem}
\begin{proof}
The no-event Hamiltonian leaves each source sector invariant and acts
on its ready/excited span by the two-by-two matrix in
\eqref{det:amplitudes}. Its propagator applied to a ready vector is
$\eta_i$. Bilinearity gives \eqref{det:full-null} for every mixed
input and reference. Application of $C_i$ annihilates every excited
vector except $A_i^*$ and maps that one to the fixed terminal $R_i$,
giving the event map. Direct differentiation gives
\[
 \frac{d}{dt}(|a_i|^2+|b_i|^2)=-\Gamma_i|b_i|^2.
\]
Thus event integration and null trace sum to one. The conditional
hazard is $w_i\Gamma_i|b_i(t)|^2/S(t)$, not the event density itself.
For equal channels,
$\langle\eta_j|\eta_i\rangle=|a|^2+\delta_{ij}|b|^2$;
tracing the receptor gives \eqref{det:dephased-null}. The normalized
event formula follows by its trace. Degenerate internal coherence is
preserved because neither operator in \eqref{det:receptor} acts on it.
\end{proof}

Both roots of
$a_i''+(\Gamma_i/2)a_i'+g_i^2a_i=0$
have negative real part. Hence the cell eventually latches with
probability one on each populated sector. Equal channels produce
integrated mark weights $w_i$, but a finite response with
\[
 b(t)=-igt+O(t^2),\qquad
 q_\Gamma(t)=\Gamma|b(t)|^2=\Gamma g^2t^2+O(t^3).
\]
The density has quadratic onset and the cumulative probability cubic
onset. No finite response is exactly the instantaneous constant-rate
detector at its start.

\section{Null recovery and a retained-excitation counterexperiment}

A receptor-only unitary cannot send the different normalized $\eta_i$
to one common ready vector: it preserves their inner products. A
source-controlled operation can do so if separately admitted. Choose
$W_i\eta_i(T)=\sqrt{n_i(T)}A_0$, where
$n_i=|a_i|^2+|b_i|^2$, and apply
$\sum_iP_i\otimes W_i$ with the latch disabled. The factors
$\sqrt{n_i}$ are forced by unitarity. For equal channels the recovered
null is a scalar multiple of the original source; unequal channels
retain sector-dependent filtering. This operation uses a shutter and
source-controlled access. It is not a source-blind reset and is not
part of the fresh-cell theorem below.

\subsection{Finite pre-latch pointer contact and two overlap scales}
\label{det:pointer-contact}

The mechanical contact of \cite[\S14]{C01} acts on the complete null
of Theorem~\ref{det:finite-receptor}. It requires a shuttered interval
with both exchange and latching disabled. Take an independent pointer
with position variance $\sigma^2>0$ and wavefunction
\[
 \varphi_0(y)=(2\pi\sigma^2)^{-1/4}e^{-y^2/(4\sigma^2)},
 \qquad \varphi_d(y)=\varphi_0(y-d).
\]
In units $\hbar=1$, the only active Hamiltonian during the pulse is
\begin{equation}
 H_Y(t)=\sum_i v_i(t)|A_i^*\rangle\langle A_i^*|\otimes P_Y,
 \qquad P_Y=-i\partial_y,\qquad d_i=\int v_i(t)\,dt.
 \label{det:pointer-pulse}
\end{equation}
The real pulse profiles have finite integrals; any pointer free
evolution is absent or compensated as part of the declared control.

\begin{theorem}[Complete contact state and distinct overlaps]
\label{det:pointer-theorem}
For every source state with an arbitrary inaccessible reference, let
$a_i=a_i(T)$, $b_i=b_i(T)$ and $S(T)>0$ be as in
Theorem~\ref{det:finite-receptor}. After the contact, its complete
unnormalized null is
\begin{equation}
 \widetilde\rho_{SRAY}
 =\sum_{i,j}\rho_{ij}\otimes|\Xi_i\rangle\langle\Xi_j|,
 \qquad
 |\Xi_i\rangle=a_i|A_0,\varphi_0\rangle
                    +b_i|A_i^*,\varphi_{d_i}\rangle.
 \label{det:pointer-full-null}
\end{equation}
Tracing only $Y$ multiplies the excited--excited receptor coherence
$|A_i^*\rangle\langle A_j^*|$ by $m_{ij}$ and each
ready--excited coherence involving $A_i^*$ by $m_i$, where
\begin{equation}
 m_{ij}=e^{-(d_i-d_j)^2/(8\sigma^2)},\qquad
 m_i=e^{-d_i^2/(8\sigma^2)}.
 \label{det:pointer-overlaps}
\end{equation}
If a pointer-position acquisition with its usual squared-amplitude
law is additionally supplied, its density conditioned on the first
null is
\begin{equation}
 p(y\mid\varnothing)=\frac{1}{S(T)}\sum_iw_i
 \left(|a_i|^2|\varphi_0(y)|^2
             +|b_i|^2|\varphi_{d_i}(y)|^2\right).
 \label{det:pointer-density}
\end{equation}
The contact by itself is unitary and supplies no acquisition or
new stochastic latch law.
\end{theorem}
\begin{proof}
The excited projectors commute and the pulse translates only their
pointer factors. Applying this unitary to \eqref{det:full-null}
gives \eqref{det:pointer-full-null}; all maps are the identity on
$R$. Completing the square in
$\int\varphi_{d_j}(y)^*\varphi_{d_i}(y)\,dy$ gives $m_{ij}$;
setting one displacement to zero gives $m_i$. Expansion of each
$|\Xi_i\rangle\langle\Xi_j|$ then gives the stated trace factors.
For the additional acquisition, replace $a_i,b_i$ in the complete
branch by $a_i\varphi_0(y),b_i\varphi_{d_i}(y)$ and take its trace.
Distinct source sectors have zero off-diagonal trace, and the ready
and excited receptor vectors are orthogonal, giving
\eqref{det:pointer-density}. Its integral is one by the definition
of $S(T)$.
\end{proof}

Orthogonal receptor labels already eliminate some interference:
tracing both $A$ and $Y$ gives the same source/reference marginal as
tracing $A$ before the contact. One must not attach $m_{ij}$ to a
source cross term that this receptor trace has already removed.
If the pointer can return in the future programme, retain
\eqref{det:pointer-full-null}, rather than only its overlap-reduced
state. A local later response with the pointer idle can be calculated
from either that full state or its exact pointer trace.

For an explicit later effect, take two equal channels with
$d_0=d_1=d\ne0$ and resume the same exchange and latch without
source control, leaving the pointer idle. Then $m_{01}=1$ but
$m_0=m_1=m=e^{-d^2/(8\sigma^2)}<1$. Define
\[
 \begin{pmatrix}u(t)&v(t)\\v(t)&z(t)\end{pmatrix}
 =\exp\left[t\begin{pmatrix}0&-ig\\-ig&-\Gamma/2\end{pmatrix}\right].
\]
The excited pointer amplitude in either populated sector is
$v(t)a\varphi_0+z(t)b\varphi_d$. Summing the two record labels,
the resumed first-event density, including the probability of the
original null, is therefore
\begin{equation}
 q_m(t)=\Gamma\left(|v(t)a|^2+|z(t)b|^2
       +2m\operatorname{Re}\{v(t)a\overline{z(t)b}\}\right).
 \label{det:pointer-resumed-density}
\end{equation}
Choose a sufficiently short original exposure $T>0$, so that
$a>0$ and $b=-i\beta$ with $\beta>0$. Since
$v(t)=-igt+O(t^2)$ and $z(t)=1-\Gamma t/2+O(t^2)$,
comparison with zero displacement ($m=1$) gives the finite-window
joint-record difference
\begin{equation}
 \int_0^\delta[q_m(t)-q_1(t)]\,dt
 =\Gamma(m-1)ga\beta\,\delta^2+O(\delta^3)\ne0
 \quad\text{for sufficiently small }\delta>0.
 \label{det:pointer-timing-gap}
\end{equation}
Dividing by $S(T)$ gives the difference conditioned on the original
null. The equal instantaneous densities at resumption do not remove
this later timing effect: equal excited displacements leave
ready--excited interference suppressed. The resumed event probabilities
still use Assumption~\ref{det:latch-postulate}; the contact calculation
does not select that event law.

\begin{counterexample}[A current event can report an old excitation]
\label{det:stale}
Start with source $|0\rangle$ and a fresh binary cell. After a null
exposure $T$ with $b(T)\ne0$, the unnormalized state is
$|0\rangle(a|A_0\rangle+b|A_0^*\rangle)$. Apply a Hadamard to the
source only. The old-excited component is now
$b|+\rangle|A_0^*\rangle$. On resuming the latch, its instantaneous
record-$0$ density is $\Gamma|b|^2$, and that contribution leaves
$|+\rangle$, not $|0\rangle$. By continuity the discrepancy persists
over a sufficiently short finite resumption window. A subsequent fresh
finite $X$-basis detector, with success probability $r>0$, declares
$+$ with conditional probability approaching $r$ on this contribution,
whereas an erroneously inserted $|0\rangle$ daughter would give
$r/2$. The corresponding unconditioned finite joint-record gap is
$\tfrac12r\Gamma|b|^2\delta+o(\delta)$ for resumption duration
$\delta$. The pre-null probability is already included by the
unnormalized $b$.
\end{counterexample}

The exact finite propagator after a control must act on the complete
source--receptor bank. An excited $A_k^*$ paired with a different
source sector has no coherent return through its stated $P_k$
coupling, but it can still latch. This is the mathematical reason that
arbitrary controls during a retained exposure are outside the simple
projective daughter claim.

\section{A quantitative finite-response instrument theorem}

Set $g=\tfrac12\sqrt{\kappa\Gamma}$, hold $\kappa>0$ fixed and put
$\varepsilon=\kappa/\Gamma\le1/8$ and $d=\sqrt{1-4\varepsilon}$.
Let $\mathfrak I_\Gamma^T$ be the finite instrument with timed events
\eqref{det:finite-events}, null \eqref{det:dephased-null}, and the
old receptor permanently excluded from later use. Its comparator is
the same output space with
\begin{equation}
 \mathfrak I_\infty(dt,i)(\rho)=\kappa e^{-\kappa t}P_i\rho P_i\,dt,
 \qquad \mathfrak I_\infty(\varnothing)(\rho)=e^{-\kappa T}\rho.
                                                        \label{det:latch-limit}
\end{equation}
Clock conventions agree exactly: both times are continuous physical
times censored at the same $T$. Known terminal receptor states can be
retained on an event in both comparators. The no-event receptor is
discarded in both; its later coherent return is excluded.

\begin{theorem}[Uniform fresh-cell error with the null retained correctly]
\label{det:thirteen}
Under the stated finite receptor and primitive latch,
\begin{equation}
 \tfrac12\|\mathfrak I_\Gamma^T-
                   \mathfrak I_\infty^T\|_\diamond
       \le\min\{1,13\kappa/\Gamma\}             \label{det:thirteen-bound}
\end{equation}
uniformly for $T\ge0$, on one unknown input with arbitrary inaccessible
reference. For at most $M$ fresh exposures, arbitrary finite durations
and the same adaptive quantum/classical controls between exposures,
the complete record/source output distance, and therefore record-history
total variation, is at most
\begin{equation}
 \min\{1,13M\kappa/\Gamma\}.                  \label{det:thirteen-network}
\end{equation}
The assertion includes stopping and censoring, but no source control
during a retained finite exposure and no return of discarded receptors.
\end{theorem}
\begin{proof}
The two amplitude decay rates are $\Gamma(1\mp d)/4$. Solving
\eqref{det:amplitudes} gives
\[
 b(t)=-i\frac{2g}{\Gamma d}
       \left(e^{-\Gamma(1-d)t/4}-e^{-\Gamma(1+d)t/4}\right).
\]
Consequently, with $\lambda_\pm=\Gamma(1\pm d)/2$,
\begin{equation}
 q_\Gamma(t)=\frac{\kappa}{d^2}
    \left(e^{-\lambda_-t}-2e^{-\Gamma t/2}
                          +e^{-\lambda_+t}\right).           \label{det:q-exact}
\end{equation}
The density is nonnegative and integrates to one by the survival
identity and complete decay. Also $\lambda_-\ge\kappa$ and
$\kappa/\lambda_-=(1+d)/2$. Split the density difference into its
leading exponential coefficient, leading exponent and two fast terms.
Integrating absolute values yields
\begin{align}
 \int_0^\infty|q_\Gamma(t)-\kappa e^{-\kappa t}|dt
 &\le(d^{-2}-1)\frac{\kappa}{\lambda_-}
       +1-\frac{\kappa}{\lambda_-}
       +\frac{\kappa}{d^2}
                     \left(\frac4\Gamma+\frac1{\lambda_+}\right).
                                                        \label{det:q-l1}
\end{align}
To check the constants explicitly, $d\ge1/\sqrt2$, so the four
terms on the right are at most
$8\varepsilon$, $2\varepsilon/(1+1/\sqrt2)$,
$8\varepsilon$ and $4\varepsilon/(1+1/\sqrt2)$, respectively.
Their sum is less than $20\varepsilon$, and in particular less than
the checkpoint's retained conservative $22\varepsilon$ bound.
Censoring a probability law cannot increase total variation. Thus
\begin{equation}
 \delta_\Gamma(T):=\tfrac12\int_0^T
      |q_\Gamma(t)-\kappa e^{-\kappa t}|dt
       +\tfrac12|S_\Gamma(T)-e^{-\kappa T}|
                 \le11\varepsilon.                         \label{det:censored}
\end{equation}

Insert an intermediate normalized instrument with the same timed
events $q_\Gamma(t)P_i\rho P_i\,dt$ but scalar null
$S_\Gamma(T)\rho$. For every referenced positive input, the trace
norm of its difference from \eqref{det:latch-limit} is exactly
$2\delta_\Gamma(T)$: the event blocks have traces summing to
$\operatorname{tr}\rho=1$ at each time, and the null block is a
scalar multiple of the same state. This proves the corresponding
half-diamond bound, since a Hermiticity-preserving channel difference
can be optimized over states with a reference.

The actual null differs from that scalar surrogate by
$|b(T)|^2(\mathcal D-\operatorname{id})(\rho)$.
From the square representation of \eqref{det:q-exact},
$q_\Gamma(T)\le\kappa/d^2\le2\kappa$, whence
$|b(T)|^2\le2\varepsilon$.
The half-diamond distance between two channels is at most one,
so this missing null term costs at most $2\varepsilon$. Adding it to
\eqref{det:censored} proves \eqref{det:thirteen-bound}.
The coefficient $13$ is retained as a simple conservative bound;
the intermediate estimates show it is not optimal.

For the adaptive programme, both local instruments are normalized CP
maps on the same retained output domain by direct calculation. The
local bound is uniform in $T$, in the input reference and in all
between-exposure controls. Pad stopping by identities, replace the
$M$ exposure slots one at a time, and contract through each common
suffix. This gives \eqref{det:thirteen-network}, including actual
null continuations at every intermediate slot. Rare conditional
outputs obey \eqref{det:rare}; they have no uniform error independent
of their success probabilities.
\end{proof}

This establishes the checkpoint's stated constants with a complete
proof and its correct scope. The convergence is integrated over
event times, not uniform pointwise in the event density at zero.
It improves useful throughput without taking $\kappa$ to zero:
$\Gamma$ and $g=\tfrac12\sqrt{\kappa\Gamma}$ increase while the
effective click rate stays fixed. The required latch strength and
coherent excitation are therefore explicit growing resources. No
theorem here obtains those resources from an unmodeled reservoir or
proves complete microscopic output convergence after permitting old
receptors to return.

\section{A three-stage null and a completed loss branch}

The distinction persists when excitation, capture and loss are
separate. For a unit bright input, let the no-latch/no-loss amplitudes
$x,y,z$ denote respectively the initial carrier, an emitted product
and a receptor excitation. Specify
\begin{equation}
 \dot x=-igy,\qquad
 \dot y=-igx-ihz-\tfrac\ell2y,\qquad
 \dot z=-ihy-\tfrac\Gamma2z,
 \quad(x,y,z)(0)=(1,0,0).                       \label{det:three-stage}
\end{equation}
The primitive latch and loss channels give record density
$\Gamma|z|^2$ and loss density $\ell|y|^2$. The norm identity
\[
 |x(T)|^2+|y(T)|^2+|z(T)|^2+
 \int_0^T(\Gamma|z(t)|^2+\ell|y(t)|^2)dt=1
\]
proves normalization, but an unobserved loss belongs in the observed
no-record branch. For a source with bright population $a$ and a
dark spectator component, put $D(T)=\int_0^T\Gamma|z|^2dt$.
The observed record survival is $1-aD(T)$, and its observed-history
hazard is
\[
 \frac{a\Gamma|z(t)|^2}{1-aD(t)}.
\]
Dividing instead by $|x|^2+|y|^2+|z|^2$ would condition on both no
latch and no loss, a different experiment. Since
$z(t)=-gh t^2/2+O(t^3)$, the physical record density is
$a\Gamma g^2h^2t^4/4+O(t^5)$. This quartic onset is an exact
finite-response prediction. A new attempt must retain the dark
component, the surviving three amplitudes and the lost branch; it
cannot restart with a newly idealized bright input.

\section{Delayed classical acquisition: a complete finite filter}

The following result consolidates the checkpoint's delayed-record
formulas, while keeping their separate statistical constitution
explicit. A native classical source generator is supplied in advance.
An event can create a pending daughter that later captures a finite
ready site or is lost. This can be a useful benchmark for delayed
recording, but it is not an admitted passive quantum current meter
without a separate material compatibility theorem.

\begin{theorem}[Finite hidden-state acquisition and continuation]
\label{det:delayed-filter}
Let $Z_t$ be a finite-state Markov process whose complete states contain
source state, pending daughters, readiness, fuel and every retained
classical memory. On each record-adapted control segment let its
generator, acting on row laws, be
\begin{equation}
 Q(t)=A(t)+\sum_{r\in\mathcal R}B_r(t).          \label{det:filter-generator}
\end{equation}
Each $B_r$ is an entrywise nonnegative kernel of transitions that
write observed mark $r$; self-transitions may represent physical
marked events. $A$ has nonnegative off-diagonal entries and
$A\mathbf1=-\sum_rB_r\mathbf1$. Its diagonal retains the killing
rate for every observed event, while all unobserved transitions,
including losses, stay in $A$. Rates are bounded and controls depend
only on admitted previous observations.

Starting from row law $\alpha$, the unnormalized no-record law solves
\begin{equation}
 \dot\alpha_t=\alpha_tA(t).                    \label{det:filter-no-event}
\end{equation}
Its trace $\alpha_t\mathbf1$ is the no-record probability. A record
$r$ at $t$ has density $\alpha_tB_r(t)\mathbf1$, and the exact
retained-state posterior is
\begin{equation}
 \frac{\alpha_tB_r(t)}{\alpha_tB_r(t)\mathbf1}.               \label{det:filter-update}
\end{equation}
The normalized update is asserted only at records of positive density,
for almost every actual event time; zero-density histories carry no
conditional-state claim. These rules iterate to every finite acquired
history, including adaptive replenishment and selective stopping, using
the actual posterior resources rather than a reset source state.
\end{theorem}
\begin{proof}
Over $dt$, retain every unobserved transition and remove every observed
inflow from the no-record law. Its balance is exactly
$\alpha_{t+dt}=\alpha_t+\alpha_tA(t)dt+o(dt)$.
For a prospective record $r$, sum the marked transition probabilities
from each incoming hidden state: the unnormalized post-event law is
$\alpha_tB_r(t)dt+o(dt)$. Its trace is the density, and ordinary
conditional probability gives \eqref{det:filter-update}. Bounded
finite rates control the multiple-event remainder and ensure the
finite propagator exists. Induct over record times; between them
propagate by the appropriate $A$, and at them multiply by the
appropriate $B_r$. The trace of this product is the joint density of
that finite history. Summing all histories and the final no-record
branch recovers total probability one from $Q\mathbf1=0$.
Adaptive control chooses subsequent matrices from the actually
observed history; the same conditional argument applies on each
branch. No quantum measurement or normalized linear extraction was
used in this classical filtering proof.
\end{proof}

If the extended physical reactions project to a native source update
with exactly its original total propensities, then for source projection
$\pi$ one has
$\mathcal L_{\mathrm{ext}}(f\circ\pi)
=(\mathcal L_{\mathrm{src}}f)\circ\pi$.
Indeed apparatus-only transitions vanish on $f\circ\pi$, and the
remaining projected transition sums agree term by term. Uniqueness
of the finite-state martingale problem gives the native projected
path law. This is a sufficient neutrality equation, not a derivation
of why an added physical daughter leaves every source propensity
unchanged. Multiplying source rates by a depleted readiness variable
instead blocks the source and changes this equality.

For a concrete complete example, let the source have transitions
$r\to q\to s$ at rates $a_1,a_2>0$, emitting labeled daughters
$X_1,X_2$. Each pending daughter captures a single ready site at rate
$\beta_R$ or is lost at rate $\beta_M$; the first capture exhausts
the site. Set $k=\beta_R+\beta_M$. A daughter's probability of no
capture by age $v$, allowing prior loss, is
\[
 A_0(v)=e^{-kv}+\int_0^v\beta_Me^{-ku}du
       =\frac{\beta_M+\beta_Re^{-kv}}{k}.
\]
For first acquired mark $R_1$ at $t$, define
\[
 h_t(s)=a_1e^{-a_1s}\beta_Re^{-k(t-s)},\qquad
 B(v)=e^{-a_2v}+\int_0^v a_2e^{-a_2u}A_0(v-u)du.
\]
Condition first on the emission of $X_1$ at $s<t$. Its survival to
capture supplies $h_t(s)ds$; $B(t-s)$ sums no second source event
and a second event whose daughter has not captured first. Hence
\begin{equation}
 f_{R_1}(t)=\int_0^th_t(s)B(t-s)ds.             \label{det:delayed-density}
\end{equation}
The actual source has already reached $s$ with conditional probability
\begin{equation}
 \frac{\displaystyle\int_0^th_t(v)
          \int_0^{t-v}a_2e^{-a_2u}A_0(t-v-u)du\,dv}
      {f_{R_1}(t)}.                             \label{det:delayed-posterior}
\end{equation}
This is strictly positive at $t>0$; an acquired $R_1$ cannot be
interpreted as a present-state projection onto its historical
destination $q$.

If a fresh site is supplied immediately without removing old daughters,
the instantaneous next-record intensity is
\begin{equation}
 \frac{\beta_R}{f_{R_1}(t)}\int_0^th_t(v)
       \int_0^{t-v}a_2e^{-a_2u}e^{-k(t-v-u)}du\,dv.            \label{det:pending-rate}
\end{equation}
The inner exponential now requires $X_2$ still to exist, rather than
merely to have avoided capture through loss. Theorem~\ref{det:delayed-filter}
derives the same expression by retaining pending daughters in the
posterior. Zero-current source holds can therefore contain new
acquired historical records without new native source events.

\section{What this part establishes for the measurement architecture}

The strongest statements now have one proof each. Joint dark-channel
and finite-order tests select held capture and its scalar rate under
their stated source premises. A conserved finite converter plus a
primitive actual diffusive reader yields a complete finite instrument,
with explicit reference, archive, resource and network bounds. The
finite receptor theorem establishes a controlled constant-rate
instrument limit only after its intrinsic latch has been supplied,
and the delayed-state theorem correctly propagates pending classical
records. These are compatible mathematical tools where their state,
output and resource hypotheses coincide.

These results do not identify three different null constitutions. A frozen
intrinsic null, a native diffusive conditional multiplier and an unobserved
finite receptor excitation are different physical states. None of the rate
constants $\gamma,\nu,\kappa$ becomes the Hamiltonian Bell intensity merely
by being an event rate. The detector theorems remain conditional on their
own actualization laws. The later pilot and massive constructions provide
separate complete measurement implementations; they do not derive every
Gaussian or intrinsic-latch primitive in this part. Their source, record
and continuation laws must be connected by the explicit state-space and
interaction comparisons stated there.

% End consolidated source: parts/05_detectors.tex

\part{Protected Transport and Complete Retained States}\label{part:protection}

% Begin consolidated source: parts/06_protection.tex
\chapter{Energy-gap protection of complete source transport}
\label{prot:gapchapter}

Exact neutrality of an aperture is stronger than conservation of its readiness charge. A charge-preserving disturbance can rotate an unknown carried state or write information into a returning memory. The construction here replaces exact neutrality of a specified coherent disturbance class by a finite energy penalty and a derived error bound \cite{M17}. It does not derive the actualization interface. Its antecedents are Hamiltonian error suppression with encoded sectors and retained environments \cite{MarvianLidar}; the contribution needed in this programme is a complete-source estimate with explicit operational boundaries.

Set $\hbar=1$ in this part. Operator norms always refer to the entire active bank. For normalized states write $D(\rho,\sigma)=\tfrac12\|\rho-\sigma\|_1$. An inaccessible reference is unrestricted and has no interaction of its own. Actual classical preparation labels and control keys are conditioned on, not averaged away to conceal their influence.

\section{Earlier algebraic protection and its statistical premise}
\label{prot:ancestral}

The precursor ownership theorem \cite{M04} protects event response by a different mechanism: conservation of noncommuting charges in an already positive marked source generator. It belongs beside the coherent protection construction because the assumptions and conclusions differ.

Let a finite provenance factor carry an irreducible spin-$j$ triple $S_k$, so $\sum_{k=1}^3S_k^2=j(j+1)I$. At fixed classical operational coordinates, assume a complete Heisenberg generator
\[
 \mathcal L^*A=i[H_0+H_I,A]
 +\sum_\alpha\left(L_\alpha^\dagger A L_\alpha
                  -\tfrac12\{L_\alpha^\dagger L_\alpha,A\}\right).
\]
The sum includes successful, failed, hidden, and unread event channels. Classical source transitions can be included when the same logical charges are identified in their incoming/outgoing sectors. $H_0$ is the specified incidence-off baseline, not a term fitted afterward to cancel disturbance. Impose the new charge-balance equations
\begin{equation}
 \mathcal L^*S_k=i[H_0,S_k],\qquad k=1,2,3.
 \label{prot:chargeward}
\end{equation}
The positive marked-generator representation is an explicit statistical input here. This theorem cannot be used to derive that representation from the charge equations.

\begin{theorem}[Noncommuting charge balance constrains every marked channel]
\label{prot:casimir}
Under these assumptions, $[L_\alpha,S_k]=[H_I,S_k]=0$ for every channel and charge. Hence $L_\alpha=\ell_\alpha I$ and $H_I=h_I I$ on the irreducible provenance factor. With an explicit additional operational Hilbert factor, the conclusion is instead membership in its commutant: $L_\alpha=I\otimes B_\alpha$.
\end{theorem}
\begin{proof}
For a self-adjoint charge define
\[
 \mathfrak D(S)=\mathcal L^*(S^2)-\mathcal L^*(S)S-S\mathcal L^*(S).
\]
Expanding each dissipator and canceling the Hamiltonian derivation gives
\[
 \mathfrak D(S_k)=\sum_\alpha[L_\alpha,S_k]^\dagger[L_\alpha,S_k]\succeq0.
\]
Unitality, the scalar Casimir and \eqref{prot:chargeward} imply
\[
 \sum_k\mathfrak D(S_k)
 =\mathcal L^*(j(j+1)I)-i[H_0,\sum_kS_k^2]=0.
\]
Every positive summand is therefore zero. Each channel commutes with each charge; its dissipator then vanishes on the charges and the balance equation forces $[H_I,S_k]=0$. Irreducibility gives the stated commutant by Schur's lemma.
\end{proof}

The theorem excludes more than unequal event intensities: $L=\sqrt b\,\sigma_z$ has scalar event effect $bI$ but dissipates transverse spin and violates \eqref{prot:chargeward}. Conversely an unprotected multiplicity register $Z$ allows $L=I_{\rm prov}\otimes\operatorname{diag}(\sqrt{b_0},\sqrt{b_1})_Z$ with unequal rates while conserving every provenance charge. Protecting only total spin or a proper subsystem therefore does not establish complete carrier neutrality.

Physical archive writing requires transported charges. Let $U_m$ be a fixed reversible append on a finite allocated bank and pointer, chosen before event rates. A concrete append increments pointer $p$ modulo capacity $K$ and adds the nonzero mark code $\eta(m)$ modulo the cell alphabet in the addressed cell. From a blank bank it preserves previous entries for at most $K$ writes. For arbitrary raw jumps $L_m$, put $R_m=U_m^\dagger L_m$. The transported balance law is
\[
 i[H-H_0,S]+\sum_m\left[
 L_m^\dagger U_m S U_m^\dagger L_m
             -\tfrac12\{L_m^\dagger L_m,S\}\right]=0.
\]
It is exactly \eqref{prot:chargeward} for the corrected operators $R_m$. Applying Theorem~\ref{prot:casimir} to a complete set of primitive bank factors yields
\[
 R_m=I_{\rm bank}\otimes B_m,\qquad
 L_m=(U_m\otimes I)(I_{\rm bank}\otimes B_m),
 \qquad L_m^\dagger L_m=I_{\rm bank}\otimes B_m^\dagger B_m.
\]
This is conservation through a supplied archive transport, not conservation of the bare written charge. The append cannot be chosen afterward to absorb an arbitrary susceptibility. With classical operational coordinates, common scalar effects give equality of successive event and null laws independently of protected bank contents, by induction over matching full histories. With an active quantum operational factor, conditioning can still steer a correlated protected state; effect factorization is not a claim of product continuation.

There is also a quantitative version. Let $E_k=\mathcal L^*S_k-i[H_0,S_k]$ and
\[
 G=\sum_{k,\alpha}[L_\alpha,S_k]^\dagger[L_\alpha,S_k]
   =-\sum_k\{S_k,E_k\},\qquad g=\|G\|.
\]
The equality follows from the same Casimir calculation without setting $E_k$ to zero. The nonnegative $G$ obeys $g\le2\sum_k\|S_k\|\|E_k\|$. Stack the channels into $V\psi=(L_\alpha\psi)_\alpha$. For a unit vector $n$, the column commutator with $n\cdot S$ has norm at most $\sqrt g$. Duhamel for $U=e^{i\theta n\cdot S}$ therefore gives
\[
 \|(I_{\rm marks}\otimes U)^\dagger VU-V\|\le|\theta|\sqrt g.
\]
Every spin conjugation has a representative rotation of angle at most $\pi$. Haar averaging gives a scalar column $V_0=(c_\alpha I)_\alpha$ satisfying
\begin{equation}
 \|V-V_0\|\le\pi\sqrt g,
 \qquad |\sqrt{\lambda_i(\rho)}-\sqrt{\gamma_i}|\le\pi\sqrt g,
 \quad\gamma_i=\sum_{\alpha\in i}|c_\alpha|^2.
 \label{prot:roothazard}
\end{equation}
The rate inequality is reverse triangle inequality applied to $V_i\rho^{1/2}$ in Hilbert--Schmidt norm, so arbitrary passive references are included. For $m$ separately protected primitive factors the same argument gives $\pi\sqrt{m g_{\rm tot}}$; no dimension-free constant for a growing bank is asserted.

To connect such bounds to actual finite exposure, define along a common stopped history
\[
 \chi=\mathbb E\int\sum_i(\sqrt{\lambda_i}-\sqrt{\gamma_i})^2dt,
 \qquad A_0=\mathbb E\int\sum_i\gamma_i dt.
\]
The identity $|a-b|\le2\sqrt b|\sqrt a-\sqrt b|+|\sqrt a-\sqrt b|^2$ and Cauchy--Schwarz give
\begin{equation}
 \mathbb E\int\sum_i|\lambda_i-\gamma_i|dt
                  \le2\sqrt{A_0\chi}+\chi.
 \label{prot:exposure}
\end{equation}
When both complete marked processes admit a common Poisson coupling on that history domain, the right side bounds unmatched-event probability before the stated stop; cutoff and preparation errors are added separately. Mean bulk neutrality alone does not supply this exposure estimate.

The exact and approximate charge results are conditional statistical protection. They do not fix the common scalar rate, derive primitive Markov chemistry, forbid an unrepresented stress register, or establish universal admission of the positive source-instrument class. The following Hamiltonian gap construction is different: it derives coherent suppression without using any event probability law in its proof.

\section{A nonempty encoding and nuisance class}

Two logical qubits $L$, possibly entangled with an inaccessible reference $R$ and old memory $E$, are encoded into four physical qubits. Define
\begin{equation}
 S_X=X_1X_2X_3X_4,\qquad S_Z=Z_1Z_2Z_3Z_4,
 \qquad P=\tfrac14(I+S_X)(I+S_Z),\quad Q=I-P.
 \label{prot:code}
\end{equation}
An isometry is
\begin{equation}
 C|a,b\rangle=\frac{|0,a,b,a\oplus b\rangle+
 |1,1\oplus a,1\oplus b,1\oplus a\oplus b\rangle}{\sqrt2}.
 \label{prot:encoder}
\end{equation}
Starting from $|0,a,b,0\rangle$, it is implemented by CNOTs $2\to4$, $3\to4$, a Hadamard on 1, and CNOTs $1\to2,3,4$. The inverse is a full unitary decoder; leakage becomes logical/syndrome amplitudes rather than being projected away. The two ready ancillas are physical independent supplies.

The protecting Hamiltonian and nuisance class are
\begin{align}
 H_{\rm pen}&=\tfrac\Delta2(I-S_X)+\tfrac\Delta2(I-S_Z),\\
 H_\Delta&=H_{\rm pen}+H_0+V,\qquad
 [H_0,P]=0,\quad\|H_0\|\le b,\\
 V&=\sum_{i=1}^4\sum_{\alpha=x,y,z}
               \sigma_i^\alpha\otimes B_{i\alpha},
 \qquad B_{i\alpha}=B_{i\alpha}^\dagger,\quad\|V\|\le v.
 \label{prot:model}
\end{align}
All operators are bounded and stationary on the exposure interval. The $B_{i\alpha}$ may act jointly on old memories, fresh archives, controller variables and live path modes, and need not commute. The bound is on their sum, not merely on each coefficient. A classical retained key may select different such generators, provided the bound is uniform in that key. The penalty has norm $2\Delta$ and complementary gap $\Delta$.

Every nonidentity one-site Pauli anticommutes with at least one stabilizer. If $S\sigma=-\sigma S$ and $SP=P$, then
\[
 P\sigma P=PS\sigma P=-P\sigma SP=-P\sigma P=0.
\]
Consequently
\begin{equation}
 PVP=0.
 \label{prot:compressionzero}
\end{equation}
The disturbing operators have not been assumed to commute with the code. They can drive transitions out of it. The code removes their first-order logical compression by an explicit algebraic identity.

The theorem below also applies to an infinite-dimensional retained bank when the relevant operators are bounded. Finite duration, finite expected energy, or a finite number of observed records does not imply these bounds. An unbounded reservoir requires a separate domain and energy-control theorem.

\section{The complete finite-gap estimate}

\begin{theorem}[Complete propagator protection]
\label{prot:gap}
Let $H_{\rm pen}P=0$, $H_{\rm pen}\ge\Delta Q$, $[H_0,P]=0$, $PVP=0$, and let all terms be bounded self-adjoint with $\|H_0\|\le b$, $\|V\|\le v$. For $\delta=\Delta-2b-v>0$, set $U_\Delta(t)=e^{-iH_\Delta t}$ and $U_0(t)=e^{-iH_0t}$. Then for every $t\ge0$,
\begin{equation}
 \|(U_\Delta(t)-U_0(t))P\|
 \le e_\Delta(t):=\min\left\{2,\frac{2v+tv^2}{\delta}\right\}.
 \label{prot:gapbound}
\end{equation}
The estimate is unchanged after tensoring any inaccessible reference. It controls all coherent output memories, not only the carried marginal.
\end{theorem}
\begin{proof}
Block the complete Hamiltonian relative to $P,Q$:
\[
 H_\Delta=H_d+W,
 \qquad H_d=\begin{pmatrix}A&0\\0&D\end{pmatrix},
 \qquad W=\begin{pmatrix}0&B^\dagger\\B&0\end{pmatrix}.
\]
Here $A=PH_0P$, $B=QVP$, $D=QH_\Delta Q$, and
$D\ge(\Delta-b-v)Q=(b+\delta)Q$, while $A\le bP$. The norm-convergent integral
\[
 X=\int_0^\infty e^{-rD}Be^{rA}dr
\]
satisfies $DX-XA=B$ and $\|X\|\le v/\delta$. To check the identity, differentiate $e^{-rD}Be^{rA}$ and integrate its vanishing boundary term. Define
\[
 S=\begin{pmatrix}0&-X^\dagger\\X&0\end{pmatrix}.
\]
Then $S^\dagger=-S$, $\|S\|=\|X\|$, and $[S,H_d]=-W$. With $f(u)=e^{uS}We^{-uS}$,
\begin{align*}
 e^SH_\Delta e^{-S}
 &=H_d+f(1)-\int_0^1f(u)du\ =H_d+R,\\
 R&=\int_0^1u e^{uS}[S,W]e^{-uS}du.
\end{align*}
Since the conjugations are unitary,
\[
 \|R\|\le\tfrac12\|[S,W]\|\le v^2/\delta,
 \qquad\|e^{\pm S}-I\|\le\|S\|\le v/\delta.
\]
Duhamel applied to $H_d+R$ and $H_d$, followed by the two changes of frame, gives
\[
 \|e^{-S}e^{-i(H_d+R)t}e^S-e^{-iH_dt}\|
 \le2v/\delta+tv^2/\delta.
\]
On $P$ the last unperturbed propagator equals $U_0(t)$. Two unitaries differ by at most two, completing \eqref{prot:gapbound}. Tensoring an identity preserves each operator norm; purification extends the associated trace-distance estimate to mixed complete inputs.
\end{proof}

For a pure complete input, output trace distance is at most $\min\{1,e_\Delta(t)\}$. This improves the ordinary $O(vt)$ bound to $O(v/\Delta+tv^2/\Delta)$ at fixed $b,v$. A useful regime is fixed finite transport time $T$ with increasing finite $\Delta$, or simultaneous scaling $v/\Delta\to0$ and $Tv^2/\Delta\to0$. The theorem is neither an all-time statement nor an assertion that the penalty is free.

\section{When the new archive is input independent}

The operator theorem allows general code-preserving $H_0$; it does not guarantee archive neutrality for every such $H_0$. For literal faithful transport choose
\begin{equation}
 H_0=I_{\rm phys}\otimes(H_E\otimes I_F+I_E\otimes H_F)
 \label{prot:split}
\end{equation}
and prepare an arbitrary old $\Psi_{LER}$ with an independent fresh $a_0\in F$. The ideal output is
\[
 (C\otimes I)(I_L\otimes e^{-iH_Et}\otimes I_R)\Psi_{LER}
                       \otimes e^{-iH_Ft}a_0.
\]
Thus the old bank follows its stated free evolution and the new archive has an input-independent state. The nuisance may couple $E,F$; Theorem~\ref{prot:gap} bounds the resulting deviation on their complete joint state. In particular two different carried inputs produce fresh-archive marginals at distance at most $2e_\Delta(t)$, by comparison with the common ideal archive. A later common coherent return unitary preserves the complete error.

The split premise is necessary. If a scalar intended interaction CNOTs a correlated old $E$ into new $F$, a Bell input on $L,E$ produces an informative new archive although that interaction is the identity on the physical code. A scalar old/new coupling defect of norm $\eta$ adds at most $\eta T$ by Duhamel; increasing $\Delta$ does not improve it. Preparation defects likewise require a complete-state bound. A correct new marginal does not establish independence from the source or the actual past.

A charge-preserving scalar valve can remain responsive:
\[
 H_{\rm valve}=\kappa I_{\rm phys}\otimes
 (|\mathrm{out}\rangle\langle\mathrm{in}|+
  |\mathrm{in}\rangle\langle\mathrm{out}|),
 \qquad T_{\rm tr}=\pi/(2\kappa).
\]
Its transfer time does not grow with $\Delta$. At earlier cuts its actual ideal path vector is
$\cos(\kappa t)|\mathrm{in}\rangle-i\sin(\kappa t)|\mathrm{out}\rangle$;
reflection and incomplete transfer have not been replaced by a completed outlet. The complete bound counts $\kappa$ in $b$.

\section{Adversarial tests and a finite-horizon obstruction}

For $H_0=0$, $V=gX_1$, a code vector and its $X_1$ image form an invariant pair with matrix
\[
 \begin{pmatrix}0&g\\g&\Delta\end{pmatrix}.
\]
Without the penalty its leakage is $\sin^2(gt)$, reaching one at $\pi/(2g)$. With the penalty leakage is
\begin{equation}
 \frac{4g^2}{\Delta^2+4g^2}
 \sin^2\!\left(\tfrac t2\sqrt{\Delta^2+4g^2}\right).
 \label{prot:singleleak}
\end{equation}
This is suppression against the same damaging interaction. Leakage alone is insufficient to certify logical fidelity, as the next exact example shows.

\begin{proposition}[Orthogonal logical return at zero leakage]
\label{prot:horizon}
For the four-qubit code, take $H_0=0$, $V=g(X_1+X_2)$ and $v=2g$. There are arbitrarily small $v/\Delta$ and finite $T$ at which an encoded vector returns to the code with zero leakage and an orthogonal logical state, while $Tv^2/\Delta\to\pi$.
\end{proposition}
\begin{proof}
The code operator $A_L=X_1X_2$ is a nontrivial logical involution. On its $-1$ eigenspace $X_2\psi=-X_1\psi$, so $V$ vanishes. On its $+1$ eigenspace the pair $\psi,X_1\psi$ has matrix
\[
 \begin{pmatrix}0&2g\\2g&\Delta\end{pmatrix},
 \quad\Omega=\sqrt{\Delta^2+16g^2}.
\]
Its code return amplitude is
\[
 a(t)=e^{-i\Delta t/2}
 [\cos(\Omega t/2)+i\Delta\sin(\Omega t/2)/\Omega].
\]
For integer $k\ge2$, choose
\[
 16g^2/\Delta^2=(2k-1)/(k-1)^2,
 \qquad T=2\pi(k-1)/\Delta.
\]
Then $\Omega/\Delta=k/(k-1)$, leakage is zero and $a(T)=-1$. An equal superposition of a dark and bright logical vector becomes orthogonal. Finally
$Tv^2/\Delta=\pi(2k-1)/[2(k-1)]\to\pi$ while $v/\Delta\to0$.
\end{proof}

A later noncommuting logical probe separates these states even though a syndrome-only inspection sees no leakage. Virtual excursions accumulate a logical phase. The $tv^2/\Delta$ term in the complete estimate therefore marks a real horizon, not just a proof artifact.

A two-body logical perturbation $\xi Z_1Z_2$ commutes with the four-qubit penalty and acts inside the code. On a logical superposition it changes a suitable later probability by $\sin^2(\xi t)$, independently of $\Delta$. A finite resonant memory is another explicit boundary: prepare $E$ excited with $H_E=\Delta|1\rangle\langle1|$ and use $V=gX_1\otimes X_E$. The states $\psi_{\rm code}|1\rangle$ and $X_1\psi_{\rm code}|0\rangle$ are degenerate in total energy and mix with amplitude $\sin(gt)$ regardless of $\Delta$. Here $b$ grows with $\Delta$, violating the required separation. A drive resonant with the gap similarly lies outside the stationary bounded-bandwidth class.

\section{Completed repair against the two-body attack}

Use the five-qubit code with commuting independent generators
\begin{equation}
 g_1=XZZXI,\quad g_2=IXZZX,\quad
 g_3=XIXZZ,\quad g_4=ZXIXZ,
 \quad P_5=\prod_{k=1}^4(I+g_k)/2.
 \label{prot:fivecode}
\end{equation}
It has rank two. Explicit nonzero logical vectors are the normalized $P_5|00000\rangle$ and its $X^{\otimes5}$ image: the initial projection has squared norm $1/16$, and the two vectors have opposite $Z^{\otimes5}$ eigenvalues.

\begin{lemma}[Detection of every weight-one and weight-two Pauli]
\label{prot:twobody}
For each nonidentity Pauli $A$ of weight at most two, $P_5AP_5=0$.
\end{lemma}
\begin{proof}
The anticommutation syndromes against $g_1,\ldots,g_4$ are
\[
\begin{array}{c|ccc}
\text{site}&X&Y&Z\\\hline
1&0001&1011&1010\\
2&1000&1101&0101\\
3&1100&1110&0010\\
4&0110&1111&1001\\
5&0011&0111&0100
\end{array}
\]
These are all fifteen nonzero four-bit strings. A Pauli acting on two different sites has the exclusive-or of two different syndromes, which is nonzero. Hence some stabilizer anticommutes with $A$. The calculation $P_5AP_5=-P_5AP_5$ proves the claim.
\end{proof}

With $H_{{\rm pen},5}=\Delta(I-P_5)$ and
\[
 V=\sum_{1\le\operatorname{wt}(A)\le2}A\otimes B_A,
 \qquad B_A=B_A^\dagger,\quad\|V\|\le v,
\]
Theorem~\ref{prot:gap} applies unchanged. This repair genuinely enlarges the admitted disturbance class. It requires a stronger supplied interaction: expanding $P_5$ uses up to fifteen commuting nonidentity Pauli products besides a scalar term. The penalty has norm $\Delta$. One unknown logical qubit and four fresh ready qubits suffice for a unitary encoding; extending the isometry to a unitary and using a Hermitian logarithm gives finite realization with norm at most $\pi/\tau$ over duration $\tau$. That establishes finite existence, not an optimized circuit.

No finite code protects against every possible logical interaction. An operator implementing $\epsilon Z$ inside its code leaves the penalty unchanged and rotates a logical $|+\rangle$ by trace distance $|\sin\epsilon t|$. This is the exact boundary of a code-based protection claim.

\chapter[Controls, keys, and operational protection]{Finite controls, retained keys, and operational protection bounds}
\label{prot:integrationchapter}

\section{Bounded-strength averaging on the complete bank}

A complementary mechanism handles time-dependent but slowly varying nuisance interactions. Let $L=\mathbb C^d$, $E$ be the complete apparatus/memory bank, and $R$ an inaccessible reference. At each actual control history use
\begin{equation}
 i\dot\Psi=[I_L\otimes B(t)+H_c(t)\otimes I_E+V(t)]\Psi.
 \label{prot:controlsource}
\end{equation}
In the interaction picture of $B$, assume $\|V_I(t)\|\le J$ and
$\|V_I(t)-V_I(s)\|\le\ell|t-s|$. These bounds are uniform in retained keys. Old correlated memories remain in $E$.

Let $X,Z$ be the $d$-dimensional Weyl operators and let $\mathcal G=\{X^aZ^b\}_{a,b=0}^{d-1}$ modulo phases. Their twirl is
\[
 \Pi(A)=I_L\otimes\operatorname{Tr}_L(A)/d.
\]
Take an Eulerian cycle in the directed Cayley graph with generators $X,Z$. It uses $2d^2$ edges. Implement an edge $\alpha$ from vertex $g$ by $u_\alpha(s)g$, $0\le s\le\tau_p$, where $u_\alpha(0)=I$, $u_\alpha(\tau_p)=\alpha$. A Hermitian logarithm implements each finite pulse at strength at most $\pi/\tau_p$. The full duration is $\tau=2d^2\tau_p$ and the control frame $g_c$ returns to the identity modulo phase. For fixed $A$,
\begin{align}
 \int_0^\tau g_c(s)^\dagger A g_c(s)ds
 &=\sum_{\alpha=X,Z}\int_0^{\tau_p}
       \sum_{g\in\mathcal G}g^\dagger u_\alpha(s)^\dagger A u_\alpha(s)g\,ds\\
 &=\tau\Pi(A).
 \label{prot:eulertwirl}
\end{align}
Thus cancellation includes finite pulse time; it is not an instantaneous-pulse approximation.

\begin{theorem}[Finite-control complete propagator estimate]
\label{prot:controlbound}
Let $U$ solve \eqref{prot:controlsource}, and let $U_{\rm sc}$ use the same $B$ with $V_I$ replaced by $\Pi(V_I)$. At $T=N\tau$,
\begin{equation}
 \|U(T)-(g_c(T)\otimes I_E)U_{\rm sc}(T)\|
          \le T\tau(J^2+\ell).
 \label{prot:cyclebound}
\end{equation}
At an arbitrary cut $t$, the corresponding bound is
$t\tau(J^2+\ell)+2J\tau$. Both hold for arbitrary inaccessible references and all retained output registers.
\end{theorem}
\begin{proof}
In the toggling and $B$ frames the generators are
$A(t)=g_c(t)^\dagger V_I(t)g_c(t)$ and $D(t)=\Pi(V_I(t))$. Freeze $V_I$ at the start of one cycle. Equation~\eqref{prot:eulertwirl} cancels the frozen integrals. The variation bounds on both remaining terms give
$\|\int(A-D)dt\|\le\ell\tau^2$.
For a self-adjoint generator $K$ with norm at most $J$, substitute its Volterra equation once:
\[
 \|U_K(\tau)-I+i\int_0^\tau K(s)ds\|\le J^2\tau^2/2.
\]
Unitarity bounds the remaining double integral. Applying this to $A,D$ makes the one-cycle error at most $(J^2+\ell)\tau^2$. Telescoping unitary products proves \eqref{prot:cyclebound}; Duhamel adds at most $2J\tau$ for the unfinished fragment.
\end{proof}

For fresh independent $E$ and scalar comparator evolution, the full output is close to $\Psi_{LR}\otimes a_T$. For an old correlated $E$ it is close to its specified comparator, which need not be independent. As an unprotected comparison, $V=JZ_L\otimes Z_E$ with $E$ in a $Z_E$ eigenstate rotates $|+\rangle$ to distance $|\sin JT|$. At $T=\pi/(2J)$ that distance is one, whereas the static protected bound is $\pi J\tau/2$.

The construction uses control strength at most $2\pi d^2/\tau$ and $2d^2T/\tau$ pulses. Complete pulse errors $\delta_k$ add at most $\sum_k\delta_k$. The classical schedule is a supplied resource; any physical clock backaction must enter the complete generator estimates. A disturbance $V_I(t)=g_c(t)V_0g_c(t)^\dagger$ becomes $V_0$ in the toggling frame and defeats averaging. Its variation grows with the control speed, so it is excluded by the uniform $\ell$ premise, not by its small norm alone.

\section{Actual event times and retained control phases}

At an event time $u$ in the middle of a cycle, stopping the controller leaves $g_c(u)\Psi$ even if $V=0$. Recording the control phase does not undo that rotation. A finite repair continues to control the routed carried bank until the next cycle boundary, adding latency at most $\tau$. When the aperture's outlet action commutes with those controls and the pre/post-event nuisance bounds remain valid, the two split fragments together last at most $\tau$. Hence a conservative complete handoff amplitude error is
\begin{equation}
 e_{\rm hand}\le T\tau(J^2+\ell)+2J\tau+\sum_k\delta_k.
 \label{prot:handoff}
\end{equation}
The ideal comparator has the same latency. If destructive capture makes the carried bank inaccessible, this repair is unavailable. Nulls, phase records, pending transport, and exhausted pulse stocks remain actual branches.

Irreducible averaging also removes a desired nonscalar sector-controlled loader. Such a loader must be performed separately, included in a protected code-preserving $H_0$, or handled by sector-wise controls with their surviving sector-dependent phases explicitly retained. Averaging cannot silently remove a nuisance while preserving an algebraically identical desired coupling.

The passive gap estimate avoids this rotating-frame problem because it is uniform at every time. Its use at an actual event still requires an outlet-adoption rule: the event must retain the actual transported row and its banks. A separately appended logical unitary at the event is an additional reaction outside the bounded pre-event Hamiltonian comparison.

\section{Coherent reversal using a physically retained history}

A different active construction is possible when an orthogonal history register is available. Let $K$ have states $|\alpha\rangle$ and retain every old correlation in $\Psi_{LKR}$. Suppose
\[
 H_{\rm err}(t)=\sum_{\alpha=1}^r|\alpha\rangle\langle\alpha|\otimes H_\alpha(t),
 \qquad U_{\rm err}=\sum_\alpha|\alpha\rangle\langle\alpha|\otimes U_\alpha.
\]
A controller with access to $K$ can apply, over duration $T_c$,
\[
 H_c(s)=-\frac T{T_c}\sum_\alpha|\alpha\rangle\langle\alpha|
             \otimes H_\alpha(T-Ts/T_c).
\]
Changing variables in the time-ordered exponential gives $U_c=U_{\rm err}^\dagger$, so the complete old state returns exactly. No syndrome is measured and no Born reset enters. The inverse uses the same integrated action, with larger strength if shortened. For an implementation error $\Delta H(t)$, Duhamel bounds the full operator error by $\int\|\Delta H(t)\|dt$, uniformly for the inaccessible reference and all old correlations.

\begin{proposition}[Distinct reversible errors require distinguishable programs]
\label{prot:programs}
If a unitary corrector using programs $|k_\alpha\rangle$ restores every unknown carried vector after errors $U_\alpha$, then programs for errors distinct up to scalar phase must be orthogonal.
\end{proposition}
\begin{proof}
Include every fresh auxiliary state in the program and every retained output in $|a_\alpha\rangle$. Exact preservation of every carried pure vector and linearity on superpositions imply
$C(U_\alpha|\psi\rangle|k_\alpha\rangle)=|\psi\rangle|a_\alpha\rangle$ with input-independent remainder. Preservation of inner products for arbitrary $\psi,\phi$ gives
\[
 \langle k_\alpha|k_\beta\rangle U_\alpha^\dagger U_\beta
                 =\langle a_\alpha|a_\beta\rangle I.
\]
Nonzero program overlap makes the relative error scalar. For $U_0=I$, $U_1=Z$, choose $|+\rangle,|-\rangle$: the two correction inputs have overlap $c=|\langle k_0|k_1\rangle|$ and their faithful outputs have orthogonal carried factors. If both complete vector errors are at most $\epsilon$, the inner-product triangle inequality gives $c\le2\epsilon$.
\end{proof}

A hidden orthogonal key in a returning memory can enable correction later while an accessible nonorthogonal proxy cannot enable it now. The physical cut and accessible couplings must therefore remain explicit. An event-safe implementation withholds readiness until $T+T_c$; an early-event intensity bounded by $\ell(t)$ costs at most $1-\exp[-\int_0^{T+T_c}\ell(t)dt]$. This is a latency/resource tradeoff. It is not protection of a payload that has already become inaccessible after destructive capture.

\section{A source-law modulus before probability selection}

Complete coherent closeness cannot automatically be propagated through an arbitrary nonlinear actualization law. The following finite theorem provides a sufficient modulus within a specified response class, without assuming density-matrix sufficiency or complete positivity.

At a source cut write $x=(c,[\psi])$, including every classical record in $c$ and every coherent reference/memory in $\psi$. Let $d(x,y)=1$ if classical data differ and otherwise
$D([\psi],[\phi])=\sqrt{1-|\langle\psi,\phi\rangle|^2}$. Use the Wasserstein distance $W_1$ induced by this bounded metric on actual source laws. Complete classical--quantum trace distance is bounded by $W_1$, but a small difference of ensemble-averaged density matrices need not make $W_1$ small.

Let $E_j$ be orthogonal outlet rows, $\sum_jE_j^\dagger E_j=I$. A stopped reader uses intensities $h(p_j)$, $p_j=\|E_j\psi\|^2$, normalized transported daughters, and common source-independent scalar marks. Assume $h(0)=0$, $0\le h$, and $\operatorname{Lip}(h)\le L$ on $[0,1]$. Held states may undergo the same prescribed coherent drive, preserving their ray distance. This is a declared family of unselected readers.

\begin{lemma}[Weighted normalization of outlets]
\label{prot:weighted}
For $a_j=E_j\psi$, $b_j=E_j\phi$, $p_j=\|a_j\|^2$, $q_j=\|b_j\|^2$,
\begin{equation}
 \sum_j\min(p_j,q_j)D([a_j],[b_j])\le D([\psi],[\phi]).
 \label{prot:weightedineq}
\end{equation}
Terms with zero branch mass vanish.
\end{lemma}
\begin{proof}
Set $c_j=|\langle a_j,b_j\rangle|$ and $z_j=\sqrt{p_jq_j-c_j^2}$. Each summand is at most $z_j$. The Euclidean triangle inequality and Cauchy--Schwarz give
\[
 \sqrt{(\sum_jz_j)^2+(\sum_jc_j)^2}
 \le\sum_j\sqrt{p_jq_j}\le1.
\]
Since $\sum_jc_j\ge|\langle\psi,\phi\rangle|$, the result follows. This is vector geometry; no outcome probability law has entered the lemma.
\end{proof}

\begin{theorem}[Complete finite response modulus]
\label{prot:modulus}
For the stopped $n$-outlet reader of duration $s$ just defined,
\begin{equation}
 W_1(\mathsf K_s(x),\mathsf K_s(y))
       \le\min\{1,[1+(n+1)Ls]d(x,y)\}.
 \label{prot:modulusbound}
\end{equation}
Every event time, null, mark, retained daughter and reference belongs to the compared output.
\end{theorem}
\begin{proof}
For matching classical data, couple channel clocks with common rates $\min\{h(p_j),h(q_j)\}$ and their excesses. Since $|p_j-q_j|\le d(x,y)$, the probability of an unmatched event is at most $nLs\,d(x,y)$. Common no-event evolution preserves the distance. At a common event the rate is at most $L\min(p_j,q_j)$ because $h(0)=0$. Lemma~\ref{prot:weighted} bounds the integrated daughter contribution by $Ls\,d(x,y)$. A common null contributes at most $d(x,y)$. Common conditional scalar marks can be coupled identically; an unmatched event costs at most one. Adding the contributions proves the theorem. Differing classical inputs use the trivial bound one.
\end{proof}

Suppose physical and comparator no-event rays differ by at most $a(t)$ uniformly, and the physical event row differs from its transported normalized row by at most $b_*$ on every nonzero common branch. Let $\zeta_s$ bound separately the full classical mark/resource mismatch. The same coupling proves
\begin{equation}
 \epsilon_s\le\min\left\{1,
 a(s)+(n+1)L\int_0^sa(t)dt+b_*+\zeta_s\right\}.
 \label{prot:slot}
\end{equation}
The gap theorem supplies $a(t)\le(2v+tv^2)/\delta$; a gate error $u_{\rm enc}$ adds to this amplitude bound. No small reduced-state error is substituted for the required complete-source comparison.

Before selection, a finite sequence with slot errors $\epsilon_k$ and proved suffix moduli $C_\ell$ has error at most
$\sum_k\epsilon_k\prod_{\ell>k}C_\ell$. This follows by replacing one slot at a time and applying each subsequent modulus. An arbitrary nonlinear future law can violate every uniform modulus; protection alone then gives no operational conclusion.

The finite rate-selection and tag arguments can use \eqref{prot:slot} only with their own nonzero response factors, interchange premises, and time regularity. Dividing a tag error by its success probability cannot be omitted. A finite event discrepancy does not bound infinitesimal Gaussian coefficients without uniform time estimates. Thus this chapter propagates physical errors into those arguments but does not derive their statistical premises.

\section{Complete stopped output after the rate has been selected}

Now retain the selected interface $\lambda_j=\gamma\|E_j\Phi\|^2$, with faithful outlet adoption. Let every live mode be registered, $\sum_jE_j^\dagger E_j=P_{\rm live}$, and let the full coherent dynamics preserve live charge. A first event terminates exposure to the nuisance; later storage and decoding are common unitaries or separately charged. This is a stronger downstream hypothesis than Theorem~\ref{prot:modulus}.

\begin{proposition}[Stopped complete-output bound]
\label{prot:stopped}
For a live code input and arbitrary inaccessible reference, the complete event/time/null/continuation output at deadline $T$ differs from the selected ideal by at most
\begin{equation}
 \epsilon_{\rm stop}(T)\le
 \min\left\{1,\frac{2v+(v^2/\gamma)(1-e^{-\gamma T})}{\delta}\right\}.
 \label{prot:stoppedbound}
\end{equation}
Total click and null probabilities are exactly $1-e^{-\gamma T}$ and $e^{-\gamma T}$. Individual outlets and their conditional states may differ.
\end{proposition}
\begin{proof}
The declared law yields subnormalized event and null vectors
\[
 K_j(u)\psi=\sqrt\gamma e^{-\gamma u/2}E_jU_\Delta(u)\psi,
 \qquad K_\varnothing\psi=e^{-\gamma T/2}U_\Delta(T)\psi.
\]
Summing $E_j^\dagger E_j$ proves normalization and constant total hazard. Append orthogonal record labels before removing their mutual coherences. That removal is an average of unitary phase conjugations and cannot increase trace norm. The event-output distance is bounded by
$\gamma e^{-\gamma u}e_\Delta(u)du$, and the null distance by $e^{-\gamma T}e_\Delta(T)$. Subsequent common unitaries preserve the complete distances. Integrate the affine gap bound and use
\[
 \int_0^T\gamma e^{-\gamma u}u\,du+Te^{-\gamma T}
            =(1-e^{-\gamma T})/\gamma.
\]
This gives \eqref{prot:stoppedbound} without discarding any environment in the proof.
\end{proof}

Stopping is physical routing out of the nuisance region, not a consequence of merely observing a click. If the event occurs during loading, readiness and incomplete-load modes also need records. For a code-preserving loader
\[
 H_{\rm load}=\kappa\sum_jQ_j\otimes(|j\rangle\langle r|+|r\rangle\langle j|),
 \qquad\tau=\pi/(2\kappa),
\]
target-$j$ density before $\tau$ is
$\gamma e^{-\gamma u}\sin^2(\kappa u)\|Q_j\Psi\|^2du$ and ready-mode failure density is $\gamma e^{-\gamma u}\cos^2(\kappa u)du$. Integrating the latter gives
\begin{equation}
 f_\tau=\frac{1-e^{-\gamma\tau}}2+
       \frac{\gamma^2(1+e^{-\gamma\tau})}{2(\gamma^2+4\kappa^2)}.
 \label{prot:loadingfailure}
\end{equation}
Holding to $s\ge\tau$ gives useful target success $r_s=1-f_\tau-e^{-\gamma s}>0$. The failed branch keeps its actual source, and a null keeps the full loaded superposition until an actual inverse is applied.

For piecewise-stationary segments of lengths $d_\ell$ starting at $t_{\ell-1}$, unitary telescoping and the same stopping integration give
\begin{equation}
 \epsilon_{\rm stop}\le\min\left\{1,
 \sum_\ell e^{-\gamma t_{\ell-1}}
 \left[\frac{2v_\ell}{\delta_\ell}
 +\frac{v_\ell^2}{\gamma\delta_\ell}(1-e^{-\gamma d_\ell})\right]\right\}.
 \label{prot:segments}
\end{equation}
The ideal prefix preserves the code; actual suffixes preserve norm. One must not assume that an actual imperfect prefix remains in the code.

\section{Reuse, references, and finite network accounting}

After an imperfect block, decode unitarily into logical $L$ and syndrome $\Sigma$, retaining all of $\Sigma$. Append fresh independently prepared ancillas $F$ and re-encode $L,F$. The new block lies exactly in the code even when $L$ is entangled with $\Sigma$ and every old bank. This restores the domain, not the correct logical state. It neither measures a syndrome nor resets the old memory. Four-qubit blocks consume two fresh ready qubits and retain two old syndrome qubits per renewal; five-qubit blocks have the corresponding four-qubit costs.

If the physical encoder differs from its ideal unitary by $u_{\rm enc}$ in operator norm, compare the output with the exact encoding of the same actual logical/syndrome input. The comparison error is at most $u_{\rm enc}$. The gap theorem applies to the exact encoded comparison input and the actual no-event unitary preserves the initial comparison error. A decoder error similarly adds its norm defect. This closes the actual-input domain needed for a finite replacement argument without assigning fictitious bad-preparation probabilities to coherent leakage.

After an ideal instrument has actually been established, a common ideal suffix contracts classical--quantum trace distance. An actual-prefix/ideal-suffix replacement then gives the complete finite-network bound
\begin{equation}
 \epsilon_{\rm net}\le\min\left\{1,
 \epsilon_{\rm prep}+p_{\rm exhausted}+\epsilon_{\rm cut}
 +\sum_k(\epsilon_{\rm aperture,k}+\epsilon_{\rm gates,k}
                 +\epsilon_{\rm rate,k}+\zeta_k)
 +\epsilon_{\rm instrument}\right\}.
 \label{prot:network}
\end{equation}
Every local bound is uniform on the actual complete input and history. Use either the preselection bound with its proved moduli or the sharper selected bound, not both for the same defect. The term $\epsilon_{\rm instrument}$ contains only the independently established loading/readout approximation. Timeout is an error only against a completed sharp measurement; a null-inclusive finite instrument already contains its timeout branch exactly.

For a growing bank, use its total $b_N,v_N$ and require $\Delta_N>2b_N+v_N$. A fixed single-memory bound cannot be reused after arbitrarily many retained resources accumulate. All error sums must tend to zero in a simultaneous limit. Readiness stock must be conditionally independent of the complete actual past or carry a quantified preparation defect; correct individual marginals are insufficient.

For an event of actual and ideal probabilities $p,q>0$, complete error $\epsilon$ implies conditional trace distance at most
\begin{equation}
 \min\{1,\epsilon/\max(p,q)\}.
 \label{prot:rare}
\end{equation}
To see this, decompose the complete output into event/complement blocks. For $p\ge q$, write the event blocks as $p\rho,q\sigma$. Then
$p\|\rho-\sigma\|_1\le\|p\rho-q\sigma\|_1+p-q$,
while the complement norm is at least $p-q$. Dividing the total trace norm by two proves the claim; exchange $p,q$ for the other case. There is no uniform perfect daughter on an event whose probability vanishes with resources.

\section{What the protection mechanism leaves fundamental}

A perfectly protected quantum carrier can coexist with the scalar intrinsic clock
\begin{equation}
 \lambda=\gamma[1+\alpha(2\operatorname{tr}\rho_L^2-1)],
 \qquad0<\alpha\le1,
 \label{prot:purityclock}
\end{equation}
with outlet marks proportional to their populations. On the equal Bell-state ensemble and equal computational-product ensemble of two qubits, both averaged density matrices are $I_4/4$, while each reduced purity is respectively $1/2$ and $1$. Their finite click gap is
\begin{equation}
 G(s)=e^{-\gamma s}-e^{-\gamma(1+\alpha)s},
 \qquad\max_sG(s)=\alpha(1+\alpha)^{-1-1/\alpha}.
 \label{prot:clockgap}
\end{equation}
Differentiation gives the maximizing time $\log(1+\alpha)/(\gamma\alpha)$. A local coherent encoding preserves nonzero reduced eigenvalues, so it does not change this gain. A printed scalar mark $z=\operatorname{tr}\rho_L^2$ similarly leaks source-ensemble information with zero coherent transport defect.

These are explicitly different stochastic laws, not Hamiltonian meters implemented by \eqref{prot:model}. They violate the population-only response or scalar-mark premises of the selected interface. If preparation labels are actively retained, they must remain in the complete input and the ensembles need not be operationally equivalent. The counterexample shows exactly why carrier protection alone cannot establish that equivalence.

The established reduction is physical suppression of nonscalar coherent nuisance interactions on a declared locality, bandwidth, horizon and resource domain. Exact carrier neutrality is replaced by a code, finite gap or finite control schedule, and a complete error estimate. Protection alone does not select outlet adoption, event statistics, fresh-resource preparation or a universal access rule. An event-only rewrite or input-sensitive scalar clock is not ruled out by an intact protected propagator. The later complete theories supply their own actual dynamics and admitted material interactions; these protection bounds apply to them only when their bounded-operator, encoding and time-domain hypotheses are verified. In particular, an unbounded massive kinetic Hamiltonian is not covered merely because its ready states have finite energy.

% End consolidated source: parts/06_protection.tex

\part{Configuration Records, Continuation, and Faithful Archives}\label{part:records}

% Begin consolidated source: parts/07_configuration.tex
\chapter{Configuration records and coherent continuation}
\label{cfg:chapter}

This chapter consolidates the configuration constitution of \cite{M18},
the finite record interactions of \cite{M30}, and the corrected coherent
continuation results of \cite{C02}. Its principal result is a complete
finite-record construction with one actual configuration and an uncollapsed
wave, as an alternative to intrinsic absorbing extraction. Here the
equivariant process and joint ready law remain explicit inputs. The later
pilot theory derives an effective Bell process from its finite microscopic
model and gives an autonomous material implementation. No argument
identifies an actual configuration jump with a global projection of the wave.

\section{Complete source and the two meanings of continuation}

Let $\mathcal H_S$ carry one unknown system, let $\mathcal H_R$ be an
inaccessible reference, and include all returning quantum memories in the
carried space $\mathcal B$. There is no control on $R$. The complete source is
$(c,\Phi,Q)$: declared classical provenance and controls $c$, a normalized
wave $\Phi$, and one actual apparatus configuration $Q$. A newly acquired
record has a physical carrier in $Q$ or in a jointly modeled configuration
register. A mathematical past trajectory is not an additional accessible
archive.

Two realizations will be used. In the continuous realization,
$\Phi\in L^2(\mathcal Q;\mathcal B)$ and the specified linear wave dynamics
have a conserved density and current,
\begin{equation}
 \rho_t(q)=\|\Phi_t(q)\|^2,\qquad
 \partial_t\rho_t+\operatorname{div}J_t=0,\qquad
 \dot Q_t=J_t(Q_t)/\rho_t(Q_t).
 \label{cfg:guidance}
\end{equation}
An admitted programme must have a unitary wave propagator and an almost-sure
conservative guidance flow. These analytic properties are checked explicitly
for the modules below; they are not asserted for arbitrary singular
coefficients. A sufficient differential expression is
\[
 H_t=-i\hbar\sum_k\left(A_k(q,t)\partial_{q_k}
             +\tfrac12\partial_{q_k}A_k(q,t)\right)+V(q,t),
 \qquad A_k=A_k^\dagger,\quad V=V^\dagger,
\]
on a domain supplying those properties, with $J_k=\Phi^\dagger A_k\Phi$.

In the discrete realization, $\{P_x:x\in\mathcal Q\}$ is a finite
orthogonal configuration resolution, $w_x=\|P_x\Phi\|^2$, and
\[
 J_{yx}=\frac2\hbar\operatorname{Im}
          \langle P_y\Phi,H P_x\Phi\rangle.
\]
The actual process is supplied by an admitted conservative equivariant law.
The minimal Bell law is one such choice, conditional on the statistical
selection developed elsewhere in this monograph. For comparison we also use
the nonexplosive Markov family
\begin{equation}
 q_{y\leftarrow x}(t)=\frac{[J_{yx}(t)]_++K_{xy}(t)}{w_x(t)},
 \qquad K_{xy}=K_{yx}\ge0,
 \label{cfg:Kfamily}
\end{equation}
where occupied-state rates are used and node behavior satisfies the
well-posedness assumptions. Finite expected integrated traffic is a sufficient
nonexplosion condition in equilibrium. The pairwise family
\eqref{cfg:Kfamily} does not exhaust every equivariant process: divergence-free
reassignments of net edge currents are additional possibilities. The present
comparison requires the displayed pairwise matching.

\begin{definition}[Relative branch and complete continuation]
For a physical record region $\Gamma_h$, the relative branch is the normalized
restriction $P_h\Phi/\|P_h\Phi\|$, when its norm is nonzero. The complete
continuation is the actual global wave, actual configuration or its conditional
law, and all retained controls and memories. A relative branch is a sufficient
replacement only on a proved future domain in which omitted branches cannot
affect the predictions being claimed.
\end{definition}

For continuous local guidance, disjoint supports preserved by future
propagation supply such a domain. For discrete processes, preservation of a
record label does not by itself ensure that rates evaluated on the global wave
equal rates evaluated on its normalized projection: an added $K$ may depend
on amplitudes in other blocks. Endpoint quantum predictions proved below do
not need that additional generator identity. Claims about an autonomous
daughter process do.

\section{A finite binary detector with exact time and null laws}

Let $Q_++Q_-=I$ be orthogonal system projectors, extended by the identity
on $\mathcal H_R$ and all carried memories. Degenerate outcomes are permitted.
Take a normalized even $H^1$ packet $\phi$, supported on $[-a,a]$, and a fresh
pointer $x\in\mathbb R$. The wave and actual ready law are
\[
 \Phi_0(x)=\phi(x)\psi,\qquad
 \mathbb P(X_0\in dx)=|\phi(x)|^2dx,
\]
conditionally on the complete actual past and independently of the unknown
$\psi$. For $u>0$, define
\begin{equation}
 H=u(Q_+-Q_-)P_x,\qquad P_x=-i\hbar\partial_x.
 \label{cfg:translationH}
\end{equation}
The exact field, density, and current are
\begin{align}
 \Phi_t(x)&=\phi(x-ut)Q_+\psi+\phi(x+ut)Q_-\psi,\\
 \rho_t(x)&=p_+ f(x-ut)+p_- f(x+ut),\\
 J_t(x)&=u[p_+f(x-ut)-p_-f(x+ut)],
 \qquad p_j=\|Q_j\psi\|^2,\quad f=|\phi|^2.
 \label{cfg:split}
\end{align}
All operators act trivially on the inaccessible reference.

\begin{lemma}[Quantile flow]\label{cfg:quantile}
Put $F_t(x)=\int_{-\infty}^x\rho_t(z)\,dz$. Almost every trajectory
has $F_t(X_t)=U$, where $U=F_0(X_0)$ is uniform on $(0,1)$.
The generalized inverse $X_t=F_t^{-1}(U)$ realizes the guidance flow
and has density $\rho_t$.
\end{lemma}
\begin{proof}
The continuity equation gives $\partial_tF_t=-J_t$.
At positive-density points, differentiating the fixed-quantile identity gives
$\dot X_t=J_t/\rho_t$. The inverse distribution transform gives its marginal
law. Plateau values correspond to a null set of uniform ranks; the smooth
interior packets used here give continuous trajectories away from those
exceptional ranks. Thus the definition supplies a nonempty almost-sure
flow without adding a second random tape.
\end{proof}

Place exit surfaces at $\pm L$, with $L>a$, and let $\tau$ be the first
exit from $(-L,L)$. The outward current at either surface has only one
nonzero packet. Define
\[
 r(t)=\int_0^t u f(L-us)\,ds.
\]
It vanishes before $(L-a)/u$ and equals one after $(L+a)/u$.

\begin{theorem}[Finite exit and retained local field]\label{cfg:exit}
The detector \eqref{cfg:translationH} obeys
\begin{equation}
 \mathbb P(j,\tau\in dt)=p_j u f(L-ut)\,dt,\qquad
 \mathbb P(j,\tau\le t)=p_jr(t),\qquad
 \mathbb P(\tau>t)=1-r(t).
 \label{cfg:exitlaw}
\end{equation}
At a nonzero-probability exit, the normalized local carried field is
$Q_j\psi/\sqrt{p_j}$ up to a global phase. At a null deadline with
$1-r(t)>0$, the actual conditional position law is
\begin{equation}
 \frac{1_{\{|x|<L\}}\rho_t(x)\,dx}{1-r(t)};
 \label{cfg:liveposition}
\end{equation}
the global field remains \eqref{cfg:split}.
\end{theorem}
\begin{proof}
At $+L$, the negative packet is absent and the outward current is
$up_+f(L-ut)$. At $-L$ the outward current is
$up_-f(-L+ut)=up_-f(L-ut)$ by evenness. Neither surface can be crossed
inward while this translation interaction remains active. Equivariance from
Lemma~\ref{cfg:quantile} therefore identifies first-exit probability with
the integrated outward current. Their sum gives the null probability and
restriction of the equilibrium density gives \eqref{cfg:liveposition}.
Only the indicated sector component of $\Phi$ is present at an exit,
so normalization yields the stated relative carried vector.
\end{proof}

This common interaction joins time, outlet, and local continuation.
Its observed hazard is
$p_jr'(t)/(1-r(t))$, conditional on the ready preparation and no observed
exit. Given the complete actual initial configuration, the exit is
deterministic. This is not the original Hamiltonian Bell conditional
jump intensity, nor an exponential clock in another notation.

A persistent position can encode the exit time. At a later cut,
\[
 R_t(x)=
 \begin{cases}
 \mathrm{null},&|x|<L,\\
 (\operatorname{sign}x,\ t-(|x|-L)/u),&|x|\ge L.
 \end{cases}
\]
Finite position bins give finite timestamp bins. If the pointer is reversed,
a lasting timestamp requires an actual copy interaction. The formula does
not provide an immutable external history.

The live relative sector populations remain $p_j$, but the coherent null
is generally different from $\psi$. Integrating over live positions gives
the unnormalized internal matrix with $(j,k)$ block
\[
 Q_j|\psi\rangle\langle\psi|Q_k
 \int_{-L}^L\phi(x-s_jut)\phi^*(x-s_kut)\,dx,
 \qquad s_\pm=\pm1.
\]
At $t>a/u$ the off-diagonal blocks vanish. A frozen-null ray from an
intrinsic extraction model therefore cannot replace this source in a return
experiment.

For an explicit finite-gradient packet,
\[
 \phi(x)=a^{-1/2}\cos(\pi x/(2a))1_{\{|x|\le a\}},
\]
direct integration gives
\begin{equation}
 \mathbb E\tau=\frac Lu,\qquad
 \operatorname{Var}(\tau)=\frac{a^2}{u^2}
          \left(\frac13-\frac2{\pi^2}\right),\qquad
 \|\phi'\|_2^2=\frac{\pi^2}{4a^2}.
 \label{cfg:timecost}
\end{equation}
Thus $u^2\operatorname{Var}(\tau)\|\phi'\|_2^2=\pi^2/12-1/2$ in this
family. This is a concrete time/gradient tradeoff, not a new universal
uncertainty relation. The transport generator is selfadjoint but unbounded
and not lower bounded; its use is an effective source sector, not a
microscopic stability theorem.

Loading errors have real branches. If a gated load prepares
\[
 \phi(x)\left[\cos\theta\,\psi|r\rangle
       -i\sin\theta\sum_jQ_j\psi|j\rangle\right]
\]
and mode $r$ is stationary while modes $j$ propagate at $s_ju$, then
$\mathbb P(j,\tau\le t)=\sin^2\theta\,p_jr(t)$.
After the traveling packets have exited, the surviving null is the stationary
ready component, with probability $\cos^2\theta$. This calculation assumes
the propagation gate is closed during loading. It does not cover simultaneous
loading and propagation by replacing their actual dynamics with these
products.

\section{Acquisition changes the packet that controls the exit}

A physical snapshot is a defined wire, not a claim that every hidden
coordinate is freely readable. Add a pointer with known packet $\chi(y)$,
jointly equilibrated with $X$, and use
\[
 H_{\rm copy}(t)=\dot\kappa(t)X P_y,\qquad
 \kappa(0)=0,\quad\kappa(\tau_c)=\kappa.
\]
Characteristics give
\begin{equation}
 \Phi_{\tau_c}(x,y)=\psi\phi(x)\chi(y-\kappa x),\qquad
 X_{\tau_c}=X_0,\quad Y_{\tau_c}=Y_0+\kappa X_0.
 \label{cfg:snapshot}
\end{equation}
The current in $x$ is zero during this pulse, while the current in $y$
is $\dot\kappa x|\Phi|^2$.

\begin{proposition}[Conditional rank after a snapshot]\label{cfg:rank}
Let
\[
 n(y)=\int|\phi(x)|^2|\chi(y-\kappa x)|^2dx,\qquad
 \phi_y(x)=\frac{\phi(x)\chi(y-\kappa x)}{\sqrt{n(y)}}.
\]
For almost every $y$ with $n(y)>0$, the conditional actual density of $X$
is $|\phi_y|^2$. Its own cumulative rank $F_y(X)$ is uniform conditionally
on $Y=y$. Consequently an estimate based on $Y$ and independent processing
keys obeys
\[
 \mathbb P(|F_Y(X)-a(Y)|\le\eta)\le2\eta,\qquad 0\le\eta\le\tfrac12.
\]
\end{proposition}
\begin{proof}
Push the initial product density through \eqref{cfg:snapshot}. It becomes
$|\phi(x)|^2|\chi(y-\kappa x)|^2dx\,dy$. Bayes' formula supplies the
normalized conditional density. Its probability integral transform is
uniform, and an interval of length at most $2\eta$ has at most that
conditional probability. Independent keys can be conditioned on and then
integrated out.
\end{proof}

The original rank and the narrowed conditional rank are different variables.
This theorem calculates the change; it does not impose numerical threshold
retirement. For real Gaussian packets of variances
$\sigma_x^2,\sigma_y^2$, set $r=\kappa^2\sigma_x^2/\sigma_y^2$.
Gaussian integration and $U^\dagger P_xU=P_x-\kappa P_y$ give
\[
 \operatorname{Var}(X\mid Y)=\frac{\sigma_x^2}{1+r},\qquad
 D_{\rm pure}(\Phi_{\tau_c},\psi\phi\chi)=\sqrt{\frac r{4+r}},
 \qquad
 \Delta\operatorname{Var}(P_x)=\frac{\hbar^2r}{4\sigma_x^2}.
\]
Indeed the wave overlap is $(1+r/4)^{-1/2}$, and the final momentum variance
adds $\kappa^2\hbar^2/(4\sigma_y^2)$.
The pure-state distance compares complete fields; the information has not
been discarded in a reduced channel.

Keeping $y$ frozen after copying and using the binary detector with $L>2a$
gives the complete finite acquisition law
\begin{equation}
 \mathbb P(j,dt,dy)=p_j u f(s_j(L-ut))
          |\chi(y-\kappa s_j(L-ut))|^2\,dt\,dy.
 \label{cfg:snapshot-time}
\end{equation}
To prove it, condition on $y$ and apply the quantile proof to $\phi_y$.
For a right exit the originating packet coordinate is $\xi=L-ut$; for a
left exit it is $\xi=ut-L$. Multiplication by $n(y)$ yields
\eqref{cfg:snapshot-time}. Integrating in time gives
$\mathbb P(j,dy)=p_j n(y)\,dy$, but the time profiles conditional on $y$
need not be equal for the two outlets. In particular, $\phi_y$ may be
asymmetric. Correct terminal weights therefore do not imply a common
conditional clock at every snapshot value.

\section{Finite wave writes, physical copies, and echo experiments}

There is a discrete finite analogue that will also be used for archives.
Let a blank register $A$ have orthogonal states $|b\rangle,|j\rangle$.
For projectors $Q_j$ on the unknown system, define
\begin{equation}
 H_M=i\hbar g\sum_jQ_j\otimes
       (|j\rangle\langle b|-|b\rangle\langle j|).
 \label{cfg:writeH}
\end{equation}
Its norm is $\hbar g$, and from $\psi|b\rangle$ the exact wave is
\begin{equation}
 \Phi_t=\sum_jQ_j\psi\otimes
             (\cos gt\,|b\rangle+\sin gt\,|j\rangle).
 \label{cfg:writewave}
\end{equation}
It preserves all within-sector amplitudes and inaccessible reference
correlations by its explicitly scalar action.

Put $e_j=\|Q_j\psi\|^2$. For the pointer-only configuration resolution,
the minimal Bell currents are $J_{jb}=2ge_j\sin gt\cos gt$.
Before $\pi/(2g)$,
\[
 q_{j\leftarrow b}=2ge_j\tan gt,\qquad
 \mathbb P(\tau>t)=\cos^2gt,\qquad
 \mathbb P(j,\tau\in dt)=e_jg\sin(2gt)\,dt.
\]
Integration of the total hazard proves survival; multiplication by the
conditional mark rate gives the density. Its divergent endpoint hazard
expels all residual blank probability without an explosion, since there
can be only one jump during this pulse.

If the fixed fundamental resolution resolves system sectors as well, use
$Q_j=P_j^S$ for this exact one-event example. Each transition
$(j,b)\to(j,j)$ has rate $2g\tan gt$ and initial equilibrium mass $e_j$.
The same mixture density results. For arbitrary projectors on a finer fixed
resolution, endpoint weights still follow from equivariance but the
pointer-only first-event formula must not be imported. Coarse flux and
path transfer require the separate conditions of
Theorem~\ref{stat:coarse}.

A second blank register $C$ can be written by
\[
 H_C=i\hbar g_C\sum_j|j\rangle\langle j|_A\otimes
           (|j\rangle\langle b|-|b\rangle\langle j|)_C
\]
for area $\pi/2$. At a completed first write,
\[
 \sum_jQ_j\psi|j\rangle_A|b\rangle_C
 \longmapsto\sum_jQ_j\psi|j\rangle_A|j\rangle_C.
\]
Reversing the first write alone produces
$\sum_jQ_j\psi|b\rangle_A|j\rangle_C$. Reversing both copies, in the
appropriate reverse order, restores the full original product wave. These
are unitary identities, not assertions that an occupied branch has changed
the global source wave.

The partial-write experiment gives a useful finite backaction calculation.
Write $c=\cos\theta,s=\sin\theta$, copy after the partial first pulse, and
reverse that pulse by $\theta$. The final wave is
\begin{equation}
 \sum_jQ_j\psi\left[
 c^2|b,b\rangle-cs|j,b\rangle
 +s^2|b,j\rangle+sc|j,j\rangle\right]_{A,C}.
 \label{cfg:partialecho}
\end{equation}
Therefore, under any admitted equivariant complete configuration law,
\[
 \mathbb P(A=b)=c^4+s^4,\qquad
 \mathbb P(A=j)=2c^2s^2e_j,\qquad
 \mathbb P(C=j)=s^2e_j.
\]
At $\theta=\pi/4$ copying reduces the return probability from one to
one half. The acquired register is retained throughout this comparison.

A later noncommuting measurement makes the continuation distinction explicit.
After a completed first write and copy, reverse $A$, apply a known unitary
$V_j$ controlled by $C=j$, and write projectors $B_k$ into a fresh register.
The final orthogonal components are
$(B_kV_jQ_j\otimes I_R)\psi\,|b,j,k\rangle$. Hence
\begin{equation}
 \mathbb P(C=j,D=k)=\|(B_kV_jQ_j\otimes I_R)\psi\|^2.
 \label{cfg:sequential}
\end{equation}
This is a coherent controlled Hamiltonian on a complete finite bank.
No external rule first samples $j$ and then supplies a daughter.
For $\psi=|+x\rangle$, a $Z$ write followed by complete erasure and an $X$
probe gives $+$ with probability one. If one copy remains, the same $X$
probe gives probability one half. A daughter-only mixture cannot reproduce
the complete-erasure experiment.

\section{What final-record equivariance actually proves}

\begin{theorem}[Complete physical endpoint law]\label{cfg:endpoint}
Fix a linear isometric apparatus preparation map $T$ independent of the unknown
input, a complete linear wave programme $U$, and an initially equivariant
actual configuration law. Suppose the actual dynamics preserve the norm
density of the complete wave. For a physically stored final record with
configuration projector $P_h$,
\begin{equation}
 \mathbb P_\psi(h)=\|P_hUT\psi\|^2.
 \label{cfg:endpointlaw}
\end{equation}
All null, opposite, failure, and exhaustion regions are included in the
resolution. The formula retains arbitrary inaccessible references and
returning quantum memories in $UT\psi$.
\end{theorem}
\begin{proof}
Equivariance identifies the actual final configuration law with the squared
norm of $UT\psi$. Summing or integrating that law over the physical record
region gives \eqref{cfg:endpointlaw}. Every step acts by the identity on the
reference. Linearity of $T$ and $U$ then gives the matrix event map
$\rho\mapsto P_hUT\rho T^\dagger U^\dagger P_h$.
This map is a consequence of the specified wave and probability laws, not
an assumption used to select those laws. Classical provenance is retained
by conditioning on its actual value and then averaging with its actual law.
\end{proof}

\begin{corollary}[Operational endpoint equivalence]\label{cfg:equivalence}
Two conservative equivariant configuration processes with the same initial
complete wave, norm-distributed configuration, global Hamiltonian programme,
and final physical record map have identical final-record distributions.
This applies to minimal Bell dynamics and any well-defined member of
\eqref{cfg:Kfamily}. It also applies to a final bank storing an entire
finite operational record string.
\end{corollary}
\begin{proof}
Both distributions equal \eqref{cfg:endpointlaw}. The record string must
be physically present in the final bank, so it is a final configuration
function; the argument does not apply to an unrecorded path functional.
\end{proof}

For a finite sequence with retained record labels, expand the complete
unitary programme into its orthogonal final branches. If their carried
coefficients are $K_{h_n}^{(n)}\cdots K_{h_1}^{(1)}\psi$, then
\[
 \mathbb P(h_1,\ldots,h_n)
 =\|K_{h_n}^{(n)}\cdots K_{h_1}^{(1)}\psi\|^2.
\]
This follows by the expansion and Theorem~\ref{cfg:endpoint}, without
independently sampling branch weights at each stage.
The $K$ operators must be the actual branch coefficients, including
coherent nulls, surviving excitations, and controlled phases.

An instructive null occurs when an active state $|e\rangle$ has a
clean-return recorder
$|e,M_0\rangle\mapsto A|e,M_0\rangle+D|e,M_1\rangle$,
$|A|^2+|D|^2=1$, while a spectator $|u,M_0\rangle$ is unchanged.
For input $\alpha|e\rangle+\beta|u\rangle$, the unnormalized null vector is
$\alpha A|e\rangle+\beta|u\rangle$. A real rotation
$|e\rangle\mapsto c|e\rangle+s|u\rangle$,
$|u\rangle\mapsto-s|e\rangle+c|u\rangle$, followed by a second
active recorder of response $R_2$, gives
\[
 \mathbb P(\mathrm{null}_1,\mathrm{record}_2)
       =|\alpha Ac-\beta s|^2 R_2.
\]
For reference vectors $R_e,R_u$, replace the square by
\[
 |\alpha Ac|^2+|\beta s|^2
 -2\operatorname{Re}\{\alpha Ac\,\beta^*s
                   \langle R_u|R_e\rangle\}.
\]
Orthogonal references remove interference but not the spectator-to-active
term. A scalar posterior for active origin is therefore insufficient.

At a separated cut and on a branch with $\|P_h\Phi\|^2>0$, conditioning on $h$ gives configuration probabilities
$\|P_xP_h\Phi\|^2/\|P_h\Phi\|^2$. For a claim that the normalized relative
wave alone generates the same microscopic future path, additionally require
\begin{equation}
 q_{y\leftarrow x}[\Phi,c]
   =q_{y\leftarrow x}[P_h\Phi/\|P_h\Phi\|,c],
 \qquad x,y\in\Gamma_h,
 \label{cfg:branchlocal}
\end{equation}
and zero transitions out of $\Gamma_h$ on the claimed horizon.
Minimal Bell rates with a block-diagonal Hamiltonian satisfy this by
cancellation of the normalization factor. An arbitrary nonlinear
$K[\Phi]$ need not. The full-wave endpoint theorem remains valid without
\eqref{cfg:branchlocal}; autonomous daughter law and endpoint operational
equivalence are different conclusions.

\section{Autonomous recording: an obstruction and a completed repair}

The shuttered write does not establish an independent memory oscillator
that restarts at each actual product entry. A configuration jump leaves
the full wave intact, so such a restart is another physical intervention.
The following exact autonomous examples from \cite{C02} keep all amplitudes.

For $e=|1,M_0\rangle$, $p=|R,M_0\rangle$, and
$m=|R,M_1\rangle$, take
\[
 H_3/\hbar=g(|p\rangle\langle e|+|e\rangle\langle p|)
          +\chi(|m\rangle\langle p|+|p\rangle\langle m|).
\]
Writing $\Omega_3=\sqrt{g^2+\chi^2}$, the initial state $e$ evolves with
\[
 a=\frac{\chi^2+g^2\cos\Omega_3t}{\Omega_3^2},\quad
 b=-i\frac g{\Omega_3}\sin\Omega_3t,\quad
 c=\frac{g\chi}{\Omega_3^2}(\cos\Omega_3t-1).
\]
Substitution into $i\dot a=gb$, $i\dot b=ga+\chi c$,
$i\dot c=\chi b$ verifies the solution and initial data. Consequently
\begin{equation}
 \sup_t\mathbb P(m,t)\le
       \frac{4g^2\chi^2}{(g^2+\chi^2)^2}\longrightarrow0
       \quad\hbox{as }\chi/g\longrightarrow\infty.
 \label{cfg:threebound}
\end{equation}
Strong memory coupling suppresses transfer; it is not automatically
arbitrarily fast faithful acquisition. The added terminal-protection
channel and its controlled weak-protection asymptotic are calculated in
Section~\ref{bench:protection}. Product-reservoir recurrences are retained
explicitly in Section~\ref{bench:reservoir}.

A scoped trapped-source obstruction is also exact. In native
$e\leftrightarrow p$ minimal dynamics, product entry occurs by $\pi/(2g)$
almost surely and the source returns to $e$ by $T=\pi/g$.
Suppose a monitored architecture traps every successfully recorded path
on the product side through $T$. Let $d$ be source-path total variation
from the native model, $E$ be monitored product entry by $T$, and
$\varepsilon=\mathbb P(E\hbox{ and no protected record})$. Then
$\mathbb P(E)\ge1-d$, the recorded probability is at least
$1-d-\varepsilon$, and its trapped product endpoint differs from the
native endpoint. Thus
\begin{equation}
 2d+\varepsilon\ge1.
 \label{cfg:trap}
\end{equation}
The inequality uses only these events and the definition of total variation.
It is not a universal recorder impossibility theorem.

A four-state reaction supplies the relevant repair. In the basis
$e0,p0,p1,e1$, take
\[
 H_4/\hbar=
 \begin{pmatrix}
 0&g&0&0\\g&0&\chi&0\\0&\chi&0&g\\0&0&g&0
 \end{pmatrix}.
\]
Put $\Omega=\sqrt{\chi^2+4g^2}$,
$C=\cos(\Omega t/2)$, $S=\sin(\Omega t/2)$,
$v=\chi t/2$, $k=2g/\Omega$, $r=\chi/\Omega$. From $e0$,
\[
 a=C\cos v+rS\sin v,\quad b=-ikS\cos v,\quad
 c=-kS\sin v,\quad d=i(-C\sin v+rS\cos v).
\]
Reflection-symmetric and antisymmetric combinations reduce $H_4$ to two
$2\times2$ matrices $\left(\begin{smallmatrix}0&g\\g&\pm\chi\end{smallmatrix}\right)$.
Their elementary exponentials give these amplitudes, proving the formula.

The source product probability and coherence are
\[
 \rho_{pp}=\frac{4g^2}{\Omega^2}\sin^2(\Omega t/2),\qquad
 \rho_{ep}=i\frac g\Omega\sin(\Omega t).
\]
Choose $\chi=2\sqrt3g$, so $\Omega=4g$.
At $t=\pi/(2g)$ the monitored source is certainly $e$, while the native
source is certainly $p$: endpoint and path total variation are one.
At $T=\pi/g$ the monitored wave is
\[
 |e\rangle\{\cos(\sqrt3\pi)|M_0\rangle
            -i\sin(\sqrt3\pi)|M_1\rangle\},
\]
and the source endpoint agrees exactly with the native one.
The written branch permits source return, so the trapping premise of
\eqref{cfg:trap} has genuinely changed. Endpoint agreement has not removed
the earlier disturbance.

At a clean factorized return $|e\rangle(a|M_0\rangle+d|M_1\rangle)$,
a subsequent source-only operation cannot reveal that isolated memory.
A real conditional contact
$\hbar\kappa(|p,M_1\rangle\langle e,M_1|+\mathrm{h.c.})$
for time $\tau$ gives instead
$\mathbb P(p)=|d|^2\sin^2(\kappa\tau)$.
Retained information matters through its correlations and later couplings,
not merely through the name archive.

\chapter{Faithful archives and complete-history error}
\label{cfg:archives}

Final record probabilities, correctness about the actual past, and
microscopic path laws are three distinct observables. This chapter expands
the archive arguments of \cite{C02,C03} into a finite theorem. It establishes
conditions sufficient for reliable sampled histories without imposing
minimality on every microscopic edge.

\section{An endpoint archive can be wrong about its own past}

Let $L_t$ be an actual sampled label and $A_T$ the label retained at the
final cut. Even exact equality $A_t=L_t$ at every current time does not
imply $A_T=L_0$.

\begin{counterexample}[Correlated label and archive turnover]
\label{cfg:turnover}
Take $\Phi=(|00\rangle+|11\rangle)/\sqrt2$, $H=0$, and symmetric
surplus $K_{00,11}=k>0$. Both conditional rates in
\eqref{cfg:Kfamily} are $2k$. The complete endpoint distribution remains
one half on each of $00,11$, and present label and archive agree always.
Nevertheless,
\begin{equation}
 \mathbb P(A_T\ne L_0)=\frac{1-e^{-4kT}}2.
 \label{cfg:turnovererror}
\end{equation}
\end{counterexample}
\begin{proof}
The number of switches is Poisson with mean $2kT$. Its odd-parity
probability is
$\frac12(1-\mathbb E(-1)^N)=\frac12(1-e^{-4kT})$.
Odd parity is exactly disagreement with the initial label. Stationarity
follows because the two rates and initial weights are equal.
\end{proof}
This rival is excluded by an additional no-transition-without-coupling
support law, not by equivariance. Multiple simultaneously changing copies
can be mutually consistent and historically unreliable.

\section{A finite write of a held actual label}

Let $\{P_\ell^L\}$ be orthogonal physical configuration projectors for a
label register. Supply a blank archive cell $|B\rangle$ orthogonal to
its outputs $|A_\ell\rangle$. The pulse is
\begin{equation}
 H_{\rm w}=\hbar\chi\sum_\ell P_\ell^L\otimes
      (|A_\ell\rangle\langle B|+|B\rangle\langle A_\ell|),
 \qquad \Delta=\frac{\pi}{2\chi}.
 \label{cfg:heldwriteH}
\end{equation}
An unknown carried vector and reference may be arbitrarily correlated with
$L$. The initial joint configuration law is equivariant. During the pulse
both the Hamiltonian and actual generator preserve $L$, and already completed
archive entries are preserved. Nonminimal traffic within fixed-label blocks
is allowed.

\begin{theorem}[Exact finite copying of a sampled actual label]
\label{cfg:heldcopy}
Under the preceding assumptions, the completed entry equals the actual
label at pulse start almost surely:
\[
 \mathbb P(A_{t+\Delta}\ne L_t)=0.
\]
The pulse uses $1+\#\{\ell\}$ archive states, duration $\Delta$, and
operator norm $\hbar\chi$. No additional copy of the unknown input is used.
\end{theorem}
\begin{proof}
Decompose the full vector as $\sum_\ell\psi_\ell|B\rangle$,
$P_\ell^L\psi_\ell=\psi_\ell$. The mutually orthogonal pulse blocks give
\[
 e^{-iH_{\rm w}\Delta/\hbar}\sum_\ell\psi_\ell|B\rangle
       =-i\sum_\ell\psi_\ell|A_\ell\rangle.
\]
The endpoint wave has zero support on every mismatch
$L=\ell,A\ne A_\ell$. Equivariance therefore makes endpoint mismatch
probability zero. Actual preservation of $L$ throughout the finite pulse
then identifies its endpoint value with $L_t$, proving the claim. The
reference and all internal correlations stay in the vectors $\psi_\ell$.
\end{proof}

The label-holding premise is physical. Freezing a native source label can
change its original event process. Copying an already stable acquired record
need not freeze the source. A shuttered write records occupancy at its
start, not all earlier first entries or every unmonitored jump.

For an imperfect pulse, let $N_{L,m}$ count changes of the sampled label
during write $m$, and let $\epsilon_m$ bound the norm difference of complete
physical and ideal endpoint waves on the actual input/reference class.
If the physical actual process is equivariant, the ideal mismatch projector
annihilates the ideal wave, giving
\begin{equation}
 \delta_{{\rm w},m}
 :=\mathbb P(A_{m,\rm end}\ne L_{t_m})
 \le \mathbb E N_{L,m}+\epsilon_m^2.
 \label{cfg:writeerror}
\end{equation}
Indeed the first disagreement requires a label change or an endpoint
mismatch. The first probability is bounded by its expected count; the
second is $\|\Pi_{\rm mis}\Phi_{\rm phys}\|^2
\le\|\Phi_{\rm phys}-\Phi_{\rm id}\|^2$.
The square is special to an ideal zero-probability event, not a general
quadratic continuity bound for all output probabilities.

\section{The archive theorem and its conditional version}

Consider a finite physical programme with $m$ writes, at prescribed or
physically controlled start times $t_i$, each with a completion cut.
Let $H_{\rm past}=(L_{t_1},\ldots,L_{t_m})$ and let $A_T$ contain
the final completed entries. On unused, rejected or exhausted branches,
put a declared symbol in the same finite output alphabet and retain the
actual source state. Let $N_{\rm corrupt}$ count transitions after completion
that alter any completed entry. A transition may alter several entries;
counting it once is enough for the following event inclusion.

\begin{theorem}[Faithful finite archive]\label{cfg:faithful}
For every well-defined process on this complete experiment,
\begin{equation}
 \delta_{\rm faith}:=\mathbb P(A_T\ne H_{\rm past})
 \le\min\left\{1,\sum_{i=1}^m\delta_{{\rm w},i}
                         +\mathbb E N_{\rm corrupt}\right\}.
 \label{cfg:faithbound}
\end{equation}
This conclusion allows adaptive controls, retained keys, returning memories,
and nonminimal internal traffic. Each write error and corruption count must
refer to those actual complete histories.
\end{theorem}
\begin{proof}
If every entry was correct when completed and no completed entry was
subsequently changed, the final bank equals the sampled history.
Therefore $\{A_T\ne H_{\rm past}\}$ is contained in the union of all write
failure events and $\{N_{\rm corrupt}\ge1\}$. The union bound and
$1_{\{N\ge1\}}\le N$ prove \eqref{cfg:faithbound}. This event argument is
valid pathwise even if corruption later repairs an error; such repairs only
make the bound conservative. Conditioning on an adaptive control history
gives the same inclusion, and averaging preserves the bound.
\end{proof}

This is stronger than endpoint correlation, but weaker than recording the
entire unmonitored trajectory. Neither correct copies nor small corruption
budgets select a unique native Bell event generator.

For a common final record event $E$ with probability $p>0$,
\[
 \mathbb P(A_T\ne H_{\rm past}\mid E)
       \le\min(1,\delta_{\rm faith}/p).
\]
A rare successful branch can thus be unreliable even when the average
error is small. Uniform conditional promises require either lower event
probabilities or a proof made separately on every admitted conditional
preparation.

\section{Physical support and quantitative crossing budgets}

Let $\mathcal C_a$ be the configurations with completed archive content $a$.
For the equilibrium family \eqref{cfg:Kfamily}, sum over unordered pairs
$\{x,y\}$ with different completed contents. The compensator of their
directed transition counts gives
\begin{equation}
 \mathbb E N_{\rm corrupt}
 =\int_0^T\sum_{\{x,y\}\ {\rm cross}}
             (|J_{yx}(t)|+2K_{xy}(t))\,dt.
 \label{cfg:traffic}
\end{equation}
Here the crossing predicate must be a function of represented complete
configuration and physical time, so its occupation-weighted average is
computed under the stated equivariant law. A deterministic schedule of
completed banks meets this condition. An adaptive completion status may
also be encoded in the complete configuration, with its own physical
currents and equilibrium law. If instead completion depends on additional
unrepresented history, the deterministic-current expression above is not
asserted: use the following predictable expectation directly. For full-history
intensities $\lambda_{y\leftarrow x}(t\mid\mathcal F_{t-})$ the exact
formula is the expectation of their predictable sum:
\[
 \mathbb E N_{\rm corrupt}
 =\mathbb E\int_0^T\sum_{y:\,{\rm cross}(X_{t-},y,t)}
            \lambda_{y\leftarrow X_{t-}}(t\mid\mathcal F_{t-})\,dt.
\]
For deterministic Markov rates and nonequilibrium occupation law $\nu_x$,
the integrand for an unordered pair is
$\nu_xq_{y\leftarrow x}+\nu_yq_{x\leftarrow y}$.
The wave-weight expression \eqref{cfg:traffic} cannot then be used.

\begin{proof}[Proof of \eqref{cfg:traffic}]
For each directed count, expected compensator equals expected count,
first with bounded stopping and then by monotone convergence.
In equilibrium the occupation-weighted directional incidences are
$[J_{yx}]_++K_{xy}$ and $[-J_{yx}]_++K_{xy}$.
Their sum is $|J_{yx}|+2K_{xy}$. Summing crossing pairs and integrating
proves the identity. No conclusion is drawn by dividing a small flux by
a possibly small sector weight.
\end{proof}

A complete support law
\begin{equation}
 H_{yx}=0\quad\Longrightarrow\quad
 q_{y\leftarrow x}=q_{x\leftarrow y}=0
 \label{cfg:support}
\end{equation}
in the absence of separately declared stochastic channels gives exact
protection if the complete Hamiltonian is block diagonal by archive content.
It is weaker than minimality: $K_{xy}=\eta|J_{yx}|$ satisfies support and
is generally nonminimal. Every controller, bath, return wire, and reset
channel belongs in the complete physical graph.

\begin{counterexample}[Support has no uniform weak-coupling modulus]
\label{cfg:discontinuity}
For $\Phi=(|0\rangle+|1\rangle)/\sqrt2$ and
$H_\varepsilon=\hbar\varepsilon\sigma_x$, the weights are constant and
$J=0$. Put $K_{01}=\gamma/2$ for $\varepsilon\ne0$ and zero for
$\varepsilon=0$. This obeys \eqref{cfg:support}, but for every nonzero
$\varepsilon$ its two rates are $\gamma$ and
\[
 \mathbb P(Q_T\ne Q_0)=\frac{1-e^{-2\gamma T}}2,
 \qquad \hbar^{-1}\int_0^T\|H_\varepsilon\|dt
       =|\varepsilon|T\longrightarrow0.
\]
\end{counterexample}
The formula follows from odd parity of a rate-$\gamma$ Poisson count.
Thus ordinary small-Hamiltonian continuity of the wave does not control
historical corruption for an unconstrained actual generator.

One sufficient quantitative replacement is
\begin{equation}
 K_{xy}\le C\,\frac{\|H_{yx}\|}{\hbar}\sqrt{w_xw_y}.
 \label{cfg:envelope}
\end{equation}
Cauchy--Schwarz gives
$|J_{yx}|\le2\|H_{yx}\|\sqrt{w_xw_y}/\hbar$.
Define the \emph{ordered} coupling budget
\[
 \Lambda_A=\int_0^T\sum_{x,y:\,{\rm cross}}
       \frac{\|H_{yx}\|}{\hbar}\sqrt{w_xw_y}\,dt.
\]
Each unordered pair occurs twice, so
\begin{equation}
 \mathbb E N_{\rm corrupt}\le(1+C)\Lambda_A.
 \label{cfg:envelopebound}
\end{equation}
The coefficient would be $2+2C$ for an unordered coupling budget.
Uniform graph bounds are needed if the number of archive configurations
grows.

The linear envelope is sufficient, not necessary. In the preceding
two-state example, $K_{01}=\frac12\sqrt{\gamma|\varepsilon|}$ gives
rates $\sqrt{\gamma|\varepsilon|}$ and final mismatch
$(1-e^{-2T\sqrt{\gamma|\varepsilon|}})/2\to0$, without a uniform
linear bound in $|\varepsilon|$. The relevant sufficient limit is the
vanishing integrated actual crossing traffic on the stated horizon.

\section{Combining wave, path, and archive comparisons}

Let two complete finite programmes use the same initial wave, fixed
configuration resolution, and physical final archive map, with
Hamiltonians $H,\widetilde H$. If
\[
 \epsilon_H=\frac1\hbar\int_0^T\|H_t-\widetilde H_t\|\,dt,
\]
Duhamel's identity and unitarity give
$\|\Phi_T-\widetilde\Phi_T\|\le\epsilon_H$.
For normalized vectors, the pure-state trace distance is at most their
vector norm difference. Measurement of the common endpoint partition
therefore gives archive-law total variation at most $\epsilon_H$.
This argument assumes equilibrium/equivariance in each compared model.

For a common projected event of ideal probability $p>0$, normalized branch
vectors satisfy
\[
 \left\|\frac{P\Phi}{\|P\Phi\|}
          -\frac{P\widetilde\Phi}{\|P\widetilde\Phi\|}\right\|
 \le \frac{2\epsilon_H}{\sqrt p}
\]
whenever the second denominator is nonzero. To prove it, add and subtract
$P\widetilde\Phi/\|P\Phi\|$ and use reverse triangle inequality for the
two norms. This is a same-projector vector estimate; it is not the
conditioning bound for arbitrary contaminated historical events.

If both models also have faithful-archive errors
$\delta_{\rm faith}$ and $\widetilde\delta_{\rm faith}$, then
\begin{equation}
 \operatorname{TV}(\mathcal L(H_{\rm past}),
                  \mathcal L(\widetilde H_{\rm past}))
 \le \min(1,\delta_{\rm faith}+\epsilon_H
                    +\widetilde\delta_{\rm faith}).
 \label{cfg:historycompare}
\end{equation}
Within each model, actual history and its archive are already coupled;
their disagreement probability bounds their law distance. Apply that
coupling inequality on both sides and insert the endpoint bound between
the two archive laws.

The fixed finite minimal Bell path-stability theorem
(Theorem~\ref{int:path-stability}) is a stronger comparison when its
node and Hamiltonian hypotheses hold. It is not needed for
Theorem~\ref{cfg:endpoint}, and it does not extend to
Counterexample~\ref{cfg:discontinuity}'s arbitrary surplus laws.

For the finite positive-background law of the statistical selection
chapter, the complete wave is identical to its zero-background Bell limit.
Final physical archive distributions are consequently identical by
Corollary~\ref{cfg:equivalence}; a path total-variation estimate is valid
but loose for those endpoint-only outputs. Positive background on
Hamiltonian-disconnected archive pairs can nevertheless corrupt actual
past labels. Exact historical protection requires the Bell limit,
an explicit small crossing budget, or a separately analyzed supported
reference graph. A change to the reference graph cannot be silently
inserted into the variational proof.

\section{Scope of the consolidated record result}

The established chain is
\[
 \begin{gathered}
 \text{linear complete wave dynamics and admitted equivariant configurations}\\
 +\ \text{joint ready resource law and physical finite writes}\\
 +\ \text{held sampled labels and bounded archive-changing traffic}\\
 \Longrightarrow\quad
 \text{quantum endpoint records and faithful sampled actual histories}.
 \end{gathered}
\]
References, nulls, failures, retained keys, and permitted coherent returns
are included when present in the modeled bank. At most a finite declared
stock is used; exhausted branches do not receive fictitious fresh cells.
This chapter reduces historical record reliability to explicit writes,
held labels and actual crossing budgets; it does not itself derive Bell
minimality or timing. The pilot theory supplies its effective event
law and proves exact monomial-copy faithfulness in the Bell comparator on
the first pass of an autonomous finite clock. Its finite-resource path
error then controls failures of that historical claim. Neither result
licenses an unchanged-source reader of every unmonitored microscopic event.

% End consolidated source: parts/07_configuration.tex


% Begin consolidated source: parts/07_stability.tex
\chapter{Complete-path stability and conditioned records}
\label{int:stability-chapter}

The finite path-stability argument in the record checkpoints \cite{C02,C03} concerns an already admitted minimal Bell generator. A small wave error alone is not a bound on pointwise jump rates near a node. The correct comparison weights a rate discrepancy by the probability that the two processes still share its origin state. The following proof supplies that step and a quantitative version on a fixed finite sector space. No lower bound on a populated sector weight is required in the estimate.

\section{A coupling estimate that retains the occupation weights}
\begin{lemma}[Overlap-weighted path coupling]\label{int:coupling}
Let $P,Q$ be nonexplosive time-inhomogeneous Markov jump laws on the same finite sector set, with rates $k_{nm},\ell_{nm}$ and marginal probabilities $p_m(t),q_m(t)$. Suppose their ordinary common-jump coupling is defined by localization on intervals of finite rates. Then
\begin{equation}\label{int:overlap-bound}
 d_{\rm TV}(P,Q)\le d_{\rm TV}(p(0),q(0))+
 \int_0^T\sum_{m\ne n}\min\{p_m(t),q_m(t)\}|k_{nm}(t)-\ell_{nm}(t)|\,dt.
\end{equation}
The comparison is on complete paths, including their jump times and null intervals.
\end{lemma}
\begin{proof}
Couple the initial states maximally. While both histories remain identical at $m$, give them simultaneous $m\to n$ jumps at rate $\min(k_{nm},\ell_{nm})$, a jump only in the first coordinate at rate $(k_{nm}-\ell_{nm})_+$, and a jump only in the second at the reverse positive difference. After the first discrepancy continue with the prescribed marginal laws. Before discrepancy the total separating rate at $m$ is $\sum_{n\ne m}|k_{nm}-\ell_{nm}|$. Write $c_m(t)$ for the probability that histories have never separated and currently agree at $m$. Each marginal dominates this event, so $c_m\le\min(p_m,q_m)$. The stopped separating-count compensator gives the integral in \eqref{int:overlap-bound}. Probability of any discrepancy bounds total variation of the path laws. Apply the construction first on compact intervals of bounded rates and then localize. The nonnegative compensator bound survives the limit. In the Bell domain below, the only possible finite endpoints are nodes or switches; zero population at a node and finite expected jump count give continuation without loss of probability, as in the node construction of Part~\ref{part:statistics}.
\end{proof}

\begin{lemma}[Cancellation at a sector node]\label{int:node-cancellation}
Let $\Psi,\Phi$ be normalized waves for the same orthogonal sectors, with $a_m=\|P_m\Psi\|$, $b_m=\|P_m\Phi\|$. Let $J,J'$ be their Hamiltonian currents for $H,G$, respectively. Put $h_* = \max(\|H\|,\|G\|)$ and set rates to zero at zero-weight origins. For every ordered pair $m\ne n$,
\begin{align}\label{int:node-ineq}
 \min(a_m^2,b_m^2)
 \left|\frac{[J_{nm}]_+}{a_m^2}-\frac{[J'_{nm}]_+}{b_m^2}\right|
 \le |J_{nm}-J'_{nm}|+
 \frac{4h_*}{\hbar}(a_n+b_n)|a_m-b_m|.
\end{align}
Consequently, with $\delta=\|\Psi-\Phi\|$ and $d$ sectors,
\begin{equation}\label{int:sumineq}
 \sum_{m\ne n}\min(a_m^2,b_m^2)|k^{\rm B}_{nm}-\ell^{\rm B}_{nm}|
 \le\frac{12d h_*\delta+2d\|H-G\|}{\hbar}.
\end{equation}
\end{lemma}
\begin{proof}
If one origin norm is zero the left side is zero. If $0<a_m\le b_m$, multiply by $a_m^2$ and add and subtract $[J'_{nm}]_+$:
\[
 a_m^2\left|\frac{[J_{nm}]_+}{a_m^2}-\frac{[J'_{nm}]_+}{b_m^2}\right|
 \le |J_{nm}-J'_{nm}|+
 \left(1-\frac{a_m^2}{b_m^2}\right)|J'_{nm}|.
\]
The current bound $|J'_{nm}|\le2\|G\|b_nb_m/\hbar$ and
$(b_m^2-a_m^2)/b_m\le2|b_m-a_m|$ prove the result with $b_n$ in the second term. Interchanging the two processes proves the other case with $a_n$, and $a_n+b_n$ is a common bound.

Write $u_m=\|P_m(\Psi-\Phi)\|$. Orthogonality gives $\sum u_m^2=\delta^2$, $\sum a_m^2=\sum b_m^2=1$, and $|a_m-b_m|\le u_m$. Expanding the bilinear current difference gives
\[
 \sum_{m\ne n}|J_{nm}-J'_{nm}|
 \le \frac{2}{\hbar}\sum_{m,n}
 \left[\|H\|(u_na_m+b_nu_m)+\|H-G\|b_nb_m\right]
 \le\frac{4d\|H\|\delta+2d\|H-G\|}{\hbar}.
\]
Also $\sum_{m,n}(a_n+b_n)|a_m-b_m|\le2d\delta$. Summing \eqref{int:node-ineq} proves \eqref{int:sumineq}. Including diagonal pairs in these upper bounds only enlarges them.
\end{proof}

\begin{theorem}[Finite Bell path stability through nodes]\label{int:path-stability}
Fix a finite orthogonal sector resolution and a physical horizon $T$. Let $H(t),G(t)$ be bounded self-adjoint Hamiltonians, each constant on finitely many time segments, on a finite-dimensional complete coherent bank. Let their normalized waves be $\Psi_t,\Phi_t$, and initialize their minimal Bell processes with respective populations $\|P_m\Psi_0\|^2,\|P_m\Phi_0\|^2$. Write
\[
 \begin{aligned}\delta_0&=\|\Psi_0-\Phi_0\|,\quad
 e(t)=\frac{\|H(t)-G(t)\|}{\hbar},\\
 D(t)&=\delta_0+\int_0^t e(s)\,ds,\quad
 h(t)=\frac{\max(\|H(t)\|,\|G(t)\|)}{\hbar}.\end{aligned}
\]
Their complete physical-time path laws satisfy
\begin{equation}\label{int:path-error}
 \boxed{\quad d_{\rm TV}(P_H,P_G)
 \le\min\left\{1,\delta_0+12d\int_0^T h(t)D(t)\,dt
                      +2d\int_0^T e(t)\,dt\right\}.\quad}
\end{equation}
The estimate is uniform over the initial normalized waves and has no inverse-weight constant. An inaccessible reference can be included in the sector fibers. The constant depends on the number of compared sectors, the integrated Hamiltonian norms and the horizon, not on a minimum sector population.
\end{theorem}
\begin{proof}
The finite-node construction in Part~\ref{part:statistics} supplies equivariant, nonexplosive minimal laws. Alternatively, apply that construction directly to $J$ and $J'$; bounded constant-Hamiltonian segments have analytic weights, finitely many isolated nodes unless a sector is identically empty, and integrable total current. Indeed
\[
\mathbb E N_{\rm all}(T)=\int_0^T\sum_{m\ne n}[J_{nm}(t)]_+dt
 \le\frac{2d}{\hbar}\int_0^T\|H(t)\|dt<\infty,
\]
and similarly for $G$. These facts justify the localization in Lemma~\ref{int:coupling}. At a node the law has zero probability in the vanishing sector; assigning a rate there has no effect on the trajectory law.

Duhamel's identity and unitarity give $\|\Psi_t-\Phi_t\|\le D(t)$. The initial population distance is at most the pure-state trace distance and hence at most $\delta_0$. Apply Lemma~\ref{int:coupling}, substitute the equivariant marginals $a_m^2,b_m^2$, and then apply Lemma~\ref{int:node-cancellation} at almost every time. Integration proves \eqref{int:path-error}. The block-norm estimates used only orthogonality of the sectors, so they are uniform over a finite inaccessible reference inside each fiber. More general time dependence is covered by the same estimate whenever equivariant nonexplosive laws and the localized coupling are separately established; it is not needed for the finite pulse theorem.
\end{proof}

For identical initial waves, $\int e\le\varepsilon$ and $\int h\le B$, this gives
\begin{equation}\label{int:simple-error}
 d_{\rm TV}(P_H,P_G)\le\min\{1,2d\varepsilon(1+6B)\}.
\end{equation}
Thus the finite checkpoint stability result does not require a node-cutoff error term. The use of the smaller occupation probability is essential; replacing it by a uniform rate norm would lose the cancellation. This is a completion of a conditional path-stability argument, not a derivation of the statistical generator. The finite sector resolution remains fixed in the limit; growing memory banks require control of $d$ as well as $B$.

\section{What a complete record comparison inherits}
\begin{corollary}[Path observables, stopping and rare records]\label{int:path-records}
For the hypotheses of Theorem~\ref{int:path-stability}, every common measurable function of the full path has output total variation bounded by \eqref{int:path-error}. This includes jump counts, first-entry times with an additional ``no entry'' value, and stopped paths. Conditioning on an event of ideal probability $p>0$ costs the factor $2/p$ in Lemma~\ref{found:errors}. A retained physical memory is covered only when it is included in the complete state and both compared record maps have the same physical meaning.
\end{corollary}
\begin{proof}
Total variation contracts under the common measurable map. The stopped output is such a map, with its stopping time and failure/null symbols included. Conditional control follows from Lemma~\ref{found:errors}. Merely naming a passive path functional does not construct a memory interaction; that additional physical claim is exactly why the last hypothesis is stated.
\end{proof}

If a finite-background variational source is used for one experiment, the triangle inequality adds its Part~\ref{part:statistics} path error $\varepsilon_{\rm bg}a_{\max}(d-1)T$. If a finite writer has error $\delta_{\rm write}$ and actual archive-crossing probability at most $\delta_{\rm corrupt}$, those are the additional errors in transferring recorded data to actual past labels. Final archive distributions alone do not provide the latter transfer.

\begin{counterexample}[The minimal-generator hypothesis is load-bearing]\label{int:nonminimal-boundary}
Take a two-sector stationary eigenwave with $w_0=w_1=1/2$ and zero current. For every nonzero $\epsilon$ let the Hamiltonian be $H_\epsilon=\hbar\epsilon\sigma_x$ and permit symmetric traffic $K_{10}=K_{01}=k/2$, $k>0$. The wave $(|0\rangle+|1\rangle)/\sqrt2$ is an eigenwave. The admitted nonminimal process jumps at rate $k$ in either direction; its probability of at least one jump by $T$ is $1-e^{-kT}$ even as $\epsilon\to0$. At exactly zero coupling an additional support rule can set $K=0$. This family obeys support at exact decoupling and all one-time weight equations, but violates uniform weak-coupling path stability. Theorem~\ref{int:path-stability} concerns the minimal Bell generator and does not exclude this different constitutive law.
\end{counterexample}

% End consolidated source: parts/07_stability.tex

\part{A Separate Massive-Configuration Completion}\label{part:massive}

% Integrated from the independently audited Event_Law_Resolution.tex.
% This fragment supplies chapters, not a new document or Part.
% Load massive_macros.tex in the book preamble and append massive_bibliography.tex
% within the book's existing thebibliography environment.

\chapter{A massive configuration constitution and its event law}
\label{mc:chapter-constitution}

Chapters~\ref{cfg:chapter} and~\ref{cfg:archives} established how an
admitted configuration law can support coherent records and faithful
archives. They did not select that law from the source/readout interface.
This chapter chooses a complete continuous constitution and then constructs
its material implementation. Its actual variables are massive positions
$Q$ guided by one uncollapsed spinor wave. Finite internal basis labels
remain amplitudes in that wave; they are not additional actual sectors.
This differs from the finite tagged configuration used by the Bell and
pilot constructions elsewhere in the book. Identical coherent gate algebra
does not identify their ontologies or microscopic path laws.

The velocity and equilibrium postulates below are standard Bohmian
ingredients, not new consequences of the current identities
\cite{mc:BohmI,mc:BohmII,mc:DGZeq,mc:DGZoperators}. In particular,
Lemma~\ref{cfg:quantile} already supplied the quantile-flow argument for a
first-order translation detector. Its generator
\eqref{cfg:translationH} is not lower bounded. The additional construction
here places the source, pointer, finite reaction stock, protection gates,
reset receivers and controller inside a single semibounded massive
Hamiltonian inventory. A controlled oscillator supplies the exact event
time; a separate derivative estimate controls whether later physical
archives remain true about their actual past.

The four chapters of this part have distinct roles. The present chapter
specifies the complete path law, proves conservative motion through nodes,
and solves the massive writer. Chapter~\ref{mc:chapter-material} supplies
the source--actuator--record chain, retained loss/null states, copies, reset
and bounded internal protection. Chapter~\ref{mc:chapter-autonomous}
includes the controller as matter and proves the retained-output and
actual-history bounds, culminating in Theorem~\ref{mc:closure}.
Chapter~\ref{mc:chapter-tests} tests the construction with a
reference-sensitive noncommuting experiment and competing actual motions.
All four chapters restore explicit $\hbar$ and use a fixed finite
nonrelativistic apparatus and observation horizon.

The result is a constitutive internal completion, with controlled
finite-resource approximation to ideal instruments. It does not derive
the original finite Bell law
$\lambda_{Y\leftarrow X}=[J_{YX}]_+/|\Psi_X|^2$ by reinterpreting a
continuous crossing. The Bell current-production and statistical-selection
results retain their own hypotheses. They are neither used to supply
random clocks here nor invalidated by this alternative.

\section{The complete material state and its initial law}
\label{mc:sec:constitution}
Use finitely many massive coordinates $q\in\mathbb R^n$, a finite internal material space $\mcH_I$ containing source, actuators, fuel, memories and spent products, and an inaccessible $R$. The complete state is
\begin{equation}\label{mc:eq:state}
 (\Psi,Q),\qquad \Psi\in L^2(\mathbb R^n;\mcH_I\otimes\mcH_R),\quad
 \norm{\Psi}=1,\quad Q\in\mathbb R^n.
\end{equation}
A quantum clock coordinate is appended in Section~\ref{mc:sec:clock}. Coordinates of otherwise passive source particles can be included in confining ground states. Internal spin is part of $\Psi$; there is no extra actual spin assignment.

\begin{assumption}[Universal material inventory]\label{mc:ax:material}
Every source contact, recorder, controller, protection device and reset receiver belongs to the same spinor Schr\"odinger inventory. Its admitted Hamiltonians are
\begin{equation}\label{mc:eq:H}
 H(t)=-\sum_{k=1}^n\frac{\hbar^2}{2m_k}\partial_k^2+V(q,t),
 \qquad V=V^\dagger,
\end{equation}
with self-adjoint semibounded realizations specified below. Every operation is identity on $R$. There is no additional classical device that reads a nonlinear function of a source ray without participating in this wave dynamics.
\end{assumption}
This last inventory statement is a physical restriction, not a theorem about every imaginable substance. It makes access and reaction compatible in this model: adding a contact changes $V$ and hence the wave controlling the actual motion. In particular a neutral ray meter from a different hybrid constitution cannot be appended without changing the theory.

\begin{assumption}[Kinetic-momentum motion]\label{mc:ax:motion}
Physical material positions have velocity given by the local real kinetic momentum per unit mass:
\begin{equation}\label{mc:eq:guidance}
 \dot Q_k(t)=v_k^\Psi(Q(t),t),\quad
 v_k^\Psi=\frac{j_k}{\rho},\quad
 \rho=\Psi^\dagger\Psi,\quad
 j_k=\frac{\hbar}{m_k}\operatorname{Im}\Psi^\dagger\partial_k\Psi.
\end{equation}
The formula is used only where $\rho>0$. There are no further random displacements or circulation terms.
\end{assumption}
This is the standard Bohmian law, with established provenance \cite{mc:BohmI,mc:BohmII,mc:DGZoperators}. Its independent physical content is differentiable material motion governed by wave kinetic momentum. For a scalar wave $\Psi=Re^{iS/\hbar}$ it says $m_kv_k=\partial_kS$. Separating the real Schr\"odinger equation gives
\[
 \partial_tS+\sum_k\frac{(\partial_kS)^2}{2m_k}+V
 -\sum_k\frac{\hbar^2}{2m_k}\frac{\partial_k^2R}{R}=0.
\]
Its gradient determines acceleration along the flow. This is a concrete mechanical postulate; it does not minimize a graph traffic functional. Symmetry or continuity alone does not force it, as Section~\ref{mc:sec:rivals} demonstrates.

\begin{assumption}[Initial complete equilibrium]\label{mc:ax:eq}
Conditional on each supplied classical preparation description $c$, the full initial configuration law is
\begin{equation}\label{mc:eq:equilibrium}
 \Prb(dQ_0\mid c)=\norm{\Psi_0^c(Q_0)}_{I,R}^2\,dQ_0.
\end{equation}
Fresh ready cells and the clock are supplied in specified product waves, independent of the unknown input. There is a finite stock. No subsequent reset is assumed to generate a fresh seed conditional on an arbitrarily exposed microscopic past.
\end{assumption}
This is the only stochastic/ensemble input in the adopted dynamics. Initial $Q$ and the deterministic flow generate every later probability. It is not derived from source/readout incompleteness, from equilibration, or from a typicality slogan. The equilibrium and conditional-wave literature distinguishes these issues \cite{mc:DGZeq}. For example a real stationary wave has $v=0$, so an initially nonequilibrium distribution is stationary too. Universal dynamical equilibration is false in this class without additional hypotheses.

\subsection{Self-adjointness, domains and nodes}
The driven library uses scalar confining quadratics, affine coordinate terms with finite Hermitian coefficients, bounded smooth matrix potentials, and finitely many smooth time windows. On each fixed finite programme, all coefficients and their required derivatives are bounded. Completing the square bounds affine forces below. Finite internal gates are bounded. The resulting oscillator operator plus infinitesimally oscillator-bounded affine terms and bounded potentials is self-adjoint on the oscillator domain. Smooth vectors are preserved on finite intervals. One can verify the last assertion by commuting $q^\alpha\partial^\beta$ through the equation: scalar quadratics keep total order fixed, affine terms lower derivative order, and bounded smooth terms contribute only lower derivatives. Finite sums of Gaussian packets used below are such vectors.

\begin{lemma}[Conservative motion through the nodal problem]\label{mc:lem:exist}
Suppose the wave is $C^2$ on each programme interval and
\begin{equation}\label{mc:eq:regularity}
 \int_0^T\left(\norm{\partial_t\Psi_t}_2+
       \sum_k\norm{\partial_k\Psi_t}_2^2\right)dt<\infty.
\end{equation}
For the smooth nonsingular inventory above, the guidance flow exists throughout $[0,T]$ for $\rho_0$-almost every initial position and pushes $\rho_0$ to $\rho_t$. It neither loses mass at a node nor reaches infinity in finite time with positive equilibrium probability.
\end{lemma}
\begin{proof}
Direct differentiation using the Hermitian potential gives $\partial_t\rho+\sum_k\partial_kj_k=0$. On compact subsets of $\rho>0$, the locally smooth velocity has a unique flow and the continuity equation gives partial equivariance up to its exit. Under this killed flow, the position distribution is dominated by $\rho_t$.

Expected distance travelled before exit is bounded by
\[
 \int_0^T\!\int |j|\,dq\,dt
 \le C\int_0^T\sum_k\norm{\partial_k\Psi_t}_2dt<\infty.
\]
Escape to infinity would require infinite distance. To control nodes, along a surviving path differentiate $\log\rho$. Its expected total variation is bounded by
\begin{align*}
 \int_0^T\!\int\left(|\partial_t\rho|
       +\frac{|j|\,|\nabla\rho|}{\rho}\right)dq\,dt
 &\le\int_0^T\left(2\norm{\partial_t\Psi_t}_2
                  +C\norm{\nabla\Psi_t}_2^2\right)dt<\infty.
\end{align*}
Here $|\nabla\rho|\le2|\Psi||\nabla\Psi|$ and $|j|\le C|\Psi||\nabla\Psi|$. Reaching a node while remaining in a bounded region would send $\log\rho$ to $-\infty$, which has probability zero by the preceding bound. Initial nodes have zero probability. Exhaust the local domains; there is no remaining loss of mass. Domination by the normalized density then becomes equality. Finitely many switches concatenate without a new random draw. This is the current-integrability argument underlying the general existence theorem of Teufel and Tumulka \cite{mc:TT}; no theorem for arbitrary singular potentials is imported.
\end{proof}
The autonomous Hamiltonian below has an additional scalar free kinetic term and smooth bounded functions of the clock multiplying affine pointer operators. The same commutator estimates, now also including clock weights and derivatives, give smooth finite-moment evolution and \eqref{mc:eq:regularity}. Its initial clock packet has finite momentum moments at each finite mass. Thus the lemma covers the \emph{complete} state, not just a reduced pointer.

\subsection{The law on complete histories}
Let $\Phi_{t,0}^\Psi$ be the almost-sure flow. The complete path measure is explicitly
\begin{equation}\label{mc:eq:pathlaw}
 \Prb(A)=\int 1_{\{(\Phi_{t,0}^\Psi(q))_{0\le t\le T}\in A\}}
                   \rho_0(q)\,dq .
\end{equation}
Continuous path space is standard Borel, so regular conditional laws for a finite or countably generated record history exist. For any past event $B$ of positive probability, the conditional future is the normalized restriction of the initial integral to its preimage under the flow. Given the complete initial state the future is deterministic. Given only coarse past records it is usually history dependent. Predictable crossing times need not have a compensator absolutely continuous in $dt$. No Bell intensity, Markov hazard or exponential threshold is asserted in this filtration.

Where the \emph{pointwise} current vanishes for an interval, positions are fixed. Current reversal reverses the corresponding instantaneous velocity, with its accumulated initial-position information retained. Nodes are handled by Lemma~\ref{mc:lem:exist}; conditioning on a zero-probability event is not assigned a normalized daughter.

\section{An exact massive detector in physical time}
\label{mc:sec:detector}
Let $A=\ket{1}\bra{1}$ act on a retained internal control label. First a finite internal unitary correlates this label with projectors $P_0,P_1$ of the unknown source, giving $\sum_a P_a\psi\ket a$. No actual internal jump is introduced by this operation. Prepare one oscillator in its ground packet
\[
 \phi_0(y)=(2\pi\sigma^2)^{-1/4}e^{-y^2/(4\sigma^2)},\qquad
 \sigma^2=\frac{\hbar}{2M\omega}.
\]
Choose a smooth centre trajectory $b(0)=0$, $b(T_w)=L$, with $\dot b=\ddot b=0$ at both ends, and set
\begin{equation}\label{mc:eq:force}
 c(t)=b(t)+\frac{\ddot b(t)}{\omega^2},\qquad
 H_w(t)=\frac{p_y^2}{2M}+\frac{M\omega^2}{2}\bigl(y-c(t)A\bigr)^2.
\end{equation}
This operator is nonnegative. No first-order unbounded-below translation is used. A convenient explicit choice is
\begin{equation}\label{mc:eq:quintic}
 b(t)=L(10s^3-15s^4+6s^5),\quad s=t/T_w,\quad
 \dot b=\frac{30L}{T_w}s^2(1-s)^2\ge0.
\end{equation}
Smooth higher-order endpoint interpolation can be used when all clock-window derivatives are required; the $C^2$ quintic suffices for the exact writer and the finite-order estimates here. The trap may overshoot the packet centre during acceleration; its finite displacement and force are resources.

\begin{proposition}[Exact wave and selected actual motion]\label{mc:prop:writer}
For this primitive, $\psi$ is an internal/reference vector, with no unresolved older spatial coordinates. Writing $p_a=\norm{P_a\psi}^2$, the wave is
\begin{align}\label{mc:eq:packet}
 \Psi_t(y)&=P_0\psi\ket0\,e^{-i\omega t/2}\phi_0(y)
 +P_1\psi\ket1\,e^{i\theta(t)}e^{iM\dot b(t)(y-b(t))/\hbar}\phi_0(y-b(t)),\\
 \dot\theta&=\frac{M\dot b^2}{2\hbar}-\frac{M\omega^2(b-c)^2}{2\hbar}-\frac\omega2.\nonumber
\end{align}
For $g_\sigma=|\phi_0|^2$,
\begin{equation}\label{mc:eq:mixture}
 \rho_t(y)=p_0g_\sigma(y)+p_1g_\sigma(y-b(t)),\qquad
 j_t(y)=p_1\dot b(t)g_\sigma(y-b(t)).
\end{equation}
Let $F_t(y)=p_0\mcNcdf(y/\sigma)+p_1\mcNcdf((y-b(t))/\sigma)$, where $\mcNcdf$ is the standard normal CDF. The complete actual pointer path is
\begin{equation}\label{mc:eq:quantile}
 Y_t=F_t^{-1}(U),\qquad U=\mcNcdf(Y_0/\sigma)\sim\operatorname{Unif}(0,1).
\end{equation}
In particular $0\le\dot Y_t\le\dot b(t)$ during the monotone write.
\end{proposition}
\begin{proof}
Substitute the Gaussian ansatz into the Schr\"odinger equation. The coefficient of $y-b$ is $M\ddot b=-M\omega^2(b-c)$ and its scalar coefficient is precisely the displayed $\dot\theta$; the Gaussian width stays at its ground value. Internal labels are orthogonal, so there are no cross terms in $\rho$ or $j$. Now $\partial_tF_t=-j_t$ and $\partial_yF_t=\rho_t>0$. Differentiating $F_t(Y_t)=U$ gives \eqref{mc:eq:guidance}. The initial inverse transform supplies the uniform rank, without a second randomness postulate.
\end{proof}

\subsection{Actual timing, false-ready tails and null continuation}
Put a threshold $h=L/2$. Let $\tau=0$ for the initially right-hand tail $Y_0\ge h$, and otherwise let $\tau$ be its first subsequent threshold crossing, with $\tau=\infty$ if none occurs before $T_w$. This convention retains the finite false-ready tail
\[
 \delta=\mcNtail\left(\frac{L}{2\sigma}\right),\qquad \Prb(\tau=0)=\delta.
\]
It is not an assertion that a preliminary check of readiness is noninvasive.
Monotonicity in Proposition~\ref{mc:prop:writer} gives the complete law
\begin{align}\label{mc:eq:timing}
 \Prb(\tau>t)&=F_t(h),\qquad
 \Prb(\tau\in dt)=p_1\dot b(t)g_\sigma(h-b(t))\,dt\quad(0<t<T_w),\\
 \Prb(\tau=\infty)&=p_0(1-\delta)+p_1\delta.\nonumber
\end{align}
The positive-time density integrates to $p_1(1-2\delta)$; together with the initial atom and final null it normalizes to one. Conditional on no pre-trigger event the survival is $F_t(h)/F_0(h)$; conditional on no crossing by $t$, its instantaneous hazard, where defined, is
\begin{equation}\label{mc:eq:hazard}
 \frac{p_1\dot b(t)g_\sigma(h-b(t))}{F_t(h)}.
\end{equation}
This is a derived physical-time formula, not a memoryless source-sector clock. Continuing a return pulse uses the same rank $U$, not a newly sampled waiting time. A dark hold with $\dot b=0$ has zero pointer current.

If an input is entangled with older spatial memories, the full velocity is evaluated before integrating those coordinates out. When they are held fixed, the quantile argument applies separately to their conditional spinor fibres, with the corresponding conditional $p_a$. The integrated $p_a$ still determines marginal density but generally does not determine an individual joint trajectory. For moving old coordinates, use the complete guidance equation rather than this one-dimensional primitive formula.

The conditional position law after a null at $t$ is
\[
 \Prb(Y_t\in dy\mid\tau>t)=\frac{1_{y<h}\rho_t(y)}{F_t(h)}\,dy.
\]
The global wave remains \eqref{mc:eq:packet}, including both packets. For an arbitrary past record event $B$, equation~\eqref{mc:eq:pathlaw} rather than a present-sector projection supplies the conditional continuation. A timestamp reader would be an additional interaction and would change this wave; equation~\eqref{mc:eq:timing} describes this specified pointer without an extra timestamp apparatus. A finite timestamp claim requires those contacts and their complete dynamics; the first-crossing formula alone does not construct that reader. Unmeasured arrivals and recorded detection times need not coincide \cite{mc:Arrival}.

\begin{figure}[tb]
\centering\includegraphics[width=\textwidth]{massive_event_mechanism_figure.pdf}
\caption{The massive writer with $L/\sigma=8$, $p_1=0.65$ and $T_w=1$.
The first panel compares the trap centre with the packet centre, the second
shows selected actual mixture-quantile trajectories and the threshold, and
the third gives the crossing-time density and cumulative probability.
These are evaluations of the derived equations. The finite initial atom
and terminal null are both retained in the timing law.}
\label{mc:fig:mechanism}
\end{figure}

\subsection{Energy and force resources}
In branch 1,
\begin{equation}\label{mc:eq:energy}
 \langle H_w(t)\rangle=\frac{\hbar\omega}{2}
       +\frac M2\dot b^2+\frac{M\omega^2}{2}(b-c)^2.
\end{equation}
It is finite for the explicit trajectory. The work in the driven description is $\int\langle\partial_tH_w\rangle dt$. It vanishes between the initial ground state and the final ground state of the shifted holding trap, but nonzero energy is borrowed and returned during the pulse. The autonomous clock below carries this exchange. Zero net work is not zero transient work or an unlimited source of reset readiness.

\chapter{Massive material records, retained resources and protection}
\label{mc:chapter-material}

The point of the material construction is to use the same Hamiltonian for
source response and physical access. The finite resource states below
describe coherent fuel, excitation and remnant amplitudes. Their algebra
can also appear in a finite Bell or pilot model, but here a declaration
occurs only through the actual position of a massive pointer. A resource
amplitude is therefore not a sampled reaction history. Every null, loss,
copy and reset receiver remains available to later coherent contacts.

The endpoint reasoning of Theorem~\ref{cfg:endpoint} is retained, while
truth about an earlier actual declaration requires a further fact:
pointwise zero archive current during its promised hold. This chapter
proves that fact for stored Gaussian records, and preserves it through the
explicit transported-trap reset. Its finite tails enter the error budget;
they are not replaced by exact compact support.

\section{A finite coherent source--actuator--resource module}
\label{mc:sec:resource}
Here is a nontrivial receptor that is physically in the same inventory. On a finite factor $D$ use orthogonal states
\[
 \ket r,\quad\ket{p_a},\quad\ket{c_a},\quad\ket{l_a}\quad(a=0,1).
\]
They include the following actual material degrees in their internal wave description:
\begin{center}\small
\begin{tabular}{@{}lcccccc@{}}\toprule
State & Production & Fuel & Site & Excitation & Memory & Remnant\\\midrule
$r$ & 1 & 1 & ready & 0 & blank & vacuum\\
$p_a$ & 0 & 1 & ready & mode $a$ & blank & vacuum\\
$c_a$ & 0 & 0 & spent & 0 & $a$ & capture $a$\\
$l_a$ & 0 & 1 & ready & 0 & blank & loss $a$\\\bottomrule
\end{tabular}
\end{center}
Assign energy $E>0$ to a production cofactor, fuel unit, excitation and loss remnant, and $2E$ to a capture remnant; the displayed labels are degenerate. Every row has total resource energy $2E$. Hence the conversion gates conserve this resource energy exactly while spending readiness and retaining energy in products. Their full tensor-factor implementation is defined to be zero outside the indicated equal-energy active subspace. Other exhausted sectors stay present.

For source projectors $P_a$, set
\begin{align}\label{mc:eq:reactionG}
 G_w&=i\sum_aP_a\otimes(\ket{p_a}\bra r-\ket r\bra{p_a}),\\
 \ket{b_a}&=\sqrt\eta\ket{c_a}+\sqrt{1-\eta}\ket{l_a},\qquad 0<\eta<1,\nonumber\\
 G_r&=i\sum_a(\ket{b_a}\bra{p_a}-\ket{p_a}\bra{b_a}).\nonumber
\end{align}
These Hermitian generators have norm one on their active subspaces. Nonoverlapping pulses $\hbar g_w(t)G_w$ and $\hbar g_r(t)G_r$ with areas $\theta,\varphi$ give exactly
\begin{equation}\label{mc:eq:resourcewave}
 \begin{split}
 \Psi_D={}&\cos\theta\,\psi\ket r+\sin\theta\sum_aP_a\psi
 \left(\cos\varphi\ket{p_a}+\sin\varphi\sqrt\eta\ket{c_a}
                    +\sin\varphi\sqrt{1-\eta}\ket{l_a}\right).
 \end{split}
\end{equation}
Indeed each generator is a two-dimensional $\sigma_y$ rotation, and the active $a$ sectors are orthogonal. There was no sampled reaction time in this calculation. Reduced excitation populations are not actual level trajectories in the adopted ontology. The actual event is a subsequent spatial registration governed by Sections~\ref{mc:sec:constitution}--\ref{mc:sec:detector}.

The finite response has four exact orthogonal status weights (and ideal resolved-pointer probabilities):
\begin{equation}\label{mc:eq:statuses}
 (P_r,P_p,P_c,P_l)=
 (\cos^2\theta,\ \sin^2\theta\cos^2\varphi,\
 \eta\sin^2\theta\sin^2\varphi,\ (1-\eta)\sin^2\theta\sin^2\varphi).
\end{equation}
The pending excitation remains a vector in the actual model at a finite cutoff. Continuing $G_r$ processes it coherently; a finite closed receptor can recur. A zero response window does not erase it. Neither an absorbing boundary nor a restart clock is imposed when a coefficient begins to populate $p_a$.

\subsection{Physical null, capture and loss}
Use three spatial readout centres $-L,0,L$ for captured $a=0$, null, and captured $a=1$. The same forced oscillator construction applies to each orthogonal control projector, using signed trajectories. Noncaptured $r,p,l$ components share the null packet. Nearest-centre cells have worst tail at most $2\delta$, with $\delta=\mcNtail(L/(2\sigma))$. The ideal orthogonal-label comparator has capture coefficient
\[
 \sqrt q\,P_a\psi\ket{c_a},\qquad q=\eta\sin^2\theta\sin^2\varphi,
\]
and complete null vector
\begin{equation}\label{mc:eq:nullvector}
 \Psi_N=\cos\theta\,\psi\ket r+
 \sin\theta\sum_aP_a\psi\left(\cos\varphi\ket{p_a}
                 +\sin\varphi\sqrt{1-\eta}\ket{l_a}\right).
\end{equation}
Only for a declared reduced comparison, tracing $D$ gives
\begin{equation}\label{mc:eq:nullmap}
 \mathcal N(\rho)=\cos^2\theta\,\rho+
 \sin^2\theta\bigl(\cos^2\varphi+(1-\eta)\sin^2\varphi\bigr)
                  \sum_aP_a\rho P_a.
\end{equation}
Equation~\eqref{mc:eq:nullvector}, along with pointer and all receiving systems, is the retained continuation. Equation~\eqref{mc:eq:nullmap} is not a global collapse rule. Finite spatial classifiers approximate these ideal labels; Section~\ref{mc:sec:complete} bounds the complete output error.

\subsection{Finite stock and exhaustion}
For a promised $m$-epoch experiment allocate $m$ ready cells, their blank archives, and receivers. Gate the active conversion only on sectors containing the required cofactor, fuel, site and blank capacity. Extend the unitary by identity on explicitly exhausted sectors, which can be spatially flagged by the same writer. A failed or null attempt does not receive a new $\ket r$ for free. The consumed-ready-cell count is bounded by the allocated finite schedule; no infinite Poisson bath is hidden in this implementation.

\section{Physical records that remain true about their past}
\label{mc:sec:archive}
After a write, retain its internal orthogonal key $K$ and hold the pointer in
\begin{equation}\label{mc:eq:store}
 H_{\rm store}=\frac{p_y^2}{2M}+\frac{M\omega^2}{2}(y-LK)^2.
\end{equation}
Known branch phases can be corrected by bounded internal potentials. A single stationary packet need not have compact support; its classification error was already accounted for.

\begin{theorem}[Exact historical storage]\label{mc:thm:store}
Suppose the wave after the write has form
\begin{equation}\label{mc:eq:storedwave}
 \Psi(y,z,t)=\sum_k\phi_0(y-Lk)\ket k_K\,\Xi_k(z,t)e^{-i\omega t/2},
\end{equation}
where $z$ denotes all other coordinates and internal factors. Future gates preserve $K$, hold $y$ as in \eqref{mc:eq:store}, and may be noncommuting on the source or act on $z$. Then $j_y=0$ \emph{pointwise} on the complete configuration space. The actual coordinate $Y$ and its finite readout label remain exactly fixed throughout that interval.
\end{theorem}
\begin{proof}
Orthogonality of the key eliminates cross terms. Each remaining contribution to $\Psi^\dagger\partial_y\Psi$ is $\phi_0(y-Lk)\phi_0'(y-Lk)\norm{\Xi_k(z,t)}^2$, which is real. Equation~\eqref{mc:eq:guidance} gives the conclusion away from the almost-sure excluded nodes.
\end{proof}
This is a path statement, not an inference from equal endpoint weights. It controls actual old declarations even when later source measurements do not commute with the first. Unknown interactions that violate its Hamiltonian conditions require their own bound.

\subsection{A genuine copy and its classification error}
Copy the key by a finite reversible unitary into a blank internal factor, amplify that factor into a fresh massive pointer, and retain both. Throughout this operation the old key is preserved, so Theorem~\ref{mc:thm:store} holds for the old actual position. If the two classifiers have worst errors $\epsilon_1,\epsilon_2$, their disagreement probability is at most $\epsilon_1+\epsilon_2$. To prove it, expand their joint squared-norm density over the orthogonal key and apply a union bound to the two conditional Gaussian tails. The key in this proof is an orthogonal expansion index, not an additional secretly actual spin variable. The copy is faithful to the first \emph{actual} declaration because the old pointer stayed fixed during the write; its finite misclassification remains in the bound.

\subsection{Reset with the receiving system retained}
Supply an identical ready factor $D'$ and let $W_{DD'}$ be SWAP. With
\begin{equation}\label{mc:eq:swap}
 H_{\rm sw}=\frac{\pi\hbar}{2\tau_s}W_{DD'},\quad
 S(t)=e^{-itH_{\rm sw}/\hbar},\quad S(\tau_s)=-iW_{DD'},
\end{equation}
an old entangled state is transferred as
\[
 \sum_j\psi_j\ket{d_j}_D\ket r_{D'}\longmapsto
 -i\ket r_D\sum_j\psi_j\ket{d_j}_{D'}.
\]
The old pending excitation, remnant, lost product and reference correlation remain in $D'$. Identical free resource Hamiltonians have $[W,H_D+H_{D'}]=0$.

There is a subtle control issue: if a stationary trap followed the old $D$ label, a bare SWAP would change its centre. The exact physical repair is
\begin{equation}\label{mc:eq:covariantreset}
 H(t)=H_{\rm sw}+S(t)H_{\rm store}S(t)^\dagger.
\end{equation}
Its propagator is $S(t)e^{-itH_{\rm store}/\hbar}$ by differentiation. Because $S$ is coordinate independent, it preserves the position density and current pointwise. The old spatial record remains held while the trap's controlling key is transferred to $D'$. The potential is a unitary conjugate of a nonnegative matrix potential plus a bounded matrix, so it stays semibounded. Expanding its square gives a scalar quadratic term and affine matrix coefficients, within the common inventory. Simply declaring a rewired trap after SWAP would omit this interaction.

If the original pointer itself is to be restored, a reverse smooth forced trap can take its branch centre back to zero while a copied key/record or receiving cell is retained. The receiving systems carry the old correlations. Future use proceeds from this full state and its conditional law; it is not assigned an independent fresh initial rank merely because a local packet now looks ready. The finite measurement theorem below uses a finite stock of fresh pointers; reset is included as a real operation and as a return test.

\subsection{Feedback from a literal spatial record}
An internal key-controlled source gate is an exact coherent operation, but it follows a finite position display only up to the display's error. Literal position feedback also belongs to \eqref{mc:eq:H}. Let $g(y)$ be smooth, $0\le g\le1$, equal to zero for $y\le h-r$ and one for $y\ge h+r$, with $0<r<L/2$. Let $B$ be a bounded Hermitian source generator and compare
\[
 H_{\rm pos}=H_{\rm store}+g(y)B,\qquad
 H_{\rm key}=H_{\rm store}+KB.
\]
The second gives the intended branch gate. Define
\begin{equation}\label{mc:eq:feedbacktail}
 \epsilon_g=\max_{k=0,1}\int |g(y)-k|^2|\phi_0(y-Lk)|^2dy
 \le\mcNtail\left(\frac{L/2-r}{\sigma}\right).
\end{equation}
Duhamel, evaluated on the exactly stationary ideal packet, gives
\begin{equation}\label{mc:eq:feedbackbound}
 \norm{\Psi_{\rm pos}(t)-\Psi_{\rm key}(t)}
 \le\frac{t\norm B}{\hbar}\sqrt{\epsilon_g}.
\end{equation}
This bound is uniform in an inaccessible reference. The physical position contact has reciprocal backaction; it is not claimed to leave the actual pointer fixed. The derivative and crossing estimate in Section~\ref{mc:sec:clock} supplies a history bound as well. Thus literal spatial feedback, rather than an idealized outside observer, has a complete implementation.

Here $g(y)$ is a multiplication operator on the wave, not a coefficient obtained by inserting the actual $Y_t$ into an externally controlled Hamiltonian.

\section{Protected coherent transport in the same inventory}
\label{mc:sec:protection}
The code and gap argument of Chapter~\ref{prot:gapchapter} apply to the
finite internal bank in the present inventory. They concern coherent
Hamiltonian transport, so their proof survives the change of actual
ontology. We give the complete bounded-bank estimate here with explicit
$\hbar$ and with the spatial-spectator condition stated before its use.
It supplies protection, not a statistical selection of guidance.

Actual archive storage and protection of unknown logical amplitudes are different tasks. The bounded internal gap construction can be implemented here without a stochastic interface. To display its nonempty domain, encode two qubits into four by
\[
 C\ket{a,b}=\frac{\ket{0,a,b,a\oplus b}+\ket{1,1\oplus a,1\oplus b,1\oplus a\oplus b}}{\sqrt2}.
\]
Let $S_X=X_1X_2X_3X_4$, $S_Z=Z_1Z_2Z_3Z_4$, $P=(I+S_X)(I+S_Z)/4$ and
\[
 H_{\rm pen}=\tfrac\Delta2(I-S_X)+\tfrac\Delta2(I-S_Z),\quad
 H_\Delta=H_{\rm pen}+H_0+V.
\]
Assume $[H_0,P]=0$, $\norm{H_0}\le b$, and
\[
 V=\sum_{i=1}^4\sum_{\alpha=x,y,z}\sigma_i^\alpha\otimes B_{i\alpha},
 \quad B_{i\alpha}=B_{i\alpha}^\dagger,\quad\norm V\le v.
\]
The $B$'s act on retained finite internal nuisance systems. During the protected exposure, the spatial holding Hamiltonian is a commuting spectator; it is factored out. We do not assert a bounded-norm theorem for arbitrary unbounded coordinate couplings. One-site Paulis anticommute with a stabilizer, so $PVP=0$. All encoding and decoding gates are finite internal unitaries.

\begin{proposition}[Retained-bank gap bound]\label{mc:prop:gap}
For $\Delta-2b-v=\gamma>0$,
\begin{equation}\label{mc:eq:gap}
 \norm{\bigl(e^{-itH_\Delta/\hbar}-e^{-itH_0/\hbar}\bigr)P}
 \le\min\left\{2,\frac{2v+tv^2/\hbar}{\gamma}\right\}.
\end{equation}
The same bound holds with every inaccessible reference and retained internal nuisance system included.
\end{proposition}
\begin{proof}
Block the Hamiltonian into $H_d=\operatorname{diag}(A,D)$ and off-diagonal $W$, with lower spectral separation $\gamma$ and $B=QVP$. The integral $X=\int_0^\infty e^{-rD}Be^{rA}dr$ solves $DX-XA=B$ and has norm at most $v/\gamma$. With $S=\left(\begin{smallmatrix}0&-X^\dagger\\X&0\end{smallmatrix}\right)$, $[S,H_d]=-W$ and
\[
 e^SH_\Delta e^{-S}=H_d+\int_0^1u e^{uS}[S,W]e^{-uS}du.
\]
The remainder is bounded by $v^2/\gamma$. Two changes of frame cost $2v/\gamma$ and Duhamel costs $tv^2/(\hbar\gamma)$. On $P$, $H_d=PH_0P$. This proves the claim.
\end{proof}
This is the book's coherent protection estimate, with its assumptions preserved; Hamiltonian error suppression has independent primary precedent \cite{MarvianLidar}. Its role here is compatibility with actual material writes and records, not selection of a noise generator.

\chapter{An autonomous controller and complete retained histories}
\label{mc:chapter-autonomous}

The driven library in Chapter~\ref{mc:chapter-material} specifies exact
contacts but still uses a pulse schedule. Here the schedule is supplied by
a massive quantum coordinate that belongs to the complete initial state.
Its recoil and correlations remain in the final comparison. The main
analytic distinction is between small wave error and reliable history:
the former controls final records, whereas the latter needs derivative
control at each actual archive decision surface.

Theorem~\ref{mc:history} supplies that additional control.
Theorem~\ref{mc:closure} then combines the material modules with all
receivers and the inaccessible reference. The result concerns complete
retained quantum outputs and specified actual declarations. It is not a
claim of total-variation continuity for every unrecorded guidance path.

\section{A massive controller within the common Hamiltonian}
\label{mc:sec:clock}
External pulse timing is a physical resource. We now include its provider and its recoil in the same Hamiltonian. This step also avoids an invalid inference from small wave error to small path error.

\subsection{The complete autonomous Hamiltonian}
Resolve a finite smooth driven programme as
\begin{equation}\label{mc:eq:program}
 H_{\rm id}(t)=H_{\rm osc}+H_{\rm const}+\sum_{i=1}^N f_i(x_0+vt)B_i(q),
 \qquad H_{\rm osc}=\sum_k\left(\frac{p_k^2}{2m_k}
                         +\frac{m_k\omega_k^2q_k^2}{2}\right).
\end{equation}
The coordinate-independent Hermitian matrix $H_{\rm const}$ is bounded. All $m_k,\omega_k$ are strictly positive. The $f_i$ are real bounded smooth profiles with bounded derivatives. The $B_i$ are affine Hermitian matrix functions of $q$, possibly plus bounded smooth matrix functions with bounded derivatives. This covers the forced trap (expand its square), finite internal rotations, transported trap controls \eqref{mc:eq:covariantreset}, and smooth position feedback. Any finite coordinate-independent internal unitary $S(t)$ while stored oscillators are present can be implemented exactly by
\[
 i\hbar\dot S(t)S(t)^\dagger+S(t)H_{\rm store}S(t)^\dagger.
\]
Its kinetic term is unchanged and its scalar quadratic term is unchanged; only finitely many bounded or affine matrix coefficients vary. This provides a direct finite gate compiler within \eqref{mc:eq:program}.

For an exactly $C^\infty$ trap programme use a flat positive bump $u(s)$ on $(0,1)$ and
\[
 b(t)=L\frac{\int_0^{t/T_w}u(s)ds}{\int_0^1u(s)ds},\qquad
 c=b+\ddot b/\omega^2.
\]
It is monotone and flat at both endpoints, so the exact writer proof is unchanged. The quintic in the figure instead gives a continuous piecewise-smooth trap profile; smoothing it has a directly bounded integrated residual. No globally smooth extension of its nonzero endpoint third derivative is presumed.

Add a massive clock coordinate $x$, prepared in a Gaussian $\chi_0$ with mean $x_0$, position deviation $s_c$, and mean momentum $M_cv$. Define
\begin{equation}\label{mc:eq:auto}
 H_{\rm aut}=\frac{P_x^2}{2M_c}+H_{\rm osc}+H_{\rm const}+
                              \sum_i f_i(x)B_i(q).
\end{equation}
This is autonomous and semibounded: bounded profile coefficients and oscillator confinement absorb each affine force by Young's inequality. It is self-adjoint on the free-clock-plus-oscillator domain, with the relative bound of the affine perturbation arbitrarily small. There is no read of the \emph{actual} clock position followed by an external switch. The quantum potential $f_i(x)B_i(q)$ is the interaction itself, and the actual clock follows \eqref{mc:eq:guidance} on the full wave.

The initial wave is $\chi_0\otimes\psi_0$, including all prepared apparatus and retained resources. Let $F_t$ be its exact evolution under \eqref{mc:eq:auto}. Compare it to
\[
 G_t=\chi_t\otimes\psi_t,\qquad
 \chi_t=e^{-itP_x^2/(2M_c\hbar)}\chi_0,\qquad
 i\hbar\dot\psi_t=H_{\rm id}(t)\psi_t.
\]
The free clock has mean $x_0+vt$ and width
\begin{equation}\label{mc:eq:clockwidth}
 s_t=\sqrt{s_c^2+\left(\frac{\hbar t}{2M_cs_c}\right)^2}.
\end{equation}
The comparator includes the clock; it is not a reduced apparatus state.

\subsection{Derivative control uniform in clock resources}
Choose fixed reference length units for the pointer coordinates. For $r\ge0$, write
\[
 W_r(F)=\sum_{|\alpha|+|\beta|\le r}\norm{q^\alpha\partial_q^\beta F}_2.
\]
Dimensional powers of those fixed length units are understood; they can equivalently be inserted term by term. The norm integrates over $x,q$ and all internal/reference indices, but differentiates only $q$.

\begin{lemma}[Uniform pointer graph norm]\label{mc:lem:graph}
For the fixed finite inventory in \eqref{mc:eq:auto},
\begin{equation}\label{mc:eq:graph}
 W_r(e^{-itH_{\rm aut}/\hbar}F)\le e^{\kappa_rt/\hbar}W_r(F)
 \quad(0\le t\le T).
\end{equation}
The constant depends on the pointer inventory and the profile bounds, but not on $M_c,s_c$, the mean clock momentum, or the reference dimension. The same estimate holds for the time-dependent ideal propagator.
\end{lemma}
\begin{proof}
For each scalar differential monomial $O=q^\alpha\partial_q^\beta$, $[O,P_x^2]=0$ and $[O,H_{\rm const}]=0$. Its commutator with the scalar oscillator is a finite sum of monomials of total order at most $r$. An affine potential removes a derivative in each nonzero commutator; multiplication by bounded smooth functions contributes bounded coefficient terms of no higher order. Thus
\[
 \sum_{|\alpha|+|\beta|\le r}\norm{[q^\alpha\partial_q^\beta,H_{\rm aut}]F}
 \le \kappa_rW_r(F).
\]
Matrices need not commute with each other: only their commutators with scalar coordinate operators have been used. Commute $O$ through the propagator, apply Duhamel and unitarity, sum, and use Gronwall. These identities hold first on the smooth core; oscillator graph-norm regularization and the same uniform estimate extend them to the displayed domain. The time-dependent case uses uniform coefficient bounds. Tensoring an identity does not change any estimate.
\end{proof}

\begin{theorem}[Complete autonomous approximation]\label{mc:thm:clock}
Assume $W_3(\psi_0)<\infty$. For $r=0,2$ define
\begin{equation}\label{mc:eq:epsclock}
 \varepsilon_r(T)=\frac1\hbar\int_0^T e^{\kappa_r(T-t)/\hbar}s_t
                         \sum_i\operatorname{Lip}(f_i)W_r(B_i\psi_t)dt.
\end{equation}
Take $\kappa_0=0$ for the $L^2$ propagation estimate, by unitarity. Then
\begin{equation}\label{mc:eq:clockbound}
 \sup_{t\le T}W_r(F_t-G_t)\le\varepsilon_r(T)
 \le A_{r,T}\left(s_c+\frac{\hbar T}{2M_cs_c}\right),
\end{equation}
where $A_{r,T}$ is finite and independent of the clock resources. All clock, fuel, receiver, record and reference factors remain in this comparison.
\end{theorem}
\begin{proof}
The defect of $G_t$ under the exact Hamiltonian is
\[
 R_t=\sum_i[f_i(x)-f_i(x_0+vt)]\chi_t(x)\otimes B_i\psi_t.
\]
The coordinate factor separates under every $q$ derivative and multiplier, so
\[
 W_r(R_t)\le s_t\sum_i\operatorname{Lip}(f_i)W_r(B_i\psi_t).
\]
Affine multiplication needs at most $W_{r+1}$ of the ideal wave; bounded smooth terms need $W_r$. Lemma~\ref{mc:lem:graph} bounds these on a fixed horizon. Apply Duhamel in the invariant $W_r$ domain and \eqref{mc:eq:graph}; $s_t\le s_c+\hbar T/(2M_cs_c)$ proves the last inequality. The exact product initial state makes the initial defect zero.
\end{proof}
For example $s_c=\sqrt{\hbar T/(2M_c)}$ gives error $O(M_c^{-1/2})$ for fixed apparatus. Its mean initial free clock energy is
\begin{equation}\label{mc:eq:clockenergy}
 E_C=\frac{M_cv^2}{2}+\frac{\hbar^2}{8M_cs_c^2}
     =\frac{M_cv^2}{2}+\frac{\hbar}{4T}.
\end{equation}
Every finite member has finite energy and normalizable resources. The ideal limit requires increasing mass/energy; Gaussian packets have unbounded support and are not claimed to have a strict energy cutoff. Total $H_{\rm aut}$ energy is conserved. The clock can recoil and entangle: \eqref{mc:eq:clockbound} bounds its complete discrepancy instead of deleting it. These are finite-horizon claims, not a perfect autonomous clock for all time.

\subsection{Why the same estimate controls actual archive history}
A small $L^2$ wave error alone does not control guidance paths. The $W_2$ estimate provides the extra information needed for a \emph{specified retained record surface}. Let $\Sigma=\{q_k=h\}$ be one such decision surface. The one-coordinate, Hilbert-valued trace estimates imply
\[
 \norm{F|_\Sigma}_2\le C_{\rm tr}W_1(F),\qquad
 \norm{\partial_kF|_\Sigma}_2\le C_{\rm tr}W_2(F).
\]
For completeness, $\norm{u(h)}^2\le2\norm u\norm{u'}$ follows by integrating the derivative of $\norm{u(s)}^2$ on a half-line; apply it also to $u'$. All other coordinates and the reference are Hilbert-valued parameters.

\begin{theorem}[Autonomous historical archive protection]\label{mc:history}
During a hold interval $I\subset[0,T]$, suppose the ideal wave has $j_k[G_t]=0$ pointwise on $\Sigma$ and $W_2(G_t)\le B$. Let $W_2(F_t-G_t)\le\varepsilon_2$. In equilibrium for the exact autonomous dynamics,
\begin{equation}\label{mc:eq:archiveerror}
 \Prb(\text{the record side of }\Sigma\text{ changes during }I)
 \le\frac{\hbar}{m_k}C_{\rm tr}^2|I|\,
                     \varepsilon_2(2B+\varepsilon_2).
\end{equation}
Sum this bound for finitely many retained record surfaces. Add their write/readout errors separately.
\end{theorem}
\begin{proof}
Write $E=F-G$. Expanding $F^\dagger\partial_kF-G^\dagger\partial_kG$ and applying the two trace estimates and Cauchy--Schwarz gives
\[
 \int_\Sigma|j_k[F]-j_k[G]|\le
 \frac\hbar{m_k}C_{\rm tr}^2\varepsilon_2(2B+\varepsilon_2).
\]
This bounds \emph{absolute} flux; cancellation of signed currents is not enough.
To avoid a hidden transversality assumption, let $s_\delta(q_k)$ smoothly approximate the indicator of one side of $\Sigma$, with $s_\delta'\ge0$ and $\int s_\delta'=1$. Along almost every complete trajectory,
\[
 \operatorname{Var}_{I}s_\delta(Q_k)\le
 \int_I |s_\delta'(Q_k)|\,|v_k(Q,t)|dt.
\]
Equivariance makes its expected right side $\int_I\int |s_\delta'(q_k)||j_k[F]|dq\,dt$. A genuine change of side contributes at least one to the limiting variation. Hilbert-valued traces make the current continuous in the normal coordinate as an $L^1$ function of the other coordinates. Fatou and the approximate-identity limit bound its probability by $\int_I\int_\Sigma|j_k[F]|$. Insert the preceding inequality and the pointwise-zero ideal current. Lemma~\ref{mc:lem:exist} already handles nodes; no positive lower density is assumed.
\end{proof}
The ideal stored wave \eqref{mc:eq:storedwave} supplies the required pointwise zero, including during noncommuting continuation on the other factors. The transported reset \eqref{mc:eq:covariantreset} also preserves that current. Finite clock tails therefore produce a quantified finite-horizon historical error, rather than being incorrectly declared harmless from endpoint equivariance.

For position feedback, keep a separate completed archive $z$ with its own immutable key. The working pointer $y$ may recoil under $g(y)B$, while $z$ still has zero ideal current by Theorem~\ref{mc:thm:store}. If preservation of $y$ is desired as well, repeat the graph-norm proof with residual $(g(y)-K)B\psi_t$. Its $W_2$ norm is computed by differentiating the known Gaussian and $g$ twice. Since $g-k$ and its derivatives are supported in the wrong half-line or central buffer, this norm is bounded by a finite polynomial in $L,\sigma^{-1},\norm{g'}_\infty,\norm{g''}_\infty$ times
\[
 \exp\left[-\frac{(L/2-r)^2}{4\sigma^2}\right].
\]
The graph propagation constant grows at most exponentially in $L$ for a fixed duration and other fixed parameters, because the affine trap coefficient is linear in $L$. Thus this derivative error, and its surface-flux budget, tend to zero as $L/\sigma\to\infty$ at fixed $\sigma,r$. This is an actual controlled limit, not a raw-path TV assertion.

\section{Complete instruments, references and finite histories}
\label{mc:sec:complete}
Let $\{K_a\}$ be a finite family on the unknown input satisfying $\sum_aK_a^\dagger K_a=I$. The map
\[
 \psi\ket{\mathrm{blank}}\longmapsto\sum_aK_a\psi\ket a
\]
is an isometry because it preserves inner products. Extend an orthonormal basis of its range to a full basis to obtain a finite unitary. A finite Hermitian logarithm supplies a bounded pulse. Alternatively the explicit resource rotations above give a fixed nontrivial family directly. The gate compiler in Section~\ref{mc:sec:clock} implements these gates while retaining all spatial storage. This argument concerns preparation-independent linear finite instruments; it does not admit arbitrary nonlinear ray maps.

\subsection{The exact finite pointer output and a strong comparator}
For one stage the physical isometry has form
\begin{equation}\label{mc:eq:W}
 W\psi=\sum_a K_a\psi\ket a_K\phi_a(q),
 \qquad\phi_a(q)=\phi_0(q-La)
\end{equation}
(or its finite multicoordinate version). Let $\Gamma_a$ be the disjoint physical readout regions and
\[
 \delta_a=\int_{\Gamma_a^c}|\phi_a|^2dq,\qquad
 \widetilde\phi_a=\frac{1_{\Gamma_a}\phi_a}{\sqrt{1-\delta_a}}.
\]
The cut packets define a \emph{comparison} isometry $\widetilde W$ with ideal disjoint records on the same retained space. They are not claimed to be physical Gaussian preparations. If a smooth comparison packet is wanted, smooth the cut in an arbitrarily narrow boundary strip and add its norm error.

\begin{lemma}[Complete retained-state record error]\label{mc:lem:instrument}
If $\delta_* =\max_a\delta_a<1$, then
\begin{equation}\label{mc:eq:isometryerror}
 \norm{W-\widetilde W}\le\sqrt{2\delta_*}.
\end{equation}
The same bound holds after tensoring any reference; it bounds the trace distance between the full pure outputs and hence the half-diamond distance of the resulting physical output channels. Any subsequent \emph{common} coherent return acting on all retained factors preserves the full-state bound.
\end{lemma}
\begin{proof}
The packet squared difference is $2(1-\sqrt{1-\delta_a})\le2\delta_a$. Orthogonality of the retained key gives
\[
 \norm{(W-\widetilde W)\psi}^2
 =\sum_a\norm{K_a\psi}^2\norm{\phi_a-\widetilde\phi_a}^2\le2\delta_*.
\]
For normalized vectors their pure-state trace distance is no greater than their norm difference. The proof is unchanged with an identity on $R$. Unitary invariance and channel contractivity prove the last claims. No continuity assertion for raw guidance paths is being used.
\end{proof}

For a source-only reduced instrument, tracing pointer and key would give weights and daughter mixtures. The comparison instead retains $K$, pointer packets, resources and $R$. The classical label channel is a representation of final physical regions: for a final wave $\Xi$, its unnormalized output is $P_{\Gamma_a}\ket\Xi\bra\Xi P_{\Gamma_a}$, optionally with a classical index. It is not a law that the global wave is physically projected at that time. When previously separated waves are to be recombined, use their complete coherent vector, not a dephased classical record representation.

\subsection{Noncommuting continuation without a fresh probability postulate}
After a first projective write $P_a$, a stored key and a copy, let $V_a$ be a source unitary and $R_b$ a noncommuting second projector family. With a fresh second pointer, the ideal complete vector is
\begin{equation}\label{mc:eq:jointwave}
 \sum_{a,b}(R_bV_aP_a\otimes I_R)\psi\,
         \ket a_K\ket b_B\phi_a(y)\chi_b(z)
         \ket{\mathrm{resources}(a,b)},
\end{equation}
with further coherent indices included if resource vectors are not single basis states. They are never discarded merely because the display says null. For exact key-controlled gates, the actual finite displayed probabilities are
\begin{equation}\label{mc:eq:noisymarks}
 \Prb(\widehat A=r,\widehat B=s)
 =\sum_{a,b}G^A_{r|a}G^B_{s|b}
                      \norm{(R_bV_aP_a\otimes I_R)\psi}^2,
\end{equation}
where $G^A_{r|a}=\int_{\Gamma_r}|\phi_a|^2$ and similarly for $B$. Since the first actual pointer is held, this is also the law of its \emph{earlier} declaration and the later declaration. Its classical history error from the ideal finite instrument is at most $\delta_A+\delta_B$. For literal position feedback add \eqref{mc:eq:feedbackbound} and keep its separate immutable archive. Conditional normalization costs the ordinary probability denominator; it is not an unqualified exact branch rule.

\begin{theorem}[Finite complete measurement-chain closure]\label{mc:closure}
Fix a finite programme of the stated resource gates, massive writes, copies, retained resets, bounded protected internal exposures, and coherent or smooth spatial feedback. Let its horizon be $T$ and let at most $m$ physical record registers be declared. Supply the complete equilibrium initial law and the finite ready stock. Then:
\begin{enumerate}
\item The autonomous model \eqref{mc:eq:auto} has a conservative complete actual path law given by \eqref{mc:eq:pathlaw}, including its clock and all returning systems.
\item For worst Gaussian classification tails $\delta_j$, total full-wave protection/gate error $\epsilon_{\rm gate}$, and the $L^2$ clock error $\varepsilon_0$, its complete final retained-state error against the ideal disjoint-record instrument is at most
\begin{equation}\label{mc:eq:totalwave}
 E_{\rm out}=\min\left\{1,\varepsilon_0+\epsilon_{\rm gate}
                                  +\sum_{j=1}^m\sqrt{2\delta_j}\right\}.
\end{equation}
The convention is trace distance for the complete retained quantum output (or half-diamond distance for its linear output channel), together with probabilities of physical records. It is not TV distance on the ontic pair $(\Psi,Q)$: different exact global waves need not be close in that much stronger sense.
\item During specified holds, let $E_{\rm arch}$ be the sum of the surface budgets \eqref{mc:eq:archiveerror}. If a declared working register is copied and then intentionally reset, also include a transfer budget $E_{\rm transfer}$: sum the pair-classification errors $\delta_{\rm old}+\delta_{\rm copy}$ and the complete endpoint comparison error at each such copy cut. The law of actual historical declarations and their retained final displays differs from the ideal record law by at most
\begin{equation}\label{mc:eq:totalhist}
 E_{\rm hist}\le\min\{1,E_{\rm out}+E_{\rm arch}+E_{\rm transfer}\}.
\end{equation}
For registers declared directly in their permanent archive and never transferred, $E_{\rm transfer}=0$. In the exact driven key-preserving library, $E_{\rm arch}=0$ and the sharper purely classical classification bound is $\sum_j\delta_j$ when every displayed occurrence is included and no other perturbation is present.
\item For every fixed finite ideal programme and tolerance $\epsilon>0$, finite pointer separations, protection gaps where used, feedback profiles and clock resources can be chosen so that these displayed bounds are below $\epsilon$. The exact physical theory remains \eqref{mc:eq:H}--\eqref{mc:eq:equilibrium}; only finite resources are adjusted.
\end{enumerate}
\end{theorem}
\begin{proof}
Conservative existence was established for the complete smooth domain in Lemma~\ref{mc:lem:exist}. Every stage is a finite unitary in that domain. Multiply the stages retaining every old key and resource, including reset receivers. First compare perturbed gates to the nominal driven key-controlled programme, using Proposition~\ref{mc:prop:gap} and feedback Duhamel on the nominal stage inputs. In particular, each protection-stage comparison is evaluated on its encoded ideal prefix in the promised subspace $P$; earlier leakage is carried by the common actual suffix, not assumed absent. Each suffix is a common unitary, so prefix discrepancies are preserved; the estimates are uniform over the unknown input and $R$ with the specified prepared material bank. Add the clock's full retained-wave error.

At the \emph{final comparison cut}, the nominal wave is an orthogonal history-key expansion with real Gaussian factors for each retained display and the complete source/resource coefficients. These displays include all replacement archive receivers; an intentionally reset original pointer is a ready factor, not a carrier of its former record. Truncate these display packets into their assigned cells only at this cut. On a branch, the norm-squared mass removed from its product of packets is at most $\sum_j\delta_j$. Orthogonality of the retained history keys then gives a full vector error at most $\sqrt{2\sum_j\delta_j}\le\sum_j\sqrt{2\delta_j}$ by the proof of Lemma~\ref{mc:lem:instrument}. This yields a disjoint-record comparator with precisely the ideal history coefficients, on the same full retained space. Pure-state trace distance and position readout contractivity prove \eqref{mc:eq:totalwave}. Cut packets are not propagated as if they were stationary oscillator ground states; the support truncation is a final comparison construction only. Subsequent common coherent returns preserve its full-state error but need not preserve its initial record separation.

For history, couple a path's earlier declared labels to its actual final archived labels on the same probability space. A held label can change only by a specified surface crossing. A transfer to a fresh copy can additionally mismatch at the copy cut: its joint endpoint probability is bounded by $\delta_{\rm old}+\delta_{\rm copy}$ in the nominal wave, plus the complete comparison error at that cut. Stop protecting the old register when its deliberate reset begins, and protect the receiving archive from then on. Theorem~\ref{mc:history} and a union bound therefore give mismatch probability at most $E_{\rm arch}+E_{\rm transfer}$. Comparing final records to the ideal law costs $E_{\rm out}$. This proves \eqref{mc:eq:totalhist} without a path-TV inference from wave closeness.

For feasibility first choose pointer separations to make Gaussian tails small. Any literal feedback derivative residual can simultaneously be made small by the Gaussian-tail estimate after Theorem~\ref{mc:history}, with a separate archive if used. Choose bounded internal protection gaps large enough for the finite sum of \eqref{mc:eq:gap}; no unbounded nuisance operators have been included. The finite apparatus inventory is then fixed. Its constants $A_{r,T},B,C_{\rm tr}$ are finite. Increase $M_c$ and choose $s_c=\sqrt{\hbar T/(2M_c)}$ so both clock-wave and archive-history bounds become arbitrarily small. All choices are finite at positive $\epsilon$. The limits are taken in this order, not uniformly over an unbounded growing graph or an infinite observation horizon.
\end{proof}

\subsection{Positive-probability conditioning and returning branches}
\label{mc:conditioning}

\begin{lemma}[Conditioning on a common retained event]
\label{mc:conditional-bound}
Let $P,Q$ be two probability laws on the same retained output or history
space, with $\TV(P,Q)\le\epsilon$. For a common event $E$, put
$p=P(E)$ and $q=Q(E)$. If \emph{both} $p>0$ and $q>0$, then
\[
 \TV\bigl(P(\,\cdot\mid E),Q(\,\cdot\mid E)\bigr)
 \le \min\{1,2\epsilon/p\}.
\]
The sufficient condition $\epsilon<p$ ensures $q>0$. For positive
unnormalized quantum outputs $\sigma,\tau$ on the same retained space,
if $\|\sigma-\tau\|_1\le d$, $p=\tr\sigma>0$ and $q=\tr\tau>0$, then
\[
 \|\sigma/p-\tau/q\|_1\le\min\{2,2d/p\}.
\]
Here $d<p$ is sufficient for positivity of the second trace.
\end{lemma}
\begin{proof}
For a measurable set $A$, add and subtract $Q(A\cap E)/p$.
The first difference is at most $\epsilon/p$, and the second is at most
$|p-q|/p\le\epsilon/p$, since $Q(A\cap E)\le q$.
Taking the supremum proves the classical estimate. For the quantum
estimate add and subtract $\tau/p$, use positivity to obtain
$\|\tau\|_1=q$, and use
$|p-q|\le\|\sigma-\tau\|_1$:
\[
 \|\sigma/p-\tau/q\|_1
 \le d/p+|p-q|/p\le2d/p.
\]
The upper bounds one and two are the maximal respective distances.
Finally $q\ge p-\epsilon$ or $q\ge p-d$ proves the positivity claims.
\end{proof}

A zero-probability event has no normalized conditional branch. Capping an
error estimate at one does not create that branch. Arbitrarily rare
events therefore have no uniform conditional guarantee.
Equation~\eqref{mc:eq:jointwave} and its resource version, rather than a
single sampled daughter, specify a complete return experiment.
The reference is never accessed; all source maps tensor $I_R$.
For a common subsequent unitary the full-state bound persists, but
record separation or historical readability need not persist after an
intentional echo. Those claims retain their separate holding and transfer
hypotheses.

\paragraph{Preparation information is retained.}
In the autonomous theory the clock's actual position is part of the initial configuration. Exposing it, or any other microscopic coordinate, changes the conditioning in \eqref{mc:eq:pathlaw}. A physically acquired coordinate record must be an extra material coupling and retain its receiver. A factorized ready packet at a fixed time does not prove independence conditional on every hypothetical unrecorded previous passage time. The theorem uses complete initial equilibrium and a finite independent stock; it does not invoke conditional nodal extraction as a nonexistent global preparation theorem.

\chapter{Noncommuting tests and rival actual motions}
\label{mc:chapter-tests}

An endpoint Born weight is too weak a test of the complete construction.
The following experiment keeps pending excitation, loss products, copied
records, a reset receiver and an inaccessible reference through a second,
noncommuting operation. The rival processes then show exactly what the
constitutive motion law selects beyond one-time equilibrium and reliable
macroscopic records.

\section{A full finite measurement and continuation example}
\label{mc:sec:example}
Use $P_a$ as $Z$ projectors and set
\[
 \theta=\pi/3,\qquad\varphi=\pi/4,\qquad\eta=2/3.
\]
The receptor weights are exactly
\begin{equation}\label{mc:eq:exampleweights}
 (P_r,P_p,P_c,P_l)=(1/4,3/8,1/4,1/8).
\end{equation}
Take one unknown source with an inaccessible two-dimensional reference:
\begin{equation}\label{mc:eq:inputexample}
 \psi=\sqrt{2/3}\ket0\ket0_R+
                       \sqrt{1/3}\ket1\ket+_R,\qquad
 \ket+_R=(\ket0_R+\ket1_R)/\sqrt2.
\end{equation}
The reduced source coherence is $\rho_{01}=1/3$. Neither independent copies nor reference control are used.

\subsection{Null followed by an incompatible measurement}
The ideal orthogonal-record null map is $\mathcal N(\rho)=\rho/4+\mathcal D_Z(\rho)/2$, with trace $3/4$; its retained coherent null is the vector \eqref{mc:eq:nullvector}. A later $X$ probe in this ideal comparator gives
\begin{equation}\label{mc:eq:nullnumbers}
 \Prb(N,X+)=\frac{11}{24},\qquad
 \Prb(X+\mid N)=\frac{11}{18},\qquad
 \sigma_R^{N,+}=\frac1{48}
               \begin{pmatrix}19&5\\5&3\end{pmatrix}.
\end{equation}
To verify, write $\bra{+_x}\psi=\sqrt{1/3}\ket0_R+\sqrt{1/6}\ket+_R$ and apply \eqref{mc:eq:nullmap} term by term, or expand \eqref{mc:eq:nullvector} and trace only the named $D$ factor. The trace of the displayed matrix is $11/24$; its normalized reference state is the matrix with entries $19,5,5,3$ divided by $22$. For $X-$ the unnormalized reference state is
\[
 \sigma_R^{N,-}=\frac1{48}\begin{pmatrix}11&1\\1&3\end{pmatrix},
 \quad\Prb(N,X-)=7/24.
\]
Their sum is $(3/4)\rho_R$, an explicit reference consistency check. A frozen original input would instead give $\Prb(X+\mid N)=5/6$; a fully $Z$-dephased daughter would give $1/2$. Both fail this finite experiment.

\subsection{Captured copy, reset, spatial feedback and a second record}
Copy the captured label, retain its spatial archive, reset $D$ by \eqref{mc:eq:covariantreset}, keeping $D'$, and use
\[
 V_0=I,\qquad V_1=e^{-i\pi\sigma_y/6}.
\]
Apply an $X$ write to a fresh pointer. The ideal branch coefficients, with all resource factors attached, are
\[
 \tfrac12(R_bV_aP_a\otimes I_R)\psi.
\]
Their probabilities are
\begin{center}
\begin{tabular}{@{}lcc@{}}\toprule
Retained first record & Second $+$ & Second $-$\\\midrule
$a=0$ & $1/12$ & $1/12$\\
$a=1$ & $(2-\sqrt3)/48$ & $(2+\sqrt3)/48$\\\bottomrule
\end{tabular}
\end{center}
They sum to capture probability $1/4$. Reference daughters are $\ket0_R$ and $\ket+_R$, respectively. Literal spatial feedback uses the smooth $g(y)B$ contact, with $B=(\pi\hbar/(6t_f))\sigma_y$, and the bound \eqref{mc:eq:feedbackbound}. The copied archive remains held while the working pointer can recoil. Thus the exact table is the ideal target with explicit finite classifier, feedback and clock errors, not an assertion that finite Gaussian records are orthogonal.

\subsection{A complete reversal remains a different experiment}
In the driven bank, if every response, source gate, copy and spatial write is coherently undone with all receiving systems, that full bank wave returns to its input, up to a known common phase. Resetting only $D$ while a copy or $D'$ survives does not achieve that return. A later incompatible probe distinguishes the two retained states. No global projection has been inserted at a declaration. In the finite autonomous realization the controller also remains in the complete state: its recoil and entanglement are bounded by the clock comparison, not claimed to be exactly undone by these apparatus inverse pulses.

An inverse need not negate a massive kinetic energy. For the piecewise constant half-period alternative to \eqref{mc:eq:quintic}, a trap centred at $aL/2$ sends a ground packet centred at zero to one centred at $aL$ in time $\pi/\omega$. Each conditional oscillator has common equally spaced spectrum, so evolution for $2\pi/\omega$ is $-I$. Evolving for a complementary positive duration realizes the inverse, up to a common phase. Finite internal gate inverses reverse a bounded matrix term. Smooth forced displacements can likewise be undone on the specified coherent packets by a reversed centre trajectory and known phase correction. A claim of inversion on an arbitrary oscillator state would require its full propagator rather than this restricted packet identity.

\subsection{Independent verification}
Exact symbolic calculations checked the full $S\otimes R\otimes D$ vector, resource unitarity, all four status weights, both null reference matrices, the four feedback probabilities, the complete SWAP export and the noncommuting transported-trap identity. All 32 exact checks passed. These finite calculations are checks of the displayed proofs, not substitutes for them.

For the exact quintic writer with $\hbar=\sigma=T_w=1$, $\omega T_w=\pi$, $M=1/(2\pi)$, $L=8$, and $p_1=0.65$, direct quadrature and inverse-CDF calculations give
\begin{center}
\begin{tabular}{@{}lr@{}}\toprule
Quantity & Value\\\midrule
Initial right-tail atom & $0.0000316712418$\\
Later threshold crossing & $0.649958827386$\\
Final no-crossing null & $0.350009501373$\\
Sum & $1.000000000000$\\\bottomrule
\end{tabular}
\end{center}
The symbolic Schr\"odinger residual is exactly zero. An independent finite-difference evaluation had relative residual below $2.0\times10^{-7}$; the quantile ODE residual was below $1.8\times10^{-9}$ and the first-passage quadrature discrepancy below $3.4\times10^{-16}$. Figure~\ref{mc:fig:mechanism} uses these parameters. No simulation evidence is used to assert a universal event-law selection.

\section{Rival actual motions and the selection obligations}
\label{mc:sec:rivals}
The strongest candidate must survive a rival that preserves more than endpoint Born weights. We give such a rival and state precisely what it defeats.

\subsection{Same local net current, mutually singular paths}
For a smooth positive density of the same complete wave, define an equivariant diffusion by
\begin{equation}\label{mc:eq:diffusion}
 dQ_t=\left(\frac j\rho+D\nabla\log\rho\right)(Q_t,t)dt
                                   +\sqrt{2D}\,dW_t,
 \qquad D>0.
\end{equation}
Its Brownian innovations are an explicit additional stochastic premise. The Fokker--Planck current is $\rho b-D\nabla\rho=j$, so it matches even the local current, not just its divergence. We only assert its global existence where checked below; this is an equivariant diffusion rival, not a claim that every axiom of Nelson's stochastic mechanics has been derived.

For a stationary harmonic ground-state pointer centred at zero, the adopted theory has $\dot Q=0$. The rival is the globally well-posed Ornstein--Uhlenbeck process
\[
 dQ=-D Q\,dt/\sigma^2+\sqrt{2D}\,dW.
\]
It has the identical invariant Gaussian density but moves at positive times. More strongly, on any $T>0$ its path law and the adopted guidance path law have TV distance one: Brownian diffusion paths have quadratic variation $2DT$, whereas the absolutely continuous guidance paths have zero. These are disjoint measurable path events. The comparison does not require an experimentally admitted passive quadratic-variation meter.

\begin{proposition}[Reliable macroscopic records do not remove the rival]\label{mc:prop:rivalarchive}
Use the same semibounded Hamiltonian
\[
 H=\frac{p^2}{2M}+\frac{M\omega^2}{2}(x-a\sigma_z)^2
\]
and the equal superposition of its two spin-labelled ground packets. Its stationary density is
\[
 \rho(x)=\tfrac12g_\sigma(x-a)+\tfrac12g_\sigma(x+a),\qquad j=0.
\]
The adopted configuration is fixed. The diffusion \eqref{mc:eq:diffusion} has smooth globally Lipschitz drift
\[
 b_D(x)=\frac D{\sigma^2}\left[-x+a\tanh\left(\frac{ax}{\sigma^2}\right)\right]
\]
and invariant law $\rho$. For $a\ge2\sigma$, the probability it changes the sign record during $[0,T]$ is no greater than
\begin{equation}\label{mc:eq:rivalbound}
 \min\left\{1,\frac{a+4DT/a}{\sqrt{2\pi}\sigma}
                         e^{-a^2/(8\sigma^2)}\right\}.
\end{equation}
\end{proposition}
\begin{proof}
Set $b=a/2$ and $f(x)=(1-|x|/b)_+$. On $(0,b)$, $\rho'\ge0$: the sign is that of $a\tanh(ax/\sigma^2)-x$, a concave function vanishing at zero and positive at $b$ for $a\ge2\sigma$. Thus $f'b_D\le0$ on both sides of the central interval. Stop at the first hit $\tau$ of zero. The It\^o--Tanaka formula gives only nonpositive interior drift and a nonpositive local-time contribution at zero; its positive contributions are $(L_{T\wedge\tau}^{-b}+L_{T\wedge\tau}^{b})/(2b)$. Since $f(Q_{T\wedge\tau})\ge1_{\{\tau\le T\}}$ and stationarity gives $\E L_T^x=2DT\rho(x)$,
\[
 \Prb(\tau\le T)\le\E f(Q_0)+\frac{DT}{b}[\rho(-b)+\rho(b)]
 \le(2b+2DT/b)\rho(b).
\]
Finally $\rho(b)\le e^{-a^2/(8\sigma^2)}/(\sqrt{2\pi}\sigma)$. Sign change requires a hit of zero, so the same bound applies. Existence and invariance follow directly from the displayed Lipschitz drift and the stationary Fokker--Planck equation.
\end{proof}
Thus arbitrarily reliable finite-horizon records can coexist with mutually singular microscopic paths. The velocity postulate selects the adopted law \emph{within the new theory}; record success does not independently force that postulate. Deterministic divergence-free changes provide further rivals: in an isotropic real two-dimensional Gaussian, $v_\Omega=\Omega(-y,x)$ preserves the same density while changing a sign record with probability $\Omega T/\pi$ for $0\le\Omega T\le\pi$. This particular rotor is a counterexample to inference from continuity, not a proposed fully symmetry-constrained replacement. More general quantum-equivalent deterministic alternatives are established in primary work \cite{mc:Deotto}.

\subsection{A finite discrete rival with a sharp complete-path discriminator}
On the book's finite graph, let
\begin{equation}\label{mc:eq:etarates}
 \lambda^{(\eta)}_{Y\leftarrow X}=
 \frac{\pos{J_{YX}}+\eta|J_{YX}|}{w_X},\qquad\eta\ge0.
\end{equation}
It preserves individual net currents, support and equilibrium and remains Markov. On a binary monotone write with $u=\sin^2(gt)$, the rates per $du$ are $(1+\eta)/(1-u)$ forward and $\eta/u$ backward. The wave weights are $1-u,u$, its final state is 1, and expected total jump count is $1+2\eta$; integrability of this count and vanishing nodal holding survival give a nonexplosive path law.

A path with exactly one jump at $u$ has density
\[
 (1-u)^{1+\eta}\frac{1+\eta}{1-u}
       \exp\left[-\int_u^1\frac\eta s ds\right]
 =(1+\eta)[u(1-u)]^\eta.
\]
The Bell law has exactly one jump, uniform in $u$. Since $(1+\eta)4^{-\eta}\le1$, the common mass of the two path measures is precisely the integral of this rival one-jump density. Therefore
\begin{equation}\label{mc:eq:exactTV}
 \TV(P_\eta,P_{\rm B})=1-(1+\eta)\mathrm B(1+\eta,1+\eta),
 \qquad\TV(P_1,P_{\rm B})=\frac23.
\end{equation}
This is a concrete surviving balanced-traffic countermodel with exact Born endpoints. It is not an independently selected new event mechanism. It confirms why endpoint agreement, even with correct individual net currents, would be insufficient to claim Bell closure.

\subsection{Obligation-by-obligation verdict in the adopted continuum}
\begin{enumerate}
\item \textbf{Individual current realization.} The kinetic-momentum postulate specifies the pointwise material current $j$. Equivariance and the flow derive expected net flux through each physical interface. There is no freely reassigned cycle current within that postulate. This does \emph{not} identify those interfaces with arbitrary finite internal Hamiltonian matrix edges in \eqref{found:bell}.
\item \textbf{Surplus traffic.} At a regular point of an interface the velocity has one sign and actual crossings realize its local direction. If an entire interface is integrated or microscopic coordinates are omitted, opposite directions on different patches give
\[
 F_+=\int_\Sigma\pos{j\cdot n},\quad F_-=\int_\Sigma\pos{-j\cdot n},\quad
 F_+-F_-=\int_\Sigma j\cdot n.
\]
In general $F_+\ne[\int_\Sigma j\cdot n]_+$. Coarse countertraffic and recrossing can survive. Neither their absence nor graph Bell minimality is falsely claimed. No extra Brownian traffic is allowed by the stated mechanical law.
\item \textbf{Conditional timing.} Equations~\eqref{mc:eq:guidance} and \eqref{mc:eq:pathlaw} give the whole path and every conditional history, including current reversals and nulls. The explicit physical-time first-passage example has a derived nonexponential law. Other waiting laws are different constitutions, not unresolved free choices inside this one. Coarse histories generally fail Markov closure and are not assigned Bell rates by projection.
\end{enumerate}
The alternative therefore closes the operational programme with a selected physical motion law and its material contacts. It gives up the original finite-sector microscopic claim rather than pretending to derive it.

\section{Dependencies and the surviving distinction}
\label{mc:dependencies}

The shared source/readout representation and finite coherent gate algebra
remain useful across constitutions. They become predictions only after a
complete motion law and an admitted initial ensemble are specified.
Within this part the dependency chain is
\[
\begin{gathered}
 \text{massive Schr\"odinger inventory, kinetic-momentum motion}\\
 +\ \text{complete initial equilibrium and finite ready resources}\\
 \Longrightarrow\ \text{conservative complete paths and physical writes}\\
 \Longrightarrow\ \text{retained output and faithful sampled histories}
 \quad\text{with the bounds of Theorem~\ref{mc:closure}.}
\end{gathered}
\]
The exact driven storage theorem and the autonomous surface-current
estimate replace an unsupported inference from final equilibrium to
archive truth. The copy-cut transfer term is needed whenever an old
declaration is moved to a receiver before its original pointer is reset.

\begin{center}
\begin{tabularx}{\textwidth}{@{}>{\raggedright\arraybackslash}p{.29\textwidth}X@{}}
\toprule
Ingredient & Status in the massive constitution\\
\midrule
Wave-current identities and coherent finite gates &
Retained, with every source, material register and reference in the
specified wave. Internal edge currents are not actual spin-jump currents.\\
Continuous configuration records &
The guidance, quantile and endpoint logic of Chapter~\ref{cfg:chapter}
is retained; positive kinetic energy, confining storage and a retained
autonomous controller supply this part's stronger material realization.\\
Preparation &
Complete initial equilibrium and a finite independent ready stock are
supplied. A conditional subsystem preparation does not imply universal
equilibrium or freshness after an exposed microscopic history.\\
Protection &
The retained-bank gap estimate applies to bounded internal perturbations
and commuting spatial spectators; arbitrary unbounded disturbances remain
outside its hypotheses.\\
Copies, nulls, loss and reset &
All coherent branches, original correlations and receiving systems are
retained. A local ready appearance is not global erasure.\\
Final outputs and actual history &
Complete retained-state trace bounds and held/transfer surface budgets
are different estimates. Neither is promoted to arbitrary raw-path TV.\\
Finite Bell or pilot constitution &
A different actual ontology and event mechanism. Shared gate matrices
do not make its microscopic path the path derived in this part.\\
\bottomrule
\end{tabularx}
\end{center}

Bohmian motion, equilibrium-conditioned measurement and the general
existence mechanism have established provenance. The construction here
uses them in a common finite semibounded resource model, with the exact
forced writer, transported-trap reset and the controller
graph-norm-to-archive-current estimate proved above. It establishes
internal closure under the stated physical axioms and controlled
finite-horizon implementation of ideal instruments. It does not establish
the universal necessity of those axioms, preparation from arbitrary initial
data, a relativistic completion, unlimited memory, or a perfect autonomous
clock for all time.

The diffusion rival demonstrates that even very reliable macroscopic
archives do not single out this microscopic motion among all equivariant
alternatives. It does not introduce an unspecified freedom inside the
adopted motion law: once the initial configuration is fixed, every event
and continuation in this constitution is fixed. This is the distinction
between a demonstrated constitutive completion and a derivation of that
constitution from the older interface alone.


\part{Physical Preparation of Apparatus Statistics}\label{part:preparation}

% Begin consolidated source: parts/08_preparation.tex
\chapter{Conditional preparation by reversible source transport}
\label{prep:deterministic}

The required apparatus law is a conditional resource statement. A pointer
with a correct marginal distribution can remain correlated with a controller
that predicts its future detector output. This chapter proves the finite
nodal extraction result of \cite{M19}, retaining its physical archive, and
separates it from global equilibration. The next chapter gives two different
statistical alternatives and the complete detector integration.

\section{The resource to be prepared}

At handoff, let $U$ be the ready configuration and let $A$ contain every
old preparation record, controller coordinate, exported cell label,
and future-active memory. Internal quantum memories and an inaccessible
reference remain in the wave. For a known normalized ready packet $\phi$,
the target is
\begin{equation}
 Q_{U,A}(du,da)=|\phi(u)|^2du\,P_A(da).
 \label{prep:target}
\end{equation}
The archive has its \emph{actual} marginal $P_A$, which may be very far
from the wave norm distribution. Equation~\eqref{prep:target} is therefore
a conditional ready-subsystem target, not a global equilibrium measure.

All comparisons use
$\operatorname{TV}(P,Q)=\sup_E|P(E)-Q(E)|=\frac12\|P-Q\|_1$.
For standard Borel conditional laws with the same archive marginal,
\begin{equation}
 \operatorname{TV}(P_{U,A},Q_{U,A})
 =\int\operatorname{TV}(P_{U|a},|\phi|^2du)\,P_A(da).
 \label{prep:averageconditional}
\end{equation}
This equality is an average conditional guarantee. It supplies no uniform
statement on arbitrary rare archive values.

\begin{proposition}[Reversible fine-grained obstruction]
\label{prep:reversible}
Let $S_t$ be a common invertible measurable guidance flow with measurable
inverse, and let its wave reference law be equivariant:
$Q_t=(S_t)_*Q_0$. For any actual law $P_t=(S_t)_*P_0$,
\[
 \operatorname{TV}(P_t,Q_t)=\operatorname{TV}(P_0,Q_0).
\]
If $P_0=fQ_0$, then
\[
 \frac{dP_t}{dQ_t}=f\circ S_t^{-1},\qquad
 \int\alpha\!\left(\frac{dP_t}{dQ_t}\right)dQ_t
      =\int\alpha(f)dQ_0
\]
for every defined integral on either side. Singular components are preserved.
\end{proposition}
\begin{proof}
For every measurable event $E$,
$P_t(E)=P_0(S_t^{-1}E)$ and likewise for $Q$.
The measurable bijection carries the full collection of events onto itself,
so taking suprema proves the total-variation identity without any
absolute-continuity assumption. In the density case, change variables in
$\int_{S_t^{-1}E}f\,dQ_0$ to obtain the displayed Radon--Nikodym derivative.
A further change of variables proves the integral identity. Null sets and
their inverse images preserve singularity.
\end{proof}

Coarse-grained relaxation, including established pilot-wave relaxation
studies such as \cite{ValentiniWestman}, does not contradict this statement.
Fine-grained information may move to unresolved scales or exported
coordinates while a restricted class of observables becomes insensitive.

\begin{lemma}[Same-wave complete future comparison]\label{prep:push}
Two preparations with the same complete wave and control programme,
and configuration-law distance at most $\delta$, have final
complete-output distance at most $\delta$ under any common measurable
source evolution and record map. The same holds for a common stochastic
kernel. Returning waves, copies, nulls, physical timestamps, and
adaptive records may be included.
\end{lemma}
\begin{proof}
A deterministic output event is a preimage under the common map.
For a kernel, its probability for an output event is a measurable function
in $[0,1]$; integrating that function against $P-Q$ has absolute value at
most $\operatorname{TV}(P,Q)$. The complete actual path itself can be the
output whenever the evolution assigns a measurable path.
\end{proof}

This lemma compares to the same-wave hybrid law
\eqref{prep:target}. It does not establish that this hybrid law predicts
Born records after arbitrary activation of a nonequilibrium archive.

\begin{lemma}[Normalization cost]\label{prep:condition}
If $\operatorname{TV}(P,Q)\le\delta$, $p=P(E)$, $q=Q(E)$ and $p,q>0$, then
\[
 \operatorname{TV}(P(\cdot|E),Q(\cdot|E))
       \le\min\{1,\delta/\max(p,q)\}.
\]
\end{lemma}
\begin{proof}
Assume $p\ge q$. The measures have common mass at least $1-\delta$.
Their common mass outside $E$ is at most $1-p$, so the common mass
inside $E$ is at least $p-\delta$. Dividing both measures by their
own probabilities leaves common conditional mass at least
$(p-\delta)/p$, since $q\le p$. Subtraction from one proves the bound.
Interchange $P$ and $Q$ for the other ordering.
\end{proof}

\section{Explicit nodal resource and admitted initial laws}

Fix $N\ge2$ and $0<\eta<1/2$. Define on $(0,1)$
\[
 a_\eta(u)=
 \begin{cases}
 \sin(\pi u/(2\eta)),&0<u<\eta,\\
 1,&\eta\le u\le1-\eta,\\
 \sin(\pi(1-u)/(2\eta)),&1-\eta<u<1,
 \end{cases}
 \qquad \phi_\eta(u)=\frac{a_\eta(u)}{\sqrt{1-\eta}},
\]
extended by zero. It is normalized, symmetric about $1/2$, and belongs to
$H^1(\mathbb R)$. Direct integration gives
\begin{equation}
 \|\phi_\eta'\|_2^2=\frac{\pi^2}{4\eta(1-\eta)},\qquad
 d_\eta:=\operatorname{TV}(|\phi_\eta|^2du,du)\le\eta.
 \label{prep:packet}
\end{equation}
For the second inequality, $\min(|\phi_\eta|^2,1)\ge a_\eta^2$,
whose integral is $1-\eta$.

The known seed wave is an array of $N$ identical cells:
\begin{equation}
 \varphi_{N,\eta}(x)=\phi_\eta(Nx-k),\qquad
 x\in I_k=(k/N,(k+1)/N).
 \label{prep:seed}
\end{equation}
It has norm one, exact nodes at cell boundaries, and
$\|\varphi_{N,\eta}'\|_2^2=N^2\|\phi_\eta'\|_2^2$.
Supplying this coherent wave is a preparation resource. The construction
does not claim to shape an arbitrary unknown wave into it without cost.

Add an internal tag with ready state $|r\rangle$ and $N$ labels
$|k\rangle$. Add an archive coordinate $y$ with known $H^1$ wave
$\chi$, supported on $[0,1]$ and positive in its interior.
Its actual configuration distribution need not be equilibrium.
Let $z$ collect its actual initial value, every old coordinate and
classical history, with a standard Borel reference measure $\mu$.
The complete initial wave is
$\varphi_{N,\eta}(x)\chi(y)|r\rangle\Xi(z_{\rm old})$,
with the unknown input and inaccessible reference inside $\Xi$.
Assume the actual law has density
\begin{equation}
 P_0(dx,dz)=f(x,z)\,dx\,\mu(dz),\quad f\ge0,\quad \int f=1,
 \qquad
 M:=\int\operatorname{Var}_x f(\cdot,z)\,\mu(dz)<\infty.
 \label{prep:BV}
\end{equation}
All actual values lie on the nonzero-wave supports needed for the stated
flow. No target wave density appears in \eqref{prep:BV}.
Multimodal distributions, steps, and correlations are allowed. An exact
old copy $Z=X$ is excluded because its conditional law is singular.
Regularity of the unconditional $X$ marginal would not suffice.

\section{The source interaction and exact actual routing}

First tag the occupied cell coherently, for time $\tau_1$:
\[
 H_{\rm tag}(x)=\frac{\pi\hbar}{2\tau_1}
       \sum_k1_{I_k}(x)(|k\rangle\langle r|+|r\rangle\langle k|).
\]
It is a bounded multiplication operator, generates no configuration current,
and sends $|r\rangle$ to $-i|k\rangle$ on cell $k$.
The sharp coefficients are applied at exact nodes, so the wave remains
$H^1$. Next use
\[
 H_y=\dot b(t)\sum_k dk\,|k\rangle\langle k|\otimes P_y,
 \qquad d>1,\quad b:0\longrightarrow1.
\]
Only the tag $k$ is present at actual $x\in I_k$, and the actual archive
position becomes $Y'=Y+dk$. The archive supports are now disjoint.

For time $\tau_3$, apply the conditional dilation
\[
 H_D=\frac12\sum_k|k\rangle\langle k|\{v_k(x),P_x\},
 \qquad v_k(x)=a(x-k/N),\qquad a=\frac{\log N}{\tau_3}.
\]
The scalar transport law
\begin{equation}
 H_v=-i\hbar(v\partial_x+v'/2),\quad J=v|\Phi|^2,\quad
 \Phi_t(x)=\sqrt{(S_t^{-1})'(x)}\,\Phi_0(S_t^{-1}x)
 \label{prep:halfdensity}
\end{equation}
follows by differentiating along characteristics. It multiplies
$x-k/N$ by $N$. A final conditional translation by $-k/N$ aligns the
dilated cells. Thus the exact actual configuration map is
\begin{equation}
 K=\lfloor NX\rfloor,\qquad U=NX-K,\qquad Y'=Y+dK.
 \label{prep:map}
\end{equation}
Although the ready coordinate packets overlap after alignment, the disjoint
$Y'$ supports identify the local tag component and hence its actual velocity.
There is no discontinuous cutting of a connected nonzero-wave flow.

The complete final field is, up to a common phase,
\begin{equation}
 \Phi_f(u,y,z_{\rm old})
   =\phi_\eta(u)\left[
        \frac1{\sqrt N}\sum_k\chi(y-dk)|k\rangle\right]
         \Xi(z_{\rm old}).
 \label{prep:finalwave}
\end{equation}
The factor $N^{-1/2}$ is the dilation Jacobian. All unoccupied archive
packets survive. On its support the retained archive
$A=(Y',z_{\rm old})$ determines $K$ and the old $Y$; this is why correlations
with $K$ cannot be discarded.

\begin{theorem}[Conditional nodal extraction]\label{prep:extractor}
The finite interaction above, applied to the initial class
\eqref{prep:BV}, succeeds with probability one and obeys
\begin{align}
 \operatorname{TV}(P_{U,A},du\,P_A)&\le\frac{M}{4N},\\
 \operatorname{TV}(P_{U,A},|\phi_\eta(u)|^2du\,P_A)
          &\le\frac{M}{4N}+d_\eta
           \le\frac{M}{4N}+\eta.
 \label{prep:extractbound}
\end{align}
The comparison retains the actual archive marginal and holds uniformly
for unknown carried inputs and inaccessible references on which the
controls act as identity.
\end{theorem}
\begin{proof}
The interaction calculation proves \eqref{prep:map} and
\eqref{prep:finalwave}. Retaining $A$ is equivalent to retaining $(K,z)$.
Their joint density after routing is
\[
 g(u,k,z)=N^{-1}f((k+u)/N,z),\qquad
 m(k,z)=\int_0^1g(u,k,z)\,du.
\]
For a bounded-variation function $h$ on $(0,1)$ with mean $\bar h$,
Jensen's inequality and its variation measure give
\begin{align*}
 \int_0^1|h(u)-\bar h|\,du
 &\le\int_0^1\int_0^1|h(u)-h(v)|\,du\,dv\\
 &\le\int_{(0,1)}2t(1-t)\,|Dh|(dt)
 \le\tfrac12\operatorname{Var}(h).
\end{align*}
The middle inequality follows by integrating the variation along intervals
between $u$ and $v$; for fixed $t$ the ordered pairs whose interval crosses
$t$ have measure $2t(1-t)$.
Apply it to $h(u)=f((k+u)/N,z)$, integrate over $z$, and sum over
open cells. Their interior variation sums to at most the total variation.
The Jacobian $1/N$ and the factor one half in total variation yield
$M/(4N)$. Replacing uniform density by $|\phi_\eta|^2$ costs $d_\eta$
with unchanged $P_A$, proving the second inequality.
\end{proof}

\section{Resources, reproducibility, and sharp failure tests}

The resource costs are $N+1$ tag states, archive extent $O(dN)$,
cell resolution $1/N$, edge resolution $\eta/N$, and integrated
dilation strain $\log N$. The exact seed gradient cost is
\begin{equation}
 \|\partial_x\varphi_{N,\eta}\|_2^2
       =\frac{N^2\pi^2}{4\eta(1-\eta)}.
 \label{prep:seedcost}
\end{equation}
Translations and dilations are unbounded selfadjoint transport generators
on their admitted domains. These statements give finite resources on the
specified support and finite-gradient fields, not a bounded operator norm
or a lower-bounded microscopic energy model. Exact nodes and the sharp
tag coefficient are ideal spatial controls. A finite-bandwidth smooth
replacement needs its own configuration-law estimate; a small wave norm
error alone does not establish that estimate for nonequilibrium inputs.

A deterministic clock can drive these pulses without an equilibrium seed.
Take a coordinate $S$ with generator $P_S$ and a known compact wave
entirely upstream of ordered pulse regions. Let
$H=P_S+\sum_rw_r(S)G_r$, with disjoint regions and fixed operators $G_r$.
Then $\dot S=1$, and each admitted initial clock position crosses the
same complete pulse integrals by a common finite deadline. The ordered
unitary is independent of that initial position. The final clock factors
and belongs to $A$. An incomplete clock traversal is a genuine failure
branch, not the ready output.

For $m$ seed coordinates with product known waves, define integrated
coordinate variations $M_i$ of their joint actual density, conditioning
on all other coordinates and old archives. Apply the preceding interaction
separately to each coordinate. Then
\begin{equation}
 \operatorname{TV}\left(P_{\mathbf U,A},
       \prod_{i=1}^m|\phi_{\eta_i}(u_i)|^2d\mathbf u\,P_A\right)
 \le\sum_{i=1}^m\left(\frac{M_i}{4N_i}+\eta_i\right).
 \label{prep:stock}
\end{equation}
To prove this, successively average the density over each rescaled
$U_i$. Each averaging is an $L^1$ contraction and does not increase the
integrated variation in another coordinate. Telescope the one-coordinate
inequality, then change each uniform factor to its ready density.
The result concerns a joint finite stock, not separate correct marginals.

For the required deformation $f(x)=1+\epsilon(2x-1)$,
$|\epsilon|\le1$, the exact routed marginal and complete archive comparison
are
\begin{equation}
 f_U(u)=1+\frac{\epsilon}{N}(2u-1),\qquad
 \operatorname{TV}(P_{U,K},du\,P_K)=\frac{|\epsilon|}{4N}.
 \label{prep:linear}
\end{equation}
Indeed $g(u,k)=N^{-1}[1+\epsilon(2(k+u)/N-1)]$,
and subtracting its $u$ average leaves
$\epsilon(2u-1)/N^2$ in each cell. Integration proves the TV identity.
For a detector prepared in $\phi_\eta$, let
$G_\eta(u)=\int_0^u|\phi_\eta|^2$ and
$u_p=G_\eta^{-1}(1-p)$. Its terminal plus probability is
\[
 1-u_p+\frac{\epsilon}{N}u_p(1-u_p).
\]
At $p=1/2$, symmetry gives exactly $1/2+\epsilon/(4N)$;
in general its Born error is at most $d_\eta+|\epsilon|/(4N)$.

The admissible oscillatory rival
$f_N(x)=1+\epsilon\sin(2\pi Nx)$ has
$M=4|\epsilon|N$. Its extracted marginal is
$1+\epsilon\sin(2\pi u)$, so the balanced terminal error remains
$-\epsilon/\pi$ for every $N$. This respects \eqref{prep:extractbound}
and disproves a uniform statement over arbitrary absolutely continuous
initial laws. Wave smoothness alone does not constrain the independent
actual density variation.

For even $N$, the retained archive also satisfies
\[
 \mathbb P(K\ge N/2)=\frac12+\frac{\epsilon}{4}.
\]
Reversing all the preparation interactions restores the original wave and
actual law exactly. A balanced reader on the returned seed therefore
recovers the unsuppressed deviation $\epsilon/4$.
The resource was prepared by exporting nonequilibrium, not by destroying
it. This return is an explicitly allowed inverse, and it identifies the
precise limit of any global equilibrium interpretation.

\chapter{Statistical alternatives and the preparation-to-record chain}
\label{prep:alternatives}

This chapter consolidates the intrinsic clock preparation of \cite{M19},
the causal quantile selection of \cite{M20}, and the SWAP preparation
calculation of \cite{C01,C02}. They have different premises. One adds
new Poisson source physics; one selects a distribution by a causal
statistical requirement; one consumes an already definite resource under
an admitted Bell law. None is a renamed proof of reversible global
equilibration.

\section{A homogeneous switching clock with exact conditional heralding}

Let a rotor have circumference $\ell$ and known flat wave
$\phi_0=\ell^{-1/2}$. Its actual initial position $Q_0$, a drive sign
$\sigma_0\in\{-1,+1\}$, and every old memory $M$ may have an arbitrary
joint law, including atoms and exact copies. The unknown input is disconnected
during preparation. Between events the wave generator is
$H_\sigma=c\sigma P_Q$, which leaves the flat wave invariant, and guidance
gives $\dot Q=c\sigma$.

The explicitly new primitive is homogeneous Poisson sign switching,
independent of the complete initial source and actual past. A physical
count register $C$ is updated at each flip. On functions of actual variables,
\begin{equation}
 \mathcal L F(q,\sigma,C,m)
 =c\sigma\partial_qF
  +\kappa\{F(q,-\sigma,C+1,m)-F(q,\sigma,C,m)\}.
 \label{prep:rotorgenerator}
\end{equation}
The wave is unchanged at the flip; the complete source is a specified
classical-controller/quantum-wave extension. It is not obtained by measuring
an auxiliary system with Born probabilities.

A block clock ends each interval at
\[
 T=\frac{\ell}{2c},\qquad \kappa T=1.
\]
The physical counter retains block counts, final sign, block number and
success/failure flags. It has no within-block flip-time register.
Accept the first completed block containing exactly one flip, with at most
$n$ blocks. This is a deterministic decision on a primitive physical count,
not a calibrated quantum measurement.

\begin{lemma}[Uniform displacement from a singleton block]
\label{prep:singleton}
Conditional on one flip in a complete block, its time $S$ is uniform
on $(0,T)$ and independent of all incoming variables. The displacement
modulo $\ell$ is uniform, independently of the incoming position and sign.
\end{lemma}
\begin{proof}
The unnormalized density for a flip at $s$ and no other flip in the block
is $\kappa e^{-\kappa T}\,ds$. Its normalization is $ds/T$.
Integrating the two constant velocities gives
\[
 Q_T=Q_0+c\sigma_0(2S-T)\pmod\ell.
\]
As $S$ traverses $(0,T)$, $2cS$ traverses one full circumference.
Translation and reflection preserve the uniform measure on the circle.
\end{proof}

\begin{theorem}[Count-heralded conditional equilibrium]
\label{prep:herald}
Retain the old bank $M$, the initial and final signs, all completed
block counts, the stopping index, and the success/failure flag.
For every accepted count history,
\[
 \mathbb P(Q_{\rm ready}\in dq\mid
         M,\sigma_0,C_1,\ldots,C_K,\sigma_K,K,\mathrm{success})
       =\frac{dq}{\ell}.
\]
The success probability is
\begin{equation}
 s_n=1-(1-e^{-1})^n.
 \label{prep:success}
\end{equation}
On failure the flat wave, actual translated position, count history and
failure flag remain; no ready configuration is substituted.
\end{theorem}
\begin{proof}
The block counts are independent Poisson$(1)$ variables. Conditional on
a complete count string, point times in distinct blocks remain independent.
The accepted block convolves its arbitrary incoming conditional position
law with a uniform displacement by Lemma~\ref{prep:singleton}.
Its starting sign is determined by the initial sign and previous count
parities. Its final sign depends on count parity, not on the new uniform
flip time. Conditioning on all listed variables therefore leaves the
ready position uniform. Rejection of earlier blocks conditions only on
their counts; it does not select a subinterval of the accepted flip time.
Each block has singleton probability $e^{-1}$, which gives
\eqref{prep:success}.
\end{proof}

There is also an unheralded theorem. Run all $n$ blocks and retain all
counts and signs. If at least one count is one, that block supplies a
uniform displacement and later independent translations preserve it.
Only count strings with no singleton can be biased. Hence
\[
 \operatorname{TV}\left(P_{Q,A},\frac{dq}{\ell}P_A\right)
       \le(1-e^{-1})^n.
\]
This conclusion retains the count archive rather than averaging it away.

The expected numbers of executed blocks and flips before the capped stopping
rule are both $[1-(1-e^{-1})^n]/e^{-1}\le e$.
For blocks this is a finite geometric sum. For flips, the indicator that a
block is executed depends only on earlier counts, independently of its
mean-one count. Thus expected preparation time is at most $e\ell/(2c)$.
Maximum time and path length are $n\ell/(2c)$ and $n\ell/2$.

A finite counter with capacity $R$ stops and rejects on the $(R+1)$st
attempted flip. Coupling to an uncapped $n$-block process gives
\[
 \mathbb P(\mathrm{success})\ge
       1-(1-e^{-1})^n-b_{n,R},\qquad
 b_{n,R}=\mathbb P(\operatorname{Poisson}(n)>R).
\]
For $R+1>n$, exponential Markov inequality optimized at
$e^t=(R+1)/n$ gives
\[
 b_{n,R}\le
 \exp\!\left[-n+(R+1)\!\left(1+\log\frac n{R+1}\right)\right].
\]
Accepted histories still have exact conditional equilibrium because the
budget conditions only on counts. The exhaustion branch retains its actual
position and stopped wave.

\section{Clock access, finite precision, and the exact interface boundary}

The conditional theorem depends on its declared count-only interface.
Take $Q_0=\ell/2$, $\sigma_0=+1$, and a singleton block. Then
$Q_T=2cS\pmod\ell$. A new physical bit
$B=1_{\{S>T/2\}}$ records exactly the half-circle of $Q_T$.
The joint law of $(Q_T,B)$ has distance $1/2$ from uniform rotor
times its unchanged bit marginal. With symmetric bit error $\zeta$,
the distance is $(1-2\zeta)/2$; including the success branch before
conditioning leaves an unconditioned distinguishing gap
$(1-2\zeta)/(2e)$.

The bit is not ruled out by wave linearity or independence of the Poisson
increments. It changes the timing-neutral update
\eqref{prep:rotorgenerator}. A mechanical version adds an archive rotor
with initial exact copy $A_Q=Q_0$ and generator
$c\sigma(P_Q+P_{A_Q})$. Both actual positions receive the same drive,
so the copy stays exact. An old copy frozen during a final independent
preparation block is harmless by Theorem~\ref{prep:herald}; a newly driven
copy is not. The completed repair is to disconnect the wire and perform
that final preparation. Continuously recording the new flip phase defeats
every such finite refresh.

Exactness has a quantitative geometric boundary. For a homogeneous
singleton block with duration $T$ and speed $c$, put
$r=2cT/\ell=m+\theta$, $m\in\mathbb N_0$, $0\le\theta<1$.
The wrapped displacement density is $(m+1)/r$ on an arc of relative length
$\theta$ and $m/r$ elsewhere. Direct integration yields
\begin{equation}
 \operatorname{TV}(P_{\rm disp},\mathrm{Uniform})
       =\frac{\theta(1-\theta)}r.
 \label{prep:wraperror}
\end{equation}
Near one wrap, $r=1+\delta$, $|\delta|<1$, this is at most $|\delta|$.
Convolution gives the same bound uniformly for singular incoming positions.
The singleton success probability is now $\kappa Te^{-\kappa T}$,
so altered throughput must be counted separately.

For a preselected positive rate schedule $\kappa(t)$, define
$A(t)=\int_0^t\kappa(s)ds$, end the block at $A(T)=1$, and set
$c(t)=\ell\kappa(t)/2$. Conditional on one flip, $A(S)$ is uniform and
the displacement is $\sigma_0\ell(A(S)-1/2)$.
Thus exact preparation survives shared preselected modulation.
It need not survive feedback from new flips: with rate one before the
first flip and two afterward in action units, a singleton flip at phase
$s$ yields physical duration $(1+s)/2$. A retained duration then reveals
the phase. Merely matching speed and rate does not remove that archive.

The new probability law has not been moved into an equilibrium bath.
Its statistical content is instead explicit in the Poisson primitive.
Its instantaneous transitions, flat wave, timing-neutral count contact and
resource independence remain constitutive inputs. A microscopic finite-band
switching derivation is not supplied by this theorem.

\section{Causal selection of an initially unspecified ready law}

Return to the continuous detector of Theorem~\ref{cfg:exit}, but initially
leave its quantile law unspecified. Let $F_0$ be the CDF of its known packet
and let $U=F_0(X_0)$ have an atomless law $\nu$ on $(0,1)$.
The guidance identity $F_t(X_t)=U$ still holds on the admitted conservative
flow. At full separation, the right packet occupies ranks $(1-p,1)$.
Therefore
\begin{equation}
 \mathbb P(+)=f(p),\qquad f(p):=\nu((1-p,1)).
 \label{prep:response}
\end{equation}
It equals $p$ for every $p$ if and only if $\nu$ is uniform.
For $\nu(du)=[1+\epsilon(2u-1)]du$,
$f(p)=p+\epsilon p(1-p)$.

The selection principle of \cite{M20} is independent of the target formula:
an optional local detector activation, with no return, communication or
feedback coupling to a distant detector, does not change that distant
detector's terminal record law. The actual resource assumption is two
independent pointer quantiles with laws $\nu_A,\nu_B$, independent of
the known coherent input and complete admitted preparation history.
This is a statistical causal premise, not a consequence of commuting wave
Hamiltonians.

Use known coherent preparations
\[
 \psi_{st}=\sqrt{st}|00\rangle+\sqrt{s(1-t)}|01\rangle
                         +\sqrt{1-s}|11\rangle,\qquad
 \chi_s=\sqrt s|01\rangle+\sqrt{1-s}|10\rangle.
\]
Alice's optional detector is completed and its actual pointer is frozen
with disjoint supports before Bob uses his unchanged local reader.
For the first input, Bob alone gives $f_B(st)$. With Alice's interaction,
her branch zero has probability $f_A(s)$, and at its occupied frozen
coordinate Bob's local population is $t$. Her other branch gives zero.
Independence of the pointer configurations leaves Bob's original quantile
law unchanged conditional on Alice's event. Thus
\begin{align}
 \mathbb P_{\rm off}(B=0)&=f_B(st),&
 \mathbb P_{\rm on}(B=0)&=f_A(s)f_B(t),\\
 \mathbb P_{\rm off}^{\chi_s}(B=0)&=f_B(1-s),&
 \mathbb P_{\rm on}^{\chi_s}(B=0)&=1-f_A(s).
 \label{prep:causaltests}
\end{align}
The local branch statement follows directly from separated packet supports
and the full wave; no Born-weighted steering measurement is assumed.
The known coherent amplitude controls are stronger preparation resources
than an arbitrary unknown input and are counted here.

\begin{theorem}[Causal quantile selection]\label{prep:causal}
Within this source and independent-resource class, neutrality in both
families \eqref{prep:causaltests} for every $s,t\in[0,1]$ holds
if and only if both pointer quantile laws are uniform.
\end{theorem}
\begin{proof}
The equalities give
$f_B(st)=f_A(s)f_B(t)$ and $f_B(1-s)=1-f_A(s)$.
Since $f_B(1)=1$, take $t=1$ to obtain $f_A=f_B=:f$.
Then $f(st)=f(s)f(t)$ and $f(s)+f(1-s)=1$.
For $x+y\le1$, $x+y>0$,
\[
 f(x)+f(y)
 =f(x+y)\left[f\!\left(\frac{x}{x+y}\right)
             +f\!\left(\frac{y}{x+y}\right)\right]
 =f(x+y).
\]
It follows that $f(k/n)=k/n$. Monotonicity of tail probabilities and
rational approximation give $f(p)=p$ on the whole interval.
Its values determine the uniform law. The converse follows by substitution
in the actual event probabilities.
\end{proof}

This is an equivalence theorem within a supplied physical class. It replaces
an explicit distribution by an independently described causal test, but it
does not derive causal neutrality, coherent access to the test preparations,
or conditional independence from the older source/readout premise.
It is not an attracting preparation dynamics.

For a common response, suppose instead
$|f(st)-f(s)f(t)|\le\epsilon$ and
$|f(s)+f(1-s)-1|\le\epsilon$, uniformly, $0\le\epsilon<1$.
Then
\begin{equation}
 D:=\sup_p|f(p)-p|\le\min\{1,5\epsilon/(1-\epsilon)\}.
 \label{prep:causalerror}
\end{equation}
Indeed put $a=f(1/2)$, so $|a-1/2|\le\epsilon/2$.
For $p\le1/2$, multiplicativity at $(1/2,2p)$ gives
$|f(p)-p|\le aD+3\epsilon/2$. Complementation adds at most
$\epsilon$ for $p\ge1/2$. Thus
$(1-a)D\le5\epsilon/2$ and $1-a\ge(1-\epsilon)/2$.
For different readers with both causal errors bounded by $\epsilon$,
the same tests first give $\|f_A-f_B\|_\infty\le\epsilon$; applying
the preceding estimate with defect $2\epsilon$ to $f_B$ gives
$D_B\le\min(1,10\epsilon/(1-2\epsilon))$ when $\epsilon<1/2$.

This is terminal interval control. It is insufficient for complete timing.
For $\nu_N(du)=[1+b\cos(2\pi Nu)]du$,
\[
 \sup_p|f_N(p)-p|\le\frac{|b|}{2\pi N},\qquad
 \operatorname{TV}(\nu_N,du)=\frac{|b|}{\pi}.
\]
For a calibrated one-outlet input, exit time is an invertible monotone
function of the initial packet rank, so its full time law retains the
latter discrepancy. Finite sufficiently fine clock bins approximate it.
Neither a finite calibration set nor arbitrarily small uniform terminal
error implies complete-history total variation.

\section{SWAP prepares a subsystem by exporting its old state}

A finite discrete construction clarifies a different preparation resource.
Let a source $S$ and ancilla $B$ have equal dimension, let $W$ swap them,
and supply a definite ancilla wave and actual configuration $|q\rangle_B$.
For
\[
 H_{\rm sw}=\hbar gW,\qquad
 U_t=\cos(gt)I-i\sin(gt)W,\qquad
 T_{\rm sw}=\frac{\pi}{2g},
\]
the endpoint unitary is $-iW$. For a full source/reference wave,
\[
 \sum_n|n\rangle_S R_n\,|q\rangle_B
 \longmapsto-i|q\rangle_S\sum_n|n\rangle_B R_n.
\]
No reference operation has been performed.

\begin{proposition}[Finite definite-source replacement]\label{prep:swap}
Assume the minimal Bell generator on the complete $(S,B)$ configuration
during this pulse. For any initial source configuration law supported on
the nonzero source wave sectors, the final source configuration is $q$
almost surely. The final ancilla configuration has the original source
configuration distribution. Its wave and reference correlations likewise
contain the old source state.
\end{proposition}
\begin{proof}
For $n\ne q$, the pair $(n,q),(q,n)$ has wave amplitudes
$\cos(gt)R_n$ and $-i\sin(gt)R_n$.
Its one-way current divided by its occupied origin weight is
$2g\tan(gt)$, independent of $R_n$. Conditional on actual initial
$(n,q)$, survival is $\cos^2(gt)$, tending to zero at $T_{\rm sw}$.
The unique transfer is $(n,q)\to(q,n)$. The $n=q$ sector stays
$(q,q)$. Thus the source is certainly $q$ and the ancilla retains the
initial actual label. The wave identity gives the correlated quantum
statement.
\end{proof}
The proof does not add an endpoint Born draw and need not assume
initial source equilibrium. It assumes Bell timing and a definite fresh
ancilla. A second SWAP returns the old wave, reference correlations, and
any old nonequilibrium to the source. This is subsystem replacement,
not a preparation of the whole closed source in equilibrium. A displayed
reset of a used ancilla does not make it fresh.

If a historical record arrives at time $t_c$, the controller responds
after latency $\delta\geq0$, and preparation lasts $T_{\rm prep}$,
the replacement statement applies only at
$t_p=t_c+\delta+T_{\rm prep}$. Any intervening accessible record must
be propagated with the actual controller, unfinished transfer and old
memories. Preparing the new source at $t_p$ cannot retrospectively
change a record already secured before that time \cite{C01}.

\section{Conditional resource banks and complete future experiments}

The correct integration target retains two kinds of data separately:
moving apparatus coordinates $W$, and an exported configuration bank $A$.
The latter may have an arbitrary actual distribution. Internal references
and returning quantum memories remain within the wave fibers.

\begin{theorem}[Conditional moving-bank closure]\label{prep:bank}
Suppose, at each actual archive value $a$, the moving coordinates have
the normalized squared-norm law of their complete initial conditional wave.
During a declared finite programme assume:
\begin{enumerate}
 \item the exported coordinate bank $A$ remains frozen and the complete
 Hamiltonian is decomposable in $a$;
 \item every moving pointer, clock, or new memory belongs to the jointly
 prepared bank, or is supplied by a preparation guarantee uniform
 conditional on the complete actual past;
 \item the complete linear wave and guidance dynamics have the admitted
 conservative flows, and all waves and memories that can return are retained.
\end{enumerate}
Then conditional equilibrium of the moving bank is preserved. Final
physical records have the conditional quantum wave probabilities averaged
with the actual law of $A$. An initial same-wave joint preparation
defect $\delta$ gives complete retained-output distance at most $\delta$
against that comparison.
\end{theorem}
\begin{proof}
At fixed $a$, write the complete wave as a scalar fiber norm times a
normalized moving/internal field. Since $A$ is frozen, its fiber norm
is conserved. That scalar cancels from all moving-coordinate guidance
velocities. The normalized fiber density and its actual conditional
distribution satisfy the same conservative transport equation.
Uniqueness of the admitted flow preserves their initial equality.
The unknown reference is only a vector factor, so no reference access
enters this argument.

A physical adaptive controller is included in the moving bank or acts as
a specified position-diagonal control in the frozen bank; the entire
programme is one conditional linear evolution. Integrating its final
record regions gives the wave probabilities at each $a$. Averaging with
the actual $P_A$ gives the asserted comparison. The initial-defect statement
is Lemma~\ref{prep:push} on the complete wave and actual bank.
\end{proof}

The class allows an old archive to return as a position-diagonal classical
control or to be read by a prepared pointer. It allows coherent returns
of moving-bank waves and arbitrary internal quantum memories already
accounted for. It excludes motion or coherent recombination of the exported
nonequilibrium coordinate itself. The explicit inverse nodal preparation
calculates a failure outside that domain. This is a theorem domain, not
a universal prohibition on archive motion.

The nodal resource transfers directly to the binary detector by the affine
map $x=2a(u-1/2)$, with packet
$(2a)^{-1/2}\phi_\eta((x+a)/(2a))$.
The actual law and its TV bound are carried by the same map.
Consequently all click, opposite-click, finite null, time, and physical-copy
statistics have error at most $M/(4N)+\eta$ under
Theorem~\ref{prep:bank}. The full null wave and every failed-loading mode
are those of Theorem~\ref{cfg:exit}, not ideal replacements.

The heralded flat rotor can be shaped with a finite nonsingular transport.
On a circle large enough to contain a nonnegative compact target packet
$\phi$, use
\[
 \rho_\xi(x)=(1-\xi)|\phi(x)|^2+\xi/\ell,\qquad
 g(q)=F_\xi^{-1}(q/\ell),\quad 0<\xi<1.
\]
Its periodic lift has
$g'=1/(\ell\rho_\xi\circ g)\le1/\xi$.
The isotopy $g_s=(1-s)\operatorname{id}+sg$ generates the
half-density flow \eqref{prep:halfdensity}, taking the flat wave to
$\sqrt{\rho_\xi}$ and uniform configurations exactly to $\rho_\xi$.
For a shaping duration $\tau_s$, its speed is at most $\ell/\tau_s$,
and its spatial derivative is bounded by
$\max\{\xi^{-1},\ell\|\rho_\xi\|_\infty\}/\tau_s$.
Both follow from differentiating
$v_s=(g-\operatorname{id})\circ g_s^{-1}/\tau_s$.
Moreover
\[
 \|\sqrt{\rho_\xi}-\phi\|_2
 \le\sqrt{2(1-\sqrt{1-\xi})}\le\sqrt{2\xi},
\]
because $\int\phi\sqrt{\rho_\xi}\ge\sqrt{1-\xi}$.
The actual floor wave, its exits, and its nulls remain present. Comparing
it to the compact detector costs this wave norm only for a common complete
physical endpoint programme in the conditional equilibrium domain.
No bare path-law estimate under different waves follows from this norm alone.

For normalized conditional equilibrium fields $\alpha(q),\beta(q)$,
\[
 \bigl\||\alpha\rangle\langle\alpha|
       -|\beta\rangle\langle\beta|\bigr\|_1
 \le(\|\alpha\|+\|\beta\|)\|\alpha-\beta\|.
\]
Integration and Cauchy--Schwarz show that half the trace norm of their
complete position/internal outputs is at most $\|\alpha-\beta\|_2$.
Grouping physical records cannot increase this distance.
This establishes the needed complete-output comparison directly from the
full fields; it is not CP contractivity for an unknown nonlinear event law.

Suppose the ideal prefixes of a finite programme enter each protected
stage in its proved code domain, and a stage of duration $t_j$ has
reference-compatible complete-field error $e_j$. With total preparation
defect $\delta_{\rm prep}$, loading norm error $d_{\rm load}$, and shaping
error $d_{\rm shape}$ when used,
\begin{equation}
 D_{\rm output}\le
 \min\left\{1,\delta_{\rm prep}+d_{\rm load}+d_{\rm shape}
                         +\sum_j e_j\right\}.
 \label{prep:completebudget}
\end{equation}
First apply Lemma~\ref{prep:push} at the same actual wave.
For the equilibrium comparator, telescope full unitary products using
ideal prefixes and actual unitary suffixes. The next stage estimate is
therefore used on an ideal state in its code domain, not on an arbitrary
leaked state. Unitary suffixes have norm one; sum the field errors and
apply the preceding integrated inequality. The bound covers actual
physical endpoint records and their retained internal states on the
stated conditional domain. Rare record conditioning requires
Lemma~\ref{prep:condition}'s denominator. For a classical--quantum
output the same bound follows directly: if the event blocks are
$p\rho$ and $q\sigma$ with $p\ge q$, then
$p\|\rho-\sigma\|_1\le\|p\rho-q\sigma\|_1+p-q$,
while the complementary record block contributes at least $p-q$ to
the total trace norm. Dividing by two proves the normalized bound.

If successive supplies each have a guarantee uniform conditional on the
complete actual past, their preparation errors add. This follows by
coupling each successive supply on its actual matching prefix; probability
of any mismatch is at most the sum of the conditional defects.
A correct individual marginal at each stage does not support that coupling.

For a fixed stock and tolerance $\delta$, one may choose
$N_i=O((1+M_i)/\delta)$, $\eta_i=O(\delta)$; seed gradient costs
are then $O(\delta^{-3})$ at fixed $M_i$.
For the stochastic route, choose $n=O(\log(1/\delta))$,
$R=Cn$ for a sufficiently large constant $C>1$, and
$\xi=O(\delta^2)$. Preparation is completed before attachment to the
unknown input. Its controls are switched off and the retained hardware
must preserve the uniform interaction bounds required by any aperture
estimate. Hardware growth with growing residual coupling is not covered.
These simultaneous finite choices keep useful detector response; they
do not imply a Gaussian-reader or original Hamiltonian Bell law.

\section{Adversarial resource tests and precise remaining assumptions}

Correct uniform marginals do not make a two-cell stock independent.
The smooth joint law
$f(u,v)=1+\epsilon(2u-1)(2v-1)$ has uniform marginals but two
balanced terminal readers give
\[
 \mathbb P(++ )=\frac14+\frac{\epsilon}{16}.
\]
The joint preparation theorem \eqref{prep:stock} controls this correlation.
A shared switching process applied to initially equal rotors preserves
their equality; independent increments conditional on the actual past
cannot be replaced by correct separate rotor marginals.

Average preparation error does not justify selective success claims.
Let an archive bit $B$ have probability $\beta$, let the ready coordinate
be uniform on $B=0$, and uniform only on its upper half on $B=1$.
Its joint defect from a uniform coordinate with the same bit law is
$\beta/2$. Conditional on $B=1$, a balanced reader has error $1/2$.
The average can vanish while the selected error stays fixed.
The exact herald theorem avoids this failure by proving equality for
each accepted count string; stopping on an informative position or phase
does not satisfy that proof.

A copy before use must be the actual interaction
\eqref{cfg:snapshot}, with its pointer in the jointly prepared bank.
Then the changed conditional packet and residual-rank theorem apply.
Using an arbitrary old nonequilibrium pointer as a calibrated probe
changes the hypotheses. Coherent erasure of every moving-bank copy is
covered by the full-wave programme; a surviving copy must remain in
the state. Returning the exported nodal archive violates the frozen-bank
premise and recovers the calculated original bias.

The assumptions reduced by this part are therefore precise.
The nodal interaction replaces exact ready-subsystem equilibrium by
conditional BV regularity, known nodal wave resources, and specified
transport, with finite joint error. The homogeneous switching mechanism
replaces that regularity by a new independent Poisson law and a timing-neutral
final preparation interface. Causal selection replaces an unspecified
quantile law by a causal neutrality requirement on a defined preparation
class. SWAP transfers a definite ready resource while exporting the old
state under an already admitted Bell generator.

The older source/readout premise selects none of those stronger physical
or statistical inputs by itself. These chapters supply conditional
preparation and a complete preparation-to-record calculation within their
respective domains. The later pilot example can start from one definite
ordinary basis configuration and deterministic carrier copies, with only
the finite gas ensemble random. This removes an unknown Born draw from
that concrete example; it does not prepare an arbitrary closed universe in
equilibrium or derive the general unknown-input ensemble postulate.

% End consolidated source: parts/08_preparation.tex


% Begin consolidated source: parts/09_integration.tex
\chapter{A limiting discrete measurement chain}
\label{int:chain-chapter}

\section{A constitution that can actually be combined}
The preceding parts contain several explicit physical theories. For a reusable limiting-model calculation, this chapter fixes the finite full-wave configuration constitution. A finite complete sector resolution is part of the model. The source has a surviving normalized wave and one actual configuration. All physical apparatus memories are coherent factors with configuration labels. An inaccessible reference remains in the carried sector fibers. Prescribed finite coherent controls act on the source and apparatus; no wave-dependent classical expectation meter is added.

For the conditional comparator in this chapter, the wave follows a specified piecewise-constant Hamiltonian in physical time. Its actual native history is either the minimal Bell process admitted directly, or the zero-background law selected by the relative-entropy principle of Part~\ref{part:statistics}. The second presentation retains expected edge-current realization, the neutral reference process and its limiting prescription. The pilot construction supplies a further microscopic route to this finite-graph Bell comparator, with explicit finite-resource error and its own complete material coupling inventory. Initial joint configurations have wave weights as a stated ensemble law or within the stronger control-stable evacuation domain proved earlier. The preparation results do not automatically prepare equilibrium of an arbitrary entire source and all returning memories.

These premises suffice for the finite construction below. They do not include CPC as an independent premise, and they do not invoke a Born-calibrated measurement of an auxiliary system. Squared amplitudes enter through the declared initial joint law and its derived equivariance. Statistical content has a named location.

\section{A finite binary write with every branch retained}
Let $Q_0+Q_1=I_S$ be any orthogonal binary resolution, with arbitrary degeneracy. Write $V\in\mathcal H_S\otimes\mathcal H_R$, $\|V\|=1$, and $V_j=(Q_j\otimes I_R)V$, $p_j=\|V_j\|^2$. A blank apparatus qubit $A$ has wave $|0\rangle_A$. During $0\le t\le t_*:=\pi/(2g)$ set
\begin{equation}\label{int:writeH}
 H_{\rm wr}=\hbar g\,Q_1\otimes I_R\otimes\sigma_y^A,
 \qquad g>0.
\end{equation}
The sector resolution includes the source block $j$ and apparatus position $a$. Its exact wave is
\begin{equation}\label{int:writewave}
 \Psi_t=V_0\otimes|0\rangle_A+
 V_1\otimes(\cos(gt)|0\rangle_A+\sin(gt)|1\rangle_A).
\end{equation}
The only nonzero inter-sector current is
\[
 J_{(1,1),(1,0)}=g p_1\sin(2gt).
\]
For $0<t<t_*$ it is positive and the Bell hazard from $(1,0)$ is $2g\tan(gt)$. Its conditional survival from the beginning of the pulse is $\cos^2(gt)$. Thus the rate diverges near the end on a vanishing population, but there is exactly one event almost surely in the $j=1$ branch. The $j=0$ branch is inert. No event at the deterministic endpoint is required.

\begin{theorem}[Complete finite write and conditional continuation]\label{int:binary}
Under the constitution just specified, observing the actual apparatus sector at any fixed $s\le t_*$ gives
\begin{align}\label{int:branches}
 \Pr(A_s=1)&=p_1\sin^2(gs),&
 \sigma_1(s)&=\sin^2(gs)|V_1\rangle\langle V_1|,\\
 \Pr(A_s=0)&=p_0+p_1\cos^2(gs),&
 \sigma_0(s)&=|V_0+\cos(gs)V_1\rangle\langle V_0+\cos(gs)V_1|.
\end{align}
Here $\sigma_a$ is the unnormalized carried conditional state for future operations that preserve the recorded apparatus block. At $s=t_*$ these are the projective instrument $V\mapsto V_j$ with weights $p_j$. At an earlier $s$, the no-click daughter is the actual state in \eqref{int:branches}, rather than either the initial state or an ideal negative daughter.

The wave \eqref{int:writewave} itself is retained globally. If later operations reconnect the $A$ sectors, continuation must use that full wave and both branches. During block-preserving future evolution, normalized linear extraction of the conditional branch follows from factorization of the Hamiltonian and cancellation of its common branch weight in the Bell ratios.
\end{theorem}
\begin{proof}
Equation~\eqref{int:writewave} follows by exponentiating $\sigma_y$, whose rotation sends $|0\rangle$ to $\cos(gt)|0\rangle+\sin(gt)|1\rangle$. The current calculation uses $\langle1|\sigma_y|0\rangle=i$. Equivariance gives the squared norms of its two apparatus components. Orthogonality of $V_0,V_1$ gives the displayed probabilities. Restricting the full vector to $A=a$ gives the unnormalized vectors in \eqref{int:branches}, establishing both the probability and carried-state expression within this constitution.

For a later $A$-block-preserving Hamiltonian, write the normalized branch wave as $\psi_a$ and its constant total weight as $r_a$. Every current and origin weight wholly inside that block are respectively $r_a J^{(a)}$ and $r_a w^{(a)}$. The common factor cancels in \eqref{found:bell}; the conditional sector law inside that block is therefore its own Bell process. All source and reference amplitudes in the block are kept. If the Hamiltonian reconnects blocks, this cancellation no longer describes the complete source; the continuing Schr\"odinger wave does.
\end{proof}

This construction includes an actual configuration event and a physical apparatus coordinate. The intermediate occupancy $A=0$ is a null, and only at the completed pulse is it a sharp negative record. Record stability requires the future Hamiltonian and actual event generator to respect the completed archive, or the quantitative crossing bound in Part~\ref{part:records}. A display label does not create that stability by itself.

A readiness variable $d\in\{0,1\}$ may be retained as active classical data. Suppose its prepared law has $\Pr(d=1)=r$, independent of the unknown input conditionally on the declared actual past. In branch $d=1$ run \eqref{int:writeH}; in branch $d=0$ apply the identity and print a failure flag. Then the complete endpoint probabilities are $rp_0,rp_1,1-r$, and the failure daughter is $V$ together with the actual failed resource state. The displayed probability $r$ is a resource premise or a previously proved preparation result, not silently declared to be one. A finite supply of $M$ blank cells supports at most $M$ fresh attempts; exhaustion is a separately recorded identity branch. Reuse requires an actual reset theorem for both the wave and its configuration law.

\section{Sequential noncommuting measurements and retained records}
After a completed write retain $A$, apply a record-controlled unitary $U_a$ to $S$, and write another resolution $(R_b)$ into a fresh cell $B$ by the same construction. In the complete coherent bank these are block-diagonal Hamiltonian pulses controlled by $A$; they do not require a second measurement postulate. The complete configuration resolution remains fixed during these controls; the later measurement changes the interaction Hamiltonian, not the declared microscopic sectors. The endpoint argument uses equivariance for that full Hamiltonian. It does not assign the first pulse's special one-jump hazard to a later noncommuting write. The final wave is
\begin{equation}\label{int:seqwave}
 \sum_{a,b}(R_bU_aQ_a\otimes I_R)V\otimes|a\rangle_A|b\rangle_B.
\end{equation}
Consequently
\begin{equation}\label{int:seqBorn}
 \Pr(A=a,B=b)=\|(R_bU_aQ_a\otimes I_R)V\|^2,
\end{equation}
and each nonzero branch continues with its displayed vector divided by its norm while the archives remain isolated. The proof is multiplication of the two finite writing unitaries and equivariance on their joint sectors, followed by the branch argument of Theorem~\ref{int:binary}. Induction proves the corresponding finite adaptive product formula. Classical records with active preparation provenance are kept rather than averaged away.

For an explicit noncommuting example, take a qubit, first $Q_a$ the $Z$ projectors and then $R_b$ the $X$ projectors, with $U_a=I$. Equation~\eqref{int:seqBorn} gives $\Pr(a,b)=p_a/2$. The conditional source is the $X$ daughter, while the inaccessible reference retains the state proportional to the original $Z$-component reference vector. If the first pulse is stopped early, $Q_0V$ in the null branch is replaced by $V_0+\cos(gs)V_1$, so subsequent $X$ probabilities contain the surviving coherence. An ideal daughter after an imperfect declaration would give a different experiment.

A later copy or coherent erase is another unitary on $S,A,B$ and any return memory. The full-wave representation supplies its actual amplitudes. The reduced branch alone is sufficient only when the required block isolation persists. This is the structural distinction between effective conditional continuation and the irreversible-extraction constitution of Part~\ref{part:detectors}.

\section{A simultaneous finite error budget}\label{int:budget}
For a fixed finite complete graph, the following errors can be combined because their state and output spaces have been specified:
\begin{itemize}
\item $\epsilon_{\rm prep}$ bounds total variation of the actual initial configuration law from the equilibrium reference law for the complete admitted preparation, including active classical keys. Applying the same subsequent transition law contracts this error.
\item $\epsilon_H=\hbar^{-1}\int_0^T\|H-G\|dt$ bounds a perturbation of the complete finite pulse programme, and $B=\hbar^{-1}\int_0^T\max(\|H\|,\|G\|)dt$. The fixed-graph Bell path bound is $2d\epsilon_H(1+6B)$ for identical initial waves.
\item A variational background of size $\epsilon_{\rm bg}$ contributes at most $\epsilon_{\rm bg}a_{\max}(d-1)T$ to the complete path comparison.
\item A physical archive used to represent earlier source labels adds the finite write error and the expected corrupting-crossing bound of Part~\ref{part:records}. These are needed for past-history meaning, even when final archive weights already agree.
\end{itemize}
Thus a common recorded path functional has the sufficient bound
\begin{equation}\label{int:budget-eq}
 \epsilon_{\rm path}\le\epsilon_{\rm prep}
   +2d\epsilon_H(1+6B)
   +\epsilon_{\rm bg}a_{\max}(d-1)T.
\end{equation}
For a complete output retaining the common global wave as well as the recorded path, add its trace-distance bound $\|\Psi_T-\Phi_T\|\le\epsilon_H$. For faithful sampled history add both writing and subsequent corruption errors. Conditioning on a rare record requires the probability denominator in Lemma~\ref{found:errors}.

One common useful regime fixes the graph, a finite number of pulses, $T$ and $B$, and takes the initial preparation error, integrated control error, variational background, finite writing error and protected crossing budget to zero. The responsiveness $g$ in \eqref{int:writeH} stays positive and the measurement time stays $\pi/(2g)$. This limit does not suppress all useful events. If resources enlarge the graph, $d$ and the relevant operator norms must be tracked explicitly. The spectral-gap protection theorem supplies a different, complete-wave estimate under its own bounded block assumptions; a small reduced-wave error alone is not substituted for $\epsilon_H$ in \eqref{int:budget-eq}.

The extraction/filtering instrument architecture has its own full-output sum of half-diamond errors in Part~\ref{part:detectors}. That is a second valid finite chain, but its Gaussian/native admission and irreversible branches are not replaced by the present configuration law without a further embedding theorem.

\section{Dependencies of the limiting-model calculation}
\begin{center}\small
\begin{longtable}{@{}>{\raggedright\arraybackslash}p{.28\textwidth}>{\raggedright\arraybackslash}p{.35\textwidth}>{\raggedright\arraybackslash}p{.29\textwidth}@{}}\toprule
Physical or mathematical input & Consequence proved in this monograph & What it does not supply\\\midrule\endhead
Source/readout descent and target completion & Exact test for lost, target-active information; predictive-current quotient & A probability measure, least-traffic law or physical access rule\\
Canonical bond action and conservative exporter & Hamiltonian current and finite signed-charge export & Actual stochastic chemistry or complete response neutrality\\
Scalar pair chemistry and calibrated carrier population & Whole tagged-path Bell limit; physical residence denominator & Universal contact admission or harmless readable microscopic histories\\
Expected edge-current realization and path-entropy principle & Selected complete finite-background history and physical-time Bell limit & Derivation of the principle or neutrality of every statistical reference\\
CPC or the explicit finite-tag/interchange premises & Shared-Wiener population and signal matching & Gaussian noise origin, positive phase lift or universal CPC\\
Responsive zero-channel tests and reference/archive conditions & Faithful intrinsic daughter within the declared class & Universal physical necessity of those tests and conditions\\
Native reader, coherent converter and finite control library & Finite instruments, actual nulls and composed error bounds & Original Hamiltonian Bell trajectories\\
Bounded block interaction and prepared code sector & Quantitative transport protection and finite horizon & Protection against unrestricted logical couplings or scalar clocks\\
Full-wave source, admitted equivariant generator and joint initial law & Actual configuration records, conditional branch predictions and coherent returns & Selection of minimal actual traffic from endpoint Born records\\
Finite correct writes and small archive-changing traffic & Reliable sampled histories & Faithfulness of every first-entry record or every microscopic event\\
Explicit regular preparation class and protocol & Conditional operational preparation within the proved future-use domain & Arbitrary global equilibrium or arbitrary nonequilibrium-memory returns\\\bottomrule
\end{longtable}
\end{center}

The monograph therefore establishes conditional realizations and controlled operational constructions. It retains the strongest incompatibility tests as results. In particular, exact readable histories of the unmodified original Bell process and a common affine complete-input law cannot be freely combined in the comparison class of Part~\ref{part:access}. Physical recording changes the complete experiment; a complete joint configuration model does not license a passive classical current writer by definition.

The statistical and interaction inputs of this limiting-model calculation remain explicit. Part~\ref{part:pilot} now supplies a deterministic microscopic approximation with controlled full-path error within P1--P4; Part~\ref{part:massive} supplies a separately postulated continuous-configuration alternative. Chapter~\ref{syn:chapter} gives the final resolution statement and its remaining foundational obligations.

% End consolidated source: parts/09_integration.tex

\chapter{Internal resolution and the remaining foundational obligations}
\label{syn:chapter}

\section{Two complete chains with distinct constitutions}
The preceding finite discrete chain isolates the familiar measurement
calculation once its event process is supplied. The pilot construction
now supplies an explicit physical approximation to that process, and
its autonomous example proves historical copies within the same
ordinary Hamiltonian. The massive construction supplies a different
exact ontology and motion law, with its own controlled autonomous
measurement implementation. The following synthesis keeps these
conclusions separate.

\begin{theorem}[Conditional internal resolution]\label{syn:resolution}
The following are two independent statements.
\begin{enumerate}
\item Fix a finite ordinary configuration graph, a bounded admitted
Hamiltonian programme with bounded-variation currents, a finite horizon,
and the pilot constitution P1--P4. Assume the initial carrier census obeys
$\mathbb E\|x^N(0)-w(0)\|_1\to0$, the predesignated tag law is dominated by
them, the gas is independently initialized as specified, packet stock is
initially empty, and all exporter and receiver budgets are supplied.
Under the simultaneous scaling in Theorem~\ref{p:main}, the complete
ordinary tagged physical-time path converges in total variation to the
minimal Bell path with the specified initial law. For the fixed
autonomous material programme of Chapter~\ref{p:chapter-records}, with
its stated calibrated or known basis-ready input, this same comparison
controls its retained ordinary outputs and first-pass monomial-copy
histories, including null/loss resources, reset receivers, feedback and
the inaccessible reference included in the declared configuration
convention.
\item Fix the massive constitution of
Chapter~\ref{mc:chapter-constitution}, including guidance, complete
initial equilibrium and its smooth finite domain. Its full dynamics
defines conservative continuous configuration paths. The finite
measurement programmes and resource choices of
Theorem~\ref{mc:closure} have the proved retained-output and
historical-archive error budgets under the autonomous implementation.
These errors can be made arbitrarily small for each fixed admitted
programme by the stated sequence of finite resource choices.
\end{enumerate}
\end{theorem}
\begin{proof}
For the first statement let $P_{\rm gas}$ be the finite spatial contact
law, $P_{\rm Pois}$ its marked Poisson comparison, and $P_{\rm exc}$ the
excess-only reaction comparison. Both contact laws are passed through
one initialized causal device, independent of the gas conditional on
the declared coherent preparation. Its overflow convention is common.
The gas comparison therefore contracts to the reaction output.
Recombination is coupled until the first discrepant matched-packet
service, with its probability bounded by the stopped compensator and
the pathwise stock budget. Finally the signed-queue tracking and
tagged-path theorem apply on the same graph. Thus
\[
 \TV(P_{\rm pilot,tag},P_{\rm Bell})\le
 \frac{(R_NT)^2}{M_N}
 +\sum_e\frac{\kappa_e\mu_N B_e}{2a_e}
 +\epsilon_{\rm kin}(N,T)\longrightarrow0.
\]
These steps are proved in Part~\ref{part:pilot} and
Theorems~\ref{kin:tracking} and \ref{kin:path}. Bell existence and
nonexplosion hold with the admitted initial domination, including nodal
endpoints. A common ordinary path functional contracts total variation.
The monomial-cut proof makes the ideal copy an actual past-label copy,
so this contraction controls the joint event of history failure as well
as endpoint outputs. It is not necessary to charge the same whole-path
error separately for every record. The first-pass restriction is part
of that proof.

For the second statement, the current-integrability existence argument
gives the complete guidance flow. The exact writer and archive
constructions establish the driven chain. Common retained-state
comparisons, the autonomous controller's pointer graph-norm estimate,
and the absolute surface-current history bound give
Theorem~\ref{mc:closure}. The resource choices fix packet separation,
feedback and protection before increasing controller resources. This
proof neither invokes nor yields the discrete Bell generator.
\end{proof}

For either chain a conditional comparison uses
Lemma~\ref{found:errors}. If the reference record has probability $p>0$
and the joint error is $\epsilon<p$, the second probability is positive
and the conditional TV bound is at most $\min\{1,2\epsilon/p\}$.
There is no uniform conclusion for arbitrarily rare records. Revealing
previously excluded pilot microcoordinates changes the output space;
it is not an application of that lemma.

\section{What the integration supersedes and what it preserves}
The older source-to-event frontier separated canonical current
production from a supplied Markov reaction law and an instantaneous
signed queue. For the enlarged pilot constitution this separation is
bridged by the finite spatial ensemble and finite recombination
construction. It is therefore no longer accurate to say that the
gas-to-reaction-to-Bell-path implication is unproved in that domain.
It remains accurate that the original source/readout premise alone
does not force this constitution.

The physical access countermodels remain useful and valid. A reciprocal
pilot-action reader, a retained ledger or a neutral native product can
carry information in a model admitting its mixed coupling. P3 excludes
that extra force from the enlarged theory. Recombination does not erase
the ledger, and the exclusion is not a theorem about every possible
wire. Similarly, a supplied minimal mean incidence does not determine
conditional timing, and an equivariant generator can contain surplus
traffic. The pilot construction addresses these freedoms inside its
specified inventory and ensemble, with finite corrections rather than
exact finite-resource minimality.

The detector, preparation and protection results are preserved under
their own assumptions. Absorbing extraction, quantum-latch filtering,
continuous affine readers and full-wave configurations have distinct
null and return rules. The apparatus comparisons do not identify them.
Finite pre-latch records, pending excitation, loss daughters, recurrence,
recovery and coherent return retain their stated mathematical content.
The historical-archive obstruction also remains: endpoint Born weights
alone cannot tell whether an archive tells the truth about its earlier
actual value.

The massive theory strengthens the continuous-configuration side by
supplying a semibounded material implementation, a transported-trap
reset and a controller estimate strong enough for history. Its diffusion
rivals show why reliable records and equivariance do not derive the
chosen guidance law. Likewise, its existence theorem cannot be imported
as a smooth realization of the pilot exporters without an embedding
proof.

\section{Distances, resources and the final boundary}
\begin{center}\small
\begin{tabularx}{\textwidth}{@{}>{\raggedright\arraybackslash}p{.25\textwidth}XX@{}}\toprule
Proved comparison & Output controlled & Required restriction\\\midrule
Pilot path TV & Complete ordinary tag path, times and common recorded functionals & Fixed graph/programme/horizon; admitted preparation and pilot scaling\\
Monomial history & Actual label at the copy crossing and retained first-pass suffix & Blank receiver, aligned fine currents, archive-preserving suffix\\
Massive state/output & Complete retained quantum state and final declared outcomes & Full receivers/reference retained; fixed input and controller domain\\
Massive history & Specified declarations, holds and transfers & Derivative and absolute-flux budgets, not wave norm alone\\
Common coherent return & Complete state-distance comparison & Same retained dynamics on both states; readability may be destroyed\\\bottomrule
\end{tabularx}
\end{center}

The remaining obligations concern physical premises and extensions,
not an omitted cancellation estimate in the established pilot chain.
First, an explanation from weaker source principles would need to
justify or replace the complete reaction catalogue and the P3 access
restriction, while retaining the known ledger and native-counter
tests. Second, the admitted statistical preparations remain resources:
independent spatial gas and general carrier calibration in the pilot
theory, and complete initial equilibrium in the massive theory. The
known basis-ready pilot example removes uncertain initial ordinary
occupancy for one explicit preparable input; it does not establish
universal unknown-input preparation.

Third, exact finite-resource Bell intensities are not established.
Finite gas depletion and finite mixed-stock service are explicit
departures, controlled by the comparison bounds. Fourth, the smooth
isolated-contact construction excludes asynchronous modifications of
its register. A single smooth Hamiltonian for the entire canonical
source, exporter, pilot and contact device remains an extension.
Finally, the finite pilot clock's recurrence and the fixed-domain constants
leave indefinite storage, arbitrary growing graphs and unrestricted
returning-memory programmes outside the present uniform claims.

Within these boundaries the measurement programme has a conditional
effective internal resolution through the pilot medium and a separate
massive operational completion. Its mathematical conclusions reach
actual events, physical records and retained noncommuting continuation.
The deeper question is whether the declared constitutions and
preparation resources follow from independently justified physics.


\appendix

% Begin consolidated source: parts/10_copies.tex
\chapter{Supplied complete-input copies and global terminal repair}
\label{copy:chapter}

This appendix preserves a stronger-resource construction from
\cite{M11}. It is not a local measurement of one unknown input with an
inaccessible reference. The apparatus is supplied with independent
copies of the entire pure input bank, including its internal reference
entanglement, and the final repair can act jointly on that entire bank.
Native probes estimate currents from those copies; a record-driven
classical pump feeds the scalar pair chemistry; further finite
tomography determines an approximate global terminal gate. These
resources give a nonempty conditional construction of original
Hamiltonian Bell paths and terminal daughters.

The distinction matters mathematically. At fixed unknown pure input
$\psi$, the initial source is
\begin{equation}
 P_\psi^{\mathrm{target}}\otimes P_\psi^{\otimes M},
 \qquad P_\psi=|\psi\rangle\langle\psi|.       \label{copy:input}
\end{equation}
The copies are supplied by an independent preparation conditional on
the actual $\psi$; they are not cloned from the target. For an actual
ensemble $\mu$, its complete preparation is
$\int P_\psi^{\otimes(M+1)}\mu(d\psi)$, which is not determined by
$\int P_\psi\mu(d\psi)$. Thus two single-copy-equivalent ensembles
need not give the same statistics for this apparatus. No single-input
affinity obstruction is evaded by silently deleting these resources.

\section{Finite native current pilots}

Fix a finite sector graph, finite physical horizon $T$ and bounded
deterministic piecewise $C^1$ Hamiltonian programme $H(t)$ on the
complete $d$-dimensional pure bank. For one orientation $e=(r,q)$ of
each of its $E$ bonds, in units $\hbar=1$ put
\begin{equation}
 A_e(t)=\frac{P_qH(t)P_r-P_rH(t)P_q}{i},\quad
 J_e(t)=\langle\psi_t,A_e(t)\psi_t\rangle,
 \quad\|A_e(t)\|\le a,\quad i\dot\psi_t=H(t)\psi_t.        \label{copy:current}
\end{equation}
The programme makes each $J_e$ of bounded variation. Assume that the
native ports $L_e=kA_e(t)$ are physically admitted on every pilot.
Their actual law is the already supplied multichannel native law,
with independent Wiener innovations for different pilots and ports:
\begin{align}
 d\psi_j&=\left[-iH-\frac{k^2}{2}\sum_e(A_e-a_{je})^2\right]
                 \psi_jdt+k\sum_e(A_e-a_{je})\psi_jdW_{je},\nonumber\\
 dY_{je}&=2ka_{je}dt+dW_{je},\qquad
 a_{je}=\langle\psi_j,A_e\psi_j\rangle.          \label{copy:pilots}
\end{align}
All pilot states, classical records and used resources are retained.
The target receives only the same coherent $H(t)$, not the diagnostic
couplings. The controller has access to $Y$, not to the separate
innovations $W$ or an unknown-wavefunction expectation wire.

Define its causal finite-bandwidth estimator by
\begin{equation}
 Z_e=\frac1{2kM}\sum_jY_{je},\qquad
 dv_e=-\omega v_e dt+\omega dZ_e,\quad  v_e(0)=0,
 \qquad\widehat J_e=\operatorname{clip}_{[-a,a]}v_e.          \label{copy:estimator}
\end{equation}
The filters, their clocks and every controller coordinate belong to
the complete classical bank. Their coefficients use known apparatus
calibrations only.

\begin{theorem}[Integrated current estimation]\label{copy:estimate}
For the independent native pilots,
\begin{equation}
 \mathbb E\int_0^T|\widehat J_e-J_e|dt
 \le\frac{|J_e(0)|+\operatorname{Var}_{[0,T]}J_e}{\omega}
   +k^2Ea^3T^2+\frac{aT}{\sqrt M}
                   +T\sqrt{\frac{\omega}{8k^2M}}.           \label{copy:estimate-bound}
\end{equation}
In particular, $k^2=M^{-1/4}$ and $\omega=M^{1/4}$ give
$b_M:=\sum_e\mathbb E\int_0^T|\widehat J_e-J_e|dt
=O(M^{-1/4})$ for the fixed programme.
\end{theorem}
\begin{proof}
The mean pilot state obeys
$\dot\rho_k=-i[H,\rho_k]+k^2\sum_e\mathcal D[A_e]\rho_k$,
where $\mathcal D[A]\rho=A\rho A-\{A^2,\rho\}/2$.
For a density matrix $\|\mathcal D[A]\rho\|_1\le2a^2$.
Variation of constants in the interaction picture therefore gives
$\|\rho_k(t)-P_{\psi_t}\|_1\le2k^2Ea^2t$ and a current bias at
most $2k^2Ea^3t$. Independence and $|a_{je}|\le a$ give
$\mathbb E|M^{-1}\sum_ja_{je}-\mathbb E a_{1e}|\le a/\sqrt M$.

Let $K_\omega f(t)=\omega\int_0^te^{-\omega(t-s)}f(s)ds$.
This convolution contracts $L^1([0,T])$. The filter is the sum of
$K_\omega J_e$, the convolutions of the bias and empirical error,
and the noise
$U_e(t)=\omega(2k\sqrt M)^{-1}
\int_0^te^{-\omega(t-s)}dB_e(s)$,
where $B_e=M^{-1/2}\sum_jW_{je}$ is Brownian. It need not be
independent of the empirical mean. Integration against the variation
measure of $J_e$, including the initial zero-filter transient and
programme switches, gives
$\int|K_\omega J_e-J_e|
\le(|J_e(0)|+\operatorname{Var}J_e)/\omega$.
The integrated bias is at most $k^2Ea^3T^2$ and the empirical term
at most $aT/\sqrt M$. It\^o isometry gives
$\mathbb E U_e(t)^2=\omega(1-e^{-2\omega t})/(8k^2M)$;
Cauchy--Schwarz and integration give the last term. Clipping cannot
increase distance to $J_e\in[-a,a]$. The stated scales balance all
dominant errors at $O(M^{-1/4})$.
\end{proof}

\section{Record-driven pumping and a robust scalar-pair bridge}

Supply a classical reservoir with
$\dot{\widehat\Pi}_e=\widehat J_e$, debiting its signed account by
the same amount. Two directional accounts of size $aT$ per edge
suffice. This is a record-driven charge pump in the classical
controller, not the forbidden passive classical expectation meter.
Its equality with the coherent norm current is approximate, whereas
its own charge accounting is exact.

Export packets of charge $1/N$ at residual thresholds $\pm1/N$ and
reset the residual after each export. The cumulative export satisfies
\begin{equation}
 A_{N,e}(t)=\int_0^t\widehat J_e(s)ds+e_{N,e}(t),
 \qquad\sup_t|e_{N,e}(t)|\le N^{-1}.            \label{copy:export}
\end{equation}
The total export count is at most $NEaT+O(E)$.
Opposite packets cancel; a packet with an empty origin waits. Every
eligible packet/carrier pair reacts at coefficient $\kappa_e\mu_N/N$
with fixed $0<\kappa_-\le\kappa_e\le\kappa_+$. This is precisely
the scalar pair model of Theorem~\ref{kin:tracking}, with its supplied
Markov chemistry and complete participation. It does not use a
separately prescribed normalized destination allocator from an earlier
model. Fixed-tag division by its origin population is therefore the
pair-counting identity \eqref{kin:tagrate}, not a new statistical
controller instruction.

\begin{theorem}[Robust current forcing and complete tagged paths]
\label{copy:robust}
Keep the fixed finite graph, bounded-variation target currents,
coherent nodal bound of \eqref{kin:nodal}, initially empty queues and
the scalar chemistry just specified. Suppose bounded predictable
$\widehat J$ obeys $b_N=\sum_e\mathbb E\int|\widehat J_e-J_e|\to0$.
Assume calibrated carrier populations $x_N(0)\to w(0)$ and compatible
fixed-tag laws $\nu_N\to\nu\le Cw(0)$. If
$\mu_N\to\infty$ and $\mu_N/N\to0$, directional flux errors,
uniform population errors and total variation of the entire tagged
path relative to the physical-time Hamiltonian Bell process tend to
zero. No condition $\mu_Nb_N\to0$ is needed.
\end{theorem}
\begin{proof}
Only the additional forcing estimates in the already-proved scalar
tracking and node localization arguments need modification. At fixed
population cutoff $\delta>0$, use their companion service
$\phi_t(u)=a_+(t)[u]_+-a_-(t)[-u]_+$ with divided-difference slopes
in $[a_0,a_1]=[\kappa_-\delta,\kappa_+]$. Each escape consumes an
export, so the variation bound on the physical population, and hence
the coefficients, remains uniform because the exports are bounded.

After removing the exporter error and reaction martingale, let
$\widehat y$ solve
$\dot{\widehat y}=\mu_N(\widehat J-\phi_t(\widehat y))$.
For the same realized coefficient history let
$\dot y_J=\mu_N(J-\phi_t(y_J))$, with matching zero starts.
Scalar monotonicity gives pathwise
\[
 \frac d{dt}|\widehat y-y_J|
 \le\mu_N|\widehat J-J|-\mu_Na_0|\widehat y-y_J|,
 \qquad
 \int_0^T|\widehat y-y_J|dt
                   \le a_0^{-1}\int_0^T|\widehat J-J|dt.
\]
This is why forcing error is not multiplied by the fast scale in the
integrated flux bound. The remaining exporter and martingale estimates
of Theorem~\ref{kin:tracking} use only the uniform bound
$\|\widehat y\|_\infty\le a/a_0$ and coefficient variation, which
still hold. Root tracking uses the bounded variation of the original
$J$, not the noisy $\widehat J$. Consequently
\begin{equation}
 R_{\delta,N}\le C_\delta\left(
     \mu_N^{-1}+\mu_N/N+\sqrt{\mu_N/N}+b_N\right).           \label{copy:robust-R}
\end{equation}

The exact population/queue balance acquires only the accumulated
forcing term:
\[
 x_N+Bz_N-w=x_N(0)-w(0)+B e_N+
                            B\int_0^t(\widehat J-J)ds.
\]
Its expected uniform size is bounded by
$\eta_N=\|x_N(0)-w(0)\|_\infty+C_B/N+C_Bb_N$.
The low-population argument of \eqref{kin:smallmass} therefore gives,
with $n$ sectors and $L=EaT+O(N^{-1})$,
\[
 D_{\delta,N}\le C_H\sqrt{T\left(n\delta T+
               \frac{L}{\kappa_-\mu_N\delta}+nT\eta_N\right)},
 \qquad\epsilon_{F,N}\le3R_{\delta,N}+2D_{\delta,N}.
\]
Taking $N\to\infty$ and then $\delta\downarrow0$ proves directional
flux convergence. The population martingale and conservation then
give $\epsilon_{x,N}\to0$ exactly as in the baseline proof.
The fixed-tag hazard is unchanged as a function of the complete
reaction state. Thus Theorem~\ref{kin:path} applies its same
regular-level node localization, with the new $\epsilon_F,\epsilon_x$,
to give full path-law convergence, including exact event times.
Initial-law error costs $\operatorname{TV}(\nu_N,\nu)$. This transfer
uses the actual scalar generator, not mean-current agreement alone.
\end{proof}

One simultaneous realization is
\begin{equation}
 M=N,\quad k^2=N^{-1/4},\quad\omega=N^{1/4},\quad\mu_N=N^{1/3}.
                                                        \label{copy:scaling}
\end{equation}
Then $R_{\delta,N}=O_\delta(N^{-1/4})$. Nodal localization gives
convergence for each fixed admitted complete input; no uniform
input-independent global path rate is asserted. The total diagnostic coupling action is bounded by
$Ea^2TN^{3/4}$; packet and reaction archives have $O(NEaT)$ entries.
An undersized bank must retain its exhaustion branch. The initial
population calibration remains a resource premise: when all admissible
inputs start in one known ready sector, $w(0)$ and the tag start are
fixed without Born sampling. For general $w(0)$ this appendix assumes
their preparation rather than deriving it from the pilots.
Continuous native pointer coordinates and calibrated clocks remain
declared ideal resources on each finite horizon; the entry count is
not a bound on their numerical recording precision.

\section{Retained pilots and a terminal continuation obstruction}

Assume the pilot bank is autonomous: no queue, tag or reaction archive
feeds back into it, and its initial state and stochastic primitives
are independent of reaction primitives conditional on $\psi$.
Conditioning on its whole actual output $Y$ then fixes a bounded
forcing $\widehat J(Y)$ without changing the reaction clocks. Applying
the preceding estimates with the conditional forcing error and
integrating proves
\begin{equation}
 \delta_{\mathrm{joint},N}:=
 \operatorname{TV}\left(\mathcal L(Q_N,Y\mid\psi),
       \mathcal L(Q^B_\psi)\otimes\mathcal L(Y\mid\psi)\right)
                       \longrightarrow0.                  \label{copy:joint}
\end{equation}
The $L^1$ forcing estimate and concavity in the low-population bound
justify the averaging. The same pilot-state conditional kernel,
including its retained quantum resources, can be attached on both
sides and preserves the bound in CQ trace distance. An ensemble
mixture keeps the common latent $\psi$ in both factors; it is not
generally the product of marginal mixture laws.

The target, however, has remained exactly $P_{\psi_t}$ conditional
on the pilot and tag history at fixed $\psi$. It has not acquired a
selected sector daughter. With Born initialization, a subsequent sharp
sector benchmark independent of that tag has repeat probability
$\sum_qw_q^2$, rather than one. Postselecting agreement of two
independent labels would produce weights proportional to $w_q^2$;
for $w=(3/4,1/4)$ the rejection probability is $3/8$ and the accepted
weights are $(9/10,1/10)$. Keeping those rejects defines a different
experiment; deleting them is not a continuation repair.

Nor does \eqref{copy:joint} automatically include microscopic queues,
other carrier histories or cancellation archives. They are not
functions of the autonomous $Y$ and no ideal common kernel for their
returns has been proved. They require a stated no-return domain or a
new benchmark-extension theorem.

\section{Finite native tomography and the global repair}

Now additionally assume that the entire target bank is controllable,
including $S+R$ if they are internally entangled. After a disclosed
embedding let $d=2^s$. Supply $(d^2-1)m$ further independent complete
input copies, independent of production conditional on $\psi$.
Allocate $m$ to each nonidentity Pauli matrix $W_\ell$. During a
known calibration hold, use a finite native monitor
$L=k_{\mathrm{cal}}W_\ell$ for time $\tau$. Its sign record has
\begin{equation}
 \mathbb P(s_{\ell j}=+1)=\frac{1+v\langle W_\ell\rangle_\psi}{2},
 \quad v=1-2\Phi(-2k_{\mathrm{cal}}\sqrt\tau)>0,
 \quad\widehat a_\ell=\frac1{mv}\sum_js_{\ell j}.             \label{copy:tomography}
\end{equation}
This follows from the actual native eigenrecord normals with means
$\pm2k_{\mathrm{cal}}\tau$ and variance $\tau$. Thus calibration
is finite and its visibility is corrected, rather than replaced by an
exact Born projective measurement. It still uses the native statistical
primitive.

\begin{lemma}[Pure-input estimation with finite records]
\label{copy:tomography-error}
Set $\widehat\rho=d^{-1}(I+\sum_\ell\widehat a_\ell W_\ell)$
and choose a top eigenvector $u$ by a fixed measurable tie rule. Then,
with $e=\min_\alpha\|u-e^{i\alpha}\psi\|$,
\begin{equation}
 \mathbb E\|\widehat\rho-P_\psi\|_F^2\le\frac d{mv^2},
 \quad\|P_u-P_\psi\|_F\le2\|\widehat\rho-P_\psi\|_F,
 \quad\mathbb E e\le2\sqrt{\frac d{mv^2}}.                  \label{copy:tomo-bound}
\end{equation}
\end{lemma}
\begin{proof}
Each visibility-corrected estimator is unbiased and has variance at
most $1/(mv^2)$. Pauli orthogonality gives mean squared Frobenius
error at most $(d^2-1)/(dmv^2)\le d/(mv^2)$.
A top eigenvector minimizes $\|\widehat\rho-P\|_F$ over pure
projectors, even when $\widehat\rho$ is not positive. Its triangle
inequality therefore gives the middle claim. For unit vectors with
phase aligned, $e^2=2(1-|\langle u,\psi\rangle|)$ and
$\|P_u-P_\psi\|_F^2=2(1-|\langle u,\psi\rangle|^2)\ge e^2$.
Cauchy--Schwarz proves the last bound. Propagating $u$ by the known
unitary programme preserves this phase distance at $T$.
\end{proof}

For terminal tag $q$, define $b_q=\|P_qu_T\|^2$. If $0<b_q<1$,
set
\begin{align}
 v_q&=P_qu_T/\sqrt{b_q},&
 z_q&=(v_q-\sqrt{b_q}u_T)/\sqrt{1-b_q},\nonumber\\
 K_q&=i(|z_q\rangle\langle u_T|-|u_T\rangle\langle z_q|),&
 V_q&=e^{-i\theta_qK_q},\quad\theta_q=\arccos\sqrt{b_q}.
                                                        \label{copy:repair-gate}
\end{align}
$u_T,z_q$ are orthonormal and direct exponentiation gives
$V_qu_T=v_q$. Use the identity at $b_q=1$ and also as a declared
fallback at $b_q=0$; no run is postselected away. The construction is
phase invariant. An admitted physical implementation has Hamiltonian
$H_q=\theta_qK_q/\tau_{\mathrm{fb}}$ of norm at most
$\pi/(2\tau_{\mathrm{fb}})$, or a finite global compiler with
uniform operator error $u_{\mathrm{impl}}$. Its coefficients depend
only on actual tomography records and the terminal tag. The gates
must respect any additional protected charges; their universal global
admission is a new resource assumption.

\begin{theorem}[Bell path and globally repaired terminal daughter]
\label{copy:terminal}
Assume the stated complete-copy resources, autonomous pilots,
Born-compatible tag initialization and global controls. For
$w_q(T)>0$ put $\phi_q=P_q\psi_T/\sqrt{w_q(T)}$. The actual path
and repaired target differ from a Bell path with terminal target
$P_{\phi_{Q_T}}$ by CQ trace distance at most
\begin{equation}
 \delta_{\mathrm{path},N}+6\sqrt{\frac d{mv^2}}+
                                               u_{\mathrm{impl}}.         \label{copy:terminal-bound}
\end{equation}
The independent tomography records and conditional states of used
tomography specimens may be retained on both sides. Including the
production pilots and their conditional quantum bank replaces
$\delta_{\mathrm{path},N}$ by $\delta_{\mathrm{joint},N}$.
\end{theorem}
\begin{proof}
Phase-align $u_T,\psi_T$ for analysis and put
$e_q=\|P_q(u_T-\psi_T)\|$. For $b_q>0$, the normalized-projection
inequality and unitarity give
\[
 D(P_{V_q\psi_T},P_{\phi_q})
       \le\min\{1,e+2e_q/\sqrt{w_q}\}.
\]
Indeed $V_qu_T=v_q$, so the first vector difference is at most $e$;
adding and subtracting $P_qu_T/\sqrt{w_q}$ bounds the two normalized
projections by $2e_q/\sqrt{w_q}$. At $b_q=0$, $e_q=\sqrt{w_q}$,
so the same bound covers the fallback. Orthogonality gives
$\sum_qe_q^2=e^2$ and hence
\[
 \sum_qw_q D(P_{V_q\psi_T},P_{\phi_q})
       \le e+2\sum_q\sqrt{w_q}e_q\le3e.
\]
Zero weights contribute nothing; no minimum-population cutoff is used.

First replace the actual path by its Bell comparator while retaining
the independent tomography record and the same controlled gates.
This costs $\delta_{\mathrm{path},N}$. Bell equivariance gives the
terminal weights $w_q$, so the preceding weighted bound costs at most
$3\mathbb E e\le6\sqrt{d/(mv^2)}$. An operator implementation
error costs at most $u_{\mathrm{impl}}$ in state trace distance.
Independent tomography specimen kernels are identical on both sides
and preserve these bounds. Use \eqref{copy:joint} for the stronger
comparison retaining autonomous production pilots.
\end{proof}

With $m=N$, fixed finite $v>0$, $u_{\mathrm{impl}}\to0$ and
\eqref{copy:scaling}, all errors vanish. The tomography stock adds
$(d^2-1)N$ complete copies and a terminal error $O(N^{-1/2})$ at fixed
$d$. Common subsequent admitted operations contract the unconditioned
CQ bound. Rare conditional branches still require inverse-probability
control. The benchmark has a terminal daughter; it does not collapse
the target after each intermediate tagged jump. The programme before
repair is deterministic. A feedback-dependent Bell theorem would need
its own conditional-current and nodal arguments, not just preallocated
copy banks.

\section{Why the resource promise cannot be hidden}

For a Bell-pair target and a tag generated independently of that target
at fixed complete $\psi$, any tag-selected trace-preserving operation
on $S$ alone leaves the conditional $R$ marginal $I/2$. A desired
selected $P_q^S$ daughter instead has $R$ marginal
$|q\rangle\langle q|$, at trace distance $1/2$. Thus the global
gates above generally cannot be relabeled as local aperture controls.
Copies of a reduced $S$ state do not supply the promised copies of its
complete entangled bank. A further exterior reference entangled with
the allegedly pure complete bank is outside the preparation promise.

The resulting implication is consequently precise: supplied complete
pure copies and an admitted native diagnostic give integrated current
estimation; a record-driven conservative pump and scalar Markov pair
chemistry give the Hamiltonian Bell tagged-path limit; further finite
native tomography and global terminal control approximate its selected
daughter. Statistical native calibration, pair-reaction completeness,
initial tag/population preparation and the stated archive-return domain
remain commitments of this stronger-resource construction. It does not
meet the main one-unknown-input, inaccessible-reference objective by
itself. The later pilot completion meets that objective with a different
material inventory, without using the copies or global reference control
assumed here.

% End consolidated source: parts/10_copies.tex


% Begin consolidated source: parts/11_chambers.tex
\chapter{Chamber transport, killing, and complete path limits}
\label{chamber:chapter}

The geometric chamber model \cite{M24} and its continuous and absorbing extensions \cite{M26,M27} use a different complete source from canonical packet exchange. There is one material tracer, squared-norm chamber volume, deterministic rectified boundary transfer, and intrinsic uniform stirring. The Bell ratio is absent from the elementary stirring generator; it emerges in the natural path limit. Its squared-norm volume, normalized face allocation, absence of extra traffic, and fresh stirring statistics are nevertheless explicit constitutive inputs. This chapter records the common proof once and states precisely which absorbing and adaptive extensions it supports.

\section{Interaction selection within the finite inventory}

Fix a finite coherent space containing source, apparatus, controller, every returning memory, and resource/failure modes. An inaccessible reference remains inside each fibre and every controlled operator acts trivially on it. Let $P_i$ be complete configuration projectors and $w_i=\|P_i z\|^2$. The actual state is $(z,Q,d,\zeta)$, with $(d,\zeta)\in[0,w_Q]\times[0,1]$. Zero-volume exceptions are assigned a fixed failure flag.

The contact-selection result is conditional on the following independent material assumptions. At fixed configuration $i$ and normalized microscopic coordinate $\xi=(d/w_i,\zeta)$, candidate energy $h_i(z,\xi;c)$ is degree two in amplitudes, continuous in $\xi$ and controls $c$, and generates a metric- and symplectic-preserving field fixing zero. Elementary stirring and native configuration transfers keep $z$ instantaneously fixed and conserve the same energy, without an additional coordinate chosen to compensate a ray-dependent energy change. Full-support stirring is admitted. Arbitrarily weak simultaneous probes connect a spanning graph of the full active inventory on nonempty open sets of rays, and energies vary continuously when these probes tend to zero.

\begin{proposition}[A common contact matrix within the specified inventory]
\label{chamber:contact}
Under these assumptions, on each connected applicable inventory,
\[
 h_i(z,\xi;c)=z^\dagger H(c)z,\qquad H(c)=H(c)^\dagger,
\]
independently of the actual configuration and microscopic coordinate.
\end{proposition}
\begin{proof}
Isoenergetic stirring gives $h_i(z,\xi';c)=h_i(z,\xi;c)$ for almost every product-uniform target $\xi'$. Full support and continuity extend equality to the whole chamber, removing microscopic dependence. The finite-cone Killing-field argument of Theorem~\ref{mat:linear}, together with degree-two energy normalization, gives $h_i(z;c)=z^\dagger H_i(c)z$. On an active transfer $i\to j$, no-work conservation gives $z^\dagger(H_j-H_i)z=0$ on an open ray set. A Hermitian quadratic form vanishing on an open sphere patch vanishes everywhere by real analyticity; polarization then gives $H_j=H_i$. Connectivity propagates equality, and continuity in simultaneous probe strength gives it at the zero-probe limit as well.
\end{proof}

The simultaneous-contact assumption covers a source-off read: testing only edges active at that exact zero-current instant would not connect the labels. The absence of a compensating coordinate is equally substantive. Adding a real work coordinate $y$ with energy $z^\dagger H_i z+y$ and transfer
\[
 (i,y)\mapsto(j,y+z^\dagger(H_i-H_j)z)
\]
preserves energy for arbitrary different $H_i,H_j$. Such a recoil changes the declared complete inventory; ordinary energy conservation alone does not exclude it. The preceding proposition is therefore a reduction within a finite material class, not universal physical admission.

\section{Volumes, portals, and the primitive stochastic law}

Let a bounded piecewise regular Hermitian $H(t)$ drive
$i\hbar\dot z=H(t)z$. Define
\[
 J_{ji}=\frac2\hbar\operatorname{Im}\langle z_j,H_{ji}z_i\rangle,
 \quad f_{ji}=[J_{ji}]_+,\quad 
 F_i=\sum_jf_{ji},\quad A_i=\sum_jf_{ij}.
\]
Then $\dot w_i=A_i-F_i$. The total chamber volume is one. Initially use density one on the union; alternatively a ready source supported in one volume-one chamber needs only two independent uniform microscopic coordinates, not an initial draw among sectors.

Between stirring and crossings,
\begin{equation}
 \dot d=-F_Q(t),\qquad\dot\zeta=0.
 \label{chamber:drift}
\end{equation}
On the lower face $d=0$, allocate an interval of transverse width $f_{ji}/F_i$ to destination $j$. Allocate its matching upper incoming interval in $j$ width $f_{ji}/A_j$. If their lower endpoints are $a_{ji},c_{ji}$, set
\begin{equation}
 Q'=j,\quad d'=w_j(t),\quad
 \zeta'=c_{ji}+(F_i/A_j)(\zeta-a_{ji}).
 \label{chamber:portal}
\end{equation}
The map satisfies $F_i\,d\zeta=A_j\,d\zeta'$. Zero-current faces are absent. This normalized flux allocation is a geometric postulate, not a consequence of conservation alone. No extra circulation or independent side jumps are allowed in this model.

The only elementary random generator is state-independent Poisson stirring at rate $\kappa>0$:
\begin{equation}
 \mathcal S_\kappa f(i,d,\zeta)=\kappa\left[
 \int_0^1\!\int_0^1f(i,w_i u,v)du\,dv-f(i,d,\zeta)\right].
 \label{chamber:stirring}
\end{equation}
Each target uses fresh product uniforms conditional on the complete past; there is no readable future seed or free stirring archive. The Poisson count on a finite horizon is finite almost surely, with mean $\kappa T$ and unbounded support. A finite consumable stirring bath has not been constructed.

\begin{lemma}[Exact volume and mean flux at finite stirring]
\label{chamber:volume}
Density one is preserved. Consequently $\Pr(Q_t=i)=w_i(t)$ and
$\mathbb E N_{ji}([s,t])=\int_s^tf_{ji}(u)du$. The characteristic construction is nonexplosive on finite horizons with integrable total flux and has no positive-probability node loss for bounded piecewise regular finite Hamiltonians.
\end{lemma}
\begin{proof}
Interior drift has zero divergence. Relative upper-boundary speed is $\dot w_i+F_i=A_i$, exactly the incoming flux supplied by \eqref{chamber:portal}; lower-face outflow is $F_i$. Resetting a mass $w_i$ uniformly in volume $w_i$ preserves density one. Localize initially to positive volumes and bounded event counts. Characteristics, flux-preserving portals and reset averaging bound the substochastic density by one. Expected counts are therefore bounded by the integrated total current, ruling out explosion as the count cutoff is removed.

After arrival at $s$, before the next event,
\[
 w_i(t)-d(t)=\int_s^tA_i(u)du.
\]
Arrival times have an absolutely continuous density bounded by $A_i(s)ds$. At almost every such arrival the right accumulated incoming flux is positive on a sufficiently short subsequent interval, by Lebesgue differentiation. The trajectory must reach its lower exit before a positive gap can coexist with $w_i=0$. Initial and refreshed depths are strictly below the upper boundary almost surely. Thus exceptional collapsing upper-boundary trajectories have zero mass. Removing the volume cutoff loses no probability. The density-one solution is normalized and its directed boundary flux is $f_{ji}$, proving both formulas.
\end{proof}

These statements do not yet determine conditional waiting times. For example under $H=\hbar g\sigma_x$, starting in sector 0, the first exit has density $g\sin(2gs)$ on $(0,\pi/(2g))$. If no stirring occurs for the next $\pi/(2g)$, the return time is exactly $s+\pi/(2g)$. This singular line in the two-event law has mass at least $e^{-\kappa\pi/(2g)}$, whereas the limiting Bell two-time law has a density. The finite chamber process is therefore genuinely different from the Bell process.

\section{The global finite-horizon path theorem}

The absorbing extension includes the unabsorbed case by setting all killing rates to zero. During one no-capture epoch let
\begin{equation}
 \dot z=-iH(t)z/\hbar-\tfrac12\Gamma(t)z,
 \qquad\Gamma=\sum_i\gamma_iP_i,
 \qquad0\le\gamma_i=\sum_c\gamma_{ci}\le\bar\gamma,
 \label{chamber:killedwave}
\end{equation}
with $\|z(0)\|=1$. Coefficients are deterministic after conditioning on the complete chemical history at the epoch start. The volumes are now unnormalized, $w_i=\|z_i\|^2$, and
\begin{equation}
 \dot w_i=A_i-F_i-\gamma_iw_i,
 \qquad\dot d=-F_i-\gamma_i d.
 \label{chamber:killeddrift}
\end{equation}
The actual chemical hazard at occupied site $i$ is $\gamma_{ci}$, independently of depth. The same portals and stirring act until capture.

The killed density equation is
\[
 \partial_t\rho_i+\partial_d[(-F_i-\gamma_i d)\rho_i]=-\gamma_i\rho_i.
\]
Thus surviving density one solves the interior equation, while relative upper-boundary speed remains $A_i$. The gap satisfies
$(w_i-d)'=A_i-\gamma_i(w_i-d)$; its positive integrating-factor form gives the preceding node argument. The reaction density is $\gamma_{ci}w_i\,dt$. Classical killing does not determine a coherent daughter; the reaction continuation below is a separate law.

Fix $n$ labels and horizon $T$. Set
\begin{align}
 M&=\sup_{i,t}F_i(t),\qquad L=\sup_{i,t}|\dot w_i(t)|,\\
 b(a)&=2na+\int_0^T\sum_{w_i\le2a}A_i\,dt
        +2\int_0^T\sum_{a\le w_i\le2a}[-\dot w_i]_+dt,
 \quad0<a<1/2.
 \label{chamber:nodebudget}
\end{align}
The limiting comparator has native rates $f_{ji}/w_i$ on occupied chambers and chemical rates $\gamma_{ci}$. Its output contains the complete native path, capture time/channel, and null or terminal failure symbols.

\begin{theorem}[Killed whole-path comparison, including nodes]
\label{chamber:path}
For $M>0$ the total-variation distance between the finite-stirring killed path law and its comparator is at most the following expression, capped at one:
\begin{equation}
 \epsilon_\kappa(a)=2e^{\bar\gamma T}b(a)
 +\frac{M(L+2M+e\bar\gamma)T}{\kappa a^2}
 +(\kappa T+1)\exp\left[-\frac{\kappa a}{4(M+\bar\gamma)}\right].
 \label{chamber:pathbound}
\end{equation}
If $M=0$, the native path is constant and the two chemical laws agree exactly. With $\bar\gamma=0$ the unabsorbed proof sharpens the final exponential denominator from $4M$ to $2M$ and removes the $e\bar\gamma$ term.
\end{theorem}
\begin{proof}
First suppress killing while retaining the compressed drift \eqref{chamber:killeddrift}. This auxiliary skeleton is only a proof device. Its density is at most $e^{\bar\gamma t}$: compression has rate at most $\bar\gamma$, portals preserve flux, and a uniform reset preserves the bound. For the comparator without killing, occupations $p_i$ obey $p_i\le e^{\bar\gamma t}w_i$. Indeed $v_i=e^{\bar\gamma t}w_i$ solves its positive forward equation with nonnegative added source $(\bar\gamma-\gamma_i)v_i$. A stopped finite-state comparison followed by removal of cutoffs proves this domination even near nodes. Both skeletons' incoming fluxes are therefore bounded by $e^{\bar\gamma T}f_{ji}$.

An occupied visit to $w_i\le a$ either starts below $2a$, arrives while $w_i\le2a$, or stays in $i$ during a downward excursion from $2a$ to $a$. Initial low-volume mass is at most $2a$ per label. The incoming visits cost their flux integral. Each completed downward excursion consumes at least $a$ of negative variation in the band; occupation at its deterministic start is at most $2ae^{\bar\gamma T}$. The union bound gives $e^{\bar\gamma T}b(a)$ for each skeleton. Stop both on such visits.

Condition on a common independent Poisson stirring-time partition. Each interval starts at $s$ with uniform depth and transverse coordinate in its current chamber. Put $G_i(s,t)=\int_s^t\gamma_i(u)du$. Until its first native exit the exact depth is
\[
 d(t)=e^{-G_i(s,t)}\left[d(s)-\int_s^te^{G_i(s,u)}F_i(u)du\right].
\]
The first-exit subdensities into $j$, before the occupied-volume stop, are
\[
 g_j(t)=e^{G_i(s,t)}f_{ji}(t)/w_i(s),
 \qquad
 g_j^*(t)=e^{-\int_s^tF_i/w_i}\,f_{ji}(t)/w_i(t).
\]
For interval length $h\le a/[4(M+\bar\gamma)]$, use
\[
 |e^{G_i}-1|\le e\bar\gamma h,\quad 
 |w_i(s)^{-1}-w_i(t)^{-1}|\le Lh/a^2,\quad 
 1-e^{-\int F_i/w_i}\le Mh/a.
\]
After summing destinations and integrating, the first-exit density difference is at most $M(L+M+e\bar\gamma a)h^2/(2a^2)$. Adding the no-exit atom costs no more than this quantity in total variation, since its discrepancy is the integral of the density difference.

After a microscopic arrival its depth is at least $ae^{-\bar\gamma h}-Mh>0$, so there cannot be a second microscopic exit within this short interval before the volume stop. The comparator's two-or-more probability is bounded by $M^2h^2/(2a^2)$, using a dominating rate-$M/a$ Poisson clock. Hence the complete interval-kernel discrepancy is at most $M(L+2M+e\bar\gamma)h^2/(2a^2)$.

For the Poisson partition, $\mathbb E\sum h^2\le2T/\kappa$. One proof integrates $e^{-\kappa|u-v|}$, the probability that $u,v$ lie in the same interval. A union bound over the initial gap and all Poisson starts bounds a gap longer than $h_0$ by $(\kappa T+1)e^{-\kappa h_0}$. Sequential conditional coupling, with the two node budgets, yields \eqref{chamber:pathbound} for the auxiliary skeletons.

Finally append the same independent unit exponential $E$ and stop each path when $\int_0^t\gamma_{Q_s}(s)ds$ reaches $E$. Select $c$ using the common probabilities $\gamma_{cQ}/\gamma_Q$. This common marking/stopping kernel contracts total variation and produces the actual killed laws. For $M=0$ native labels never change and the same construction couples the whole law exactly. When $\bar\gamma=0$, a fresh depth at least $a$ cannot be swept twice if $h<a/(2M)$; the sharper unabsorbed estimate follows by the same proof.
\end{proof}

If $\sup_t\|H(t)\|\le\hbar\Omega$, let $C=2n\Omega$. Then $|J_{ji}|\le2\Omega\sqrt{w_iw_j}$ and $\sum_iw_i\le1$ imply
\begin{equation}
 M\le C,\quad L\le C+\bar\gamma,\quad
 b(a)\le2na+3nCT\sqrt{2a}+4n\bar\gamma aT.
 \label{chamber:explicitnodes}
\end{equation}
The incoming and outgoing flux at a label are bounded by $C\sqrt{w_i}$; insert that bound into \eqref{chamber:nodebudget}, using $[-\dot w_i]_+\le F_i+\gamma_iw_i$. These finite integrals justify cutoff removal. At fixed graph, horizon and reaction bounds, $a=\kappa^{-2/5}$ in fixed units gives a coarse $O(\kappa^{-1/5})$ path error. Growing resources require the actual right-hand side, including $e^{\bar\gamma T}$, to vanish in the same limit.

\section{Conditional capture and finite adaptive epochs}

For a product mark $c$ that does not reveal the native site, define
\[
 D_c=\sum_i\gamma_{ci}P_i,
 \qquad r_c=\|\sqrt{D_c}z\|^2.
\]
The declared coherent reaction is $z\mapsto\sqrt{D_c}z/\sqrt{r_c}$, followed by its specified fibre-preserving outgoing isometry. This assumption fixes more than the scalar hazard. Density-one capture makes the old-site posterior $\gamma_{ci}w_i/r_c$ and leaves depth uniform in each site. The map
\begin{equation}
 i'=f_c(i),\qquad d'=\gamma_{ci}d/r_c,\qquad\zeta'=\zeta
 \label{chamber:capturereset}
\end{equation}
therefore produces density one in the normalized daughter chambers, when $f_c$ is the declared injective site embedding. If channels merge sites, the outgoing multiplicity must remain in the retained configuration or the construction requires an explicitly specified measure-preserving merge. Zero-rate branches are absent.

This readiness is conditional on the complete chemical history, not necessarily on the earlier native path. To concatenate at most $m$ adaptive epochs, require either an explicitly fresh carrier conditional on the complete past or a quiet interval $s_\ell>0$ before the next active epoch, with $H=0$, all killing off, and stirring retained. Conditional on at least one stir in that interval, the final normalized depth and transverse coordinate are independent fresh uniforms given the earlier macroscopic path and its final site. Failure costs $e^{-\kappa s_\ell}$.

\begin{corollary}[Complete finite adaptive path budget]
\label{chamber:adaptive}
Suppose the coherent daughter and resource maps are as just specified, the coefficients of each epoch are deterministic conditional on its complete chemical history, all adaptive controls use retained physical records, and the epoch bounds are uniform on the admitted complete histories. Then
\begin{equation}
 \operatorname{TV}(P_\kappa,P_*)\le
 \min\left\{1,\sum_{\ell=1}^m
       \sup_{\text{chemical histories}}\epsilon_{\kappa,\ell}(a_\ell)
       +\sum_{\text{required refreshes}}e^{-\kappa s_\ell}\right\}.
 \label{chamber:adaptivebound}
\end{equation}
The common output includes all native paths, event times, product marks, nulls, failures, exhaustion, copied memories and permitted resource returns represented by the specified subsequent maps.
\end{corollary}
\begin{proof}
On matching chemical events, exact event times and prior retained outputs, the normalized source and next bounded programme coincide. Successful quiet refresh supplies identical independent microscopic readiness in the matched occupied site. The interval comparison in Theorem~\ref{chamber:path} is uniform conditional on that site. Node costs need not be uniform conditional on the whole prior native path: average them under each model's actual epoch law conditional on chemical history, where density one gives the stated site distribution, and bound the failure events by these marginal probabilities. Sequential maximal couplings of the short-interval kernels, the union bound for the two node events in each epoch, and the refresh-failure union bound prove \eqref{chamber:adaptivebound}. A common subsequent record map contracts the bound.
\end{proof}

An arbitrary event-dependent change of wave generator does not satisfy these hypotheses. Start in sector 0 under $H=\hbar g\sigma_x$ and freeze $H$ at the first actual native exit $S$. Its one-way density is $g\sin(2gS)$ on $(0,\pi/(2g))$. After freezing, $Q=1$ certainly, but the wave stays at $w_1=\sin^2(gS)$ and
\[
 \mathbb E w_1=\int_0^{\pi/(2g)}\sin^2(gs)g\sin(2gs)ds=1/2.
\]
Thus $\Pr(Q=1)=1\ne\mathbb Ew_1$. The reaction continuation and conditional readiness assumptions cannot be omitted merely because the policy is causal.

\section{Continuous coordinates and the finite split extension}

Let $y$ be a material work coordinate, with local weights $w_i(y,t)$ and currents obeying
\[
 \partial_tw_i+\partial_y(w_iv_i)=A_i-F_i-\gamma_iw_i.
\]
The density-one microscopic region is $0<d<w_i(y,t)$, $0<\zeta<1$, with dynamics
\begin{equation}
 \dot y=v_i,\qquad
 \dot d=-F_i-d\partial_yv_i-\gamma_i d,
 \qquad\dot\zeta=0.
 \label{chamber:continuum}
\end{equation}
Interior divergence is $-\gamma_i$, exactly balanced by killing. The moving upper boundary has relative incoming speed $A_i$. Thus the same portals preserve the correct flux, wherever the stated flow is defined. A global completion is available for the following finite split class; it is not inferred for arbitrary simultaneous singular transport and reactions.

During a no-event translation block, $f=\gamma=0$ and sector mass $m_i=\int w_i(y,t)dy$ is fixed. Set
\[
 F_{i,t}(y)=m_i^{-1}\int_{-\infty}^yw_i(x,t)dx,
 \quad y(t)=F_{i,t}^{-1}(F_{i,s}(y(s))).
\]
Use generalized inverses and retain $d/w_i,\zeta$ between stirs. The change of $w_i(y,t)dy$ cancels the depth scaling, preserving complete volume. If $|v_i|\le V(t)$, integrating the continuity equation against monotone cutoffs gives $|y(t)-y(s)|\le\int_s^tV(u)du$ for almost every rank. Vacuum gaps do not require division by zero.

During a gate block take $v_i=0$ and first omit killing. Then $\rho(y)=\sum_iw_i(y,t)$ is fixed. Conditional on $y$ with $\rho(y)>0$, define $\widehat w_i=w_i/\rho$, $\widehat f_{ji}=f_{ji}/\rho$ and normalized depth $d/\rho$. The unabsorbed theorem applies on this fibre. Its error is
\begin{equation}
 2\int\rho(y)b_y(a)dy+
 \frac{M(L+2M)T}{\kappa a^2}
 +(\kappa T+1)e^{-\kappa a/(2M)},
 \label{chamber:fibrebound}
\end{equation}
where $M,L$ are essential uniform bounds for the normalized fibre rates and weight derivatives, and $b_y$ is \eqref{chamber:nodebudget} on that fibre. Bounded local gate matrices make these estimates uniform; zero-$\rho$ fibres have no probability. A chemical fixed-$y$ block instead uses Theorem~\ref{chamber:path} conditioned on its complete initial fibre law and the corresponding uniform bounds.

Alternating finitely many exact translations with such gates or chemical blocks, and supplying quiet refreshes wherever conditional depth is not fresh, yields the sum of the local bounds and $e^{-\kappa s_\ell}$ failures. The proof couples matching labels and exact event times at each gate; intervening identical quantile maps append identical continuous work paths. The common path space is
\[
 D([0,T_h],\mathcal I\cup\{\dagger\})\times C([0,T_h],\mathbb R)
\]
with finite retained record labels incorporated in $\mathcal I$. Resource parameters and the unknown complete input are fixed for each law; mixing admitted preparations averages their error bounds. A physically recorded derivative of $y$ requires another admitted contact and is not supplied by placing the mathematical path in this comparison space.

\section{Filtration, resources, and target status}

The limiting intensities are conditional on the complete input parameters, the natural native/chemical path history, and $y$ when that coordinate is included. Omitting unresolved preparation parameters or an unrecorded $y$ requires posterior averaging of those rates. In the microscopic filtration revealing $d,\zeta$ and all stirring draws, the next boundary crossing is predictable between stirs. The Bell compensator is a natural-filtration limit, not a formula valid after arbitrary microscopic archive exposure.

On isolated $\Gamma=0$ epochs the limit uses the specified Hamiltonian current in physical time. During a non-Hermitian no-capture epoch or a changed acquisition Hamiltonian, the wave and its current differ from the isolated experiment. The theorem does not equate those paths. Common measurable stopping and output maps contract path total variation. Conditioning on an event of probability at least $q>\epsilon$ costs at most $2\epsilon/q$; no uniform rare-event claim follows if $q$ vanishes.

The technical result derives the complete path law from finite geometric residence and intrinsic stirring; it is stronger than mean-flux matching. Norm volumes, normalized rectified faces, no extra traffic, intrinsic Poisson stirring and fresh conditional microscopic readiness remain premises of this chamber model. Capture continuation and adaptive renewal require their additional assumptions. The contact-selection theorem restricts its specified material inventory; it does not prove that every substance belongs to that inventory. The later finite-gas pilot construction is a different microscopic realization: it neither derives this chamber geometry nor converts the chamber's primitive stirring into a theorem retrospectively. Both must be distinguished from a universal derivation using only the older source principles.

% End consolidated source: parts/11_chambers.tex


% Begin consolidated source: parts/12_benchmarks.tex
\chapter{Additional finite-model benchmarks and rejected shortcuts}
\label{bench:chapter}

The main chapters prove the general instrument, historical-record and
traffic statements. The following calculations retain distinct physical
examples from the corrected checkpoints \cite{C01,C02,C03}. Each example
states its own dynamics and statistical input. They are regression tests
for continuation and acquisition claims, rather than an additional
selection of the event law. The Gaussian pre-latch contact is treated
alongside its finite receptor in Chapter~\ref{det:chapter-response}.


% Begin consolidated source: parts/12_history_benchmark.tex
\section{A generated first record and a finite-delay history separator}
\label{bench:history}

This example makes the readable-history comparison of \cite{C01}
explicit. Its source constitution is the minimal Bell process
\eqref{found:bell} for the declared projectors
$P_n=|n\rangle\langle n|\otimes I_R$. The inaccessible reference is
retained coherently inside each sector. Resolving reference states as
additional actual configuration labels would specify a different
process; coarse-graining such a process need not recover these minimal
sector rates.

The apparatus is a stipulated \emph{neutral actual-event marker}:
a monitored native jump produces a daughter without changing the source
wave or its prescribed rates. Downstream capture writes a physical
memory. This is a mathematical coupling assumption of the kind used in
Theorem~\ref{acc:reporter} and Theorem~\ref{det:delayed-filter}.
Those results do not identify it with a universally admitted quantum
instrument or a coherent source--product interaction.

\subsection{One-way transport and a finite response}

Fix an acquired earlier history $H$ and the subsequent control schedule.
Suppose the only active source edge on $[0,\tau]$ is $0\to1$, with
$w_0(u)>0$, $J_{10}(u)=-\dot w_0(u)\ge0$ and other sectors inert.
Put $p=\Pr(Q_0=0\mid H)$. There is at most one subsequent native jump.
A fresh detector supplies one daughter-production cofactor, one ready
site, finite capture fuel and a blank persistent memory. The native
jump consumes the cofactor and produces a daughter $X$.
Independently of subsequent source evolution,
\[
 X+A_{\rm ready}+M_{\rm blank}
 \xrightarrow{\ \beta_R\ }D_R+A_{\rm spent}+M_R,\qquad
 X\xrightarrow{\ \beta_M\ }D_L.
\]
The first arrow includes the supplied capture fuel consumption.
The cofactor, fuel, site, pending daughter or reaction remnant, and
memory remain in the extended state. Loss writes no record and
restores no cofactor. An exhausted production channel leaves the
source jump available without another daughter, preserving the
stipulated source generator. No exhausted repeat is needed here.
For $\beta_R>0$, $\beta_M\ge0$, $k=\beta_R+\beta_M$, the response is
\begin{equation}
 F_{\rm det}(v)=\frac{\beta_R}{k}(1-e^{-kv}).
 \label{bench:history-response}
\end{equation}
The mean time to capture or loss is $1/k$, as is the mean delay
conditional on capture. This model has no additional propagation lag.

\begin{proposition}[A one-way finite-delay separator]
\label{bench:history-separator}
Assume no competing old daughters. The actual source law conditioned
only on $H$, and its emission density, are
\begin{equation}
 \nu_u(0\mid H)=\frac{p}{w_0(0)}w_0(u),\qquad
 f_{\rm emit}(u\mid H)=\frac{p}{w_0(0)}J_{10}(u).
 \label{bench:history-transport}
\end{equation}
Thus the acquired-record probability and its difference from the
equilibrium benchmark for the unchanged wave are
\begin{align}
 P_H(\tau)&=\frac{p}{w_0(0)}
       \int_0^\tau J_{10}(u)F_{\rm det}(\tau-u)\,du,
       \label{bench:history-record}\\
 P_H(\tau)-P_{\rm eq}(\tau)&=
 \left(\frac{p}{w_0(0)}-1\right)
 \int_0^\tau J_{10}(u)F_{\rm det}(\tau-u)\,du.
 \label{bench:history-gap}
\end{align}
The difference is nonzero if the prefactor is nonzero and positive
current overlaps positive response on a set of positive measure.
\end{proposition}
\begin{proof}
Integrating the occupied-source hazard
$\lambda_{1\leftarrow0}=-\dot w_0/w_0$ gives survival
$w_0(u)/w_0(0)$. This proves \eqref{bench:history-transport}.
Condition on the unique possible emission time and multiply its
density by the response at its remaining age. For the equilibrium
benchmark replace $p$ by $w_0(0)$ and subtract.
\end{proof}

Here $\nu_u$ averages over later detector outcomes. It is not the
posterior additionally conditioned on every subsequent observed null.
For \eqref{bench:history-response}, the capture density is
\[
 f_R(t\mid H)=\int_0^t f_{\rm emit}(u\mid H)
                          \beta_Re^{-k(t-u)}\,du.
\]
The hazard conditioned on no new record is $f_R(t\mid H)/(1-P_H(t))$;
its denominator includes no emission, pending daughter and loss.
Multiple daughters or shared sites require the full retained-state
filter of Theorem~\ref{det:delayed-filter}. A general signed integrand
can cancel, and a propagation delay longer than the observation
window can eliminate the overlap required for strict positivity.

A useful source-transport check takes
$\Psi_0=\cos\phi|0\rangle+\sin\phi|1\rangle$ and
$H_c=\hbar\omega\sigma_y$. Then
\[
 \Psi_u=\cos(\phi+\omega u)|0\rangle+\sin(\phi+\omega u)|1\rangle,
 \qquad\lambda_{1\leftarrow0}(u)=2\omega\tan(\phi+\omega u).
\]
For $\phi=\pi/4$, $p=1$, $\omega u_*=\pi/12$, the wave weight is
$1/4$ while the actual probability is $1/2$. During a further
$0<\tau<\pi/(6\omega)$, the departure probability from occupied
sector 0 is $A_\tau=1-4\cos^2(\pi/3+\omega\tau)$.
The two ensemble emission probabilities are $A_\tau/2$ and
$A_\tau/4$. Squared unitary entries are not the Bell transition
kernel; $1-e^{-2\sqrt3\omega\tau}$ freezes the hazard and is only a
short-window approximation. Starting instead at $\phi=\pi/4$, the
finite response gives, for $0<\tau<\pi/(4\omega)$,
\[
 P_H(\tau)-P_{\rm eq}(\tau)\ge
 \omega\cos(2\omega\tau)\frac{\beta_R}{k}
 \left[\tau-\frac{1-e^{-k\tau}}k\right]>0,
\]
because $J_{10}(u)=\omega\cos(2\omega u)$. Its short-window value is
$\omega\beta_R\tau^2/2+O(\tau^3)$.

\subsection{Generating the first record}

Prepare
\begin{equation}
 |\Psi(0)\rangle=\sqrt a\,|r,R_r\rangle+\sqrt b\,|1,R_1\rangle,
 \qquad a=\frac23,\quad b=\frac13,\quad
 \langle R_r|R_1\rangle=0,
 \label{bench:history-initial}
\end{equation}
with normalized references and initial actual equilibrium.
Write $\mathcal R_1$ for the first acquired record event, distinct
from the reference vector $R_1$. Apply
$H_1=\hbar\Omega(|0\rangle\langle r|+|r\rangle\langle0|)\otimes I_R$
until $T_1=\pi/(3\Omega)$. The complete source wave is
\begin{equation}
 |\Psi(t)\rangle=
 \sqrt a\cos(\Omega t)|r,R_r\rangle
 -i\sqrt a\sin(\Omega t)|0,R_r\rangle
 +\sqrt b|1,R_1\rangle.
 \label{bench:history-first-wave}
\end{equation}
Consequently $J^{(1)}_{0r}=a\Omega\sin(2\Omega t)$ and
$\lambda_{0\leftarrow r}=2\Omega\tan(\Omega t)$.
Sector 0 is absorbing throughout this monotone first interval.
A type-1 detector initially has no daughter and has the response
$F_1(v)=\beta_{R,1}(1-e^{-k_1v})/k_1$,
where $k_1=\beta_{R,1}+\beta_{M,1}$. Therefore
\begin{equation}
 \mathcal K=\int_0^{T_1}\Omega\sin(2\Omega s)F_1(T_1-s)\,ds,
 \qquad
 \Pr(\mathcal R_1)=a\mathcal K,\qquad
 \Pr(Q_{T_1}=0\mid\mathcal R_1)=1.
 \label{bench:history-first-record}
\end{equation}
The event $\mathcal R_1$ means acquisition by the fixed time $T_1$,
not a later capture of an old daughter.

For the null define
\[
 E_1=\int_0^{T_1}\Omega\sin(2\Omega s)e^{-k_1(T_1-s)}\,ds,\qquad
 L_1=\frac{\beta_{M,1}}{k_1}
                 \{\sin^2(\Omega T_1)-E_1\}.
\]
The first-stage endpoint alternatives have unnormalized weights
\begin{center}
\begin{tabular}{@{}lll@{}}
\toprule
Actual source & Retained alternative & Probability\\
\midrule
$r$ & No emission; blank memory & $1/6$\\
$1$ & No eligible emission; blank memory & $1/3$\\
$0$ & Pending daughter; blank memory & $aE_1$\\
$0$ & Lost daughter; blank memory & $aL_1$\\
$0$ & Captured daughter; acquired memory & $a\mathcal K$\\
\bottomrule
\end{tabular}
\end{center}
They sum to one since $E_1+L_1+\mathcal K=\sin^2(\Omega T_1)$.
The first null comprises the first four rows, normalized by
$1-a\mathcal K$. Every row retains its apparatus resources and the
same stipulated source wave. Pending daughters can still acquire
late first records; the acquired memory remains a physical register.
In particular,
\begin{equation}
 |\Psi(T_1)\rangle=|r\rangle v_r+|0\rangle v_0+|1\rangle v_1,
 \quad v_r=\frac{R_r}{\sqrt6},\quad
 v_0=-\frac{iR_r}{\sqrt2},\quad v_1=\frac{R_1}{\sqrt3}.
 \label{bench:history-endpoint}
\end{equation}
The wave weights are $(W_r,W_0,W_1)=(1/6,1/2,1/3)$ even on
$\mathcal R_1$, whose actual law is $\delta_0$. Neither the residual
$r$ amplitude nor the relative $-i$ phase has been removed.

\subsection{Three continuations, with finite positive numbers}

At $T_1$ switch to $H_2=\hbar\omega\sigma_y\otimes I_R$ on 0,1,
leaving $r$ inert. With $u=t-T_1$,
\[
 v_r(u)=v_r,\qquad
 v_0(u)=\cos(\omega u)v_0-\sin(\omega u)v_1,\qquad
 v_1(u)=\sin(\omega u)v_0+\cos(\omega u)v_1.
\]
For general reference-valued amplitudes, put
$c=\operatorname{Re}\langle v_0,v_1\rangle$. Direct differentiation
gives
\[
 W_0(u)=W_0(0)\cos^2(\omega u)+W_1(0)\sin^2(\omega u)
                                      -c\sin(2\omega u),
\]
\[
 J^{(2)}_{10}(u)=\omega(W_0(0)-W_1(0))\sin(2\omega u)
                                      +2\omega c\cos(2\omega u).
\]
For \eqref{bench:history-endpoint}, $c=0$, so
\begin{equation}
 W_0(u)=\frac12\cos^2(\omega u)+\frac13\sin^2(\omega u),
 \qquad J^{(2)}_{10}(u)=\frac{\omega}{6}\sin(2\omega u).
 \label{bench:history-second-current}
\end{equation}
This current includes the inaccessible reference. Use
$0<\tau<\pi/(2\omega)$, so the interval is one-way with finite rates.
A fresh type-2 detector with response $F_2$ records only $0\to1$.
Its separate site cannot capture type-1 daughters. The first detector,
memory and any old daughter persist. The same fixed source schedule
therefore also defines first-null continuations, including late
type-1 captures, without competition for the second site. At most
two monitored native emissions occur in this programme; two finite
production cofactors and two finite capture sites suffice.

Define
\begin{equation}
 I_2=\int_0^\tau\omega\sin(2\omega u)F_2(\tau-u)\,du.
 \label{bench:history-I2}
\end{equation}
On $\mathcal R_1$, \eqref{bench:history-transport} gives
$\nu_u(0\mid\mathcal R_1)=2W_0(u)$ and hence
\begin{equation}
 \Pr_{\rm marker}(\mathcal R_2\mid\mathcal R_1)=I_2/3,\qquad
 \Pr_{\rm marker}(\mathcal R_1,\mathcal R_2)=a\mathcal K I_2/3.
 \label{bench:history-joint}
\end{equation}
Here $\mathcal R_2$ means capture by $T_1+\tau$. On this branch the
no-second-emission weight is $1-\sin^2(\omega\tau)/3$. Among the
emission branches, the pending weight is
$\frac13\int_0^\tau\omega\sin(2\omega u)e^{-k_2(\tau-u)}\,du$;
the lost and captured weights use the responses
$\beta_{M,2}(1-e^{-k_2v})/k_2$ and $F_2(v)$, respectively.
The four alternatives sum to one. Thus the second observed null is
$1-I_2/3$, including pending and lost daughters.

\begin{center}
\begin{tabularx}{\textwidth}{@{}Xcc@{}}
\toprule
Continuation after the first record
& $\Pr(\mathcal R_2\mid\mathcal R_1)$ & Joint probability\\
\midrule
Neutral marker: full wave \eqref{bench:history-endpoint},
actual law $\delta_0$
& $I_2/3$ & $a\mathcal K I_2/3$\\
Equilibrium replacement for the unchanged full wave:
actual law $(1/6,1/2,1/3)$
& $I_2/6$ & $a\mathcal K I_2/6$\\
New source preparation in $|0\rangle$:
wave weight and actual probability one in sector 0
& $I_2$ & $a\mathcal K I_2$\\
\bottomrule
\end{tabularx}
\end{center}
The last row instead has $w_0(u)=\cos^2(\omega u)$ and
$J_{10}(u)=\omega\sin(2\omega u)$. The common joint factor
$a\mathcal K$ stipulates the same first stage followed by each
replacement on every first-record branch. It does not assert that
either replacement has been implemented. The equilibrium replacement
is not an ordinary-quantum prediction for the original first
apparatus: that claim requires a physical instrument with its actual
conditional states. These are three different continuations.

For generic frequency $g$ and rates $\beta_R,k$, elementary
exponential-trigonometric integration gives
\begin{align}
 E(g,k,t)&=\int_0^t g\sin(2gu)e^{-k(t-u)}\,du
 =\frac{g\{k\sin(2gt)-2g\cos(2gt)+2ge^{-kt}\}}{k^2+4g^2},
 \nonumber\\
 C(g,k,\beta_R,t)&=\frac{\beta_R}{k}
                       \{\sin^2(gt)-E(g,k,t)\}.
 \label{bench:history-closed-integral}
\end{align}
The second expression is the same integral with the capture response
in place of the exponential. In one chosen time unit set
$\Omega=\omega=\beta_{R,1}=\beta_{M,1}=\beta_{R,2}=\beta_{M,2}=1$,
$T_1=\pi/3$ and $\tau=\pi/4$. Both detectors have positive loss,
eventual capture probability $1/2$ and mean resolution delay $1/2$.
Then
\[
 \begin{aligned}
 \mathcal K&=\frac{5-\sqrt3-2e^{-2\pi/3}}{16}
            \simeq0.188853736,\\
 I_2&=\frac{1-e^{-\pi/2}}8\simeq0.099015053,
 \qquad a\mathcal K\simeq0.125902490.
 \end{aligned}
\]
\begin{center}
\begin{tabular}{@{}lrr@{}}
\toprule
Continuation & Second record, conditional & Both records\\
\midrule
Neutral marker & $0.033005018$ & $0.004155414$\\
Equilibrium replacement & $0.016502509$ & $0.002077707$\\
New $|0\rangle$ preparation & $0.099015053$ & $0.012466242$\\
\bottomrule
\end{tabular}
\end{center}
The event that both records are acquired within their windows
separates the first two programmes by
$a\mathcal K I_2/6\simeq0.002077707$ with finite positive delay and
loss. For short second windows,
$I_2=\beta_{R,2}\omega^2\tau^3/3+O(\tau^4)$, so their conditional
difference starts as $\beta_{R,2}\omega^2\tau^3/18+O(\tau^4)$.

\subsection{The rejected endpoint identification}

The proposed operators
\[
 M_R=\sqrt{\mathcal K}|0\rangle\langle r|,\qquad
 M_N=\sqrt{1-\mathcal K}|r\rangle\langle r|
                              +|1\rangle\langle1|
\]
form an abstract instrument on
$\operatorname{span}\{|r\rangle,|1\rangle\}$ and reproduce
$a\mathcal K$. They do not reproduce the neutral marker's
continuation. Disabling acquisition makes $\mathcal K=0$ and
$M_N=I$ on the input subspace, although the source still undergoes
\eqref{bench:history-first-wave}. At finite inefficiency, the actual
null additionally retains pending and lost daughters. Removing those
states requires a physical recovery of the source and all
information-bearing apparatus; agreement of one endpoint effect
does not supply it. The preparation in Proposition~\ref{prep:swap} is a
different operation, with its completion time and retained old-state
ancilla included before using the new source as the third benchmark.
Later preparation cannot change earlier acquired records.

% End consolidated source: parts/12_history_benchmark.tex


% Begin consolidated source: parts/13_reservoir_benchmark.tex
\section{Explicit record and loss products: the memory kernel}
\label{bench:reservoir}

The product-field construction of \cite[\S6]{C02} supplies a finite or
spectral Hamiltonian behind the distinction between pending excitation,
record products and hidden loss. It is a different realization from the
directed contact in \eqref{mat:chiralH}. The calculation below retains the
products; eliminating their amplitudes is an algebraic reduction, not a
physical deletion or an actualization rule.

\paragraph{A complete Hamiltonian on the admitted sector.}
Fix a real coupling $g\ge0$ and orthogonal states
\[
 |0\rangle=|r,A_0,\mathrm{vac},M_0\rangle,\qquad
 |1\rangle=|q,A^*,\mathrm{vac},M_0\rangle,
 \qquad |a,j\rangle=|q,A_0,a_j,M_0\rangle,
 \quad a\in\{R,L\}.
\]
Here $M_0$ is an unchanged blank memory. The $R$ and $L$ products occupy
orthogonal field sectors, both orthogonal to the vacuum. For finitely
many modes, the entire Hamiltonian on their span is
\begin{equation}
 \frac{H_N}{\hbar}
 =g(|1\rangle\langle0|+|0\rangle\langle1|)
 +\sum_{a,j}\omega_{aj}|a,j\rangle\langle a,j|
 +\sum_{a,j}\bigl(\kappa_{aj}|a,j\rangle\langle1|
                  +\overline{\kappa}_{aj}|1\rangle\langle a,j|\bigr).
 \label{bench:reservoir-finiteH}
\end{equation}
The real $\omega_{aj}$ are detunings in a rotating frame; ready and
excited energies have been set to zero. There are no further interactions
in this model. This subspace is invariant; an unused orthogonal complement
may be given any specified decoupled self-adjoint Hamiltonian.

For a continuum replace each mode space by $L^2(I_a,d\omega)$, with
$I_a\subseteq\mathbb R$, and assume $\kappa_a\in L^2(I_a)$.
The complete Hilbert space and Hamiltonian are then
\begin{align}
 \mathcal H&=\mathbb C^2\oplus L^2(I_R)\oplus L^2(I_L),\notag\\
 \frac{H}{\hbar}
 &=g(|1\rangle\langle0|+|0\rangle\langle1|)
   +\sum_{a=R,L}\int_{I_a}\omega|a,\omega\rangle
                          \langle a,\omega|\,d\omega\notag\\
 &\quad+\sum_{a=R,L}\int_{I_a}
  \bigl(\kappa_a(\omega)|a,\omega\rangle\langle1|
       +\overline{\kappa_a(\omega)}|1\rangle\langle a,\omega|\bigr)
       \,d\omega .
 \label{bench:reservoir-continuumH}
\end{align}
The multiplication operator by $\omega$ has its usual domain
$\{f:\omega f\in L^2\}$. The displayed coupling is a bounded finite-rank
perturbation, so this specifies a self-adjoint Hamiltonian and unitary
evolution. The continuum kets denote the corresponding spectral
representation, not normalizable additional vectors.

Start with $|0\rangle$ and empty product sectors. In the finite model write
\begin{equation}
 |\Psi(t)\rangle=x(t)|0\rangle+y(t)|1\rangle
                  +\sum_{a,j}z_{aj}(t)|a,j\rangle,
 \quad x(0)=1,\quad y(0)=z_{aj}(0)=0.
 \label{bench:reservoir-wave}
\end{equation}
The continuum expression replaces the sums by integrals. Schr\"odinger's
equation gives
\[
 \dot x=-igy,\qquad
 \dot y=-igx-i\sum_{a,j}\overline{\kappa}_{aj}z_{aj},\qquad
 \dot z_{aj}=-i\omega_{aj}z_{aj}-i\kappa_{aj}y.
\]
Solving the last equation with its stated initial condition yields
\begin{equation}
 z_{aj}(t)=-i\kappa_{aj}\int_0^t
            e^{-i\omega_{aj}(t-s)}y(s)\,ds,
 \qquad
 \dot y(t)=-igx(t)-\int_0^t\Sigma(t-s)y(s)\,ds,
 \label{bench:reservoir-kernel}
\end{equation}
where
\begin{equation}
 \Sigma(v)=\sum_{a,j}|\kappa_{aj}|^2e^{-i\omega_{aj}v},
 \qquad
 \Sigma(v)=\sum_{a=R,L}\int_{I_a}|\kappa_a(\omega)|^2
                                      e^{-i\omega v}\,d\omega
 \quad\hbox{in the continuum}.
 \label{bench:reservoir-spectrum}
\end{equation}
The continuum formula follows by the same variation-of-constants
argument. The $L^2$ coupling assumption makes $|\kappa_a|^2$ integrable
and justifies these finite-time integrals.

Put $p_a(t)=\sum_j|z_{aj}(t)|^2$, or its continuum integral. Unitarity
gives $|x|^2+|y|^2+p_R+p_L=1$. With $\Sigma_a$ denoting one channel's
kernel, its exact flux is
\begin{equation}
 \dot p_a(t)=2\operatorname{Re}\left[
  \overline{y(t)}\int_0^t\Sigma_a(t-s)y(s)\,ds\right].
 \label{bench:reservoir-flux}
\end{equation}
It need not be positive: products can return. A finite Hamiltonian has
recurrent unitary evolution, and neither an exponential survival law nor
an irreversible acquisition clock follows from
\eqref{bench:reservoir-kernel}.

\paragraph{The retained branch and the operational null.}
Let $|Z_a(t)\rangle$ denote the complete product wave in channel $a$,
including its mode amplitudes, and suppress the unchanged factor $M_0$.
The record-product projection is exactly
\begin{equation}
 \Pi_R\Psi(t)=|q,A_0\rangle\otimes|Z_R(t)\rangle.
 \label{bench:reservoir-recordbranch}
\end{equation}
Thus an admitted physical configuration readout of this sector, together
with the complete equilibrium/equivariance premises, selects source
factor $q$ at this time. The unitary Hamiltonian alone does not select an
actual sector. It also has not written $M_0$: product occupation is not
automatically a protected acquired memory. Neither this endpoint
factorization nor its label guarantees source factor $q$ after later
source interactions; their full continuation must be propagated as
qualified below.

For clarity, now admit an endpoint readout at time $t$ that resolves the
$R$ sector against its complement. This access premise defines the
following ``no readable record'' outcome $N$; it does not assert that no
record-product entry ever occurred. Its complete unnormalized projected
component is
\begin{equation}
 |\Phi_N(t)\rangle
 =|\phi_0(t)\rangle\otimes|\mathrm{vac}\rangle
       +|q,A_0\rangle\otimes|Z_L(t)\rangle,
 \qquad |\phi_0(t)\rangle=x(t)|r,A_0\rangle+y(t)|q,A^*\rangle.
 \label{bench:reservoir-nullcomplete}
\end{equation}
Tracing the unobserved field gives the generally mixed
source--receptor state
\begin{equation}
 \widetilde\rho_N^{SA}(t)
 =|\phi_0(t)\rangle\langle\phi_0(t)|
       +p_L(t)|q,A_0\rangle\langle q,A_0|,
 \qquad
 \rho_N^{SA}(t)=\frac{\widetilde\rho_N^{SA}(t)}{1-p_R(t)}
 \quad (p_R(t)<1).
 \label{bench:reservoir-nullmixed}
\end{equation}
The vacuum--loss cross terms vanish in this partial trace by field
orthogonality; they remain in \eqref{bench:reservoir-nullcomplete}.
These projected components are autonomous conditional instrument states
only with an admitted projective extraction or dynamically separated
physical pointer outcomes. Conditioning an actual configuration alone
does not remove the unoccupied global wave. If $R$ and $N$ can later
recombine, propagate the full original wave together with the actual
conditioning. Even within a separated $N$ continuation, future return of
the hidden loss field requires the complete component
\eqref{bench:reservoir-nullcomplete}, not merely its partial trace. The
pure no-product amplitude $\phi_0$ cannot replace an operational null
containing hidden loss. A null defined by absence of retained memory
acquisition is a different event and requires the corresponding memory
dynamics.

\paragraph{A controlled convolution limit, with its scope.}
An explicit continuum idealization makes one Markov limit provable.
Take $I_R=I_L=\mathbb R$, bandwidth $\Lambda>0$, and rates
$\gamma_R=\Gamma\ge0$, $\gamma_L=\ell\ge0$, with
\begin{equation}
 |\kappa_{a,\Lambda}(\omega)|^2
   =\frac{\gamma_a}{2\pi}\frac{\Lambda^2}{\omega^2+\Lambda^2},
 \qquad
 \Sigma_{a,\Lambda}(v)=\frac{\gamma_a\Lambda}{2}e^{-\Lambda v}
 \quad(v\ge0).
 \label{bench:reservoir-lorentzian}
\end{equation}
The second identity is the Fourier transform of the displayed
Lorentzian. Every finite $\Lambda$ has an $L^2$ form factor and the
self-adjoint Hamiltonian above. Its two-sided detuning spectrum is
unbounded below in this rotating-frame description. It is an explicit
wide-band mathematical idealization, not a lower-bounded material bath
construction.

Let $k=\Gamma+\ell$ and let $y_\Lambda$ be the exact continuum solution.
Norm conservation and the nonnegative exponential kernel give
$|\dot y_\Lambda|\le g+k/2$. Integration by parts, using
$y_\Lambda(0)=0$, therefore proves
\begin{align}
 r_{a,\Lambda}(t)
 &:=\int_0^t\Sigma_{a,\Lambda}(t-s)y_\Lambda(s)\,ds
                          -\frac{\gamma_a}{2}y_\Lambda(t)\notag\\
 &=-\frac{\gamma_a}{2}\int_0^t
              e^{-\Lambda(t-s)}\dot y_\Lambda(s)\,ds,
 \qquad
 |r_{a,\Lambda}(t)|\le
       \frac{\gamma_a(g+k/2)}{2\Lambda}.
 \label{bench:reservoir-residual}
\end{align}
Let $(\bar x,\bar y)$ solve the Markov equations
\[
 \dot{\bar x}=-ig\bar y,\qquad
 \dot{\bar y}=-ig\bar x-\frac{k}{2}\bar y,
 \qquad (\bar x(0),\bar y(0))=(1,0).
\]
Their $2\times2$ propagator is a contraction since the squared norm has
derivative $-k|\bar y|^2$. Duhamel's formula and
\eqref{bench:reservoir-residual} imply the finite-horizon bound
\begin{equation}
 \sup_{0\le t\le T}
 \left\|(x_\Lambda(t),y_\Lambda(t))-(\bar x(t),\bar y(t))\right\|_2
 \le \frac{kT(g+k/2)}{2\Lambda}.
 \label{bench:reservoir-markovbound}
\end{equation}
This controls the two no-product amplitudes. It does not compare complete
emitted field states or derive a microscopic event generator. For other
spectra, a claimed limit
\[
 \int_0^t\Sigma(t-s)y(s)\,ds
       \longrightarrow \left(\frac{\Gamma+\ell}{2}+i\Delta\right)y(t)
\]
requires its own approximation theorem or an explicit premise controlling
the integrated convolution residual. The same contraction argument then
applies when $\Delta$ is real. Stating the convolution itself fixes the
normalization without a half-delta convention.

\section{Weak terminal protection of the three-state recorder}
\label{bench:protection}

The three-state calculation behind \eqref{cfg:threebound} has a useful
quantitative protection benchmark \cite[\S12]{C02}. Let
\[
 e=|1,M_0\rangle,\qquad p=|R,M_0\rangle,\qquad m=|R,M_1\rangle,
 \qquad
 \frac{H_3}{\hbar}
 =g(|p\rangle\langle e|+|e\rangle\langle p|)
  +\chi(|m\rangle\langle p|+|p\rangle\langle m|),
\]
where $g,\chi\ge0$ and $\Omega_3=\sqrt{g^2+\chi^2}>0$.
Starting from $e$, the unprotected amplitudes are
\begin{equation}
 a_0(t)=\frac{\chi^2+g^2\cos\Omega_3t}{\Omega_3^2},\qquad
 b_0(t)=-i\frac{g}{\Omega_3}\sin\Omega_3t,\qquad
 c_0(t)=\frac{g\chi}{\Omega_3^2}(\cos\Omega_3t-1).
 \label{bench:protection-unperturbed}
\end{equation}

Admit an additional absorbing protection instrument with jump operator
$C=\sqrt\gamma|M\rangle\langle m|$, $\gamma\ge0$, where $M$ is an
orthogonal terminal state with no outgoing channel. This stochastic
instrument, including its conditional wave law, is a new primitive in
this benchmark. It is not derived by merely adding an unitarily coupled
product mode, and no claim is made here to derive it from the preceding
reservoir Hamiltonian.

The no-protection wave $v_\gamma=a_\gamma e+b_\gamma p+c_\gamma m$
obeys
\begin{equation}
 \dot a_\gamma=-igb_\gamma,\qquad
 \dot b_\gamma=-iga_\gamma-i\chi c_\gamma,\qquad
 \dot c_\gamma=-i\chi b_\gamma-\frac\gamma2c_\gamma,
 \qquad v_\gamma(0)=e.
 \label{bench:protection-null}
\end{equation}
The admitted instrument gives exactly
\begin{equation}
 P_M(t)=\gamma\int_0^t|c_\gamma(s)|^2\,ds
       =1-\|v_\gamma(t)\|^2,
 \qquad
 \rho(t)=|v_\gamma(t)\rangle\langle v_\gamma(t)|
                          +P_M(t)|M\rangle\langle M|.
 \label{bench:protection-exact}
\end{equation}
The norm identity follows directly from
\eqref{bench:protection-null}. It displays both the unfinished coherent
branch and the terminal branch. When interpreted by actual local clocks,
the same expression requires matching killed occupations and the
compatible damped wave law; the clock $\gamma\mathbf1_{\{Q=m\}}$
alone does not establish that compatibility.

\begin{proposition}[One-cycle protection with a relative error bound]
\label{bench:protection-bound}
Put $T_3=2\pi/\Omega_3$ and
$P_0=3\pi\gamma g^2\chi^2/\Omega_3^5$. For the admitted absorbing
instrument,
\begin{equation}
 P_M(T_3)=P_0+R,\qquad
 |R|\le P_0\left[
  \frac{2\pi}{3}\frac\gamma{\Omega_3}
  +\left(\frac{2\pi^2}{9}+\frac5{12}\right)
                        \left(\frac\gamma{\Omega_3}\right)^2\right].
 \label{bench:protection-relative}
\end{equation}
The bound holds for every $\gamma\ge0$; its first-order use requires
$\gamma/\Omega_3\ll1$. If $\gamma g\chi=0$, both probabilities in
the comparison vanish exactly.
\end{proposition}

\begin{proof}
Write $B=g\chi/\Omega_3^2$ and
$V_\gamma(t)=\exp[(-iH_3/\hbar-\gamma|m\rangle\langle m|/2)t]$.
The norm derivative in \eqref{bench:protection-exact}, applied to any
initial vector, proves $\|V_\gamma(t)\|\le1$. Duhamel's formula in the
order using the damped propagator on the left gives
\[
 c_\gamma(t)-c_0(t)
 =-\frac\gamma2\int_0^t
     \langle m|V_\gamma(t-s)|m\rangle c_0(s)\,ds.
\]
Since $|c_0(s)|=B(1-\cos\Omega_3s)$, it follows that
\begin{equation}
 |c_\gamma(t)-c_0(t)|\le\frac{\gamma B}{2}F(t),
 \qquad F(t)=t-\frac{\sin\Omega_3t}{\Omega_3}.
 \label{bench:protection-duhamel}
\end{equation}
In particular this estimate retains the factor $B$ suppressed by strong
memory coupling. The unprotected integral is
\[
 \gamma\int_0^{T_3}|c_0(t)|^2\,dt
  =\frac{\gamma B^2}{\Omega_3}
        \int_0^{2\pi}(1-\cos u)^2\,du
  =\frac{3\pi\gamma B^2}{\Omega_3}=P_0.
\]
Using $\bigl||c_\gamma|^2-|c_0|^2\bigr|
\le2|c_0||c_\gamma-c_0|+|c_\gamma-c_0|^2$ and
\eqref{bench:protection-duhamel} gives
\[
 |R|\le\frac{\gamma^2B^2T_3^2}{2}
       +\frac{\gamma^3B^2}{4}
            \left(\frac{T_3^3}{3}+\frac{5T_3}{2\Omega_3^2}\right).
\]
Here the first integral is
$\int_0^{T_3}(1-\cos\Omega_3t)F(t)\,dt=T_3^2/2$, since
$F'=1-\cos\Omega_3t$, and direct integration gives
$\int_0^{T_3}F(t)^2dt=T_3^3/3+5T_3/(2\Omega_3^2)$.
Substitution of $T_3=2\pi/\Omega_3$ proves
\eqref{bench:protection-relative}, including its zero cases without
division by $P_0$.
\end{proof}

For fixed $g>0$ and $\gamma>0$, the proposition proves the genuine
large-$\chi$ asymptotic
\begin{equation}
 P_M(2\pi/\Omega_3)
 =\frac{3\pi\gamma g^2}{\chi^3}
   \left[1+O\!\left(\frac\gamma\chi+\frac{g^2}{\chi^2}\right)\right],
 \qquad \chi\longrightarrow\infty.
 \label{bench:protection-asymptotic}
\end{equation}
The absolute damping error is $O(\chi^{-4})$ at these fixed parameters,
so it cannot overwhelm the $\chi^{-3}$ leading term. This is a
one-undamped-cycle horizon, which itself decreases as the coupling
increases. It is neither an exact use of $c_0$ in the protected model nor
a uniform assertion over arbitrary simultaneous scalings of protection,
coupling and observation time. The trapping tradeoff in
\eqref{cfg:trap} remains conditional on its separate architecture
premises.

% End consolidated source: parts/13_reservoir_benchmark.tex


% Begin consolidated source: parts/14_scalar_benchmarks.tex
\section{Missed absorption and a later rotated probe}
\label{bench:inefficient}

This finite scalar benchmark \cite{C01} illustrates why a null must retain
unobserved transitions. It assumes the amplitude-damping jump instrument;
it does not derive its statistical law. Let
$L=\sqrt\gamma\,|g\rangle\langle e|$, let $0<\eta<1$ be the fixed
recording efficiency, and start in $|e\rangle$. A jump is recorded with
probability $\eta$ or transferred to a distinct unobserved loss register
with probability $1-\eta$. No further source drive acts during an exposure.
All recording sites are fresh and loss products cannot return during the
specified two-exposure test.

Put $u=e^{-\gamma T}$. The no-jump, missed-jump and recorded-jump
contributions after duration $T$ are respectively
\[
 u|e\rangle\langle e|,\qquad
 (1-\eta)(1-u)|g\rangle\langle g|,\qquad
 \eta(1-u)|g\rangle\langle g|.
\]
These follow by integrating the first-decay density $\gamma e^{-\gamma t}$
and assigning the two admitted acquisition channels. Hence the
unnormalized no-readable-record source state is
\begin{equation}
 \widetilde\rho_{\varnothing}
 =u|e\rangle\langle e|+(1-\eta)(1-u)|g\rangle\langle g|,
 \qquad p_{\varnothing}=1-\eta(1-u).
 \label{bench:inefficient-null}
\end{equation}
In the complete retained description the two null contributions carry
different vacuum/loss flags, and a returning loss register must not be
discarded. Equation~\eqref{bench:inefficient-null} is their source marginal.

Next use a supplied unitary with
$U|e\rangle=c|e\rangle+s|g\rangle$ and
$U|g\rangle=-s|e\rangle+c|g\rangle$, where $c,s$ are real and
$c^2+s^2=1$. The unnormalized excited weight becomes
$uc^2+(1-\eta)(1-u)s^2$. A second fresh exposure of duration $\tau$
therefore gives the joint history probability
\begin{equation}
 P(\varnothing_1,R_2)=\eta(1-e^{-\gamma\tau})
 \bigl[uc^2+(1-\eta)(1-u)s^2\bigr].
 \label{bench:inefficient-second}
\end{equation}
Division by $p_{\varnothing}$ gives its conditional version; the second
null has joint mass $p_{\varnothing}-P(\varnothing_1,R_2)$ and retains
both old loss and new unresolved/missed branches. For $c=0$, the entire
second-click contribution comes from the formerly missed decays. Replacing
the first null by an attenuated $|e\rangle$ would predict zero instead.

\section{A binary pulse with correct endpoints and excess actual jumps}
\label{bench:binary-surplus}

The stationary examples in Chapter~\ref{stat:current-chapter} already prove
the stronger general distinctions. The following nonstationary two-state
calculation preserves a useful exact checkpoint test \cite{C02,C03}.
Let $H=\hbar\chi\sigma_x$, $\chi>0$, and start with wave and actual
configuration $|0\rangle$. Up to $T=\pi/(2\chi)$,
\[
 \psi_t=\cos(\chi t)|0\rangle-i\sin(\chi t)|1\rangle,
 \qquad J_{10}(t)=\chi\sin(2\chi t)\geq0.
\]
Choose a fixed surplus parameter $\zeta\geq0$ and set
$K_{01}=\zeta J_{10}$. On the open interval $(0,T)$ the Markov rates are
\begin{equation}
 \lambda_{10}=2(1+\zeta)\chi\tan(\chi t),\qquad
 \lambda_{01}=2\zeta\chi\cot(\chi t).
 \label{bench:binary-rates}
\end{equation}
Their forward equation is solved by
$p_1(t)=\sin^2(\chi t)$: the net inflow is
$(1+\zeta)J_{10}-\zeta J_{10}=J_{10}=\dot p_1$.
Despite the nodal conditional rates, the occupation-weighted total
activity is $(1+2\zeta)J_{10}$, and
\begin{equation}
 \mathbb E N_{[0,T]}=1+2\zeta<\infty.
 \label{bench:binary-count}
\end{equation}
For completeness, construct from the definite sector $0$ at time zero
using its locally integrable outward rate, then use the regular jump
construction between each pair of interior times. The first jump occurs
strictly after zero, so there is no accumulation of jumps at zero.
The difference of two solutions of the scalar forward equation on
$(0,T)$ is
$C\cos^{2(1+\zeta)}(\chi t)\sin^{-2\zeta}(\chi t)$.
Boundedness at zero forces $C=0$ when $\zeta>0$, and the initial value
does so when $\zeta=0$. Thus the displayed population solution is the
entrance law of this construction. There is no interior explosion
because the rates are bounded on each compact subinterval. Taking
limits in the expected compensated jump counts gives
\eqref{bench:binary-count}; finite expected activity excludes infinitely
many jumps accumulating at $T$. The limiting state at $T$ is $1$ almost
surely. Rates assigned to unoccupied endpoint nodes have no effect.

Until the first jump the actual state is $0$, so the exact survival is
\begin{equation}
 P(\tau_1>t)=\exp\!\left[-\int_0^t
                 2(1+\zeta)\chi\tan(\chi s)ds\right]
             =\cos^{2(1+\zeta)}(\chi t).
 \label{bench:binary-first}
\end{equation}
Thus the endpoint law is independent of $\zeta$ while the first event
and the expected number of events are not. These are native path
predictions. Access to that first-event time still requires a physical
reporter; an endpoint pointer alone does not measure it. Conversely, an
admitted neutral reporter with finite positive response must be treated
using the delay and competition calculation in
Section~\ref{bench:history}, rather than identifying its latch with the
native jump.

% End consolidated source: parts/14_scalar_benchmarks.tex


% End consolidated source: parts/12_benchmarks.tex

\chapter{Smooth autonomous realization of isolated pilot contacts}\label{p:smooth}
\label{p:mech-section}

This chapter supplies an explicit Hamiltonian for a finite family of reversible
classical contact updates. Its scope is the contact module: packet births,
continuous pilot-field evolution and any asynchronous exporter remain separate
laws unless their simultaneous coupling is independently supplied. No Bell
intensity is used in this Hamiltonian.

\section{Register cells and an explicit permutation compiler}

Let the finite register have values $i\in\{1,\ldots,K\}$ and let each contact mark
$c\in\mathcal C$ specify a permutation $\pi_c$ of those values. A register value
may encode a finite product of carrier, packet-slot, eligibility, resource and
retained-record registers. An irreversible update must first be extended
injectively by retaining its overwritten value and a finite ready stock.
The permutations of that enlarged register are the input to this construction.

Use canonical coordinates $q,p\in\mathbb R^2$, put $Q_i=(id,0)$, and prepare
\begin{equation}\label{p:mech-cells}
 |q-Q_i|\le r_0,\qquad |p|\le p_*.
\end{equation}
For each $c$, choose smooth tracks $\Gamma_{c,i}:[0,1]\to\mathbb R^2$ with
\begin{equation}\label{p:mech-tracks}
 \Gamma_{c,i}(0)=Q_i,\qquad
 \Gamma_{c,i}(1)=Q_{\pi_c(i)},\qquad
 |\Gamma_{c,i}(s)-\Gamma_{c,j}(s)|\ge a>0\quad(i\ne j).
\end{equation}
There is an elementary explicit compiler. In the first third of the programme,
lift $(id,0)$ to $(id,ih)$; in the middle third translate it to
$(\pi_c(i)d,ih)$; in the last third lower it to $(\pi_c(i)d,0)$.
Use a smooth increasing interpolation flat to all orders at the phase endpoints.
The horizontal coordinates give separation at least $d$ during the lifts and
descents, and the distinct heights give separation at least $h$ during the
horizontal transport. Thus $a=\min(d,h)$ works. Extend the tracks constantly
outside $[0,1]$.

Choose
\[
 0<r_0<r_{\rm core}<r_{\rm supp}<a/2
\]
and a smooth cutoff $\chi$ equal to one on $|z|\le r_{\rm core}$ and zero on
$|z|\ge r_{\rm supp}$. For prepared beam speed $v>0$ and gate duration
$0<\delta<T$, define
\begin{equation}\label{p:mech-field}
 b_c(x,q)=\frac1{v\delta}\sum_{i=1}^K
 \Gamma'_{c,i}\!\left(\frac{x}{v\delta}\right)
 \chi\!\left(q-\Gamma_{c,i}\!\left(\frac{x}{v\delta}\right)\right).
\end{equation}
This field is smooth and vanishes outside $0<x<v\delta$. At each phase its
cutoff supports are disjoint. In the core of track $i$,
$b_c=\Gamma'_{c,i}/(v\delta)$ and $\partial_qb_c=0$.

\section{Positive kinetic coupling and exact passage time}

First regard $c_j$ as particle $j$'s fixed channel. Let $x_j,P_j$ be its
longitudinal canonical coordinates, $m_b>0$ its finite mass parameter, and
$\mu>0$ the register mass parameter. Take
\begin{equation}\label{p:mech-H}
 H_{\rm gate}=
 \sum_{j=1}^M\frac{[P_j+b_{c_j}(x_j,q)\cdot p]^2}{2m_b}
 +\frac{|p|^2}{2\mu}+V(q).
\end{equation}
Here $V\ge0$ is smooth and zero on a region containing every track and its
support. One may take $V=0$ globally for the finite-horizon theorem.
All coefficients are fixed functions of position; there is no external
time-dependent drive.

\begin{proposition}[Autonomous nonnegative contact Hamiltonian]
\label{p:mech-hamiltonian}
For every finite parameter choice, \eqref{p:mech-H} is smooth, autonomous,
nonnegative, and has a global classical Hamiltonian flow when $V=0$.
It retains the register and all incoming and outgoing beam particles.
Its interaction is an explicitly declared positive position-dependent kinetic
metric, rather than an ordinary scalar-potential impact.
\end{proposition}
\begin{proof}
Write $k_j=P_j+b_{c_j}\cdot p$. At each configuration the triangular map
$(P_1,\ldots,P_M,p)\mapsto(k_1,\ldots,k_M,p)$ is invertible.
The fields and their derivatives are globally bounded at fixed parameters,
so this positive quadratic kinetic form is uniformly positive definite for
that parameter choice. Its ellipticity constant may depend on $\delta,M$
and the programme.

Conserved energy bounds $p$, every $k_j$, and hence every $P_j$.
Hamilton's velocities and forces are then bounded at that energy because
the field derivatives are bounded. Local smooth flow therefore cannot
escape to infinity in finite time. Nonnegativity follows directly from the
squares. No particle coordinate has been removed.
\end{proof}

\begin{lemma}[Exact separated gate on a robust core]
\label{p:mech-exactgate}
Suppose particle $j$ enters at time $t_{\rm in}$ with $x_j=0$,
$P_j=m_bv$, mark $c$, and the register in the core of track $i$.
Suppose no other particle is in an interaction region.
While the core condition holds, its exit time and register evolution are
\begin{align}
 x_j(t)&=v(t-t_{\rm in}),\qquad
 t_{\rm out}=t_{\rm in}+\delta,\label{p:mech-time}\\
 p(t)&=p_{\rm in},\nonumber\\
 q(t)&=\Gamma_{c,i}\!\left(\frac{t-t_{\rm in}}{\delta}\right)
       +(q_{\rm in}-Q_i)
       +\frac{p_{\rm in}}{\mu}(t-t_{\rm in}).\label{p:mech-transport}
\end{align}
At exit $P_j=m_bv$ and the register is in the output cell
$Q_{\pi_c(i)}$ up to the displayed offset and free drift.
\end{lemma}
\begin{proof}
On the core, $\partial_qb_c=0$ and $\nabla V=0$, so $\dot p=0$.
For the active particle, Hamilton's equations give
\[
 \dot x_j=k_j/m_b,\qquad
 \dot P_j=-(k_j/m_b)(\partial_xb_c)\cdot p,\qquad
 \dot q=p/\mu+(k_j/m_b)b_c.
\]
Since $\dot p=0$ and $\partial_qb_c=0$,
\[
 \frac{d}{dt}(b_c\cdot p)
       =(k_j/m_b)(\partial_xb_c)\cdot p.
\]
Hence $\dot k_j=0$. At entry $b_c=0$, so $k_j=m_bv$ throughout.
This proves \eqref{p:mech-time}; inserting \eqref{p:mech-field} into the
register equation gives \eqref{p:mech-transport}.
At exit the field vanishes again, so $P_j=k_j=m_bv$.
\end{proof}

\begin{corollary}[Uniform composition through separated contacts]
\label{p:mech-robust}
Assume \eqref{p:mech-cells} and
\begin{equation}\label{p:mech-margin}
 r_0+\frac{p_*(T+\delta)}{\mu}<r_{\rm core}.
\end{equation}
Every sequence of nonoverlapping contacts implements the specified register
permutations correctly through time $T$, including every intermediate track
stage. This conclusion is uniform over the compact preparation domain and
does not require a large incident mass.
\end{corollary}
\begin{proof}
Each successful gate preserves $p$ and translates the cell offset without
rotating or amplifying it. Idle motion contributes the same free drift.
Thus the offset at time $t$ from the prescribed idle cell or active track
is at most $r_0+p_*t/\mu$. Inequality \eqref{p:mech-margin} prevents a first
exit from a track core and closes the bootstrap used in
Lemma~\ref{p:mech-exactgate}.
\end{proof}

The energy before and after every separated gate is
\begin{equation}\label{p:mech-energy}
 E=\frac{Mm_bv^2}{2}+\frac{|p_{\rm in}|^2}{2\mu}<\infty.
\end{equation}
During a gate the canonical momentum changes by $-b_c\cdot p$ while its
kinetic momentum remains $m_bv$. This is the backreaction in the specified
inertial interaction. The outgoing particles continue freely and are retained;
on the good domain a finite receiver region of length greater than $vT$
suffices for the horizon. The incoming spatial stock is not reset.
The metric, its growing derivatives as $\delta\to0$, the finite geometry and
the register mass are declared resources. In particular, this is not a
fixed-resource limit or a derivation of these interactions from a
diagonal-mass scalar-potential collision model.

\section{Present-state logical decoding and path topology}

During a finite gate, a quantizer of the bare coordinate $q$ may report a
transit value or several intermediate values. Those paths need not approach
a direct discrete jump in the $J_1$ Skorokhod topology. A nonzero-duration
excursion through a third discrete value cannot be erased by a continuous
time change.

The output used below is the \emph{committed register value}, defined as a
fixed function of the present enlarged physical state. When idle, decode the
unique nearby $Q_i$. When one particle of mark $c$ is inside, its position
gives $s=x/(v\delta)$; decode the unique nearby track $\Gamma_{c,i}(s)$ and
report input value $i$ until exit. At $x=v\delta$, decode the output idle
cell and report $\pi_c(i)$. The separation and margin conditions make this
definition unambiguous. It uses the current busy-particle coordinate, not an
external clock or an unretained past. A separate persistent logical record,
if required, must be included in the finite register and permutation library.
It is not supplied merely by naming the decoder.

For path comparison, at the first overlap or transition-strip encounter
stop this decoder and set its value to an absorbing cemetery symbol;
stop the completed-contact log there as well, retaining a failure flag.
Before that time there are only finitely many successful contacts, so
the extended logical path is c\`adl\`ag. This is a convention for the
compared output on exceptional trajectories, not a change to the
Hamiltonian flow, which continues. An arbitrary Borel decoder on bad
phase-space states would not by itself ensure a c\`adl\`ag path.

On separated contacts the resulting path is exactly the specified register
circuit with each contact delayed by $\delta$. If incident times are $t_j$,
there is no contact in $(T-\delta,T]$, and the order is unchanged, a
piecewise linear time change aligning $t_j$ to $t_j+\delta$ gives
\begin{equation}\label{p:mech-J1}
 d_{J_1}(S^{\rm committed},S^{\rm instantaneous})\le\delta.
\end{equation}
One may align null contacts as well. No analogous claim is made for the
unprocessed continuous-coordinate path.

\section{Smooth channel plates and the projected history bound}

Prepare independent uniform incoming positions on $[-L,0]$, where
\[
 L=Mv/R,\qquad M/R>T,
\]
with longitudinal momenta $m_bv$. Independently prepare uniform transverse
coordinates in a finite aperture, with zero transverse momentum.
Partition the aperture into mark areas of fractions $\pi_c$.
Choose smooth nonnegative channel functions $\zeta_c(y)$ equal to one on
the respective interior plateaus, zero on other plateaus, and satisfying
$\sum_c\zeta_c\le1$. Let the total transition-strip area fraction be at most
$\eta$. Replace $b_{c_j}$ by
\[
 b(x_j,y_j,q)=\sum_c\zeta_c(y_j)b_c(x_j,q)
\]
and add free transverse kinetic energies. On a plateau the transverse
derivatives vanish, so the mark stays fixed; globally the Hamiltonian is
still smooth and nonnegative. A strip encounter is assigned to the failure
event rather than asserted to perform an ideal gate.

Before entry the field is zero, so incident times are exactly
$T_j=-x_j(0)/v$, independent uniform variables on $[0,M/R]$.
Their independent marks come from the transverse geometry. Initial register
state and offset are independent of the beam ensemble. No draw occurs at
a contact.

\begin{theorem}[Full projected logical-history comparison for the contact module]
\label{p:mech-history}
Assume the finite permutation library, independent incoming beam ensemble,
smooth plateau construction, and margin condition
\eqref{p:mech-margin}. Assume that no other interaction changes the register
during this module. Let $\widehat S^\delta$ be the committed register path
and $\widehat\Xi^\delta$ its completed-contact history, using the absorbing
cemetery and stopped-log convention after a first overlap or strip encounter.
Let $(S^{\rm P},\Xi)$ be the same permutation circuit driven instantaneously
by a marked Poisson process of intensity $R\dd t\,\pi_c$ on $[0,T]$.
Then
\begin{equation}\label{p:mech-TV}
 \TV\!\left(
 \pLaw(\widehat S^\delta,\widehat\Xi^\delta),
 \pLaw(S^{\rm P},\Xi)\right)
 \le
 \frac{(RT)^2}{M}+R^2T\delta+R\delta+RT\eta.
\end{equation}
The compared space consists of projected logical/contact histories, including
any finite logical archives contained in the permutation library.
It does not contain all microscopic beam coordinates or continuous
phase-space paths.
\end{theorem}
\begin{proof}
For an unordered pair of independent incident times on $[0,M/R]$, the
area in which both lie in $[0,T]$ and differ by at most $\delta$ is at most
$2T\delta$. A union bound therefore gives
\[
 \mathbb P(\text{overlap})\le
 {M\choose2}\frac{2T\delta}{(M/R)^2}\le R^2T\delta.
\]
The expected number of strip encounters is $RT\eta$, so their probability
is at most that quantity. A first-failure argument is sufficient: until
the first such event, all previous gates are exact and the still-incoming
particles are free.

Define the fictitious completed process
\[
 \Xi_M^\delta=\sum_{j:T_j+\delta\le T}
                 \delta_{(T_j+\delta,C_j)}.
\]
Its count is $\operatorname{Bin}(M,R(T-\delta)/M)$.
Conditional on the count, the times are independent uniform on $[\delta,T]$
and the marks have law $\pi$. This is the same conditional kernel as the
Poisson process restricted to $[\delta,T]$.
The elementary Bernoulli--Poisson coupling gives distance at most
$[R(T-\delta)]^2/M$. Adding the independent Poisson points in $[0,\delta)$
costs at most $1-e^{-R\delta}\le R\delta$.

On the good event, Lemma~\ref{p:mech-exactgate} and
Corollary~\ref{p:mech-robust} identify the actual committed circuit with the
same causal permutation circuit applied to $\Xi_M^\delta$.
A common measurable circuit cannot increase total variation.
Add the two failure probabilities and bound $T-\delta\le T$ to obtain
\eqref{p:mech-TV}.
\end{proof}

For fixed programme and horizon, the bound tends to zero if
\[
 (RT)^2/M\to0,\qquad
 R^2T\delta\to0,\qquad R\delta\to0,\qquad RT\eta\to0.
\]
These conditions also state the required scales if an outer construction
uses a varying attempt rate $R=R_N$. Every finite member has a finite
particle stock, finite energy and a smooth nonnegative Hamiltonian.
Conditioning on all initial particle positions instead makes the history
deterministic; the theorem does not assert a Poisson intensity under that
enlarged microscopic filtration.

\begin{remark}[Exporters and other asynchronous changes remain separate]
\label{p:mech-scope}
Theorem~\ref{p:mech-history} is a supplementary smooth realization of the
finite contact module. A packet exporter that creates, cancels or changes
a queue while a cell is being transported is not covered by its hypotheses.
An arbitrary finite transition can be permutation-lifted with a retained
receiver, but that algebraic fact does not establish that independently
scheduled or state-dependent transitions can interleave inside this gate
without changing its motion. If an overall theory retains such exports
as deterministic hybrid laws, they must be stated as such and given their
own composition analysis. No additional full-TV bound for interleaved
exporter events is claimed here.

The static field \eqref{p:mech-field} encodes the supplied reaction
permutations. It does not select directional packet exposure, exclude
other physical reaction channels, or derive their rates from a coherent
Hamiltonian edge current. Nor does the classical gate proof establish
quantum source/readout admission through an inaccessible reference.
Those obligations are logically distinct from the mechanical theorem
proved in this section.
\end{remark}




\backmatter

% Begin consolidated source: bibliography.tex
\begin{thebibliography}{P-Computer}

\addcontentsline{toc}{chapter}{Bibliography and source editions}

\bibitem[Pilot]{Pilot}
Jeremy Rodgers.
\emph{A Deterministic Pilot Medium: Bell Path Selection and Autonomous Material Records}.
Version 2, 15 September 2026. Zenodo preprint.
\href{https://doi.org/10.5281/zenodo.22774634}{doi:10.5281/zenodo.22774634}.

\bibitem[Massive]{MassiveCompletion}
Jeremy Rodgers.
\emph{A Massive Configuration Completion of the Quantum Measurement Programme}.
Version 2, 15 September 2026. Zenodo preprint.
\href{https://doi.org/10.5281/zenodo.22774739}{doi:10.5281/zenodo.22774739}.

\bibitem[M01]{M01}
Shadow Theory research programme.
\emph{Canonical Source Exchange and Bell Incidence A common charge interaction, a binary reaction mechanism, and a global path-law limit}.
Internal research manuscript, supplied corpus edition.
Source file: \path{Shadow_Source_Canonical_Exchange.tex}.

\bibitem[M02]{M02}
Shadow Theory research programme.
\emph{Selecting Source Incidence A cancellation mechanism, a global path-law limit, and the remaining physical freedom}.
Internal research manuscript, supplied corpus edition.
Source file: \path{Shadow_Source_Incidence_Selection.tex}.

\bibitem[M03]{M03}
Shadow Theory research programme.
\emph{Carrier Neutrality}.
Internal research manuscript, supplied corpus edition.
Source file: \path{Shadow_Source_Carrier_Neutrality.tex}.

\bibitem[M04]{M04}
Shadow Theory research programme.
\emph{Aperture Ownership}.
Internal research manuscript, supplied corpus edition.
Source file: \path{Shadow_Source_Aperture_Ownership.tex}.

\bibitem[M05]{M05}
Shadow Theory research programme.
\emph{Interface and Event Closure}.
Internal research manuscript, supplied corpus edition.
Source file: \path{Shadow_Source_Interface_and_Event_Closure.tex}.

\bibitem[M06]{M06}
Shadow Theory research programme.
\emph{Source Charge Balance and Actual Calibration A dynamical unification, finite-exposure limits, and the remaining causal constitution}.
Internal research manuscript, supplied corpus edition.
Source file: \path{Shadow_Source_Calibration_Dynamics_and_Constitutive_Core.tex}.

\bibitem[M07]{M07}
Shadow Theory research programme.
\emph{Attachment and Continuation}.
Internal research manuscript, supplied corpus edition.
Source file: \path{Shadow_Source_Attachment_and_Continuation.tex}.

\bibitem[M08]{M08}
Shadow Theory research programme.
\emph{Record Transduction and Innovation Taps A constructive output law, its unavoidable backaction, and the remaining source constitution}.
Internal research manuscript, supplied corpus edition.
Source file: \path{Shadow_Source_Record_Transduction_and_Innovation_Taps.tex}.

\bibitem[M09]{M09}
Shadow Theory research programme.
\emph{First Writing and Finite Causal Closure Source-contact regularity, latent records, and auxiliary-load stability}.
Internal research manuscript, supplied corpus edition.
Source file: \path{Shadow_Source_First_Writing_and_Finite_Causal_Closure.tex}.

\bibitem[M10]{M10}
Shadow Theory research programme.
\emph{Admission and Controlled Measurement Closure}.
Internal research manuscript, supplied corpus edition.
Source file: \path{Shadow_Source_Admission_and_Controlled_Measurement_Closure.tex}.

\bibitem[M11]{M11}
Shadow Theory research programme.
\emph{Shadow source production and measurement compatibility}.
Internal research manuscript, supplied corpus edition.
Source file: \path{Shadow_Source_Production_and_Measurement_Compatibility.tex}.

\bibitem[M12]{M12}
Shadow Theory research programme.
\emph{An integrated Shadow source measurement theory}.
Internal research manuscript, supplied corpus edition.
Source file: \path{Shadow_Source_Integrated_Measurement_Theory.tex}.

\bibitem[M13]{M13}
Shadow Theory research programme.
\emph{Selection of measurement currents from preparation consistency A conditional source/readout construction and an output-access bound}.
Internal research manuscript, supplied corpus edition.
Source file: \path{Shadow_Record_Law_Selection.tex}.

\bibitem[M14]{M14}
Shadow Theory research programme.
\emph{Finite absorbing tags and causal selection of preparation consistency A source interaction, a finite record theorem, and its admission boundary}.
Internal research manuscript, supplied corpus edition.
Source file: \path{Shadow_CPC_Physical_Selection.tex}.

\bibitem[M15]{M15}
Shadow Theory research programme.
\emph{Intrinsic outlet actualization: capture, access and causal rate selection A constructive replacement for the absorbing-threshold constitution}.
Internal research manuscript, supplied corpus edition.
Source file: \path{Shadow_Detector_Constitution.tex}.

\bibitem[M16]{M16}
Shadow Theory research programme.
\emph{Dark channels and joint event--continuation selection Source apertures, stochastic daughters and retained memories}.
Internal research manuscript, supplied corpus edition.
Source file: \path{Shadow_Joint_Event_Continuation.tex}.

\bibitem[M17]{M17}
Shadow Theory research programme.
\emph{Protecting the Shadow source aperture Finite-gap transport, retained memories and preselection stability}.
Internal research manuscript, supplied corpus edition.
Source file: \path{Shadow_Aperture_Protection.tex}.

\bibitem[M18]{M18}
Shadow Theory research programme.
\emph{A common interaction for Shadow detector records Finite arrival, conditional transport and physical archives}.
Internal research manuscript, supplied corpus edition.
Source file: \path{Shadow_Common_Source_Interaction.tex}.

\bibitem[M19]{M19}
Shadow Theory research programme.
\emph{Conditional preparation of Shadow detector statistics Deterministic extraction, exact clock heralding and returning-memory limits}.
Internal research manuscript, supplied corpus edition.
Source file: \path{Shadow_Apparatus_Equilibrium_Preparation.tex}.

\bibitem[M20]{M20}
Shadow Theory research programme.
\emph{Source-law selection across Shadow measurement architectures}.
Internal research manuscript, supplied corpus edition.
Source file: \path{Shadow_Source_Law_Reconciliation.tex}.

\bibitem[M21]{M21}
Shadow Theory research programme.
\emph{Dynamic carrier neutrality for canonical source events Complete-history bounds and a reciprocal readout obstruction Shadow Theory: Event-Law Bridge, Round 001}.
Internal research manuscript, supplied corpus edition.
Source file: \path{Shadow_Event_Law_Bridge_Round_001.tex}.

\bibitem[M22]{M22}
Shadow Theory research programme.
\emph{One-packet records and a coherent-ledger repair Access--reaction compatibility for the original Hamiltonian current}.
Internal research manuscript, supplied corpus edition.
Source file: \path{Shadow_Event_Law_Bridge_Round_002.tex}.

\bibitem[M23]{M23}
Shadow Theory research programme.
\emph{Reaction ownership and physical record access A native-reporter obstruction and a coherent contact derived from source chemistry}.
Internal research manuscript, supplied corpus edition.
Two distinct Round 003 sources are retained; this is the original reaction-reporter manuscript.
Source file: \path{Shadow_Event_Law_Bridge_Round_003.tex}.

\bibitem[M24]{M24}
Shadow Theory research programme.
\emph{A common material interface for source reactions and records Isoenergetic contact selection, a stirred-volume Bell limit, and the finite acquisition experiment}.
Internal research manuscript, supplied corpus edition.
This is the separate Physical Port manuscript.
Source file: \path{Shadow_Event_Law_Bridge_Round_003_Physical_Port.tex}.

\bibitem[M25]{M25}
Shadow Theory research programme.
\emph{An accounted continuous work port and an autonomous copied-return obstruction Canonical recoil, actual finite records, and the surviving event-law boundary}.
Internal research manuscript, supplied corpus edition.
Source file: \path{Shadow_Event_Law_Bridge_Round_004_Accounted_Work.tex}.

\bibitem[M26]{M26}
Shadow Theory research programme.
\emph{Event Bridge 005: Material Interaction}.
Internal research manuscript, supplied corpus edition.
Source file: \path{Shadow_Event_Law_Bridge_Round_005_Material_Interaction.tex}.

\bibitem[M27]{M27}
Shadow Theory research programme.
\emph{Record contacts from source-charge conversion Finite reaction closure, physical acquisition, and the surviving admission boundary Shadow Theory --- Event-Law Bridge, Round 006}.
Internal research manuscript, supplied corpus edition.
Source file: \path{Shadow_Event_Law_Bridge_Round_006_Contact_Closure.tex}.

\bibitem[M28]{M28}
Shadow Theory research programme.
\emph{Common exchange response and a retained null-record obstruction Shadow Theory --- Event-Law Bridge, Round 007}.
Internal research manuscript, supplied corpus edition.
Source file: \path{Shadow_Event_Law_Bridge_Round_007_Response_Selection.tex}.

\bibitem[M29]{M29}
Shadow Theory research programme.
\emph{Event Bridge 008: Common Interaction}.
Internal research manuscript, supplied corpus edition.
Source file: \path{Shadow_Event_Law_Bridge_Round_008_Common_Interaction.tex}.

\bibitem[M30]{M30}
Shadow Theory research programme.
\emph{Strengthening the Shadow Event-Law Proposal: Predictive Completion, Statistical Current Realization, and a Variational Selection of Bell Histories}.
Internal research manuscript, supplied corpus edition.
Source file: \path{Shadow_Event_Law_Foundational_Completion.tex}.

\bibitem[C01]{C01}
Shadow Theory research programme.
\emph{Checkpoint 01: finite response and delayed acquisition}.
Corrected cumulative research checkpoint, supplied corpus edition.
Source file: \path{Shadow_QM_Born_Event_Record_Physics_Checkpoint_2026-09-14.md}.

\bibitem[C02]{C02}
Shadow Theory research programme.
\emph{Checkpoint 02: coherent records and faithful copying}.
Corrected cumulative research checkpoint, supplied corpus edition.
Source file: \path{Shadow_QM_Born_Event_Record_Physics_Checkpoint_02_2026-09-14.md}.

\bibitem[C03]{C03}
Shadow Theory research programme.
\emph{Checkpoint 03: protection, coarse currents and MPBT}.
Corrected cumulative research checkpoint, supplied corpus edition.
Source file: \path{Shadow_QM_Born_Event_Record_Physics_Checkpoint_03_2026-09-15.md}.

\bibitem[F01]{F01}
Shadow Theory research programme.
\emph{Source--Readout Non-Equivalence: Descent and Equivariant Reconstruction Obstructions}.
Foundational Shadow Theory paper, supplied corpus edition.
Source file: \path{shadow_theory_paper_01.tex}.

\bibitem[F02]{F02}
Shadow Theory research programme.
\emph{Target-Relative Necessity of Completion: When Readout Loss Obstructs, and What a Sufficient Extension Must Retain}.
Foundational Shadow Theory paper, supplied corpus edition.
Source file: \path{paper02_final_v4.tex}.

\bibitem[F03]{F03}
Shadow Theory research programme.
\emph{Canonical Minimal Source Completion: The Coarsest Readout Extension on which a Nominated Family of Source Relations Becomes Well Defined}.
Foundational Shadow Theory paper, supplied corpus edition.
Source file: \path{shadow_theory_paper_03.tex}.

\bibitem[F04]{F04}
Shadow Theory research programme.
\emph{Geometric realization and source field equations}.
Foundational Shadow Theory paper, supplied corpus edition.
Source file: \path{shadow_theory_paper_04.tex}.

\bibitem[F05]{F05}
Shadow Theory research programme.
\emph{Observable Quotients and Exact Projected Dynamics: Closure, Memory, Minimal Dynamical Completion, and Effective Field Operators}.
Foundational Shadow Theory paper, supplied corpus edition.
Source file: \path{SHADOW_THEORY_PAPER_05.tex}.

\bibitem[F06]{F06}
Shadow Theory research programme.
\emph{Non-Source Projection and Internal Identifiability}.
Foundational Shadow Theory paper, supplied corpus edition.
Source file: \path{paper06.tex}.

\bibitem[F07]{F07}
Shadow Theory research programme.
\emph{Bulk-to-Brane Projection, Dynamical Nonclosure, and Observable Residues in Randall--Sundrum Gravity}.
Foundational Shadow Theory paper, supplied corpus edition.
Source file: \path{Paper 7 rs2_projection.tex}.

\bibitem[DGGTZ]{BellQFT}
D. D\"urr, S. Goldstein, R. Tumulka and N. Zangh\`i.
\emph{Bell-Type Quantum Field Theories}.
Journal of Physics A 38, R1--R43 (2005).
\url{https://arxiv.org/abs/quant-ph/0407116}.

\bibitem[BvHJ]{Filtering}
L. Bouten, R. van Handel and M. R. James.
\emph{An Introduction to Quantum Filtering}.
SIAM Journal on Control and Optimization 46, 2199--2241 (2007).
\url{https://arxiv.org/abs/math/0601741}.

\bibitem[BFG]{BFG}
L. Bertini, A. Faggionato and D. Gabrielli.
\emph{Flows, Currents, and Cycles for Markov Chains: Large Deviation Asymptotics}.
\url{https://arxiv.org/abs/1408.5477}.

\bibitem[ML]{MarvianLidar}
M. Marvian and D. A. Lidar.
\emph{Error Suppression for Hamiltonian-Based Quantum Computation Using Subsystem Codes}.
\url{https://arxiv.org/abs/1606.03795}.

\bibitem[VW]{ValentiniWestman}
A. Valentini and H. Westman.
\emph{Dynamical Origin of Quantum Probabilities}.
\url{https://arxiv.org/abs/quant-ph/0403034}.

\bibitem[P-Clock]{pChristandl2004}
M. Christandl, N. Datta, A. Ekert and A. J. Landahl,
``Perfect state transfer in quantum spin networks,''
\emph{Physical Review Letters} \textbf{92}, 187902 (2004).
\href{https://arxiv.org/abs/quant-ph/0309131}{arXiv:quant-ph/0309131};
\href{https://doi.org/10.1103/PhysRevLett.92.187902}{doi:10.1103/PhysRevLett.92.187902}.
The specific engineered spin-chain identification is equations (13)--(15).

\bibitem[P-Computer]{pFeynman1986}
R. P. Feynman, ``Quantum mechanical computers,''
\emph{Foundations of Physics} \textbf{16}, 507--531 (1986).
\href{https://doi.org/10.1007/BF01886518}{doi:10.1007/BF01886518}.
Earlier version in \emph{Optics News}, February 1985;
\href{https://www.cs.princeton.edu/courses/archive/fall05/frs119/papers/feynman85_optics_letters.pdf}{primary-paper scan}.

\bibitem[P-Audit]{pIndependentAudit}
\emph{Independent audit of the pilot-medium and massive-configuration
completions}, supplied AI-generated working audit, 15 September 2026.
Source: \path{Shadow_Event_Law_Independent_Audit.md}.
Section 6.6 is provenance for the conservative fixed-circuit cutoff
bookkeeping. This is supporting working material, not external validation;
the present text supplies the proof and its scope restrictions.

\bibitem[Bohm52a]{mc:BohmI}
D. Bohm,
A suggested interpretation of the quantum theory in terms of ``hidden'' variables. I,
\emph{Physical Review} \textbf{85}, 166--179 (1952).
\href{https://doi.org/10.1103/PhysRev.85.166}{doi:10.1103/PhysRev.85.166}.

\bibitem[Bohm52b]{mc:BohmII}
D. Bohm,
A suggested interpretation of the quantum theory in terms of ``hidden'' variables. II,
\emph{Physical Review} \textbf{85}, 180--193 (1952).
\href{https://doi.org/10.1103/PhysRev.85.180}{doi:10.1103/PhysRev.85.180}.

\bibitem[DGZ92]{mc:DGZeq}
D. D\"urr, S. Goldstein and N. Zangh\`i,
Quantum equilibrium and the origin of absolute uncertainty,
\emph{Journal of Statistical Physics} \textbf{67}, 843--907 (1992).
\href{https://arxiv.org/abs/quant-ph/0308039}{arXiv:quant-ph/0308039}.
In particular, the complete initial equilibrium measure and conditional-probability analysis;
no claim of global relaxation is imported.

\bibitem[DGZ04]{mc:DGZoperators}
D. D\"urr, S. Goldstein and N. Zangh\`i,
Quantum equilibrium and the role of operators as observables in quantum theory,
\emph{Journal of Statistical Physics} \textbf{116}, 959--1055 (2004).
\href{https://doi.org/10.1023/B:JOSS.0000037234.80916.d0}{doi:10.1023/B:JOSS.0000037234.80916.d0};
\href{https://arxiv.org/abs/quant-ph/0308038}{arXiv:quant-ph/0308038}.

\bibitem[TT05]{mc:TT}
S. Teufel and R. Tumulka,
Simple proof for global existence of Bohmian trajectories,
\emph{Communications in Mathematical Physics} \textbf{258}, 349--365 (2005).
\href{https://doi.org/10.1007/s00220-005-1302-0}{doi:10.1007/s00220-005-1302-0};
\href{https://arxiv.org/abs/math-ph/0406030}{arXiv:math-ph/0406030}.
The current-integrability assumptions and the spinor theorem are checked in
Chapter~\ref{mc:chapter-constitution}.

\bibitem[DG98]{mc:Deotto}
E. Deotto and G. C. Ghirardi,
Bohmian mechanics revisited,
\emph{Foundations of Physics} \textbf{28}, 1--30 (1998).
\href{https://doi.org/10.1023/A:1018752202576}{doi:10.1023/A:1018752202576};
\href{https://arxiv.org/abs/quant-ph/9704021}{arXiv:quant-ph/9704021}.

\bibitem[GTZ24]{mc:Arrival}
S. Goldstein, R. Tumulka and N. Zangh\`i,
Arrival times versus detection times,
\emph{Foundations of Physics} \textbf{54}, 63 (2024).
\href{https://doi.org/10.1007/s10701-024-00798-y}{doi:10.1007/s10701-024-00798-y};
\href{https://arxiv.org/abs/2405.04607}{arXiv:2405.04607}.

\end{thebibliography}

% End consolidated source: bibliography.tex

\end{document}
